\documentclass[uplatex, dvipdfmx, 11pt]{jsbook}

\usepackage[utf8]{inputenc}
\usepackage[margin=35truemm]{geometry} % 余白を35mmに設定
\usepackage{tikz-feynhand}
% --- 数学関連のパッケージ ---
\usepackage{amsmath}      % 様々な数式環境
\usepackage{amsthm}       % 定理環境
\usepackage{amsfonts}     % 数学用フォント
\usepackage{physics}      % 物理学の数式用マクロ (bra, ket, dvなど)
\usepackage{bm}           % 太字の数式 \bm{}
\usepackage{showkeys}
% --- 図や色、レイアウト関連のパッケージ ---
\usepackage{graphicx, color} % 図の挿入や色の利用
\usepackage{tikz}         % 図形描画
\usetikzlibrary{intersections,calc,arrows.meta,decorations.markings,decorations.pathreplacing}
% --- 図用のスタイル（辺の同一視の矢印など） ---
\tikzset{
  mid arrow/.style={postaction={decorate},decoration={markings,
      mark=at position .55 with {\arrow{Stealth[length=2.2mm]}}}},
  mid arrows/.style={postaction={decorate},decoration={markings,
      mark=at position .48 with {\arrow{Stealth[length=2.2mm]}},
      mark=at position .60 with {\arrow{Stealth[length=2.2mm]}}}},
  vtx/.style={circle,fill,inner sep=1.1pt},
  hvtx/.style={circle,draw,fill=white,inner sep=1.3pt},
}
% 図の下に付ける説明（solution 環境の中でも使えるよう float は使わない）
\newcommand{\zucap}[1]{\par\nopagebreak{\centering\footnotesize #1\par}\medskip}
\usepackage{float}        % 図表の位置を調整 [H]など
\usepackage{siunitx}      % 単位をきれいに表示 \SI{100}{\kilo\gram}など
\usepackage{ascmac}       % itemboxなどの囲み枠
\usepackage{amsmath}
\usepackage{amssymb}      % 数学記号 (\becauseなど)
\numberwithin{equation}{chapter}
% --- その他 ---
\parindent = 0pt % 全体の段落開始時のインデントをなくす
\newcommand{\intinf}{\int_{-\infty}^\infty}
\newcommand{\pathint}{\int \mathcal{D}\varphi}
\newenvironment{note}{
    \vspace{1ex} % 上部の余白（お好みで調整）
    \par\footnotesize\noindent\textbf{【注】}
}{
    \par\vspace{1ex} % 下部の余白
}
% --- ハイパーリンク関連 (できるだけ最後に読み込む) ---
\usepackage[
    dvipdfmx,
    bookmarks=true,
    bookmarksnumbered=true,
    bookmarkstype=toc,
    colorlinks=true,
    linkcolor=blue,
    citecolor=blue,
    urlcolor=blue
]{hyperref}
\usepackage{pxjahyper} % hyperrefの日本語文字化け対策
\usepackage{amsmath, amssymb}
\usepackage[many]{tcolorbox} % manyオプションでskinsやtheoremsなどを一括読み込み

% tcolorboxを用いた新しい定理環境の定義
\tcbset{
    myboxstyle/.style={
        enhanced,
        colback=white,           % メインボックスの背景色（白）
        boxrule=0.8pt,           % 枠線の太さ
        arc=1.5mm,               % 角の丸み
        coltitle=white,          % タイトル文字の色（白）
        attach boxed title to top left={xshift=5mm, yshift=-\tcboxedtitleheight/2},
        description font=\normalfont,
        separator sign={:}       % 番号の後のコロン
    }
}

% ==========================================
% 2. 各環境の定義（色のみ個別設定）
% ==========================================

% --- 定理 (Theorem) : グレー系 ---
\newtcbtheorem[number within=section]{mythm}{Theorem}{
    myboxstyle,
    colframe=red!75!black,           % 枠線の色（赤系に合わせて変更）
    boxed title style={
        colback=red!50,              % タイトル背景色をグレーから赤に変更
        colframe=red!75!black,       % タイトル枠線の色（赤系に合わせて変更）
        boxrule=0.8pt,
        arc=1.5mm
    }
}{thm}

% --- 定義 (Definition) : ブルー系 ---
\newtcbtheorem[number within=section]{mydef}{Definition}{
    myboxstyle,
    colframe=blue!70!black,      % 枠線の色（青70% + 黒30%）
    boxed title style={
        colback=blue!60!gray!80, % タイトル背景色（少しグレーを混ぜた落ち着いた青）
        colframe=blue!70!black,
        boxrule=0.8pt,
        arc=1.5mm
    }
}{def}

% --- 補題 (Lemma) : グリーン系 ---
\newtcbtheorem[number within=section]{mylem}{Lemma}{
    myboxstyle,
    colframe=green!50!black,     % 枠線の色（緑50% + 黒50%）
    boxed title style={
        colback=green!50!gray,   % タイトル背景色（渋めの緑）
        colframe=green!50!black,
        boxrule=0.8pt,
        arc=1.5mm
    }
}{lem}

% --- 問題と解答 ---
\theoremstyle{definition}
\newtheorem{problem}{問題}[chapter]       % 「問題 1.1」のように章ごとに番号付け
% solution 環境の中身はその場では表示せず、\printsolutions で最後にまとめて出力する
\usepackage{environ}
\newcommand{\allsolutions}{}            % 集めた解答
\newcommand{\currentchaptertitle}{}     % 現在の章タイトル（解答編の見出し用）
\let\origchaptermark\chaptermark
\renewcommand{\chaptermark}[1]{\gdef\currentchaptertitle{#1}\origchaptermark{#1}}
\newcounter{solchapter}                 % 解答編で最後に見出しを出した章
\NewEnviron{solution}{%
    \ifnum\value{chapter}>\value{solchapter}%
        \setcounter{solchapter}{\value{chapter}}%
        \xdef\allsolutions{\unexpanded\expandafter{\allsolutions}%
            \noexpand\solutionchapter{\thechapter}{\unexpanded\expandafter{\currentchaptertitle}}}%
    \fi
    \xdef\allsolutions{\unexpanded\expandafter{\allsolutions}%
        \noexpand\solutionitem{\theproblem}{\unexpanded\expandafter{\BODY}}}%
}
\newcommand{\solutionchapter}[2]{\section*{第#1章\quad #2}}
\newcommand{\solutionitem}[2]{\par\medskip\noindent\textbf{問題 #1 の解答}\par\nopagebreak #2\qed\par}
\newcommand{\printsolutions}{%
    \chapter*{解答}
    \addcontentsline{toc}{chapter}{解答}
    \allsolutions
}
% 各章の冒頭に置く「この章で用いる事項」の枠（ページをまたいでも分割される）
\newtcolorbox{summary}{breakable, enhanced, colback=gray!5, colframe=gray!60!black,
    boxrule=0.6pt, arc=1.5mm, fonttitle=\bfseries, title=この章で用いる事項}
\newcommand{\id}{\mathrm{id}}  % 恒等写像
\newcommand{\Z}{\mathbb{Z}}    % 整数
\newcommand{\R}{\mathbb{R}}    % 実数

% --- 以前のテスト用の本文（不要なら \iffalse から \fi まで削除してよい） ---
\iffalse
\begin{document}

\title{テスト文書}
\author{小武}
\date{\today}
\maketitle
\section{はじめに}


\maketitle
\tableofcontents
\clearpage

% ---------------------------------------------------
% 問題編
% ---------------------------------------------------
\section{ホモトピーと基本群}

\begin{problem}
位相空間 $X, Y$ と基点 $x_0 \in X, y_0 \in Y$ について、直積空間 $X \times Y$ の基本群が
\[
\pi_1(X \times Y, (x_0, y_0)) \cong \pi_1(X, x_0) \times \pi_1(Y, y_0)
\]
と同型になることを示せ。
\end{problem}

\begin{problem}
$n \ge 2$ のとき、球面 $S^n$ の基本群 $\pi_1(S^n)$ が自明群（$\{1\}$）であることを示せ。
\end{problem}

\section{ファン・カンペンの定理}

\begin{problem}
8の字空間 $X = S^1 \vee S^1$ の基本群 $\pi_1(X)$ が、2つの生成元を持つ自由群 $F_2$ となることを、ファン・カンペンの定理を用いて示せ。
\end{problem}

\begin{problem}
2次元トーラス $T^2 = S^1 \times S^1$ を、1点の周りに2つの円周を付着させた空間（1-スケルトン）に2次元セル（円板）を1つ付着させたCW複体とみなすことで、その基本群 $\pi_1(T^2)$ がアーベル群 $\mathbb{Z} \times \mathbb{Z}$ と同型になることを示せ。
\end{problem}

\section{被覆空間}

\begin{problem}
実射影平面 $\mathbb{RP}^2$ の普遍被覆空間を構成し、その基本群 $\pi_1(\mathbb{RP}^2)$ を計算せよ。
\end{problem}

\begin{problem}
円周 $S^1$ の連結な被覆空間は、同相を除いて $p_n: S^1 \to S^1, z \mapsto z^n$ ($n \in \mathbb{N}$) または指数写像 $p: \mathbb{R} \to S^1, t \mapsto e^{2\pi i t}$ に限られることを示せ。
\end{problem}

\clearpage
\fi

%================================================
% Document
%================================================
\title{
  代数トポロジー問題集\\
  \large 学部4年程度・全200問
}
\author{}
\date{}

\begin{document}

\frontmatter
\maketitle
\tableofcontents

\chapter*{はじめに}
本書は学部4年程度の代数トポロジーを対象とする
演習問題集である。

証明問題と計算問題の両方を扱い、
定義・定理の理解から具体的な計算、
理論の応用までを段階的に学習する。

各問題には詳細な解答を付す。

各問題には難易度を A〜D の4段階で付した。
\begin{itemize}
  \item A：定義の確認や、ほぼ1行で済む基本問題
  \item B：標準的な定理を1つ用いる計算・証明
  \item C：複数の道具を組み合わせる、またはやや長い計算・証明
  \item D：発展的な内容や、長い証明・工夫を要する問題
\end{itemize}

\mainmatter

%================================================
\chapter{位相空間の基礎とホモトピー}
% 20問
%================================================

\begin{summary}
\begin{itemize}
  \item \textbf{ホモトピー}：連続写像 $f,g\colon X\to Y$ がホモトピック（$f\simeq g$）とは、連続写像 $H\colon X\times[0,1]\to Y$ で $H(\cdot,0)=f$、$H(\cdot,1)=g$ となるものが存在すること。$\simeq$ は同値関係であり、合成と両立する（$f\simeq f'$, $g\simeq g'$ なら $g\circ f\simeq g'\circ f'$）。
  \item \textbf{ホモトピー同値}：$f\colon X\to Y$ に対し $g\colon Y\to X$ で $g\circ f\simeq\id_X$、$f\circ g\simeq\id_Y$ となるものがあるとき $X\simeq Y$ と書く。一点とホモトピー同値な空間を\textbf{可縮}という。
  \item \textbf{レトラクト}：$A\subset X$ と連続写像 $r\colon X\to A$ で $r|_A=\id_A$ となるものがあるとき、$A$ は $X$ のレトラクト。さらに $\id_X$ と $i\circ r$ が（$A$ を動かさずに）ホモトピックなら\textbf{（強）変位レトラクト}といい、このとき $A\simeq X$。
  \item \textbf{構成}：$f\colon X\to Y$ の写像柱 $M_f=(X\times[0,1])\sqcup Y/(x,1)\sim f(x)$、錐 $CX=X\times[0,1]/X\times\{1\}$、懸垂 $SX=X\times[0,1]/(X\times\{0\},\,X\times\{1\}\text{ をそれぞれ1点に})$。$M_f$ は $Y$ に変位レトラクトし、$CX$ は可縮、$SS^n\cong S^{n+1}$。
  \item \textbf{CW複体}：$(n-1)$ 骨格 $X^{n-1}$ に $n$ 胞体 $D^n$ を接着写像 $\varphi\colon S^{n-1}\to X^{n-1}$ で貼って $X^n$ を作り、和集合に弱位相を入れたもの。有限CW複体はコンパクト。
  \item \textbf{貼り合わせの補題}：$X=A\cup B$（$A,B$ 閉）で $f|_A$, $f|_B$ が連続なら $f$ は連続。ホモトピーや写像を区分的に定義するときに常に用いる。
\end{itemize}
\end{summary}

\begin{problem}[難易度 A]
区間 $[0,1]$ が一点とホモトピー同値であることを示せ。
\end{problem}

\begin{solution}
写像
\[
f:[0,1]\longrightarrow \{*\},
\qquad
g:\{*\}\longrightarrow [0,1]
\]
を考える。

$f\circ g=\id_{\{*\}}$ である。

一方、
\[
g\circ f(x)=0
\]
であるから、恒等写像との間に
\[
H(x,t)=(1-t)x
\]
というホモトピーが存在する。

したがって
\[
[0,1]\simeq \{*\}.
\]
\end{solution}

\begin{problem}[難易度 A]
円周 $S^1$ が可縮でないことを示せ。
\end{problem}

\begin{solution}
もし $S^1$ が可縮ならば、
その基本群は自明でなければならない。

しかし
\[
\pi_1(S^1)\cong\Z
\]
である。

よって $S^1$ は可縮ではない。
\end{solution}

\begin{problem}[難易度 A]
位相空間 $X, Y$ に対し、連続写像 $f, g\colon X\to Y$ がホモトピック（$f\simeq g$）であるとは、連続写像 $H\colon X\times[0,1]\to Y$ で $H(x,0)=f(x)$、$H(x,1)=g(x)$ をみたすものが存在することをいう。関係 $\simeq$ が $X$ から $Y$ への連続写像全体の集合上の同値関係であることを示せ。
\end{problem}

\begin{solution}
$I=[0,1]$ と書く。

反射律：$H(x,t)=f(x)$ は連続であり、$f\simeq f$ を与える。

対称律：$H$ が $f\simeq g$ を与えるとき、$H'(x,t)=H(x,1-t)$ は連続写像 $X\times I\to X\times I$、$(x,t)\mapsto(x,1-t)$ と $H$ の合成なので連続で、$H'(x,0)=g(x)$、$H'(x,1)=f(x)$ である。よって $g\simeq f$。

推移律：$H$ が $f\simeq g$ を、$K$ が $g\simeq h$ を与えるとする。
\[
L(x,t)=\begin{cases} H(x,2t) & (0\le t\le 1/2),\\ K(x,2t-1) & (1/2\le t\le 1)\end{cases}
\]
と定める。$t=1/2$ では $H(x,1)=g(x)=K(x,0)$ なので $L$ は矛盾なく定義される。$X\times[0,1/2]$ と $X\times[1/2,1]$ は $X\times I$ の閉集合で、その和は $X\times I$ であり、$L$ は各々の上で連続である。よって貼り合わせの補題により $L$ は連続である。$L(x,0)=f(x)$、$L(x,1)=h(x)$ なので $f\simeq h$ である。
\end{solution}

\begin{problem}[難易度 A]
連続写像 $f, f'\colon X\to Y$、$g, g'\colon Y\to Z$ が $f\simeq f'$、$g\simeq g'$ をみたすとき、$g\circ f\simeq g'\circ f'$ であることを示せ。
\end{problem}

\begin{solution}
$H\colon X\times I\to Y$ を $f\simeq f'$ を与えるホモトピー、$K\colon Y\times I\to Z$ を $g\simeq g'$ を与えるホモトピーとする。

まず $g\circ H\colon X\times I\to Z$ は連続であり、$(g\circ H)(x,0)=g(f(x))$、$(g\circ H)(x,1)=g(f'(x))$ をみたす。よって $g\circ f\simeq g\circ f'$ である。

次に $K'(x,t)=K(f'(x),t)$ とおく。これは連続写像 $f'\times\id_I\colon X\times I\to Y\times I$ と $K$ の合成なので連続であり、$K'(x,0)=g(f'(x))$、$K'(x,1)=g'(f'(x))$ である。よって $g\circ f'\simeq g'\circ f'$ である。

ホモトピーは同値関係（特に推移的）であるから、$g\circ f\simeq g\circ f'\simeq g'\circ f'$ より $g\circ f\simeq g'\circ f'$ を得る。
\end{solution}

\begin{problem}[難易度 A]
$C\subset\R^n$ を空でない凸集合とする。任意の位相空間 $X$ と任意の連続写像 $f, g\colon X\to C$ に対して $f\simeq g$ であることを示せ。また、$C$ が可縮であることを示せ。
\end{problem}

\begin{solution}
$H\colon X\times I\to C$ を
\[
H(x,t)=(1-t)f(x)+t\,g(x)
\]
で定める（直線ホモトピー）。$f(x), g(x)\in C$ で $C$ は凸なので、$0\le t\le 1$ に対して $H(x,t)\in C$ である。$H$ は $\R^n$ のベクトル演算と連続写像の組合せなので連続であり、$H(x,0)=f(x)$、$H(x,1)=g(x)$ をみたす。よって $f\simeq g$ である。

$C$ が可縮であるとは、恒等写像 $\id_C$ がある定値写像とホモトピックであることである。点 $c_0\in C$ をとり、$X=C$、$f=\id_C$、$g$ を値 $c_0$ の定値写像として上の結果を適用すれば $\id_C\simeq g$ を得る。よって $C$ は可縮である。
\end{solution}

\begin{problem}[難易度 A]
位相空間 $X, Y$ がホモトピー同値であるとは、連続写像 $f\colon X\to Y$、$g\colon Y\to X$ で $g\circ f\simeq\id_X$、$f\circ g\simeq\id_Y$ となるものが存在することをいう。$X$ と $Y$、$Y$ と $Z$ がそれぞれホモトピー同値ならば、$X$ と $Z$ もホモトピー同値であることを示せ。
\end{problem}

\begin{solution}
$f\colon X\to Y$、$g\colon Y\to X$ を $g\circ f\simeq\id_X$、$f\circ g\simeq\id_Y$ をみたす写像、$h\colon Y\to Z$、$k\colon Z\to Y$ を $k\circ h\simeq\id_Y$、$h\circ k\simeq\id_Z$ をみたす写像とする。$h\circ f\colon X\to Z$ と $g\circ k\colon Z\to X$ を考える。

ホモトピーは合成と両立する（$a\simeq a'$、$b\simeq b'$ ならば $b\circ a\simeq b'\circ a'$）。これを用いると
\[
(g\circ k)\circ(h\circ f)=g\circ(k\circ h)\circ f\simeq g\circ\id_Y\circ f=g\circ f\simeq\id_X
\]
であり、同様に
\[
(h\circ f)\circ(g\circ k)=h\circ(f\circ g)\circ k\simeq h\circ k\simeq\id_Z
\]
である。よって $h\circ f$ はホモトピー同値写像であり、$X$ と $Z$ はホモトピー同値である。
\end{solution}

\begin{problem}[難易度 A]
空でない可縮な空間 $X$ は弧状連結であることを示せ。
\end{problem}

\begin{solution}
$X$ は可縮なので、点 $x_0\in X$ と連続写像 $H\colon X\times I\to X$ で $H(x,0)=x$、$H(x,1)=x_0$（すべての $x\in X$）をみたすものが存在する。

任意の $x\in X$ に対し、$\gamma_x(t)=H(x,t)$ とおく。$\gamma_x$ は連続写像 $I\to X\times I$、$t\mapsto(x,t)$ と $H$ の合成なので連続であり、$\gamma_x(0)=x$、$\gamma_x(1)=x_0$ をみたす。すなわち $\gamma_x$ は $x$ から $x_0$ への道である。

$x, y\in X$ を任意にとると、$\gamma_x$ に続けて $\gamma_y$ を逆向きにたどる道が $x$ と $y$ を結ぶ。よって $X$ は弧状連結である。
\end{solution}

\begin{problem}[難易度 B]
位相空間 $X, Y$ について、次を示せ。
\begin{enumerate}
\item $Y$ が可縮ならば、任意の連続写像 $f, g\colon X\to Y$ はホモトピックである。
\item $X$ が可縮で $Y$ が弧状連結ならば、任意の連続写像 $f, g\colon X\to Y$ はホモトピックである。
\end{enumerate}
\end{problem}

\begin{solution}
写像の合成とホモトピーが両立すること、およびホモトピーが同値関係であることを用いる。

(1) $\id_Y\simeq c$（$c$ は値 $y_0$ の定値写像）とする。このとき
\[
f=\id_Y\circ f\simeq c\circ f,\qquad g=\id_Y\circ g\simeq c\circ g
\]
であり、$c\circ f$ と $c\circ g$ はともに値 $y_0$ の定値写像 $X\to Y$ で等しい。よって $f\simeq g$ である。

(2) $\id_X\simeq c'$（$c'$ は値 $x_0$ の定値写像）とする。すると $f=f\circ\id_X\simeq f\circ c'$ であり、$f\circ c'$ は値 $f(x_0)$ の定値写像である。同様に $g$ は値 $g(x_0)$ の定値写像とホモトピックである。$Y$ は弧状連結なので、$f(x_0)$ から $g(x_0)$ への道 $\gamma$ がある。$H(x,t)=\gamma(t)$ は連続で、値 $f(x_0)$ の定値写像と値 $g(x_0)$ の定値写像の間のホモトピーを与える。以上をつなげて $f\simeq g$ を得る。
\end{solution}

\begin{problem}[難易度 B]
$A\subset X$ が $X$ の強変位レトラクトであるとは、連続写像 $H\colon X\times I\to X$ で、$H(x,0)=x$、$H(x,1)\in A$（$x\in X$）、$H(a,t)=a$（$a\in A$、$t\in I$）をみたすものが存在することをいう。$n\ge1$ のとき、$S^{n-1}$ が $\R^n\setminus\{0\}$ の強変位レトラクトであることを示し、$\R^n\setminus\{0\}$ と $S^{n-1}$ がホモトピー同値であることを示せ。
\end{problem}

\begin{solution}
$X=\R^n\setminus\{0\}$ とし、$H\colon X\times I\to\R^n$ を
\[
H(x,t)=(1-t)x+t\,\frac{x}{\lvert x\rvert}=\Bigl((1-t)+\frac{t}{\lvert x\rvert}\Bigr)x
\]
で定める。$x\ne0$ なので $\lvert x\rvert$ は連続かつ正であり、$H$ は連続である。係数 $(1-t)+t/\lvert x\rvert$ は、非負の二数 $1-t$ と $t/\lvert x\rvert$ の少なくとも一方が正なので正であり、 $H(x,t)\ne0$、すなわち $H$ は $X$ に値をとる。

さらに $H(x,0)=x$、$H(x,1)=x/\lvert x\rvert\in S^{n-1}$ であり、$\lvert a\rvert=1$ ならば $H(a,t)=a$ である。よって $S^{n-1}$ は $X$ の強変位レトラクトである。

$r(x)=x/\lvert x\rvert\colon X\to S^{n-1}$、$i\colon S^{n-1}\to X$ を包含とすると、$r\circ i=\id_{S^{n-1}}$ であり、$H$ は $i\circ r\simeq\id_X$ を与える。よって $X$ と $S^{n-1}$ はホモトピー同値である。
\end{solution}

\begin{problem}[難易度 B]
$X$ を位相空間、$f, g\colon X\to S^n$ を連続写像とし、すべての $x\in X$ に対して $f(x)\ne -g(x)$ であるとする。このとき $f\simeq g$ であることを示せ。さらに、不動点をもたない連続写像 $f\colon S^n\to S^n$ は対蹠写像 $a(x)=-x$ とホモトピックであることを示せ。
\end{problem}

\begin{solution}
$(x,t)\in X\times I$ に対し $v(x,t)=(1-t)f(x)+t\,g(x)\in\R^{n+1}$ とおく。$v(x,t)\ne0$ であることを示す。もし $v(x,t)=0$ なら $(1-t)f(x)=-t\,g(x)$ であり、両辺のノルムをとると $1-t=t$、すなわち $t=1/2$ となる。このとき $f(x)=-g(x)$ となり仮定に反する。

よって
\[
H(x,t)=\frac{(1-t)f(x)+t\,g(x)}{\lvert(1-t)f(x)+t\,g(x)\rvert}
\]
は連続写像 $X\times I\to S^n$ として定義でき、$H(x,0)=f(x)$、$H(x,1)=g(x)$ である。したがって $f\simeq g$ である。

後半：$f$ が不動点をもたないとき、すべての $x$ で $f(x)\ne x=-a(x)$ である。よって前半を $g=a$ に適用して $f\simeq a$ を得る。
\end{solution}

\begin{problem}[難易度 B]
連続写像 $f\colon S^n\to Y$ が定値写像とホモトピックであるための必要十分条件は、$f$ が連続写像 $F\colon D^{n+1}\to Y$ に拡張できることであることを示せ。ここで $D^{n+1}=\{x\in\R^{n+1}:\lvert x\rvert\le1\}$、$S^n=\partial D^{n+1}$ である。
\end{problem}

\begin{solution}
（十分性）$F|_{S^n}=f$ とする。$H(x,t)=F((1-t)x)$（$x\in S^n$）は連続で、$H(x,0)=f(x)$、$H(x,1)=F(0)$ である。よって $f$ は値 $F(0)$ の定値写像とホモトピックである。

（必要性）$H\colon S^n\times I\to Y$ を $H(x,0)=f(x)$、$H(x,1)=y_0$ をみたすホモトピーとする。写像
\[
q\colon S^n\times I\to D^{n+1},\qquad q(x,s)=(1-s)x
\]
は連続な全射である。$S^n\times I$ はコンパクト、$D^{n+1}$ はハウスドルフなので $q$ は閉写像であり、特に商写像である。$q$ のファイバーは、$0$ 上では $S^n\times\{1\}$、それ以外の点上では一点である。$H$ は $S^n\times\{1\}$ 上で値 $y_0$ をとるので、$H$ は $q$ のファイバー上で一定である。よって商写像の普遍性から、連続写像 $F\colon D^{n+1}\to Y$ で $F\circ q=H$ となるものが存在する。$x\in S^n$ に対し $F(x)=F(q(x,0))=H(x,0)=f(x)$ なので、$F$ は $f$ の拡張である。
\end{solution}

\begin{problem}[難易度 B]
メビウスの帯を $M=[0,1]\times[-1,1]/{\sim}$（$(0,y)\sim(1,-y)$）で定める。中心線 $C=\{[s,0]:s\in[0,1]\}$ が $M$ の強変位レトラクトであることを示し、$M$ が $S^1$ とホモトピー同値であることを示せ。
\begin{center}
\begin{tikzpicture}[scale=1.1]
  \fill[blue!6] (0,-1) rectangle (4,1);
  \draw (0,-1) -- (4,-1) (0,1) -- (4,1);
  \draw[very thick,red!70!black,mid arrow] (0,-1) -- (0,1);
  \draw[very thick,red!70!black,mid arrow] (4,1) -- (4,-1);
  \draw[thick,dashed] (0,0) -- (4,0) node[right] {$C$};
  \node[below left] at (0,-1) {$(0,-1)$};
  \node[above left] at (0,1) {$(0,1)$};
  \node[below right] at (4,-1) {$(1,-1)$};
  \node[above right] at (4,1) {$(1,1)$};
  \node at (2,-1.35) {$s$};
\end{tikzpicture}
\zucap{左右の辺を矢印の向きがそろうように（ひねって）貼り合わせるとメビウスの帯になる。破線が中心線 $C$。}
\end{center}
\end{problem}

\begin{solution}
\begin{center}
\begin{tikzpicture}[scale=1.0]
  \fill[blue!6] (0,-1) rectangle (4,1);
  \draw (0,-1) -- (4,-1) (0,1) -- (4,1);
  \draw[thick,red!70!black,mid arrow] (0,-1) -- (0,1);
  \draw[thick,red!70!black,mid arrow] (4,1) -- (4,-1);
  \draw[thick] (0,0) -- (4,0);
  \foreach \x in {0.5,1.25,2,2.75,3.5}{
    \draw[-{Stealth[length=1.8mm]},gray] (\x,0.9) -- (\x,0.15);
    \draw[-{Stealth[length=1.8mm]},gray] (\x,-0.9) -- (\x,-0.15);
  }
\end{tikzpicture}
\zucap{各縦線分（ファイバー）を中心に向かって一様に縮める。$(0,y)$ と $(1,-y)$ の行き先も再び同一視されるので、この変形は商空間 $M$ 上で矛盾なく定義される。}
\end{center}
$q\colon[0,1]\times[-1,1]\to M$ を商写像とする。$\tilde H\colon([0,1]\times[-1,1])\times I\to[0,1]\times[-1,1]$ を $\tilde H((s,y),t)=(s,(1-t)y)$ で定める。$(0,y)$ と $(1,-y)$ の行き先 $(0,(1-t)y)$ と $(1,-(1-t)y)$ は再び $\sim$ で同一視されるので、$q\circ\tilde H$ は $q\times\id_I$ のファイバー上で一定である。

$I$ は局所コンパクト・ハウスドルフなので、商写像 $q$ に対して $q\times\id_I$ も商写像である。よって $H\colon M\times I\to M$、$H([s,y],t)=[s,(1-t)y]$ は矛盾なく定義され連続である。$H(\cdot,0)=\id_M$、$H([s,y],1)=[s,0]\in C$、$H([s,0],t)=[s,0]$ であるから、$C$ は $M$ の強変位レトラクトである。

最後に、$[0,1]$ において $0$ と $1$ のみを同一視した空間 $[0,1]/\{0,1\}$ は $S^1$ と同相であり（$s\mapsto e^{2\pi is}$ が誘導する連続全単射はコンパクト空間からハウスドルフ空間への写像なので同相）、$s\mapsto[s,0]$ は同様に $[0,1]/\{0,1\}$ から $C$ への同相を誘導する。よって $C\cong S^1$ であり、強変位レトラクトの包含はホモトピー同値なので $M\simeq S^1$ である。
\end{solution}

\begin{problem}[難易度 A]
位相空間 $X$ の錐を $CX=X\times I/(X\times\{1\})$ で定める。$X\ne\emptyset$ のとき、$CX$ が可縮であることを示せ。
\end{problem}

\begin{solution}
$q\colon X\times I\to CX$ を商写像とし、$v=q(X\times\{1\})$ を錐の頂点とする。$\tilde H\colon (X\times I)\times I\to X\times I$ を
\[
\tilde H((x,s),t)=(x,\,(1-t)s+t)
\]
で定める。$s=1$ ならば $(1-t)s+t=1$ なので、$q\circ\tilde H$ は $X\times\{1\}$ の点を常に頂点 $v$ に送る。したがって $q\circ\tilde H$ は $q\times\id_I$ のファイバー上で一定である。

$I$ は局所コンパクト・ハウスドルフなので $q\times\id_I$ は商写像であり、$H([x,s],t)=[x,(1-t)s+t]$ は連続写像 $CX\times I\to CX$ を定める。$H(\cdot,0)=\id_{CX}$ であり、$H([x,s],1)=[x,1]=v$ である。よって $\id_{CX}$ は値 $v$ の定値写像とホモトピックであり、$CX$ は可縮である。
\end{solution}

\begin{problem}[難易度 B]
連続写像 $f\colon X\to Y$ の写像柱を
\[
M_f=\bigl((X\times I)\sqcup Y\bigr)/{\sim},\qquad (x,1)\sim f(x)
\]
で定める。$Y$（の像）が $M_f$ の強変位レトラクトであることを示せ。また、$j\colon X\to M_f$、$j(x)=[x,0]$ と変位レトラクション $r\colon M_f\to Y$ に対して $r\circ j=f$ となることを示せ。
\end{problem}

\begin{solution}
$q\colon(X\times I)\sqcup Y\to M_f$ を商写像とする。$q$ の $Y$ への制限は単射で、像を $Y$ と同一視する。

$(X\times I)\sqcup Y$ 上のホモトピー $\tilde H$ を、$X\times I$ 上では $\tilde H((x,s),t)=(x,(1-t)s+t)$、$Y$ 上では $\tilde H(y,t)=y$ で定める。$s=1$ のとき $\tilde H((x,1),t)=(x,1)\sim f(x)=\tilde H(f(x),t)$ なので、$q\circ\tilde H$ は $q\times\id_I$ のファイバー上で一定である。$q\times\id_I$ は商写像（$I$ が局所コンパクト・ハウスドルフ）なので、
\[
H([x,s],t)=[x,(1-t)s+t],\qquad H(y,t)=y
\]
は連続写像 $M_f\times I\to M_f$ を与える。$H(\cdot,0)=\id$、$H([x,s],1)=[x,1]=f(x)\in Y$、$Y$ の点は動かないので、$Y$ は $M_f$ の強変位レトラクトである。

このとき変位レトラクションは $r=H(\cdot,1)$、すなわち $r([x,s])=f(x)$、$r(y)=y$ である。よって $r(j(x))=r([x,0])=f(x)$ となり、$r\circ j=f$ である。特に任意の写像は、包含 $j$ とホモトピー同値写像 $r$ の合成に分解される。
\end{solution}

\begin{problem}[難易度 C]
位相空間 $X$ の懸垂を $SX=X\times[-1,1]/{\sim}$ で定める。ここで $\sim$ は $X\times\{1\}$ を一点に、$X\times\{-1\}$ を別の一点につぶす同値関係である。$SS^n$ が $S^{n+1}$ と同相であることを示せ。
\end{problem}

\begin{solution}
$\varphi\colon S^n\times[-1,1]\to S^{n+1}\subset\R^{n+1}\times\R$ を
\[
\varphi(x,t)=\bigl(\sqrt{1-t^2}\,x,\ t\bigr)
\]
で定める。$\lvert\varphi(x,t)\rvert^2=(1-t^2)+t^2=1$ なので $\varphi$ は $S^{n+1}$ に値をとる連続写像である。

$\varphi(S^n\times\{1\})=\{(0,1)\}$、$\varphi(S^n\times\{-1\})=\{(0,-1)\}$ なので、$\varphi$ は連続写像 $\bar\varphi\colon SS^n\to S^{n+1}$ を誘導する。$\bar\varphi$ が全単射であることを示す。$(y,t)\in S^{n+1}$ をとると $\lvert y\rvert^2=1-t^2$ である。$\lvert t\rvert<1$ ならば $y\ne0$ で、$x=y/\lvert y\rvert$ は $\varphi(x,t)=(y,t)$ をみたし、このような $x$ は唯一つである。$t=\pm1$ ならば $(y,t)=(0,\pm1)$ で、その逆像は $SS^n$ のつぶされた一点である。よって $\bar\varphi$ は全単射である。

$SS^n$ はコンパクト空間 $S^n\times[-1,1]$ の連続像なのでコンパクトであり、$S^{n+1}$ はハウスドルフである。コンパクト空間からハウスドルフ空間への連続全単射は同相写像なので、$SS^n\cong S^{n+1}$ である。
\end{solution}

\begin{problem}[難易度 C]
$O(n)$ が一般線形群 $GL_n(\R)$（$\R^{n^2}$ の部分空間）の強変位レトラクトであることを示せ。
\end{problem}

\begin{solution}
$A\in GL_n(\R)$ の列ベクトル $a_1,\dots,a_n$ は一次独立である。これにグラム・シュミットの直交化を施す：
\[
u_k=a_k-\sum_{j<k}\langle a_k,q_j\rangle q_j,\qquad q_k=\frac{u_k}{\lvert u_k\rvert}.
\]
一次独立性より $u_k\ne0$ であり、$q_k$ は $A$ の成分の連続関数である（和・積・商・平方根を、分母が $0$ にならない範囲で組み合わせたもの）。$Q(A)=(q_1,\dots,q_n)\in O(n)$ とおき、$R(A)=Q(A)^{\mathsf T}A$ とおくと $A=Q(A)R(A)$ であり、$R(A)$ は対角成分 $\lvert u_k\rvert>0$ をもつ上三角行列である。$Q, R$ は連続である。また $A\in O(n)$ ならば $u_k=a_k$、$q_k=a_k$ となるので $Q(A)=A$、$R(A)=E$（単位行列）である。

そこで
\[
H(A,t)=Q(A)\bigl((1-t)R(A)+tE\bigr)
\]
とおく。$(1-t)R(A)+tE$ は対角成分が正の上三角行列なので正則であり、$H(A,t)\in GL_n(\R)$ である。$H$ は連続で、$H(A,0)=A$、$H(A,1)=Q(A)\in O(n)$、$A\in O(n)$ ならば $H(A,t)=A$ である。よって $O(n)$ は $GL_n(\R)$ の強変位レトラクトである。
\end{solution}

\begin{problem}[難易度 B]
有限個の胞体からなる CW 複体（有限 CW 複体）はコンパクトであることを示せ。
\end{problem}

\begin{solution}
$X$ を有限 CW 複体とし、その $k$ 骨格を $X^k$ とする。定義により $X^0$ は離散空間、$X^k$ は $X^{k-1}$ に $k$ 胞体 $D^k_\alpha$ を接着写像 $\varphi_\alpha\colon S^{k-1}\to X^{k-1}$ で貼り付けた空間
\[
X^k=\Bigl(X^{k-1}\sqcup\bigsqcup_\alpha D^k_\alpha\Bigr)\Big/(x\sim\varphi_\alpha(x),\ x\in\partial D^k_\alpha)
\]
である。胞体は有限個なので、各 $X^k$ に貼る胞体は有限個であり、ある $N$ について $X=X^N$ である。

$X^k$ がコンパクトであることを $k$ についての帰納法で示す。$X^0$ は有限離散空間なのでコンパクトである。$X^{k-1}$ がコンパクトならば、$X^{k-1}\sqcup\bigsqcup_\alpha D^k_\alpha$ はコンパクト空間の有限個の非交和なのでコンパクトであり、$X^k$ はその商写像（連続全射）による像なのでコンパクトである。

よって $X=X^N$ はコンパクトである。
\end{solution}

\begin{problem}[難易度 C]
トーラス $T^2$ を正方形 $I^2$ の辺を $(0,y)\sim(1,y)$、$(x,0)\sim(x,1)$ で同一視した空間とし、$p\in T^2$ を正方形の中心 $c=(1/2,1/2)$ の像とする。$T^2\setminus\{p\}$ が $S^1\vee S^1$ と同相な部分空間に強変位レトラクトすることを示せ。
\end{problem}

\begin{solution}
\begin{center}
\begin{tikzpicture}[scale=2.2]
  \fill[blue!6] (0,0) rectangle (1,1);
  \draw[thick,red!70!black,mid arrow] (0,0) -- (1,0);
  \draw[thick,red!70!black,mid arrow] (0,1) -- (1,1);
  \draw[thick,blue!70!black,mid arrows] (0,0) -- (0,1);
  \draw[thick,blue!70!black,mid arrows] (1,0) -- (1,1);
  \foreach \a in {0,30,...,330}{
    \draw[-{Stealth[length=1.6mm]},gray] ($(0.5,0.5)+(\a:0.08)$) -- ($(0.5,0.5)+(\a:0.3)$);
  }
  \node[hvtx] at (0.5,0.5) {};
  \node[above right] at (0.52,0.5) {\footnotesize $c$};
  \node[below] at (0.5,0) {$a$};
  \node[left] at (0,0.5) {$b$};
\end{tikzpicture}
\qquad
\begin{tikzpicture}[scale=1.2]
  \draw[thick,red!70!black,mid arrow] (0,0) .. controls (-1.6,1.2) and (-1.6,-1.2) .. (0,0);
  \draw[thick,blue!70!black,mid arrows] (0,0) .. controls (1.6,1.2) and (1.6,-1.2) .. (0,0);
  \node[vtx] at (0,0) {};
  \node at (-1.4,0) {$a$};
  \node at (1.4,0) {$b$};
  \node at (0,-1.1) {$A=q(\partial I^2)\cong S^1\vee S^1$};
\end{tikzpicture}
\zucap{中心 $c$ を除いた正方形を、中心からの放射状の線分に沿って境界へ押し出す（左）。境界は4辺が $a,b$ の2本に同一視され、4頂点はすべて1点になるので、8の字 $S^1\vee S^1$ になる（右）。}
\end{center}
$q\colon I^2\to T^2$ を商写像とし、$Q=I^2\setminus\{c\}$ とおく。$\lVert z\rVert_\infty=\max(\lvert z_1\rvert,\lvert z_2\rvert)$ とすると、$I^2=\{z:\lVert z-c\rVert_\infty\le1/2\}$ であり、$\partial I^2$ は $\lVert z-c\rVert_\infty=1/2$ の部分である。$z\in Q$ に対し
\[
\tilde H(z,t)=c+\Bigl((1-t)+\frac{t}{2\lVert z-c\rVert_\infty}\Bigr)(z-c)
\]
とおく。係数は正なので $\tilde H(z,t)\ne c$ であり、$\lVert\tilde H(z,t)-c\rVert_\infty=(1-t)\lVert z-c\rVert_\infty+t/2\le1/2$ なので $\tilde H(z,t)\in Q$ である。$\tilde H(z,0)=z$、$\tilde H(z,1)\in\partial I^2$ であり、$z\in\partial I^2$ ならば $\tilde H(z,t)=z$ である。

$\sim$ で同一視されるのは $\partial I^2$ の点のみで、それらは $\tilde H$ で動かないから、$q\circ\tilde H$ は $q\times\id_I$ のファイバー上で一定である。$Q=q^{-1}(T^2\setminus\{p\})$ は飽和な開集合なので $q|_Q\colon Q\to T^2\setminus\{p\}$ は商写像であり、$I$ が局所コンパクト・ハウスドルフであることから $q|_Q\times\id_I$ も商写像である。よって $\tilde H$ は連続写像 $H\colon(T^2\setminus\{p\})\times I\to T^2\setminus\{p\}$ を誘導し、$A=q(\partial I^2)$ は $T^2\setminus\{p\}$ の強変位レトラクトである。

最後に、$A$ は $q([0,1]\times\{0\})$ と $q(\{0\}\times[0,1])$ の和である。各々は両端点を同一視した閉区間なので $S^1$ と同相であり、両者は一点 $q(0,0)$ のみを共有する。よって $A\cong S^1\vee S^1$ である。
\end{solution}

\begin{problem}[難易度 C]
$D^n\times\{0\}\cup S^{n-1}\times I$ が $D^n\times I$ のレトラクトであることを示せ。これを用いて、任意の連続写像 $F\colon D^n\times\{0\}\cup S^{n-1}\times I\to Y$ が $D^n\times I$ 上の連続写像に拡張できることを示せ。
\end{problem}

\begin{solution}
$A=D^n\times\{0\}\cup S^{n-1}\times I$ とおく。点 $(0,2)$ から $(x,t)\in D^n\times I$ を通る半直線と $A$ の交点を $r(x,t)$ とする。具体的に計算すると
\[
r(x,t)=\begin{cases}
\Bigl(\dfrac{2x}{2-t},\ 0\Bigr) & (2\lvert x\rvert\le 2-t),\\[2ex]
\Bigl(\dfrac{x}{\lvert x\rvert},\ 2-\dfrac{2-t}{\lvert x\rvert}\Bigr) & (2\lvert x\rvert\ge 2-t)
\end{cases}
\]
である。$2-t\ge1>0$ なので第1式は連続であり、第2式の領域では $\lvert x\rvert\ge(2-t)/2\ge1/2$ なので第2式も連続である。$2\lvert x\rvert=2-t$ のときは両式とも $(x/\lvert x\rvert,0)$ となり一致する。二つの領域は閉集合なので、貼り合わせの補題により $r$ は連続である。

値域について：第1式では $\lvert 2x/(2-t)\rvert\le1$ なので $D^n\times\{0\}$ に入る。第2式では高さは $2\lvert x\rvert\ge2-t$ より $\ge0$、$\lvert x\rvert\le1$ より $\le t\le1$ なので $S^{n-1}\times I$ に入る。$A$ 上で恒等であること：$t=0$ なら $2\lvert x\rvert\le2$ で第1式より $r(x,0)=(x,0)$、$\lvert x\rvert=1$ なら第2式より $r(x,t)=(x,t)$ である。よって $r$ はレトラクションである。

拡張は $F\circ r\colon D^n\times I\to Y$ で与えられる（$A$ 上で $F\circ r=F$）。
\end{solution}

\begin{problem}[難易度 C]
実射影空間 $\mathbb{RP}^n=S^n/(x\sim-x)$ は、各 $k=0,1,\dots,n$ に対してちょうど一つの $k$ 胞体をもつ CW 複体の構造をもつことを示せ。ただし $\mathbb{RP}^n$ がハウスドルフであることは用いてよい。
\end{problem}

\begin{solution}
$S^{n-1}=\{x\in S^n:x_{n+1}=0\}$ とみなすと、包含は連続単射 $\mathbb{RP}^{n-1}\to\mathbb{RP}^n$ を誘導する。$\varphi\colon S^{n-1}\to\mathbb{RP}^{n-1}$ を商写像とし、$\mathbb{RP}^n\cong\mathbb{RP}^{n-1}\cup_\varphi D^n$ を示せば、$\mathbb{RP}^0$ が一点であることと合わせ、帰納的に主張が従う。

$h\colon D^n\to S^n$、$h(x)=(x,\sqrt{1-\lvert x\rvert^2})$ は閉上半球への同相である。$\Phi\colon\mathbb{RP}^{n-1}\sqcup D^n\to\mathbb{RP}^n$ を、$\mathbb{RP}^{n-1}$ 上では上の包含、$D^n$ 上では $x\mapsto[h(x)]$ で定める。$\Phi$ は連続である。

$S^n$ の各点 $y$ は $y$ か $-y$ の一方が閉上半球に入るので、$\Phi$ は全射である。$\lvert x\rvert<1$ ならば $h(x)$ の第 $n+1$ 座標は正で、$-h(x)$ は閉上半球に入らないから、$[h(x)]$ を値にとる点は $x$ のみである。$\lvert x\rvert=1$ ならば $h(x)=(x,0)$ で、$\Phi(x)=[x]=\Phi(\varphi(x))$ である。したがって $\Phi$ による同一視は、ちょうど $x\in\partial D^n$ と $\varphi(x)\in\mathbb{RP}^{n-1}$ の同一視である。

よって $\Phi$ は連続全単射 $\bar\Phi\colon\mathbb{RP}^{n-1}\cup_\varphi D^n\to\mathbb{RP}^n$ を誘導する。左辺はコンパクト空間の商なのでコンパクト、右辺はハウスドルフなので、$\bar\Phi$ は同相写像である。
\end{solution}

\begin{problem}[難易度 B]
下図のように、アルファベットの大文字を（黒丸を頂点とする）有限グラフ、すなわち1次元CW複体とみなす。
\begin{center}
\begin{tikzpicture}[scale=0.62,line width=1.1pt]
  \begin{scope}[xshift=0cm] % A
    \draw (0,0) -- (0.2,0.64) -- (0.5,1.6) -- (0.8,0.64) -- (1,0) (0.2,0.64) -- (0.8,0.64);
    \foreach \p in {(0,0),(0.2,0.64),(0.5,1.6),(0.8,0.64),(1,0)} \node[vtx] at \p {};
  \end{scope}
  \begin{scope}[xshift=1.8cm] % B
    \draw (0,0) -- (0,1.6);
    \draw (0,1.6) .. controls (0.9,1.6) and (0.9,0.8) .. (0,0.8);
    \draw (0,0.8) .. controls (1.05,0.8) and (1.05,0) .. (0,0);
    \foreach \p in {(0,0),(0,0.8),(0,1.6)} \node[vtx] at \p {};
  \end{scope}
  \begin{scope}[xshift=3.4cm] % C
    \draw (1.03,1.33) arc (45:315:0.75);
    \foreach \p in {(1.03,1.33),(1.03,0.27)} \node[vtx] at \p {};
  \end{scope}
  \begin{scope}[xshift=5.2cm] % D
    \draw (0,0) -- (0,1.6) .. controls (1.2,1.6) and (1.2,0) .. (0,0);
    \foreach \p in {(0,0),(0,1.6)} \node[vtx] at \p {};
  \end{scope}
  \begin{scope}[xshift=6.9cm] % E
    \draw (0.9,1.6) -- (0,1.6) -- (0,0) -- (0.9,0) (0,0.8) -- (0.7,0.8);
    \foreach \p in {(0,0),(0,0.8),(0,1.6),(0.9,1.6),(0.7,0.8),(0.9,0)} \node[vtx] at \p {};
  \end{scope}
  \begin{scope}[xshift=8.5cm] % O
    \draw (0.5,0.8) ellipse (0.5 and 0.8);
    \foreach \p in {(0.5,0),(0.5,1.6)} \node[vtx] at \p {};
  \end{scope}
  \begin{scope}[xshift=10.1cm] % P
    \draw (0,0) -- (0,1.6) .. controls (0.9,1.6) and (0.9,0.8) .. (0,0.8);
    \foreach \p in {(0,0),(0,0.8),(0,1.6)} \node[vtx] at \p {};
  \end{scope}
  \begin{scope}[xshift=11.6cm] % Q
    \draw (0.5,0.8) ellipse (0.5 and 0.8);
    \draw (0.65,0.45) -- (1.06,0.02);
    \foreach \p in {(0.5,1.6),(0.854,0.234),(0.65,0.45),(1.06,0.02)} \node[vtx] at \p {};
  \end{scope}
  \begin{scope}[xshift=13.3cm] % R
    \draw (0,0) -- (0,1.6) .. controls (0.9,1.6) and (0.9,0.8) .. (0,0.8) -- (0.9,0);
    \foreach \p in {(0,0),(0,0.8),(0,1.6),(0.9,0)} \node[vtx] at \p {};
  \end{scope}
  \begin{scope}[xshift=14.9cm] % X
    \draw (0,0) -- (1,1.6) (0,1.6) -- (1,0);
    \foreach \p in {(0,0),(1,1.6),(0,1.6),(1,0),(0.5,0.8)} \node[vtx] at \p {};
  \end{scope}
\end{tikzpicture}
\zucap{A, B, C, D, E, O, P, Q, R, X をグラフとみなしたもの（Q の尾は円周と1点で交わる）。}
\end{center}
\begin{enumerate}
\item 連結な有限グラフ $G$ の頂点数を $V$、辺数を $E$ とすると、$G$ は $E-V+1$ 個の円周の楔和 $\bigvee^{E-V+1}S^1$ とホモトピー同値であることを示せ（$E-V+1=0$ のときは一点、すなわち可縮）。ただし「CW対 $(X,A)$ で $A$ が可縮ならば商写像 $X\to X/A$ はホモトピー同値である」こと、および「有限グラフの木（閉路をもたない連結な部分グラフ）は可縮である」ことは用いてよい。
\item 上の10文字をホモトピー型によって分類せよ。
\end{enumerate}
\end{problem}

\begin{solution}
(1) $G$ の極大な木 $T$（閉路をもたない連結部分グラフで、包含について極大なもの）をとる。$T$ は $G$ のすべての頂点を含む。実際、$T$ に含まれない頂点 $v$ があれば、$G$ が連結なので $T$ の頂点と $T$ 外の頂点を結ぶ辺 $e$ があり、$T\cup e$（と $e$ の端点）は閉路をもたない連結部分グラフとなって極大性に反する。

$V$ 頂点の木の辺数は $V-1$ である（葉を1つずつ取り除く帰納法：有限の木には次数 $1$ の頂点があり、それと接続する辺を除いても木である）。$T$ は $G$ の部分複体で可縮なので、商写像 $G\to G/T$ はホモトピー同値である。$G/T$ は $T$ がつぶれた1点を唯一の頂点とし、$T$ に含まれない $E-(V-1)$ 本の辺がすべてその点を両端とするループになったグラフ、すなわち $\bigvee^{E-V+1}S^1$ である。
\begin{center}
\begin{tikzpicture}[scale=0.9,line width=0.9pt]
  \draw (0,0) -- (0.2,0.64) -- (0.5,1.6) -- (0.8,0.64) -- (1,0);
  \draw[densely dotted] (0.2,0.64) -- (0.8,0.64);
  \foreach \p in {(0,0),(0.2,0.64),(0.5,1.6),(0.8,0.64),(1,0)} \node[vtx] at \p {};
  \draw[-{Stealth[length=2mm]}] (1.5,0.8) -- (2.6,0.8) node[midway,above] {\footnotesize $T$ をつぶす};
  \draw[densely dotted] (3.4,0.8) .. controls (4.9,2.0) and (4.9,-0.4) .. (3.4,0.8);
  \node[vtx] at (3.4,0.8) {};
  \node at (5.4,0.8) {$\cong S^1$};
\end{tikzpicture}
\zucap{A の場合。実線が極大な木 $T$（4辺）、点線が木に含まれない辺（1本）。$T$ を1点につぶすと点線の辺がループになる。}
\end{center}

(2) 各文字の頂点数・辺数を数えると次のようになる。
\begin{center}
\begin{tabular}{c|cccccccccc}
文字 & A & B & C & D & E & O & P & Q & R & X\\ \hline
$V$ & 5 & 3 & 2 & 2 & 6 & 2 & 3 & 4 & 4 & 5\\
$E$ & 5 & 4 & 1 & 2 & 5 & 2 & 3 & 4 & 4 & 4\\
$E-V+1$ & 1 & 2 & 0 & 1 & 0 & 1 & 1 & 1 & 1 & 0
\end{tabular}
\end{center}
（Q では円周が2頂点で2本の弧に分かれ、尾が交点で2本の辺に分かれる。）よって (1) から
\begin{itemize}
\item 可縮：C, E, X
\item $S^1$ とホモトピー同値：A, D, O, P, Q, R
\item $S^1\vee S^1$ とホモトピー同値：B
\end{itemize}
であり、同じ類の文字どうしはホモトピー同値である。

異なる類どうしがホモトピー同値でないことは、ホモトピー不変量で区別する必要がある。基本群は一点で自明、$S^1$ で $\Z$（第2章）、$S^1\vee S^1$ で非可換な自由群 $F_2$（第2章、ファン・カンペンの定理）であり、互いに同型でないので3つの類は異なる。（あるいは第3章の $H_1$ の階数 $E-V+1$ が $0,1,2$ と異なることからもわかる。）
\end{solution}

%================================================
\chapter{基本群と被覆空間}
% 25問
%================================================

\begin{summary}
\begin{itemize}
  \item \textbf{基本群}：$x_0$ を基点とするループの、端点を止めたホモトピー類全体 $\pi_1(X,x_0)$ は、道の積 $[f][g]=[f\cdot g]$ で群になる。単位元は定値ループ、逆元は逆向きの道 $\bar f$。
  \item \textbf{基点の取り替え}：$x_0$ から $x_1$ への道 $h$ に対し $\beta_h\colon\pi_1(X,x_1)\to\pi_1(X,x_0)$, $[f]\mapsto[h\cdot f\cdot\bar h]$ は同型。弧状連結なら $\pi_1(X)$ は基点によらず同型。
  \item \textbf{誘導準同型}：$f\colon(X,x_0)\to(Y,y_0)$ は $f_*[\gamma]=[f\circ\gamma]$ を誘導し、$(g\circ f)_*=g_*\circ f_*$、$(\id)_*=\id$。ホモトピー同値写像は $\pi_1$ の同型を誘導する。
  \item \textbf{基本的な計算}：$\pi_1(S^1)\cong\Z$（$n$ 回まわるループが $n$ に対応）、$\pi_1(X\times Y)\cong\pi_1(X)\times\pi_1(Y)$、凸集合や可縮空間の $\pi_1$ は自明。
  \item \textbf{ファン・カンペンの定理}：$X=U\cup V$（$U,V,U\cap V$ は弧状連結な開集合、$x_0\in U\cap V$）なら
  \[
  \pi_1(X)\cong\pi_1(U)*_{\pi_1(U\cap V)}\pi_1(V).
  \]
  特に、2胞体を語 $w$ に沿って貼ると、基本群は生成元に関係式 $w=1$ を加えた群になる。
  \item \textbf{被覆空間}：$p\colon\tilde X\to X$ は各点に均等に覆われた近傍をもつ全射。道とホモトピーは始点を決めると一意に持ち上がり、$p_*\colon\pi_1(\tilde X)\to\pi_1(X)$ は単射。葉数は指数 $[\pi_1(X):p_*\pi_1(\tilde X)]$ に等しい。
  \item \textbf{持ち上げ判定法}：$Y$ が弧状連結かつ局所弧状連結のとき、$f\colon Y\to X$ が $\tilde X$ に持ち上がる $\iff$ $f_*\pi_1(Y)\subset p_*\pi_1(\tilde X)$。
  \item \textbf{分類}：（良い空間 $X$ について）連結被覆の同型類は $\pi_1(X)$ の部分群の共役類と1対1に対応する。単連結な被覆が普遍被覆で、その被覆変換群は $\pi_1(X)$ と同型。単連結な $Y$ に群 $G$ が被覆空間作用するとき $\pi_1(Y/G)\cong G$。
\end{itemize}
\end{summary}

\begin{problem}[難易度 B]
円周 $S^1$ の基本群を計算せよ。
\end{problem}

\begin{solution}
普遍被覆
\[
p:\R\longrightarrow S^1,
\qquad
p(t)=e^{2\pi it}
\]
を考える。$\R$ を螺旋として描き、真下の円周へ射影したものが $p$ である。
\begin{center}
\begin{tikzpicture}[scale=1.0]
  % 螺旋 R（t in [-0.25, 2.25]）
  \draw[thick,gray,domain=-0.25:2.25,samples=200,smooth]
    plot ({1.6*cos(360*\x)}, {0.45*sin(360*\x)+1.1*\x+1.2});
  % 0 から 2 への持ち上げ（赤）
  \draw[very thick,red!75!black,domain=0:2,samples=160,smooth]
    plot ({1.6*cos(360*\x)}, {0.45*sin(360*\x)+1.1*\x+1.2});
  \foreach \n in {0,1,2}{
    \node[vtx] at (1.6,{1.1*\n+1.2}) {};
    \node[right] at (1.65,{1.1*\n+1.2}) {\footnotesize $\n$};
  }
  \node[left] at (-1.7,3.0) {$\R$};
  % 円周
  \draw[thick] (0,-1.3) ellipse (1.6 and 0.45);
  \draw[very thick,red!75!black,mid arrow] (1.6,-1.3) arc (0:350:1.6 and 0.45);
  \node[vtx] at (1.6,-1.3) {};
  \node[right] at (1.65,-1.3) {\footnotesize $1=p(n)$};
  \node[left] at (-1.7,-1.3) {$S^1$};
  \draw[-{Stealth[length=2.2mm]}] (2.6,1.6) -- (2.6,-0.9) node[midway,right] {$p$};
\end{tikzpicture}
\zucap{$S^1$ を2周するループ $\gamma_2(s)=e^{4\pi is}$ の、$0$ を始点とする持ち上げ $\tilde\gamma_2(s)=2s$（赤）。終点 $2$ が「周回数」を表す。}
\end{center}

基点 $1\in S^1$ のループ $\gamma$ に対し、$\tilde\gamma(0)=0$ となる持ち上げ $\tilde\gamma\colon[0,1]\to\R$ がただ1つ存在し（道の持ち上げ）、$p(\tilde\gamma(1))=1$ より $\tilde\gamma(1)\in p^{-1}(1)=\Z$ である。そこで $\Phi([\gamma])=\tilde\gamma(1)$ とおく。
\begin{itemize}
\item well-defined：$\gamma\simeq\gamma'$（端点を止めたホモトピー $H$）ならば、ホモトピーの持ち上げ $\tilde H$ は各 $s$ ごとに $\tilde H(1,s)\in\Z$ をとり連続なので一定であり、$\tilde\gamma(1)=\tilde\gamma'(1)$。
\item 準同型：$\tilde\gamma(1)=m$ とすると、$\gamma'$ の持ち上げを $m$ だけずらした $\tilde\gamma'+m$ は $\gamma'$ の $m$ を始点とする持ち上げであり、$\tilde\gamma$ に続けてこれをたどると $\gamma\cdot\gamma'$ の持ち上げになる。その終点は $m+\tilde\gamma'(1)$ なので $\Phi([\gamma][\gamma'])=\Phi([\gamma])+\Phi([\gamma'])$。ここで $\tilde\gamma'+m$ が持ち上げになるのは、被覆変換 $T_m(t)=t+m$ が $p\circ T_m=p$ をみたすからである。
\item 全射：$\gamma_n(s)=e^{2\pi ins}$ の持ち上げは $\tilde\gamma_n(s)=ns$ で、$\Phi([\gamma_n])=n$。
\item 単射：$\Phi([\gamma])=0$ ならば $\tilde\gamma$ は $\R$ 内のループで、$\R$ は凸なので直線ホモトピー $\tilde H(s,t)=(1-t)\tilde\gamma(s)$ で定値ループとホモトピック。$p\circ\tilde H$ が $\gamma$ と定値ループの間のホモトピーを与える。
\end{itemize}
したがって
\[
\pi_1(S^1)\cong\Z,\qquad [\gamma]\mapsto\tilde\gamma(1)\ (\text{周回数})
\]
であり、生成元は1周するループ $[\gamma_1]$ である。
\end{solution}

\begin{problem}[難易度 A]
$C\subset\R^n$ を凸集合、$x_0\in C$ とする。$\pi_1(C,x_0)$ が自明群であることを示せ。
\end{problem}

\begin{solution}
$\gamma\colon I\to C$ を $x_0$ を基点とするループとする。$F\colon I\times I\to\R^n$ を
\[
F(s,t)=(1-t)\gamma(s)+t\,x_0
\]
で定める。$\gamma(s), x_0\in C$ で $C$ は凸なので $F(s,t)\in C$ であり、$F$ は連続である。

$F(s,0)=\gamma(s)$、$F(s,1)=x_0$ であり、さらに $\gamma(0)=\gamma(1)=x_0$ より
\[
F(0,t)=F(1,t)=(1-t)x_0+t\,x_0=x_0
\]
である。よって $F$ は $\gamma$ から定値ループ $c_{x_0}$ への端点を固定したホモトピーであり、$[\gamma]=[c_{x_0}]$ である。$\gamma$ は任意なので $\pi_1(C,x_0)$ は自明である。
\end{solution}

\begin{problem}[難易度 A]
$h$ を空間 $X$ の $x_0$ から $x_1$ への道とし、$\bar h(s)=h(1-s)$ とする。$\beta_h\colon\pi_1(X,x_1)\to\pi_1(X,x_0)$ を $\beta_h[f]=[h\cdot f\cdot\bar h]$ で定めるとき、$\beta_h$ が矛盾なく定義された群同型であることを示せ。また、$\beta_h$ は $h$ の（端点を固定した）ホモトピー類のみによることを示せ。
\end{problem}

\begin{solution}
道の積は端点を固定したホモトピー類の上で矛盾なく定義され、結合的で、定値道が単位元、$\bar h$ が逆元としてはたらく（$h\cdot\bar h\simeq c_{x_0}$、$\bar h\cdot h\simeq c_{x_1}$）。以下 $\simeq$ は端点を固定したホモトピーを表す。

well-defined 性：$f\simeq f'$ ならば $h\cdot f\cdot\bar h\simeq h\cdot f'\cdot\bar h$ である（$h$ と $\bar h$ の部分は定値ホモトピーにすればよい）。同様に $h\simeq h'$ ならば $\bar h\simeq\bar h'$ なので $\beta_h=\beta_{h'}$ であり、後半の主張も従う。

準同型であること：
\[
\beta_h([f][g])=[h\cdot f\cdot g\cdot\bar h]=[h\cdot f\cdot\bar h\cdot h\cdot g\cdot\bar h]=\beta_h[f]\,\beta_h[g].
\]
ここで $\bar h\cdot h\simeq c_{x_1}$ を挿入した。

同型であること：$\beta_{\bar h}\colon\pi_1(X,x_0)\to\pi_1(X,x_1)$ を考えると
\[
\beta_h\beta_{\bar h}[f]=[h\cdot\bar h\cdot f\cdot h\cdot\bar h]=[f]
\]
であり、同様に $\beta_{\bar h}\beta_h=\id$ である。よって $\beta_h$ は同型である。
\end{solution}

\begin{problem}[難易度 B]
$X$ を弧状連結な空間、$x_0\in X$ とする。$\pi_1(X,x_0)$ がアーベル群であるための必要十分条件は、任意の $x_1\in X$ と $x_0$ から $x_1$ への任意の道 $h, h'$ に対して $\beta_h=\beta_{h'}$ となることであることを示せ。
\end{problem}

\begin{solution}
（必要性）$\pi_1(X,x_0)$ がアーベル群であるとする。$g=h'\cdot\bar h$ は $x_0$ を基点とするループである。任意の $[f]\in\pi_1(X,x_1)$ に対し、$\bar h\cdot h\simeq c_{x_1}$ を挿入して
\[
\beta_{h'}[f]=[h'\cdot\bar h\cdot h\cdot f\cdot\bar h\cdot h\cdot\bar h']=[g]\,\beta_h[f]\,[g]^{-1}
\]
を得る。$\pi_1(X,x_0)$ はアーベル群なので右辺は $\beta_h[f]$ に等しい。よって $\beta_{h'}=\beta_h$ である。

（十分性）$x_1=x_0$、$h'=c_{x_0}$（定値道）とし、$h$ を $x_0$ を基点とする任意のループとする。$\beta_{c_{x_0}}=\id$ なので、仮定より任意の $[f]\in\pi_1(X,x_0)$ に対し
\[
[h][f][h]^{-1}=\beta_h[f]=[f]
\]
である。すなわち $[h][f]=[f][h]$ であり、$\pi_1(X,x_0)$ はアーベル群である。
\end{solution}

\begin{problem}[難易度 A]
$A\subset X$ がレトラクトである（連続写像 $r\colon X\to A$ で $r|_A=\id_A$ となるものが存在する）とき、$a_0\in A$ に対し、包含 $i\colon A\to X$ が誘導する $i_*\colon\pi_1(A,a_0)\to\pi_1(X,a_0)$ は単射であり、$r_*\colon\pi_1(X,a_0)\to\pi_1(A,a_0)$ は全射であることを示せ。
\end{problem}

\begin{solution}
誘導準同型は関手的である。すなわち基点を保つ連続写像 $\varphi\colon(X,x_0)\to(Y,y_0)$、$\psi\colon(Y,y_0)\to(Z,z_0)$ に対して
\[
(\psi\circ\varphi)_*=\psi_*\circ\varphi_*,\qquad (\id_X)_*=\id
\]
が成り立つ。実際、$(\psi\circ\varphi)_*[\gamma]=[\psi\circ\varphi\circ\gamma]=\psi_*[\varphi\circ\gamma]=\psi_*\varphi_*[\gamma]$ であり、後者は明らかである。

$r\circ i=\id_A$ で、$r(a_0)=a_0$ なので $r$ は基点を保つ。よって
\[
r_*\circ i_*=(r\circ i)_*=(\id_A)_*=\id_{\pi_1(A,a_0)}
\]
である。$r_*\circ i_*$ が単射なので $i_*$ は単射であり、$r_*\circ i_*$ が全射なので $r_*$ は全射である。
\end{solution}

\begin{problem}[難易度 B]
$x_0\in X$、$y_0\in Y$ に対し、$\pi_1(X\times Y,(x_0,y_0))\cong\pi_1(X,x_0)\times\pi_1(Y,y_0)$ であることを示せ。これを用いて $n$ 次元トーラス $T^n=(S^1)^n$ の基本群を計算せよ。
\end{problem}

\begin{solution}
$p_X, p_Y$ を射影とし、$\Phi=(p_{X*},p_{Y*})$、すなわち
\[
\Phi[\gamma]=([p_X\circ\gamma],[p_Y\circ\gamma])
\]
とおく。各成分は誘導準同型なので $\Phi$ は準同型である。

全射性：$x_0$ を基点とするループ $\alpha$ と $y_0$ を基点とするループ $\beta$ に対し、$\gamma(s)=(\alpha(s),\beta(s))$ は $X\times Y$ のループで、$\Phi[\gamma]=([\alpha],[\beta])$ である。

単射性：$\Phi[\gamma]$ が単位元とする。$F$ を $p_X\circ\gamma$ から定値ループへの、$G$ を $p_Y\circ\gamma$ から定値ループへの端点を固定したホモトピーとすると、$H(s,t)=(F(s,t),G(s,t))$ は連続で、$\gamma$ から定値ループへの端点を固定したホモトピーである。よって $[\gamma]$ は単位元である。

以上より $\Phi$ は同型である。帰納的に $\pi_1(X_1\times\dots\times X_n)\cong\prod_k\pi_1(X_k)$ であり、$\pi_1(S^1)\cong\Z$ より $\pi_1(T^n)\cong\Z^n$ である。
\end{solution}

\begin{problem}[難易度 B]
$X$ を弧状連結な空間とする。$X$ が単連結であるための必要十分条件は、$X$ の任意の 2 点 $x_0, x_1$ と $x_0$ から $x_1$ への任意の道 $\gamma, \delta$ が、端点を固定してホモトピックであることであることを示せ。
\end{problem}

\begin{solution}
以下 $\simeq$ は端点を固定したホモトピーを表す。

（必要性）$X$ が単連結ならば、$\gamma\cdot\bar\delta$ は $x_0$ を基点とするループなので $\gamma\cdot\bar\delta\simeq c_{x_0}$ である。よって
\[
\gamma\simeq\gamma\cdot c_{x_1}\simeq\gamma\cdot\bar\delta\cdot\delta\simeq c_{x_0}\cdot\delta\simeq\delta
\]
を得る。

（十分性）$x_1=x_0$ とし、$x_0$ を基点とする任意のループ $\gamma$ と定値ループ $\delta=c_{x_0}$ に仮定を適用すれば $\gamma\simeq c_{x_0}$、すなわち $\pi_1(X,x_0)=1$ である。$X$ は弧状連結なので基本群は基点によらず、$X$ は単連結である。
\end{solution}

\begin{problem}[難易度 C]
連続写像 $f, g\colon X\to Y$ がホモトピー $H\colon X\times I\to Y$（$H(\cdot,0)=f$、$H(\cdot,1)=g$）でホモトピックであるとし、$x_0\in X$、$h(t)=H(x_0,t)$ とおく。このとき
\[
f_*=\beta_h\circ g_*\colon\pi_1(X,x_0)\to\pi_1(Y,f(x_0))
\]
であることを示せ。ここで $\beta_h[\delta]=[h\cdot\delta\cdot\bar h]$ である。
\end{problem}

\begin{solution}
$h$ は $f(x_0)$ から $g(x_0)$ への道なので、$\beta_h\colon\pi_1(Y,g(x_0))\to\pi_1(Y,f(x_0))$ である。$\gamma$ を $x_0$ を基点とするループとし、$F\colon I\times I\to Y$、$F(s,t)=H(\gamma(s),t)$ とおく。

正方形 $I\times I$ 内の次の二つの道を考える：下辺 $b(s)=(s,0)$ と、左辺 $(0,t)$ を上り、上辺 $(s,1)$ を進み、右辺 $(1,t)$ を下る道 $c$。どちらも $(0,0)$ から $(1,0)$ への道であり、$I\times I$ は凸なので、直線ホモトピーにより $b\simeq c$（端点固定）である。これを $F$ で写すと $F\circ b\simeq F\circ c$ となる。

ここで $\gamma(0)=\gamma(1)=x_0$ なので、
\[
F(s,0)=f(\gamma(s)),\quad F(0,t)=F(1,t)=h(t),\quad F(s,1)=g(\gamma(s))
\]
であり、$F\circ b=f\circ\gamma$、$F\circ c=h\cdot(g\circ\gamma)\cdot\bar h$ である。よって
\[
f_*[\gamma]=[f\circ\gamma]=[h\cdot(g\circ\gamma)\cdot\bar h]=\beta_h(g_*[\gamma])
\]
である。
\end{solution}

\begin{problem}[難易度 C]
$f\colon X\to Y$ がホモトピー同値写像ならば、任意の $x_0\in X$ に対して $f_*\colon\pi_1(X,x_0)\to\pi_1(Y,f(x_0))$ は同型であることを示せ。特に、可縮な空間は単連結であることを示せ。
\end{problem}

\begin{solution}
$g\colon Y\to X$ を $g\circ f\simeq\id_X$、$f\circ g\simeq\id_Y$ をみたす写像とする。ホモトピックな写像 $\varphi\simeq\psi$ について、$\varphi_*$ は $\psi_*$ と基点の取り替え同型 $\beta_h$ の合成に等しい（前問）。基点の取り替えは同型なので、$\varphi_*$ が同型であることと $\psi_*$ が同型であることは同値である。

$g\circ f\simeq\id_X$ より、合成
\[
\pi_1(X,x_0)\xrightarrow{f_*}\pi_1(Y,f(x_0))\xrightarrow{g_*}\pi_1(X,gf(x_0))
\]
は同型である。よって $f_*$ は単射で、この $g_*$ は全射である。次に $f\circ g\simeq\id_Y$ を基点 $f(x_0)$ で用いると、合成
\[
\pi_1(Y,f(x_0))\xrightarrow{g_*}\pi_1(X,gf(x_0))\xrightarrow{f_*}\pi_1(Y,fgf(x_0))
\]
は同型であり、この $g_*$ は単射である。以上よりこの $g_*$ は同型であり、$f_*=g_*^{-1}\circ(g_*\circ f_*)$ も同型である。

可縮な空間 $X$ は一点とホモトピー同値なので（一点への写像がホモトピー同値写像）、各基点で $\pi_1(X,x_0)$ は一点の基本群、すなわち自明群と同型である。また可縮な空間は弧状連結なので、$X$ は単連結である。
\end{solution}

\begin{problem}[難易度 B]
$S^1=\{z\in\mathbb{C}:\lvert z\rvert=1\}$、$n\in\Z$ とし、$f_n\colon S^1\to S^1$、$f_n(z)=z^n$ とする。同型 $\pi_1(S^1,1)\cong\Z$ のもとで、$(f_n)_*$ が $n$ 倍写像 $m\mapsto nm$ であることを示せ。
\end{problem}

\begin{solution}
普遍被覆 $p\colon\R\to S^1$、$p(s)=e^{2\pi is}$ を用いる。同型 $\Phi\colon\pi_1(S^1,1)\to\Z$ は、ループ $\gamma$ の $0$ を始点とする持ち上げ $\tilde\gamma$ に対し $\Phi[\gamma]=\tilde\gamma(1)$ で与えられる。

$m\in\Z$ に対し $\omega_m(s)=e^{2\pi ims}$ とおくと、$\tilde\omega_m(s)=ms$ はその持ち上げで、$\tilde\omega_m(1)=m$ なので $\Phi[\omega_m]=m$ である。$\Phi$ は全単射なので、$\pi_1(S^1,1)$ の元はすべて $[\omega_m]$ の形である。

$f_n(1)=1$ であり、
\[
f_n\circ\omega_m(s)=\bigl(e^{2\pi ims}\bigr)^n=e^{2\pi inms}=\omega_{nm}(s)
\]
である。よって $(f_n)_*[\omega_m]=[\omega_{nm}]$ であり、$\Phi$ で移せば $(f_n)_*$ は $m\mapsto nm$ である。
\end{solution}

\begin{problem}[難易度 B]
$D^2$ から $S^1=\partial D^2$ へのレトラクションが存在しないことを示せ。これを用いて、任意の連続写像 $f\colon D^2\to D^2$ が不動点をもつこと（2 次元ブラウアーの不動点定理）を示せ。
\end{problem}

\begin{solution}
レトラクション $r\colon D^2\to S^1$ があったとする。包含 $i\colon S^1\to D^2$ について $r\circ i=\id_{S^1}$ なので、基点 $1\in S^1$ に関して $r_*\circ i_*=\id_{\pi_1(S^1,1)}$ である。しかし $D^2$ は凸なので $\pi_1(D^2,1)=1$ であり、$\id_\Z$ が自明群を経由することになって矛盾する。

次に $f$ が不動点をもたないとする。$x\in D^2$ に対し $u=x-f(x)\ne0$ とし、$f(x)$ から $x$ に向かう半直線が $S^1$ と交わる点を $r(x)=x+\lambda u$（$\lambda\ge0$）とする。$\lvert x+\lambda u\rvert^2=1$ を解くと
\[
\lambda(x)=\frac{-\langle x,u\rangle+\sqrt{\langle x,u\rangle^2+\lvert u\rvert^2(1-\lvert x\rvert^2)}}{\lvert u\rvert^2}
\]
である（根号の中は非負で、$\lambda\ge0$ となるのはこの根）。これは $x$ について連続なので、$r\colon D^2\to S^1$ は連続である。

$\lvert x\rvert=1$ ならば、コーシー・シュワルツの不等式より $\langle x,u\rangle=1-\langle x,f(x)\rangle\ge0$ であり、$1-\lvert x\rvert^2=0$ なので $\lambda(x)=(-\langle x,u\rangle+\lvert\langle x,u\rangle\rvert)/\lvert u\rvert^2=0$、すなわち $r(x)=x$ である。よって $r$ はレトラクションとなり、前半に矛盾する。したがって $f$ は不動点をもつ。
\end{solution}

\begin{problem}[難易度 C]
メビウスの帯 $M=[0,1]\times[-1,1]/{\sim}$（$(0,y)\sim(1,-y)$）の境界 $\partial M=\{[s,\pm1]\}$ は $M$ のレトラクトでないことを示せ。
\end{problem}

\begin{solution}
中心線 $C=\{[s,0]\}$ は $S^1$ と同相で、$r\colon M\to C$、$r[s,y]=[s,0]$ は変位レトラクション（$H([s,y],t)=[s,(1-t)y]$）である。ループ $c(s)=[s,0]$ は $\pi_1(C,[0,0])\cong\Z$ の生成元を代表する。$r$ はホモトピー同値写像なので、境界上の点 $[0,1]$ を基点として $r_*\colon\pi_1(M,[0,1])\to\pi_1(C,[0,0])\cong\Z$ は同型である。

境界は次の道 $\beta\colon[0,2]\to\partial M$ で一周される：
\[
\beta(u)=[u,1]\ (0\le u\le1),\qquad \beta(u)=[u-1,-1]\ (1\le u\le2).
\]
$\beta(1)=[1,1]=[0,-1]$、$\beta(2)=[1,-1]=[0,1]=\beta(0)$ であり、$\beta$ は $[0,2]/\{0,2\}\cong S^1$ から $\partial M$ への連続全単射、したがって同相を誘導する。よって $[\beta]$ は $\pi_1(\partial M,[0,1])\cong\Z$ の生成元である。$r\circ\beta$ は $[u,0]$ $(0\le u\le1)$ に続いて $[u-1,0]$ $(1\le u\le2)$ をたどるので、$c\cdot c$ と同じループであり、$r_*i_*[\beta]=[c]^2$ である（$i\colon\partial M\to M$ は包含）。したがって $r_*$ により $\pi_1(M,[0,1])\cong\Z$ と同一視すると、$i_*\colon\pi_1(\partial M,[0,1])\cong\Z\to\Z$ は $2$ 倍写像である。

レトラクション $\rho\colon M\to\partial M$ があれば、基点 $[0,1]$ について $\rho_*\circ i_*=\id_\Z$ である。しかし $i_*$ の像は $2\Z$ に含まれるので、$\rho_*i_*(1)=\rho_*(2)=2\rho_*(1)\in2\Z$ となり、$1$ に等しくなりえない。よって $\partial M$ はレトラクトでない。
\end{solution}

\begin{problem}[難易度 B]
$n\ge2$ のとき、$S^n$ が単連結であることをファン・カンペンの定理を用いて示せ。
\end{problem}

\begin{solution}
$N=(0,\dots,0,1)$、$S=-N$ とし、$U=S^n\setminus\{N\}$、$V=S^n\setminus\{S\}$ とおく。これらは開集合で $U\cup V=S^n$ である。立体射影により $U\cong\R^n$、$V\cong\R^n$ なので、$U, V$ は単連結である。

$U\cap V=S^n\setminus\{N,S\}$ は立体射影により $\R^n\setminus\{0\}$ と同相である。$n\ge2$ のとき $\R^n\setminus\{0\}$ は弧状連結である（$x\mapsto x/\lvert x\rvert$ により各点は $S^{n-1}$ に道で結ばれ、$S^{n-1}$ は $n-1\ge1$ より弧状連結）。

基点 $x_0\in U\cap V$ をとる。ファン・カンペンの定理（の全射性の部分）により、$\pi_1(U,x_0)$ と $\pi_1(V,x_0)$ の像は $\pi_1(S^n,x_0)$ を生成する。どちらも自明群なので $\pi_1(S^n,x_0)=1$ であり、$S^n$ は弧状連結なので単連結である。
\end{solution}

\begin{problem}[難易度 B]
$n\ge0$、$n\ne2$ ならば $\R^2$ と $\R^n$ は同相でないことを示せ。
\end{problem}

\begin{solution}
$n=0$ のとき $\R^0$ は一点なので、$\R^2$ と同相でない。$n\ge1$ とし、同相写像 $h\colon\R^2\to\R^n$ があったとする。$h$ は同相 $\R^2\setminus\{0\}\cong\R^n\setminus\{h(0)\}$ を誘導し、平行移動により $\R^n\setminus\{h(0)\}\cong\R^n\setminus\{0\}$ である。

$n=1$ のとき：$\R\setminus\{0\}$ は連結でないが、$\R^2\setminus\{0\}$ は弧状連結（$S^1$ への強変位レトラクトがあり、$S^1$ は弧状連結）なので連結である。矛盾。

$n\ge3$ のとき：$\R^n\setminus\{0\}$ は $S^{n-1}$ とホモトピー同値で、$n-1\ge2$ なので $S^{n-1}$ は単連結である。一方 $\R^2\setminus\{0\}$ は $S^1$ とホモトピー同値で、$\pi_1\cong\Z$ である。ホモトピー同値写像（特に同相写像）は基本群の同型を誘導するので、$\Z\cong1$ となり矛盾する。

よって $n\ne2$ ならば $\R^2\not\cong\R^n$ である。
\end{solution}

\begin{problem}[難易度 C]
すべての成分が正である $3\times3$ 実行列 $A$ は正の固有値をもつことを、ブラウアーの不動点定理を用いて示せ。
\end{problem}

\begin{solution}
$T=\{x\in\R^3:\lvert x\rvert=1,\ x_1,x_2,x_3\ge0\}$ とおく。$\Delta=\{y\in\R^3:y_i\ge0,\ y_1+y_2+y_3=1\}$ とすると、$x\mapsto x/(x_1+x_2+x_3)$ と $y\mapsto y/\lvert y\rvert$ は互いに逆な連続写像で、$T\cong\Delta$ である。$\Delta$ は三角形なので $D^2$ と同相であり（重心からの放射状の拡大縮小）、ブラウアーの不動点定理は $T$ に対しても成り立つ。

$x\in T$ に対し、$Ax$ の第 $i$ 成分は $\sum_j a_{ij}x_j$ である。$a_{ij}>0$ で、$x_j\ge0$ かつ少なくとも一つの $x_j$ は正なので、$Ax$ の各成分は正である。特に $Ax\ne0$ であり、
\[
f(x)=\frac{Ax}{\lvert Ax\rvert}
\]
は連続写像 $T\to T$ を定める。

ブラウアーの不動点定理より $f(x)=x$ となる $x\in T$ が存在する。このとき $Ax=\lvert Ax\rvert x$ であり、$\lambda=\lvert Ax\rvert>0$ は $A$ の固有値（固有ベクトル $x\ne0$）である。
\end{solution}

\begin{problem}[難易度 B]
$X=S^1\vee S^1$ の基本群が階数 2 の自由群 $F_2=\Z*\Z$ であることを、ファン・カンペンの定理を用いて示せ。
\end{problem}

\begin{solution}
$X=A\cup B$、$A\cap B=\{x_0\}$（$A, B$ は円周）とし、$a\in A$、$b\in B$ を $x_0$ と異なる点とする。$U=X\setminus\{b\}$、$V=X\setminus\{a\}$ とおくと、これらは $X$ の開集合で $U\cup V=X$、$x_0\in U\cap V$ である。

$B\setminus\{b\}$ は $x_0$ を内部に含む開区間と同相なので、$x_0$ を固定して $x_0$ に強変位レトラクトする。このホモトピーを $A$ 上で恒等として延長すれば（$A$ と $B\setminus\{b\}$ は $U$ の閉集合で共通部分は $\{x_0\}$）、$U$ は $A$ に強変位レトラクトする。よって $\pi_1(U,x_0)\cong\pi_1(A,x_0)\cong\Z$ であり、同様に $\pi_1(V,x_0)\cong\Z$ である。また $U\cap V=(A\setminus\{a\})\cup(B\setminus\{b\})$ は二つの開区間を一点 $x_0$ で貼ったもので、同様にして $x_0$ に強変位レトラクトするから単連結である。$U, V, U\cap V$ はいずれも弧状連結である。

ファン・カンペンの定理より $\pi_1(X,x_0)\cong\pi_1(U,x_0)*\pi_1(V,x_0)\cong\Z*\Z$ であり、生成元は $A$ を一周するループと $B$ を一周するループである。
\end{solution}

\begin{problem}[難易度 B]
$p=(-1,0)$、$q=(1,0)$ とする。$\pi_1(\R^2\setminus\{p,q\})$ を計算せよ。
\end{problem}

\begin{solution}
$X=\R^2\setminus\{p,q\}$ とし、
\[
U=\{(x,y):x<1/2\}\setminus\{p\},\qquad V=\{(x,y):x>-1/2\}\setminus\{q\}
\]
とおく。これらは $X$ の開集合で $U\cup V=X$ である。$U\cap V=\{(x,y):-1/2<x<1/2\}$（$p, q$ はこの帯に含まれない）は凸なので弧状連結かつ単連結である。基点を $(0,0)$ にとる。

$\phi(x)=-\log(1/2-x)$ は $(-\infty,1/2)$ から $\R$ への同相なので、$(x,y)\mapsto(\phi(x),y)$ は半平面 $\{x<1/2\}$ から $\R^2$ への同相であり、$U\cong\R^2\setminus\{(\phi(-1),0)\}$ となる。一点を除いた平面は $S^1$ とホモトピー同値なので $\pi_1(U)\cong\Z$ である。同様に $\pi_1(V)\cong\Z$ である。$U, V$ は弧状連結である。

ファン・カンペンの定理により、$\pi_1(U\cap V)$ が自明であるから
\[
\pi_1(X,(0,0))\cong\pi_1(U)*\pi_1(V)\cong\Z*\Z=F_2
\]
である。生成元は $p$ のまわりを一周するループと $q$ のまわりを一周するループである。
\end{solution}

\begin{problem}[難易度 B]
$X=\R^2\setminus\{p,q\}$（$p=(-1,0)$、$q=(1,0)$）、基点を $x_0=(0,-1)$ とする。$a$ を $p$ のまわりを反時計回りに1周するループ、$b$ を $q$ のまわりを反時計回りに1周するループとし（下図左）、前問より $\pi_1(X,x_0)$ は $a,b$ を生成元とする自由群であることを用いてよい。道の積 $f\cdot g$ は $f$ を先にたどるものとする。
\begin{center}
\begin{tikzpicture}[scale=0.85]
  \begin{scope}
    \draw[thick,red!75!black,mid arrow] (0,-1) .. controls (-0.1,0.3) and (-0.5,0.9) .. (-1,0.9)
       .. controls (-1.7,0.9) and (-2,0) .. (-1.7,-0.6) .. controls (-1.3,-1.1) and (-0.4,-1.1) .. (0,-1);
    \draw[thick,blue!75!black,mid arrow] (0,-1) .. controls (0.4,-1.1) and (1.3,-1.1) .. (1.7,-0.6)
       .. controls (2,0) and (1.7,0.9) .. (1,0.9) .. controls (0.5,0.9) and (0.1,0.3) .. (0,-1);
    \node[red!75!black] at (-2.0,0.6) {$a$};
    \node[blue!75!black] at (2.0,0.6) {$b$};
    \node[hvtx] at (-1,0) {}; \node[below] at (-1,0) {\footnotesize $p$};
    \node[hvtx] at (1,0) {};  \node[below] at (1,0) {\footnotesize $q$};
    \node[vtx] at (0,-1) {};  \node[below] at (0,-1) {\footnotesize $x_0$};
  \end{scope}
  \begin{scope}[xshift=5.2cm]
    \draw[thick,mid arrow] (0,-1) .. controls (1.6,-1.05) and (2.2,-0.6) .. (2.2,0)
       .. controls (2.2,0.8) and (1.2,1.1) .. (0,1.1) .. controls (-1.2,1.1) and (-2.2,0.8) .. (-2.2,0)
       .. controls (-2.2,-0.6) and (-1.6,-1.05) .. (0,-1);
    \node at (0,1.4) {$\gamma_1$};
    \node[hvtx] at (-1,0) {}; \node[below] at (-1,0) {\footnotesize $p$};
    \node[hvtx] at (1,0) {};  \node[below] at (1,0) {\footnotesize $q$};
    \node[vtx] at (0,-1) {};  \node[below] at (0,-1) {\footnotesize $x_0$};
  \end{scope}
  \begin{scope}[xshift=10.4cm]
    \draw[thick,mid arrow] (0,-1) .. controls (-0.1,0.3) and (-0.5,0.9) .. (-1,0.9)
       .. controls (-1.7,0.9) and (-2,0) .. (-1.7,-0.6) .. controls (-1.3,-1.1) and (-0.4,-1.1) .. (0,-1);
    \draw[thick,mid arrow] (0,-1) .. controls (0.1,0.3) and (0.5,0.9) .. (1,0.9)
       .. controls (1.7,0.9) and (2,0) .. (1.7,-0.6) .. controls (1.3,-1.1) and (0.4,-1.1) .. (0,-1);
    \node at (0,1.4) {$\gamma_2$（左の輪が先）};
    \node[hvtx] at (-1,0) {}; \node[below] at (-1,0) {\footnotesize $p$};
    \node[hvtx] at (1,0) {};  \node[below] at (1,0) {\footnotesize $q$};
    \node[vtx] at (0,-1) {};  \node[below] at (0,-1) {\footnotesize $x_0$};
  \end{scope}
\end{tikzpicture}
\zucap{左：生成元 $a,b$。中：両方の点を囲む反時計回りのループ $\gamma_1$。右：8の字のループ $\gamma_2$（左の輪をたどってから右の輪をたどる）。}
\end{center}
\begin{enumerate}
\item $[\gamma_1]$、$[\gamma_2]$ を $a,b$ の語で表せ。
\item ループ $\gamma$ の $p$ のまわりの回転数を $w_p(\gamma)$、$q$ のまわりの回転数を $w_q(\gamma)$ とする（$w_p(\gamma)$ は、$\gamma(s)-p=r(s)e^{i\theta(s)}$ と偏角を連続に選んだときの $(\theta(1)-\theta(0))/2\pi$）。$w=(w_p,w_q)\colon\pi_1(X,x_0)\to\Z^2$ が well-defined な準同型であることを認めて、$w(\gamma_1)$、$w(\gamma_2)$ を求めよ。
\item 交換子 $c=aba^{-1}b^{-1}$ は $w(c)=(0,0)$ をみたすが、定値ループとホモトピックでないことを示せ。
\end{enumerate}
\end{problem}

\begin{solution}
(1) $\gamma_1$ の上半分を、$p$ と $q$ の間の点 $(0,0)$ を通るように引き下げる（$p,q$ を横切らないのでホモトピー）。
\begin{center}
\begin{tikzpicture}[scale=0.8]
  \draw[thick,blue!75!black,mid arrow] (0,-1) .. controls (0.4,-1.1) and (1.3,-1.1) .. (1.7,-0.6)
     .. controls (2,0) and (1.7,0.9) .. (1,0.9) .. controls (0.5,0.9) and (0.15,0.4) .. (0.08,-0.95);
  \draw[thick,red!75!black,mid arrow] (-0.08,-0.95) .. controls (-0.15,0.4) and (-0.5,0.9) .. (-1,0.9)
     .. controls (-1.7,0.9) and (-2,0) .. (-1.7,-0.6) .. controls (-1.3,-1.1) and (-0.4,-1.1) .. (0,-1);
  \node[hvtx] at (-1,0) {}; \node[hvtx] at (1,0) {};
  \node[vtx] at (0,-1) {};
  \node[blue!75!black] at (2.1,0.5) {$b$};
  \node[red!75!black] at (-2.1,0.5) {$a$};
  \node at (3.6,0) {$\gamma_1\simeq b\cdot a$};
\end{tikzpicture}
\zucap{中央の縦の往復（下りと上り）は打ち消し合うので、$\gamma_1$ はまず $b$、次に $a$ をたどるループとホモトピック。}
\end{center}
$\gamma_1$ は $x_0$ から右へ出て $q$ の右側を上り、上を左へ通って $p$ の左側を下り $x_0$ に戻る。上の図のように中央に「下って上る」往復を挿入すると、前半は $q$ を反時計回りに1周する $b$、後半は $p$ を反時計回りに1周する $a$ になる。往復 $\delta\cdot\bar\delta$ は定値の道とホモトピックなので
\[
[\gamma_1]=ba .
\]
$\gamma_2$ の左の輪は $a$ そのものである。右の輪は $x_0$ から中央を上り、$q$ の上を右へ通って右側を下り、下を通って戻る。これは $q$ のまわりを時計回りに1周するループで、$b$ を逆向きにたどったもの $\bar b$ に等しい。よって
\[
[\gamma_2]=ab^{-1}.
\]

(2) $w$ は準同型で、$w(a)=(1,0)$、$w(b)=(0,1)$ である（$a$ は半平面 $\{x<1\}$ に含まれ、そこでは $q$ から見た偏角を $(\pi/2,3\pi/2)$ の範囲で連続に選べるので $w_q(a)=0$。$w_p(b)=0$ も同様）。よって
\[
w(\gamma_1)=w(b)+w(a)=(1,1),\qquad w(\gamma_2)=w(a)-w(b)=(1,-1).
\]

(3) $\Z^2$ は可換なので $w(c)=w(a)+w(b)-w(a)-w(b)=(0,0)$。一方、$\pi_1(X,x_0)$ は $a,b$ 上の自由群であり、自由群の元は既約な語で一意に表される。$aba^{-1}b^{-1}$ は隣り合う文字が互いに逆元でない既約語で、空語ではないから $c\neq1$。すなわち $c$ は定値ループとホモトピックでない。

回転数（あとで見る $H_1$ のクラス）は $\pi_1$ を可換化した情報しかもたないので、$c$ のような交換子は検出できない（第4章のフレヴィッチの定理 $H_1\cong\pi_1^{\mathrm{ab}}$ を参照）。
\end{solution}

\begin{problem}[難易度 C]
種数 $g\ge1$ の向き付け可能閉曲面 $M_g$ は、一つの 0 胞体、$2g$ 個の 1 胞体 $a_1,b_1,\dots,a_g,b_g$、および語 $[a_1,b_1]\cdots[a_g,b_g]$（$[a,b]=aba^{-1}b^{-1}$）に沿って貼られた一つの 2 胞体からなる CW 構造をもつ。$\pi_1(M_g)$ とその可換化を計算し、$g\ne h$ ならば $M_g$ と $M_h$ はホモトピー同値でないことを示せ。
\end{problem}

\begin{solution}
1 骨格 $X^1$ は $2g$ 個の円周の一点和であり、ファン・カンペンの定理を繰り返し用いれば（$S^1\vee S^1$ の場合と同様）、$\pi_1(X^1)$ は $a_1,b_1,\dots,a_g,b_g$ で生成される自由群である。弧状連結な空間 $X$ に接着写像 $\varphi\colon S^1\to X$ で 2 胞体を貼ると、$\pi_1(X\cup_\varphi D^2)\cong\pi_1(X)/N$（$N$ は $[\varphi]$ の生成する正規部分群）となる（ファン・カンペンの定理の帰結）。よって
\[
\pi_1(M_g)\cong\langle a_1,b_1,\dots,a_g,b_g\mid[a_1,b_1]\cdots[a_g,b_g]\rangle
\]
である。

可換化は、生成元を同じにして関係式を可換群の中で課したもの、すなわち $\Z^{2g}$ を関係子の像で割ったものである。可換群では各交換子 $[a_i,b_i]$ は $a_i+b_i-a_i-b_i=0$ に写るので関係子の像は $0$ であり、
\[
\pi_1(M_g)^{\mathrm{ab}}\cong\Z^{2g}
\]
である。

$M_g\simeq M_h$ ならば基本群が同型、したがって可換化も同型で $\Z^{2g}\cong\Z^{2h}$ である。$\Z^m/2\Z^m$ の位数は $2^m$ なので $2g=2h$、すなわち $g=h$ である。（$g=0$ の $M_0=S^2$ は単連結で、可換化 $0=\Z^0$ として同じ議論に含まれる。）
\end{solution}

\begin{problem}[難易度 C]
クラインの壺 $K$ は、一つの 0 胞体、二つの 1 胞体 $a, b$、語 $abab^{-1}$ に沿って貼られた一つの 2 胞体からなる CW 構造をもつ。$\pi_1(K)$ がアーベル群でないことを示せ。
\end{problem}

\begin{solution}
前問と同様に、1 骨格 $S^1\vee S^1$ の基本群は $a, b$ を生成元とする自由群で、2 胞体を貼ることにより
\[
G=\pi_1(K)\cong\langle a,b\mid abab^{-1}\rangle
\]
である。関係 $abab^{-1}=1$ は $bab^{-1}=a^{-1}$ と同値である。

$\R$ の等長変換群において、$t(x)=x+1$、$\rho(x)=-x$ とおくと
\[
\rho t\rho^{-1}(x)=-((-x)+1)=x-1=t^{-1}(x)
\]
なので $t\rho t\rho^{-1}=\id$ である。よって $a\mapsto t$、$b\mapsto\rho$ は準同型 $\psi\colon G\to\operatorname{Isom}(\R)$ を定める（関係子が単位元に写るため）。

もし $G$ がアーベル群なら $bab^{-1}=a$ なので $a=a^{-1}$、すなわち $a^2=1$ であり、$t^2=\psi(a^2)=\id$ となる。しかし $t^2(x)=x+2\ne x$ なので矛盾する。よって $ab\ne ba$ であり、$\pi_1(K)$ はアーベル群でない。
\end{solution}

\begin{problem}[難易度 B]
$p\colon\tilde X\to X$ を被覆写像、$\tilde x_0\in\tilde X$、$x_0=p(\tilde x_0)$ とする。$p_*\colon\pi_1(\tilde X,\tilde x_0)\to\pi_1(X,x_0)$ が単射であることを示せ。また、$x_0$ を基点とするループ $\gamma$ について、$[\gamma]\in p_*\pi_1(\tilde X,\tilde x_0)$ であることと、$\gamma$ の $\tilde x_0$ を始点とする持ち上げ $\tilde\gamma$ がループであることが同値であることを示せ。
\end{problem}

\begin{solution}
被覆写像の道の持ち上げの一意性と、ホモトピー持ち上げ性質（$F\colon I\times I\to X$ と $F(\cdot,0)$ の持ち上げが与えられれば、$F$ の持ち上げ $\tilde F$ が一意に存在する）を用いる。また定値道の持ち上げは定値道である（一意性による）。

主張「$G\colon I\times I\to X$ が端点を固定したホモトピーで、$\tilde G$ がその持ち上げならば、$\tilde G$ も端点を固定したホモトピーで、$\tilde G(\cdot,1)$ は $G(\cdot,1)$ の $\tilde G(0,0)$ を始点とする持ち上げである」：$\tilde G(0,\cdot)$、$\tilde G(1,\cdot)$ は定値道の持ち上げなので定値である。これから従う。

単射性：$\tilde\gamma$ を $\tilde x_0$ を基点とするループで $p\circ\tilde\gamma\simeq c_{x_0}$ とし、そのホモトピー $G$ を $\tilde\gamma$ から持ち上げる。主張より $\tilde\gamma$ は $c_{x_0}$ の持ち上げ、すなわち定値ループ $c_{\tilde x_0}$ と端点固定でホモトピックであり、$[\tilde\gamma]=1$ である。

後半：$\tilde\gamma$ がループならば $[\gamma]=p_*[\tilde\gamma]$ である。逆に $[\gamma]=p_*[\tilde\delta]$（$\tilde\delta$ は $\tilde x_0$ を基点とするループ）とし、$p\circ\tilde\delta$ から $\gamma$ への端点固定ホモトピーを $\tilde\delta$ から持ち上げる。主張より終わりの道は $\tilde x_0$ を始点とする $\gamma$ の持ち上げなので $\tilde\gamma$ に等しく、その終点は $\tilde\delta(1)=\tilde x_0$ である。よって $\tilde\gamma$ はループである。
\end{solution}

\begin{problem}[難易度 C]
$p\colon\tilde X\to X$ を被覆写像とし、$\tilde X, X$ は弧状連結とする。$\tilde x_0\in\tilde X$、$x_0=p(\tilde x_0)$、$H=p_*\pi_1(\tilde X,\tilde x_0)$ とおく。被覆の葉の数 $\lvert p^{-1}(x_0)\rvert$ が指数 $[\pi_1(X,x_0):H]$ に等しいことを示せ。
\end{problem}

\begin{solution}
$G=\pi_1(X,x_0)$ とし、$\Phi\colon G\to p^{-1}(x_0)$ を $\Phi[\gamma]=\tilde\gamma(1)$（$\tilde\gamma$ は $\gamma$ の $\tilde x_0$ を始点とする持ち上げ）で定める。端点を固定してホモトピックな道の持ち上げは同じ終点をもつ（ホモトピー持ち上げ性質）ので、$\Phi$ は well-defined である。

全射性：$\tilde x\in p^{-1}(x_0)$ に対し、$\tilde X$ は弧状連結なので $\tilde x_0$ から $\tilde x$ への道 $\tilde\gamma$ がある。$\gamma=p\circ\tilde\gamma$ はループで $\Phi[\gamma]=\tilde x$ である。

ファイバー：$\Phi[\gamma]=\Phi[\delta]\iff H[\gamma]=H[\delta]$ を示す。$\tilde\gamma(1)=\tilde\delta(1)$ ならば、$\tilde\gamma\cdot\bar{\tilde\delta}$ は $\gamma\cdot\bar\delta$ の持ち上げで $\tilde x_0$ を基点とするループなので、$[\gamma][\delta]^{-1}\in H$ である。逆に $[\gamma\cdot\bar\delta]\in H$ ならば、前問よりその持ち上げはループである。この持ち上げは $\tilde\gamma$ に続いて $\bar\delta$ の $\tilde\gamma(1)$ を始点とする持ち上げ $\eta$ をたどるもので、$\eta$ の終点は $\tilde x_0$ である。$\bar\eta$ は $\delta$ の $\tilde x_0$ を始点とする持ち上げなので $\bar\eta=\tilde\delta$ であり、$\tilde\delta(1)=\eta(0)=\tilde\gamma(1)$ である。

$[\gamma][\delta]^{-1}\in H$ は右剰余類が一致すること $H[\gamma]=H[\delta]$ と同値なので、$\Phi$ は右剰余類の集合 $H\backslash G$ から $p^{-1}(x_0)$ への全単射を誘導する。よって $\lvert p^{-1}(x_0)\rvert=[G:H]$ である。
\end{solution}

\begin{problem}[難易度 D]
群 $G$ が単連結かつ弧状連結な空間 $Y$ に被覆空間作用（各 $y\in Y$ が開近傍 $U$ をもち、$e\ne g\in G$ ならば $gU\cap U=\emptyset$）で作用しているとき、$\pi_1(Y/G)\cong G$ であることを示せ。ただし商写像 $p\colon Y\to Y/G$ が被覆写像であることは用いてよい。これを用いて $n\ge2$ のとき $\pi_1(\mathbb{RP}^n)$ を計算せよ。
\end{problem}

\begin{solution}
$y_0\in Y$、$x_0=p(y_0)$ とする。$g\in G$ に対し $y_0$ から $gy_0$ への道 $\gamma_g$ をとり、$\Phi(g)=[p\circ\gamma_g]\in\pi_1(Y/G,x_0)$ とおく。$Y$ は単連結なので $\gamma_g$ は端点固定ホモトピー類として一意であり、$\Phi$ は well-defined である。

準同型：$\gamma_g\cdot(g\circ\gamma_h)$ は $y_0$ から $gy_0$、さらに $ghy_0$ に至る道なので $\gamma_{gh}$ として使える。$p\circ g=p$ より $p\circ(g\circ\gamma_h)=p\circ\gamma_h$ なので、$\Phi(gh)=\Phi(g)\Phi(h)$ である。

全射：$x_0$ を基点とするループ $\alpha$ の $y_0$ を始点とする持ち上げ $\tilde\alpha$ の終点は $p^{-1}(x_0)=Gy_0$ に属するので $\tilde\alpha(1)=gy_0$ となる $g$ がある。$\tilde\alpha$ を $\gamma_g$ にとれば $\Phi(g)=[\alpha]$ である。

単射：$\Phi(g)=1$ ならば、前々問より $p\circ\gamma_g$ の $y_0$ を始点とする持ち上げ $\gamma_g$ はループで、$gy_0=y_0$ である。被覆空間作用の条件より $g\ne e$ ならば $gy_0\ne y_0$（$gy_0\in gU\cap U$ となるため）なので、$g=e$ である。以上より $\Phi$ は同型である。

$\mathbb{RP}^n=S^n/\{\pm1\}$ では、$U=\{z\in S^n:\langle z,y\rangle>0\}$ が $(-U)\cap U=\emptyset$ をみたすので $\{\pm1\}\cong\Z/2$ の作用は被覆空間作用であり、$n\ge2$ で $S^n$ は単連結である。よって $\pi_1(\mathbb{RP}^n)\cong\Z/2$ である。
\end{solution}

\begin{problem}[難易度 C]
任意の連続写像 $f\colon\mathbb{RP}^2\to S^1$ は定値写像とホモトピックであることを示せ。また、任意の連続写像 $S^2\to S^1$ についても同じことを示せ。
\end{problem}

\begin{solution}
持ち上げ判定法を用いる：$Y$ が弧状連結かつ局所弧状連結、$p\colon\tilde X\to X$ が被覆写像のとき、$f\colon(Y,y_0)\to(X,x_0)$ が $(\tilde X,\tilde x_0)$ への持ち上げをもつための必要十分条件は $f_*\pi_1(Y,y_0)\subset p_*\pi_1(\tilde X,\tilde x_0)$ である。

$\mathbb{RP}^2$ は多様体なので弧状連結かつ局所弧状連結であり、$\pi_1(\mathbb{RP}^2)\cong\Z/2$ である。$f_*\colon\Z/2\to\pi_1(S^1)\cong\Z$ は準同型であり、$\Z$ には位数 2 の元がないので $f_*$ は自明である。よって $f_*\pi_1(\mathbb{RP}^2)=1\subset p_*\pi_1(\R)$ であり、普遍被覆 $p\colon\R\to S^1$ に関して $f=p\circ\tilde f$ となる連続写像 $\tilde f\colon\mathbb{RP}^2\to\R$ が存在する。

$\R$ は凸なので $H(y,t)=(1-t)\tilde f(y)$ は $\tilde f$ から値 $0$ の定値写像へのホモトピーであり、$p\circ H$ は $f$ から値 $p(0)=1$ の定値写像へのホモトピーである。

$S^2$ の場合は、$S^2$ が単連結かつ局所弧状連結であることから $f_*$ は自明であり、同じ議論が適用できる。
\end{solution}

\begin{problem}[難易度 C]
連結な被覆空間 $p\colon\tilde X\to S^1$ を同型を除いてすべて求めよ。また $n\ge1$ に対し、被覆 $p_n\colon S^1\to S^1$、$p_n(z)=z^n$ の被覆変換群を求め、$p_n$ が正規被覆であることを示せ。
\end{problem}

\begin{solution}
$S^1$ は弧状連結・局所弧状連結・半局所単連結なので、被覆空間の分類定理により、連結な被覆空間の同型類と $\pi_1(S^1)\cong\Z$ の部分群の共役類が $p\mapsto p_*\pi_1(\tilde X)$ で一対一に対応する。$\Z$ はアーベル群なので共役類は部分群そのもので、部分群は $n\Z$（$n\ge0$）で尽くされる。

$n\ge1$ のとき、$p_n$ は被覆写像であり（$S^1$ の真の開弧の逆像は $n$ 個の開弧に同相に分かれる）、$(p_n)_*$ は $n$ 倍写像なので像は $n\Z$ である。$n=0$ には普遍被覆 $\R\to S^1$、$s\mapsto e^{2\pi is}$ が対応する（$\R$ は単連結）。よって連結な被覆は $p_n$（$n\ge1$）と $\R\to S^1$ で尽くされ、これらは互いに同型でない。

被覆変換群：$\zeta^n=1$ のとき回転 $\tau_\zeta(z)=\zeta z$ は $p_n\circ\tau_\zeta=p_n$ をみたす同相なので被覆変換である。逆に被覆変換 $\tau$ に対し $\zeta=\tau(1)\in p_n^{-1}(1)$ とおくと、$\tau$ と $\tau_\zeta$ はともに $p_n$ の持ち上げで点 $1$ で一致するので、持ち上げの一意性（$S^1$ は連結）より $\tau=\tau_\zeta$ である。よって被覆変換群は 1 の $n$ 乗根の群 $\mu_n\cong\Z/n$ である。

これはファイバー $p_n^{-1}(1)=\mu_n$ に推移的に作用するので $p_n$ は正規被覆である（あるいは $\Z$ がアーベル群なので $n\Z$ は正規部分群である）。
\end{solution}

\begin{problem}[難易度 D]
$X=S^1\vee S^1$ の連結な 2 重被覆（葉の数が 2 の被覆）は同型を除いていくつあるか。その数を求め、各被覆を具体的に記述せよ。
\end{problem}

\begin{solution}
$X$ は有限グラフなので被覆空間の分類定理が使え、連結な被覆の同型類は $\pi_1(X)=F_2=\langle a,b\rangle$ の部分群の共役類と対応し、葉の数は部分群の指数に等しい。よって指数 2 の部分群の共役類を数えればよい。

指数 2 の部分群 $H$ は正規である（$g\notin H$ ならば $gH$ と $Hg$ はともに $H$ の補集合）。したがって共役類は一元からなり、$H$ は全射準同型 $F_2\to F_2/H\cong\Z/2$ の核である。逆に全射準同型 $\psi\colon F_2\to\Z/2$ の核は指数 2 であり、$\Z/2$ への全射は核によって決まる（核の元を $0$、それ以外を $1$ に送る）。$F_2$ は自由群なので $\psi$ は $(\psi(a),\psi(b))\in(\Z/2)^2$ で決まり、全射であるのは $(1,0),(0,1),(1,1)$ の 3 通りである。よって連結な 2 重被覆はちょうど 3 個ある。

具体的には、$\tilde X$ は 2 頂点 $v_0, v_1$ をもつグラフで、$v_i$ を基点 $x_0$ の上に、各辺を $a$ または $b$ の上に写すものである。
\begin{itemize}
\item $(1,0)$：$v_0$ と $v_1$ を結ぶ 2 本の $a$ 辺（合わせて円周）と、各頂点での $b$ ループ。
\item $(0,1)$：上で $a$ と $b$ の役割を入れ替えたもの。
\item $(1,1)$：$v_0$ と $v_1$ を結ぶ 2 本の $a$ 辺と 2 本の $b$ 辺。
\end{itemize}
\begin{center}
\begin{tikzpicture}[scale=0.9,thick]
  \tikzset{aedge/.style={red!75!black,mid arrow},bedge/.style={blue!75!black,dashed,mid arrow}}
  \begin{scope}
    \draw[aedge] (0,0) to[bend left=40] (2,0);
    \draw[aedge] (2,0) to[bend left=40] (0,0);
    \draw[bedge] (0,0) .. controls (-1.3,0.9) and (-1.3,-0.9) .. (0,0);
    \draw[bedge] (2,0) .. controls (3.3,-0.9) and (3.3,0.9) .. (2,0);
    \node[vtx] at (0,0) {}; \node[vtx] at (2,0) {};
    \node[below] at (0,-0.1) {\footnotesize $v_0$}; \node[below] at (2,-0.1) {\footnotesize $v_1$};
    \node at (1,-1.2) {$(\psi(a),\psi(b))=(1,0)$};
  \end{scope}
  \begin{scope}[xshift=5cm]
    \draw[bedge] (0,0) to[bend left=40] (2,0);
    \draw[bedge] (2,0) to[bend left=40] (0,0);
    \draw[aedge] (0,0) .. controls (-1.3,0.9) and (-1.3,-0.9) .. (0,0);
    \draw[aedge] (2,0) .. controls (3.3,-0.9) and (3.3,0.9) .. (2,0);
    \node[vtx] at (0,0) {}; \node[vtx] at (2,0) {};
    \node[below] at (0,-0.1) {\footnotesize $v_0$}; \node[below] at (2,-0.1) {\footnotesize $v_1$};
    \node at (1,-1.2) {$(0,1)$};
  \end{scope}
  \begin{scope}[xshift=10cm]
    \draw[aedge] (0,0) to[bend left=30] (2,0);
    \draw[aedge] (2,0) to[bend left=30] (0,0);
    \draw[bedge] (0,0) to[bend left=80] (2,0);
    \draw[bedge] (2,0) to[bend left=80] (0,0);
    \node[vtx] at (0,0) {}; \node[vtx] at (2,0) {};
    \node[left] at (-0.05,0) {\footnotesize $v_0$}; \node[right] at (2.05,0) {\footnotesize $v_1$};
    \node at (1,-1.2) {$(1,1)$};
  \end{scope}
\end{tikzpicture}
\zucap{連結な2重被覆の3つ。実線（赤）は $a$ の上に、破線（青）は $b$ の上に矢印の向きを保って写る。各頂点から出る $a$ 辺・入る $a$ 辺、出る $b$ 辺・入る $b$ 辺がちょうど1本ずつあることが、被覆であるための条件である。}
\end{center}
各場合、$\psi$ の値が $1$ の文字に沿う辺だけが頂点を入れ替えるので、ループ $\gamma$ の $v_0$ を始点とする持ち上げが閉じるのは $\psi[\gamma]=0$ のときである。持ち上げが閉じることと $[\gamma]\in p_*\pi_1(\tilde X,v_0)$ は同値なので、$p_*\pi_1(\tilde X,v_0)=\ker\psi$ が確かめられる。
\end{solution}

\begin{problem}[難易度 C]
$X=S^1\vee S^1$ の2つの円周を $a,b$ とし、$\pi_1(X,x_0)=F_2=\langle a,b\rangle$ とする。下図のグラフ $\tilde X$（頂点 $v_0,v_1,v_2$）から $X$ への写像 $p$ を、頂点をすべて $x_0$ に、実線（赤）の辺を $a$ に、破線（青）の辺を $b$ に、矢印の向きを保って写すものとして定める。
\begin{center}
\begin{tikzpicture}[scale=1.0,thick]
  \tikzset{aedge/.style={red!75!black,mid arrow},bedge/.style={blue!75!black,dashed,mid arrow}}
  \coordinate (V0) at (0,0); \coordinate (V1) at (3,0); \coordinate (V2) at (1.5,2.3);
  \draw[aedge] (V0) to[bend right=15] (V1);
  \draw[aedge] (V1) to[bend right=15] (V2);
  \draw[aedge] (V2) to[bend right=15] (V0);
  \draw[bedge] (V1) to[bend right=55] (V2);
  \draw[bedge] (V2) to[bend right=25] (V1);
  \draw[bedge] (V0) .. controls (-1.5,0.8) and (-1.5,-0.8) .. (V0);
  \node[vtx] at (V0) {}; \node[vtx] at (V1) {}; \node[vtx] at (V2) {};
  \node[below] at (0,-0.1) {$v_0$}; \node[below right] at (3,-0.05) {$v_1$}; \node[above] at (1.5,2.4) {$v_2$};
\end{tikzpicture}
\zucap{$a$ 辺：$v_0\to v_1\to v_2\to v_0$。$b$ 辺：$v_0$ でのループ、$v_1\to v_2$、$v_2\to v_1$。}
\end{center}
\begin{enumerate}
\item $p$ が3重被覆であることを確かめよ。
\item $H=p_*\pi_1(\tilde X,v_0)$ の自由な生成系を求めよ。
\item $p$ は正規被覆か。また被覆変換群を求めよ。
\end{enumerate}
\end{problem}

\begin{solution}
(1) 各頂点で、出ていく $a$ 辺・入ってくる $a$ 辺・出ていく $b$ 辺・入ってくる $b$ 辺がちょうど1本ずつある（$v_0$ の $b$ ループは出る辺と入る辺を1本ずつ与える）。したがって $x_0$ の小さな近傍（$x_0$ から4方向に出る短い線分の和）の逆像は、各頂点のまわりの同じ形の近傍3つの非交和に同相に写り、辺の内部の点の近傍も同様である。よって $p$ は3重被覆である。

(2) 道の持ち上げは「各頂点で、読んでいる文字の辺を矢印の向き（逆元なら逆向き）にたどる」ことで得られる。グラフの基本群の生成系は、極大な木 $T$ を選び、$T$ に含まれない各辺 $e$（$u$ から $u'$ へ）に対して「$T$ 内で $v_0$ から $u$ へ行き、$e$ を通り、$T$ 内で $u'$ から $v_0$ へ戻る」ループをとれば得られる（$\tilde X/T$ は円周の楔和であり、第1章の問題と同じ）。$T$ として $a$ 辺 $v_0\to v_1$、$v_1\to v_2$ をとると、$T$ 外の辺は4本で
\begin{center}
\begin{tabular}{lll}
$T$ 外の辺 & 対応するループ & $H$ の元\\ \hline
$a\colon v_2\to v_0$ & $v_0\xrightarrow{a}v_1\xrightarrow{a}v_2\xrightarrow{a}v_0$ & $a^3$\\
$b\colon v_0\to v_0$ & $v_0\xrightarrow{b}v_0$ & $b$\\
$b\colon v_1\to v_2$ & $v_0\xrightarrow{a}v_1\xrightarrow{b}v_2\xrightarrow{a^{-1}}v_1\xrightarrow{a^{-1}}v_0$ & $aba^{-2}$\\
$b\colon v_2\to v_1$ & $v_0\xrightarrow{a}v_1\xrightarrow{a}v_2\xrightarrow{b}v_1\xrightarrow{a^{-1}}v_0$ & $a^2ba^{-1}$
\end{tabular}
\end{center}
となる。$p_*$ は単射なので
\[
H=\langle a^3,\ b,\ aba^{-2},\ a^2ba^{-1}\rangle
\]
で、これらは $H$ の自由な生成系である（階数 $E-V+1=6-3+1=4$）。シュライアーの公式 $\operatorname{rank}H=1+[F_2:H](\operatorname{rank}F_2-1)=1+3\cdot1=4$ とも一致する。

(3) 正規被覆ではない。実際 $b\in H$ だが、$aba^{-1}$ の $v_0$ からの持ち上げは $v_0\xrightarrow{a}v_1\xrightarrow{b}v_2\xrightarrow{a^{-1}}v_1$ で閉じないので $aba^{-1}\notin H$、すなわち $aH a^{-1}\neq H$ である。図形的には、$v_0$ ではループ $b$ の持ち上げが閉じるのに、$v_1,v_2$ では閉じない、という「頂点ごとの見え方の違い」が非正規性を表している。

被覆変換 $\tau$ は $v_0$ の像で決まり、$b$ ループを $b$ ループに写さなければならない。$b$ ループをもつ頂点は $v_0$ だけなので $\tau(v_0)=v_0$、よって $\tau=\id$ である。被覆変換群は自明群であり、これは $N(H)/H$（$N(H)$ は正規化群）とも整合する：$H$ は指数3で正規でないので $N(H)=H$ となる（$H\subset N(H)\subset F_2$ で指数は3の約数）。
\end{solution}

%================================================
\chapter{単体的ホモロジー}
% 25問
%================================================

\begin{summary}
\begin{itemize}
  \item \textbf{単体と向き}：$n$ 単体 $[v_0,\dots,v_n]$ の向きは頂点の順序で決まり、偶置換で同じ向き、奇置換で逆向き（$-1$ 倍）とみなす。
  \item \textbf{鎖群と境界}：$C_n(K)$ は向き付き $n$ 単体で生成される自由アーベル群。境界作用素は
  \[
  \partial[v_0,\dots,v_n]=\sum_{i=0}^n(-1)^i[v_0,\dots,\widehat{v_i},\dots,v_n]
  \]
  で、$\partial\circ\partial=0$ をみたす。
  \item \textbf{ホモロジー}：$H_n(K)=\ker\partial_n/\operatorname{im}\partial_{n+1}$（サイクル／境界）。$H_0(K)\cong\Z^{(\text{連結成分の数})}$。添加写像 $\varepsilon\colon C_0\to\Z$ を付けたものが簡約ホモロジー $\tilde H_n$。
  \item \textbf{計算の道具}：境界行列をスミス標準形 $\operatorname{diag}(d_1,\dots,d_r,0,\dots)$ に変形すると、$\operatorname{rank}\operatorname{im}\partial_n=r$ であり、$H_{n-1}$ のねじれ部分は $\bigoplus\Z/d_i$ になる。有限生成アーベル群の基本定理 $A\cong\Z^b\oplus\bigoplus\Z/d_i$ も用いる。
  \item \textbf{単体写像}：単体写像は鎖写像 $f_\sharp$（退化した単体は $0$ に送る）を与え、準同型 $f_*$ を誘導する。鎖ホモトピー $\partial D+D\partial=f_\sharp-g_\sharp$ があれば $f_*=g_*$。
  \item \textbf{オイラー標数}：$c_n$ を $n$ 単体の個数とすると
  \[
  \chi(K)=\sum_n(-1)^n c_n=\sum_n(-1)^n\operatorname{rank}H_n(K).
  \]
  \item \textbf{相対ホモロジー}：部分複体 $L\subset K$ に対し $C_n(K,L)=C_n(K)/C_n(L)$ から $H_n(K,L)$ を定める。
\end{itemize}
\end{summary}

\begin{problem}[難易度 A]
$n$ 単体 $\Delta^n=[v_0,v_1,\dots,v_n]$ の $k$ 次元の面（$0\le k\le n$）の個数は $\binom{n+1}{k+1}$ であることを示せ。また、$\Delta^n$ の面（$\Delta^n$ 自身を含む）の総数は $2^{n+1}-1$ であることを示せ。
\end{problem}

\begin{solution}
$\Delta^n$ の面とは、頂点集合 $\{v_0,\dots,v_n\}$ の空でない部分集合が張る単体のことであり、面と空でない部分集合は一対一に対応する。$k$ 次元の面は $k+1$ 個の頂点で張られるから、その個数は $n+1$ 個の頂点から $k+1$ 個を選ぶ方法の数 $\binom{n+1}{k+1}$ に等しい。

面の総数は、二項定理 $\sum_{j=0}^{n+1}\binom{n+1}{j}=(1+1)^{n+1}=2^{n+1}$ より
\[
\sum_{k=0}^{n}\binom{n+1}{k+1}=\sum_{j=1}^{n+1}\binom{n+1}{j}=2^{n+1}-\binom{n+1}{0}=2^{n+1}-1
\]
である。これは $\{v_0,\dots,v_n\}$ の部分集合 $2^{n+1}$ 個から空集合を除いた個数と言っても同じである。例えば $\Delta^2$ の面は頂点 3 個、辺 3 個、2 単体 1 個の計 $7=2^3-1$ 個、$\Delta^3$ の面は $4+6+4+1=15=2^4-1$ 個である。
\end{solution}

\begin{problem}[難易度 A]
向き付けられた $n$ 単体 $[v_0,\dots,v_n]$（$n\ge 1$）の境界を
\[
\partial[v_0,\dots,v_n]=\sum_{i=0}^{n}(-1)^i[v_0,\dots,\widehat{v_i},\dots,v_n]
\]
で定める。隣接する 2 頂点 $v_j,v_{j+1}$ を入れ替えた単体について
\[
\partial[v_0,\dots,v_{j+1},v_j,\dots,v_n]=-\partial[v_0,\dots,v_j,v_{j+1},\dots,v_n]
\]
が成り立つことを示せ。これにより、$\partial$ が向き付き単体の鎖群（$[\dots]$ の頂点を奇置換すると $-1$ 倍になるという関係を課したもの）の上で矛盾なく定義されることを説明せよ。
\end{problem}

\begin{solution}
$\sigma=[v_0,\dots,v_n]$、$\sigma'=[w_0,\dots,w_n]$ を $v_j$ と $v_{j+1}$ を入れ替えたもの（$w_j=v_{j+1}$、$w_{j+1}=v_j$、他は $w_i=v_i$）とする。$\partial\sigma'=\sum_i(-1)^i[w_0,\dots,\widehat{w_i},\dots,w_n]$ の各項を $\partial\sigma$ の項と比べる。

(i) $i\ne j,j+1$ のとき：$\sigma'$ の第 $i$ 面は $\sigma$ の第 $i$ 面の中で $v_j,v_{j+1}$ を入れ替えたものなので、$(n-1)$ 単体として $-1$ 倍である。符号 $(-1)^i$ は共通なので、この項は $\partial\sigma$ の対応する項の $-1$ 倍である。

(ii) $i=j$ と $i=j+1$ のとき：$\sigma'$ から $w_j=v_{j+1}$ を除くと $[v_0,\dots,v_j,\widehat{v_{j+1}},\dots,v_n]$ が符号 $(-1)^j$ で現れる。$\partial\sigma$ ではこれは $i=j+1$ の項として符号 $(-1)^{j+1}$ で現れる。同様に $\sigma'$ から $w_{j+1}=v_j$ を除いた項は符号 $(-1)^{j+1}$、$\partial\sigma$ では符号 $(-1)^j$ である。よってこの 2 項の和も $\partial\sigma$ の対応する 2 項の和の $-1$ 倍である。

以上で $\partial\sigma'=-\partial\sigma$ が示された。任意の置換は隣接互換の積であり、偶置換・奇置換はその個数の偶奇で決まるので、頂点の置換 $\pi$ に対して $\partial[v_{\pi(0)},\dots,v_{\pi(n)}]=\operatorname{sgn}(\pi)\,\partial[v_0,\dots,v_n]$ となる。したがって $\partial$ は関係式
\[
[v_{\pi(0)},\dots,v_{\pi(n)}]=\operatorname{sgn}(\pi)[v_0,\dots,v_n]
\]
と両立し、向き付き単体の鎖群上の準同型として矛盾なく定まる。
\end{solution}

\begin{problem}[難易度 A]
3 単体 $\sigma=[v_0,v_1,v_2,v_3]$ について $\partial\sigma$ と $\partial(\partial\sigma)$ を具体的に計算し、$\partial\circ\partial(\sigma)=0$ を確かめよ。
\end{problem}

\begin{solution}
定義より
\[
\partial\sigma=[v_1,v_2,v_3]-[v_0,v_2,v_3]+[v_0,v_1,v_3]-[v_0,v_1,v_2]
\]
である。各項の境界は
\begin{align*}
\partial[v_1,v_2,v_3]&=[v_2,v_3]-[v_1,v_3]+[v_1,v_2],\\
\partial[v_0,v_2,v_3]&=[v_2,v_3]-[v_0,v_3]+[v_0,v_2],\\
\partial[v_0,v_1,v_3]&=[v_1,v_3]-[v_0,v_3]+[v_0,v_1],\\
\partial[v_0,v_1,v_2]&=[v_1,v_2]-[v_0,v_2]+[v_0,v_1]
\end{align*}
である。符号 $+,-,+,-$ を付けて足すと、各辺の係数は
\begin{align*}
[v_2,v_3]&:\ 1-1=0, & [v_1,v_3]&:\ -1+1=0, & [v_1,v_2]&:\ 1-1=0,\\
[v_0,v_3]&:\ 1-1=0, & [v_0,v_2]&:\ -1+1=0, & [v_0,v_1]&:\ 1-1=0
\end{align*}
となる。よって $\partial(\partial\sigma)=0$ である。各辺 $[v_a,v_b]$ はちょうど 2 つの 2 次元面に含まれ、その 2 つの寄与が逆符号で打ち消し合っていることに注意する。
\end{solution}

\begin{problem}[難易度 B]
任意の単体複体 $K$ と $n\ge 2$ について、$\partial_{n-1}\circ\partial_n=0\colon C_n(K)\to C_{n-2}(K)$ を示せ。
\end{problem}

\begin{solution}
$C_n(K)$ は向き付き $n$ 単体で生成されるので、$\sigma=[v_0,\dots,v_n]$ について示せばよい。
\[
\partial\partial\sigma=\sum_{i=0}^{n}(-1)^i\,\partial[v_0,\dots,\widehat{v_i},\dots,v_n].
\]
$[v_0,\dots,\widehat{v_i},\dots,v_n]$ において、$j<i$ の頂点 $v_j$ は第 $j$ 番目、$j>i$ の頂点 $v_j$ は第 $j-1$ 番目にある。したがって
\[
\partial\partial\sigma
=\sum_{j<i}(-1)^{i}(-1)^{j}[\dots,\widehat{v_j},\dots,\widehat{v_i},\dots]
+\sum_{j>i}(-1)^{i}(-1)^{j-1}[\dots,\widehat{v_i},\dots,\widehat{v_j},\dots]
\]
となる。$v_a,v_b$（$a<b$）を除いた $(n-2)$ 単体 $[\dots,\widehat{v_a},\dots,\widehat{v_b},\dots]$ は、第 1 の和では $(j,i)=(a,b)$ として係数 $(-1)^{a+b}$ で、第 2 の和では $(i,j)=(a,b)$ として係数 $(-1)^{a+b-1}$ でちょうど 1 回ずつ現れる。両者は打ち消し合うので、$\partial\partial\sigma=0$ である。線形に拡張して $\partial_{n-1}\circ\partial_n=0$ を得る。特に $\operatorname{im}\partial_{n+1}\subset\ker\partial_n$ であり、ホモロジー群 $H_n(K)=\ker\partial_n/\operatorname{im}\partial_{n+1}$ が定義できる。
\end{solution}

\begin{problem}[難易度 A]
(1) 1 点 $\{v\}$ からなる単体複体 $P$ のホモロジー群を計算せよ。\\
(2) 1 単体 $\Delta^1=[v_0,v_1]$ とその面全体からなる単体複体のホモロジー群を、境界作用素の行列を書いて計算せよ。
\end{problem}

\begin{solution}
(1) $C_0(P)=\Z v$、$C_n(P)=0$（$n\ne 0$）であり、$\partial_0=0$、$\partial_1=0$ である。よって $H_0(P)=\ker\partial_0/\operatorname{im}\partial_1=\Z v\cong\Z$、$H_n(P)=0$（$n\ne 0$）である。

(2) 鎖群は $C_0=\Z v_0\oplus\Z v_1$、$C_1=\Z e$（$e=[v_0,v_1]$）、$C_n=0$（$n\ge 2$）である。$\partial_1 e=v_1-v_0$ なので、基底 $e$ と $v_0,v_1$ に関する $\partial_1$ の行列は
\[
\partial_1=\begin{pmatrix}-1\\ 1\end{pmatrix}
\]
である。$\partial_1(me)=m(v_1-v_0)=0$ なら $m=0$ なので $\ker\partial_1=0$、よって $H_1=0$。また $\ker\partial_0=C_0$、$\operatorname{im}\partial_1=\Z(v_1-v_0)$ である。準同型 $\varepsilon\colon C_0\to\Z$、$\varepsilon(av_0+bv_1)=a+b$ は全射で、$\ker\varepsilon=\{a v_0+bv_1: a+b=0\}=\Z(v_1-v_0)=\operatorname{im}\partial_1$ である。よって
\[
H_0=C_0/\operatorname{im}\partial_1\cong\Z,\qquad H_1=0,\qquad H_n=0\ (n\ge 2)
\]
であり、$H_0$ は $[v_0]=[v_1]$ で生成される。
\end{solution}

\begin{problem}[難易度 B]
$K$ を $d$ 次元の有限単体複体とする。$n>d$ に対して $H_n(K)=0$ であること、および $H_d(K)$ が自由アーベル群であることを示せ。
\end{problem}

\begin{solution}
$n>d$ のとき $K$ には $n$ 単体が存在しないので $C_n(K)=0$ であり、$H_n(K)=\ker\partial_n/\operatorname{im}\partial_{n+1}$ は $0$ の部分商なので $0$ である。

$n=d$ のときは $C_{d+1}(K)=0$ より $\operatorname{im}\partial_{d+1}=0$ なので
\[
H_d(K)=\ker\partial_d\subset C_d(K)
\]
である。$C_d(K)$ は $d$ 単体（に向きを 1 つずつ選んだもの）を基底とする有限生成自由アーベル群であり、有限生成自由アーベル群の部分群は自由アーベル群である（単項イデアル整域上の自由加群の部分加群は自由）。よって $H_d(K)$ は自由アーベル群である。

一方 $d-1$ 次以下のホモロジー群はねじれを持ちうる。例えば実射影平面の単体分割（2 次元）では $H_1\cong\Z/2$ となる（本章の後の問題を参照）。
\end{solution}

\begin{problem}[難易度 A]
2 単体 $[v_0,v_1,v_2]$ の境界（3 頂点と 3 辺からなる 1 次元単体複体）$L=\partial\Delta^2$ は円周 $S^1$ の単体分割である。鎖群の基底と境界作用素の行列を書き、$H_n(L)$ をすべて計算せよ。
\end{problem}

\begin{solution}
$C_0=\Z v_0\oplus\Z v_1\oplus\Z v_2$、$C_1=\Z e_{01}\oplus\Z e_{02}\oplus\Z e_{12}$（$e_{ij}=[v_i,v_j]$）、$C_n=0$（$n\ge2$）とする。$\partial e_{ij}=v_j-v_i$ より、列を $e_{01},e_{02},e_{12}$、行を $v_0,v_1,v_2$ として
\[
\partial_1=\begin{pmatrix}-1&-1&0\\ 1&0&-1\\ 0&1&1\end{pmatrix}.
\]
$c=ae_{01}+be_{02}+ce_{12}$ に対し $\partial c=0$ は $-a-b=0$、$a-c=0$、$b+c=0$ と同値で、解は $b=-a$、$c=a$ である。よって
\[
H_1(L)=\ker\partial_1=\Z z,\qquad z=e_{01}+e_{12}-e_{02}=[v_0,v_1]+[v_1,v_2]+[v_2,v_0]
\]
であり $H_1(L)\cong\Z$（$C_2=0$ なので割る必要はない）。

次に $\operatorname{im}\partial_1$ は $v_1-v_0,\ v_2-v_0,\ v_2-v_1$ で生成され、これは $\varepsilon(\sum a_iv_i)=\sum a_i$ の核 $\{\sum a_iv_i:\sum a_i=0\}$ に一致する（$\sum a_iv_i=a_1(v_1-v_0)+a_2(v_2-v_0)$ と書ける）。$\varepsilon$ は全射なので $H_0(L)=C_0/\ker\varepsilon\cong\Z$。まとめると $H_0\cong\Z$、$H_1\cong\Z$、$H_n=0$（$n\ge 2$）である。
\end{solution}

\begin{problem}[難易度 A]
2 単体 $\Delta^2=[v_0,v_1,v_2]$ とそのすべての面からなる単体複体 $K$ のホモロジー群を計算せよ。
\end{problem}

\begin{solution}
$K$ は前問の $L=\partial\Delta^2$ に 2 単体 $t=[v_0,v_1,v_2]$ を加えたものなので、$C_0,C_1,\partial_1$ は前問と同じで、さらに $C_2=\Z t$ である。
\[
\partial_2 t=[v_1,v_2]-[v_0,v_2]+[v_0,v_1]=e_{12}-e_{02}+e_{01},
\qquad
\partial_2=\begin{pmatrix}1\\-1\\1\end{pmatrix}
\]
（行は $e_{01},e_{02},e_{12}$）。$\partial_2(mt)=0$ なら $m=0$ なので $\ker\partial_2=0$、よって $H_2(K)=0$。

前問より $\ker\partial_1=\Z z$、$z=e_{01}+e_{12}-e_{02}$ であり、$\partial_2 t=z$ なので $\operatorname{im}\partial_2=\Z z=\ker\partial_1$。よって $H_1(K)=0$。$H_0$ は $C_0,C_1,\partial_1$ だけで決まるので前問と同じく $H_0(K)\cong\Z$。まとめると
\[
H_0(K)\cong\Z,\qquad H_n(K)=0\ (n\ge 1).
\]
円周 $L$ の 1 次元のサイクル $z$ が、2 単体を加えると境界になって消えることが分かる。
\end{solution}

\begin{problem}[難易度 B]
単体複体 $K$ の 2 頂点 $u,v$ が、辺の列 $[u=w_0,w_1]$, $[w_1,w_2]$, $\dots$, $[w_{m-1},w_m=v]$ で結べるとき $u\sim v$ と書く（$u\sim u$ とする）。これは同値関係であり、その同値類の個数 $r$ は $\lvert K\rvert$ の連結成分の個数に一致する。$H_0(K)\cong\Z^r$ であることを示せ。
\end{problem}

\begin{solution}
同値類を $A_1,\dots,A_r$ とし、準同型 $\varphi\colon C_0(K)\to\Z^r$ を、$v\in A_j$ に対して $\varphi(v)=\mathbf{e}_j$（第 $j$ 標準基底）と定める。$\varphi$ は明らかに全射である。$H_0(K)=C_0(K)/\operatorname{im}\partial_1$ なので、$\ker\varphi=\operatorname{im}\partial_1$ を示せばよい。

($\supset$) 辺 $[u,v]$ の両端は同じ類に属するので $\varphi(\partial[u,v])=\varphi(v)-\varphi(u)=0$。

($\subset$) 各 $A_j$ に代表元 $p_j$ を選ぶ。$u\in A_j$ なら辺の列で $p_j$ と結べるので
\[
u-p_j=\partial\bigl([p_j,w_1]+[w_1,w_2]+\dots+[w_{m-1},u]\bigr)\in\operatorname{im}\partial_1 .
\]
$x=\sum_u a_u u\in\ker\varphi$ とすると、各 $j$ について $\sum_{u\in A_j}a_u=0$ である。よって
\[
x=\sum_{j=1}^r\sum_{u\in A_j}a_u u=\sum_{j=1}^r\sum_{u\in A_j}a_u(u-p_j)\in\operatorname{im}\partial_1 .
\]
以上より $H_0(K)\cong\Z^r$ であり、各連結成分の 1 頂点のホモロジー類が基底を与える。
\end{solution}

\begin{problem}[難易度 B]
空でない単体複体 $K$ に対し、添加写像 $\varepsilon\colon C_0(K)\to\Z$、$\varepsilon(\sum a_v v)=\sum a_v$ を用いて、鎖複体
\[
\cdots\to C_1(K)\xrightarrow{\partial_1}C_0(K)\xrightarrow{\varepsilon}\Z\to 0
\]
のホモロジー（$C_n$ の位置、$n\ge 0$）を簡約ホモロジー $\tilde H_n(K)$ と定める。これが鎖複体であることを確かめ、$H_0(K)\cong\tilde H_0(K)\oplus\Z$、および $n\ge 1$ で $\tilde H_n(K)=H_n(K)$ であることを示せ。
\end{problem}

\begin{solution}
$\varepsilon(\partial_1[u,v])=\varepsilon(v-u)=1-1=0$ より $\varepsilon\circ\partial_1=0$ であり、他の合成は通常の鎖複体と同じく $0$ なので、これは鎖複体である。

$n\ge 1$ では $C_n$ の前後の写像が通常の鎖複体と同じなので $\tilde H_n(K)=H_n(K)$ である。$n=0$ では
\[
\tilde H_0(K)=\ker\varepsilon/\operatorname{im}\partial_1\ \subset\ C_0(K)/\operatorname{im}\partial_1=H_0(K).
\]
$\varepsilon\circ\partial_1=0$ より $\varepsilon$ は準同型 $\bar\varepsilon\colon H_0(K)\to\Z$ を誘導し、$K$ が頂点 $v$ を持つので $\bar\varepsilon([v])=1$ より全射である。その核は $\ker\varepsilon/\operatorname{im}\partial_1=\tilde H_0(K)$ なので、完全列
\[
0\to\tilde H_0(K)\to H_0(K)\xrightarrow{\bar\varepsilon}\Z\to 0
\]
を得る。$\Z$ は自由なので $1\mapsto[v]$ で分裂し、$H_0(K)\cong\tilde H_0(K)\oplus\Z$ である。特に前問より $\tilde H_0(K)\cong\Z^{r-1}$（$r$ は連結成分の個数）であり、$K$ が連結であることと $\tilde H_0(K)=0$ は同値である。
\end{solution}

\begin{problem}[難易度 C]
$K$ を単体複体、$w$ を $K$ の頂点でない点とし、錐 $CK=w*K$ を、$K$ の単体、頂点 $w$、および $K$ の各単体 $[v_0,\dots,v_k]$ に対する $[w,v_0,\dots,v_k]$ 全体からなる単体複体とする。すべての $n$ で $\tilde H_n(CK)=0$ であることを示せ。これを用いて、$n$ 単体 $\Delta^n$（とその面全体）について $\tilde H_*(\Delta^n)=0$ を示せ。
\end{problem}

\begin{solution}
準同型 $c\colon C_k(CK)\to C_{k+1}(CK)$ を、$w$ を含まない単体 $\sigma=[v_0,\dots,v_k]$ には $c\sigma=[w,v_0,\dots,v_k]$、$w$ を含む単体には $c\sigma=0$ と定める（頂点の置換と両立するので矛盾なく定まる）。

(a) $k\ge1$、$\sigma$ が $w$ を含まないとき：
\[
\partial c\sigma=[v_0,\dots,v_k]-\sum_{i}(-1)^{i}[w,v_0,\dots,\widehat{v_i},\dots,v_k]=\sigma-c\,\partial\sigma .
\]
(b) $k\ge1$、$\sigma=[w,v_1,\dots,v_k]$ のとき：$c\sigma=0$ であり、$\partial\sigma=[v_1,\dots,v_k]+(w$ を含む項$)$ なので $c\,\partial\sigma=[w,v_1,\dots,v_k]=\sigma$。

よって $k\ge 1$ では $\partial c+c\partial=\id$ である。$k=0$ では $v\ne w$ に対し $\partial c(v)=v-w$、$\partial c(w)=0$ なので、$x\in C_0$ に対し $\partial c(x)=x-\varepsilon(x)w$ である。

したがって $k\ge 1$ のサイクル $z$ は $z=\partial(cz)+c(\partial z)=\partial(cz)$ と境界になり $\tilde H_k(CK)=H_k(CK)=0$。また $x\in\ker\varepsilon$ なら $x=\partial(cx)$ なので $\tilde H_0(CK)=0$。

$n\ge1$ のとき $\Delta^n=[v_0,\dots,v_n]$ は、面 $[v_1,\dots,v_n]$ とその面全体のなす複体 $K$ の $w=v_0$ による錐 $CK$ に等しい。$n=0$ なら 1 点で、$\tilde H_*=0$ は直接分かる。よって $\tilde H_*(\Delta^n)=0$ である。
\end{solution}

\begin{problem}[難易度 B]
3 単体 $[v_0,v_1,v_2,v_3]$ の境界 $L=\partial\Delta^3$（4 頂点、6 辺、4 つの 2 単体からなる 2 次元単体複体）は 2 次元球面 $S^2$ の単体分割である。境界作用素の行列を書き、$H_n(L)$ を直接計算せよ。
\end{problem}

\begin{solution}
$e_{ij}=[v_i,v_j]$（$i<j$）、$t_0=[v_1,v_2,v_3]$、$t_1=[v_0,v_2,v_3]$、$t_2=[v_0,v_1,v_3]$、$t_3=[v_0,v_1,v_2]$ とする。行を $e_{01},e_{02},e_{03},e_{12},e_{13},e_{23}$、列を $t_0,\dots,t_3$（$\partial_2$）、行を $v_0,\dots,v_3$、列を $e_{01},\dots,e_{23}$（$\partial_1$）として
\[
\partial_2=\begin{pmatrix}0&0&1&1\\0&1&0&-1\\0&-1&-1&0\\1&0&0&1\\-1&0&1&0\\1&1&0&0\end{pmatrix},\quad
\partial_1=\begin{pmatrix}-1&-1&-1&0&0&0\\1&0&0&-1&-1&0\\0&1&0&1&0&-1\\0&0&1&0&1&1\end{pmatrix}.
\]
$H_2$：$\partial_2(at_0+bt_1+ct_2+dt_3)=0$ は $c+d=0$、$b-d=0$、$-b-c=0$、$a+d=0$、$-a+c=0$、$a+b=0$ と同値で、解は $(a,b,c,d)=a(1,-1,1,-1)$。$C_3=0$ より $H_2(L)=\Z(t_0-t_1+t_2-t_3)\cong\Z$。

$H_1$：$z\in\ker\partial_1$ の $e_{12},e_{13},e_{23}$ の係数を $\alpha,\beta,\gamma$ とする。$\partial t_3,\partial t_2,\partial t_1$ はそれぞれ $e_{12},e_{13},e_{23}$ のみを係数 $1$ で含む（残り 2 本は含まない）ので、$z'=z-\alpha\,\partial t_3-\beta\,\partial t_2-\gamma\,\partial t_1$ は $z'=xe_{01}+ye_{02}+we_{03}$ の形のサイクルである。$\partial z'$ の $v_1,v_2,v_3$ の係数は $x,y,w$ なので $x=y=w=0$、すなわち $z=\partial(\gamma t_1+\beta t_2+\alpha t_3)$。よって $\ker\partial_1=\operatorname{im}\partial_2$ で $H_1(L)=0$。

$H_0$：$L$ は連結なので $H_0(L)\cong\Z$。以上より $H_0\cong\Z$、$H_1=0$、$H_2\cong\Z$、$H_n=0$（$n\ge3$）。
\end{solution}

\begin{problem}[難易度 B]
$n\ge 1$ とし、$(n+1)$ 単体の境界 $L=\partial\Delta^{n+1}$（$\Delta^{n+1}$ の $n$ 次元以下の面全体）を考える。これは $S^n$ の単体分割である。錐のホモロジーに関する問題の結果を用いて
\[
H_0(L)\cong\Z,\qquad H_n(L)\cong\Z,\qquad H_k(L)=0\ (k\ne 0,n)
\]
を示せ。
\end{problem}

\begin{solution}
$K=\Delta^{n+1}$（すべての面を含む）とする。$k\le n$ では $C_k(L)=C_k(K)$ で境界作用素も等しく、$C_{n+1}(L)=0$、$C_{n+1}(K)=\Z\sigma$（$\sigma=[v_0,\dots,v_{n+1}]$）である。また錐の問題より $\tilde H_*(K)=0$ である。

(i) $k\le n-1$ のとき：$H_k$ は $C_{k+1}\to C_k\to C_{k-1}$ だけで決まり、$k+1\le n$ なので $H_k(L)=H_k(K)$。よって $H_0(L)\cong\Z$、$1\le k\le n-1$ で $H_k(L)=0$。

(ii) $k=n$ のとき：$C_{n+1}(L)=0$ より
\[
H_n(L)=\ker\partial_n=Z_n(K)=\operatorname{im}\bigl(\partial_{n+1}\colon C_{n+1}(K)\to C_n(K)\bigr)=\Z\,\partial\sigma .
\]
ここで $n\ge 1$ より $H_n(K)=0$ を使った。$\partial\sigma$ は $n$ 単体の基底で係数 $\pm1$ を持つ $0$ でない元であり、$C_n$ はねじれを持たないので $\Z\,\partial\sigma\cong\Z$。

(iii) $k>n$ では $C_k(L)=0$ より $H_k(L)=0$。

以上で主張が示された。$H_n(L)$ の生成元は $\partial\sigma=\sum_i(-1)^i[v_0,\dots,\widehat{v_i},\dots,v_{n+1}]$ であり、$n=2$ の場合は前問の $t_0-t_1+t_2-t_3$ と一致する。$n=1$ の場合は (i) の範囲が $k=0$ のみとなり、円周の計算と一致する。
\end{solution}

\begin{problem}[難易度 B]
$G$ を連結な 1 次元有限単体複体（グラフ）とし、頂点数を $V$、辺数を $E$ とする。$H_1(G)$ は階数 $E-V+1$ の自由アーベル群であることを示せ。また、$G$ が木（$H_1(G)=0$）であることと $E=V-1$ は同値であることを示せ。
\end{problem}

\begin{solution}
$C_2(G)=0$ より $H_1(G)=\ker\partial_1\subset C_1(G)\cong\Z^E$ であり、自由アーベル群の部分群なので自由である。

階数を求める。$\partial_1\colon C_1\to C_0$ について、準同型定理より $C_1/\ker\partial_1\cong\operatorname{im}\partial_1$ なので
\[
\operatorname{rank}\ker\partial_1=E-\operatorname{rank}\operatorname{im}\partial_1 .
\]
$G$ は連結なので $H_0(G)=C_0/\operatorname{im}\partial_1\cong\Z$ であり、$\operatorname{rank}\operatorname{im}\partial_1=V-1$。よって $\operatorname{rank}H_1(G)=E-(V-1)=E-V+1$。

後半：$H_1(G)$ は自由なので、$H_1(G)=0$ と $\operatorname{rank}H_1(G)=0$ は同値であり、これは $E=V-1$ と同値である。（ここでの「木」は $H_1(G)=0$ で定めたが、これはサイクルを含まない連結グラフという通常の定義と一致する：単純閉路 $w_0,w_1,\dots,w_m=w_0$ があれば $\sum[w_{i},w_{i+1}]$ は $0$ でないサイクルであり、逆に $0$ でないサイクルの台には各頂点の次数が $2$ 以上の部分グラフが含まれ、閉路を持つ。）
\end{solution}

\begin{problem}[難易度 B]
2 つの円周を 1 点で貼り合わせた 8 の字 $S^1\vee S^1$ を、頂点 $p,a_1,a_2,b_1,b_2$ と 6 本の辺
\[
[p,a_1],\ [a_1,a_2],\ [a_2,p],\ [p,b_1],\ [b_1,b_2],\ [b_2,p]
\]
からなる単体複体 $X$ で表す。$H_n(X)$ を計算し、$H_1(X)$ の基底を具体的に与えよ。
\end{problem}

\begin{solution}
$X$ は連結な 1 次元複体なので $H_0(X)\cong\Z$、$H_n(X)=0$（$n\ge 2$）、$H_1(X)=\ker\partial_1$ である。

1 鎖 $c=x_1[p,a_1]+x_2[a_1,a_2]+x_3[a_2,p]+y_1[p,b_1]+y_2[b_1,b_2]+y_3[b_2,p]$ の境界は
\[
\partial c=(x_1-x_2)a_1+(x_2-x_3)a_2+(y_1-y_2)b_1+(y_2-y_3)b_2+(x_3-x_1+y_3-y_1)p
\]
である。$\partial c=0$ は $x_1=x_2=x_3$、$y_1=y_2=y_3$ と同値（このとき $p$ の係数は自動的に $0$）。よって
\[
H_1(X)=\ker\partial_1=\Z z_a\oplus\Z z_b\cong\Z^2,
\]
\[
z_a=[p,a_1]+[a_1,a_2]+[a_2,p],\qquad z_b=[p,b_1]+[b_1,b_2]+[b_2,p].
\]
これは前問の公式 $E-V+1=6-5+1=2$ とも一致する。$z_a,z_b$ はそれぞれ一方の円周を 1 周するサイクルである。
\end{solution}

\begin{problem}[難易度 B]
単体複体 $K,L$ が $K\cap L=\{v\}$（共通部分は 1 頂点のみ）を満たすとき、$K\vee L=K\cup L$ とおく。すべての $n\ge 0$ について
\[
\tilde H_n(K\vee L)\cong\tilde H_n(K)\oplus\tilde H_n(L)
\]
を示せ。
\end{problem}

\begin{solution}
$n\ge1$ の単体は $K$ と $L$ の一方にのみ属するので、$n\ge 1$ で $C_n(K\vee L)=C_n(K)\oplus C_n(L)$ であり、$\partial$ はこの分解を保つ（$n\ge 2$ なら値も $C_{n-1}(K)\oplus C_{n-1}(L)$ で分解する）。

$n\ge2$：$\ker\partial_n$、$\operatorname{im}\partial_{n+1}$ ともに直和に分解するので $H_n(K\vee L)\cong H_n(K)\oplus H_n(L)$。

$n=1$：$c=a+b$（$a\in C_1(K)$、$b\in C_1(L)$）が $\partial c=0$ を満たすとする。$\partial a=-\partial b\in C_0(K)\cap C_0(L)=\Z v$ なので $\partial a=tv$ と書ける。$\varepsilon\circ\partial=0$ より $t=\varepsilon(\partial a)=0$。よって $\partial a=\partial b=0$ で、$\ker\partial_1=\ker\partial_1^K\oplus\ker\partial_1^L$。$\operatorname{im}\partial_2=\partial C_2(K)\oplus\partial C_2(L)$ と合わせて $H_1(K\vee L)\cong H_1(K)\oplus H_1(L)$。$n\ge 1$ では $\tilde H_n=H_n$ である。

$n=0$：連結成分の個数を $r_K,r_L$ とすると、$v$ を含む成分同士が 1 つになるので $K\vee L$ の成分数は $r_K+r_L-1$。$\tilde H_0\cong\Z^{(\text{成分数})-1}$ より
\[
\tilde H_0(K\vee L)\cong\Z^{r_K+r_L-2}\cong\Z^{r_K-1}\oplus\Z^{r_L-1}\cong\tilde H_0(K)\oplus\tilde H_0(L).
\]
例えば円周 2 つの楔和では $\tilde H_1\cong\Z^2$ となり前問と一致し、$S^2\vee S^1$（$\partial\Delta^3$ と $\partial\Delta^2$ を 1 頂点で貼る）では $\tilde H_1\cong\Z$、$\tilde H_2\cong\Z$、他は $0$ である。
\end{solution}

\begin{problem}[難易度 B]
有限単体複体 $K$ の $k$ 単体の個数を $\alpha_k$ とし、オイラー標数を $\chi(K)=\sum_k(-1)^k\alpha_k$ と定める。ベッチ数 $b_k=\operatorname{rank}H_k(K)$ を用いて
\[
\chi(K)=\sum_k(-1)^k b_k
\]
が成り立つことを示せ。
\end{problem}

\begin{solution}
$Z_k=\ker\partial_k$、$B_k=\operatorname{im}\partial_{k+1}$ とおく。有限生成アーベル群の階数（$\otimes\mathbb{Q}$ した次元）は短完全列に関して加法的である。

完全列 $0\to Z_k\to C_k\xrightarrow{\partial_k}B_{k-1}\to 0$ より
\[
\alpha_k=\operatorname{rank}C_k=\operatorname{rank}Z_k+\operatorname{rank}B_{k-1}.
\]
完全列 $0\to B_k\to Z_k\to H_k(K)\to 0$ より $\operatorname{rank}Z_k=\operatorname{rank}B_k+b_k$。よって
\[
\alpha_k=b_k+\operatorname{rank}B_k+\operatorname{rank}B_{k-1}.
\]
$(-1)^k$ を掛けて $k$ について和をとると（$K$ は有限なので有限和で、$B_{-1}=0$、十分大きい $k$ で $B_k=0$）、$\operatorname{rank}B_k$ の項は $(-1)^k$ と $(-1)^{k+1}$ で 1 回ずつ現れて打ち消し合う。したがって
\[
\chi(K)=\sum_k(-1)^k\alpha_k=\sum_k(-1)^kb_k .
\]
例：$\partial\Delta^3$ では $4-6+4=2=1-0+1$、8 の字では $5-6=-1=1-2$ である。右辺はホモロジーのみで決まるので、$\chi$ は単体分割の取り方によらない量の候補となる。
\end{solution}

\begin{problem}[難易度 C]
単体複体の間の単体写像 $f\colon K\to L$ に対し、$f_\sharp\colon C_n(K)\to C_n(L)$ を
\[
f_\sharp[v_0,\dots,v_n]=\begin{cases}[f(v_0),\dots,f(v_n)] & f(v_0),\dots,f(v_n)\text{ が相異なるとき},\\ 0 & \text{それ以外}\end{cases}
\]
で定める。これが well-defined であり、$\partial\circ f_\sharp=f_\sharp\circ\partial$ を満たすこと（鎖写像であること）を示せ。したがって $f$ は準同型 $f_*\colon H_n(K)\to H_n(L)$ を誘導する。
\end{problem}

\begin{solution}
well-defined：頂点を置換 $\pi$ で並べ替えると像の列も同じ $\pi$ で並べ替わるので、像が相異なれば両辺とも $\operatorname{sgn}(\pi)$ 倍、そうでなければ両辺とも $0$ である。

$\sigma=[v_0,\dots,v_n]$ とし、$f(v_0),\dots,f(v_n)$ の中の相異なる値の個数で場合分けする。

(i) $n+1$ 個のとき：$[f(v_0),\dots,f(v_n)]$ の面は $\sigma$ の面の像なので、両辺は明らかに一致する。

(ii) $n-1$ 個以下のとき：$f_\sharp\sigma=0$。どの面 $[\dots,\widehat{v_i},\dots]$ の像にも重複が残るので $f_\sharp\partial\sigma=0$。

(iii) ちょうど $n$ 個のとき：$f(v_j)=f(v_k)=u$（$j<k$）となる組がただ 1 つある。$f_\sharp\sigma=0$ であり、$f_\sharp\partial\sigma$ で $0$ にならないのは $i=j,k$ の 2 項である：
\[
(-1)^j[f(v_0),\dots,\widehat{f(v_j)},\dots,f(v_n)]+(-1)^k[f(v_0),\dots,\widehat{f(v_k)},\dots,f(v_n)].
\]
2 つの列は $u$ の位置だけが異なり、第 1 の列では $u$ は $k-1$ 番目、第 2 の列では $j$ 番目にある（0 番目から数える）。$u$ を $k-1$ 番目から $j$ 番目へ移すのは $k-1-j$ 回の隣接互換なので、第 1 項 $=(-1)^j(-1)^{k-1-j}X=(-1)^{k-1}X$（$X$ は第 2 の列）となり、第 2 項 $(-1)^kX$ と打ち消し合う。よって $f_\sharp\partial\sigma=0=\partial f_\sharp\sigma$。

鎖写像はサイクルをサイクルに（$\partial f_\sharp z=f_\sharp\partial z=0$）、境界を境界に（$f_\sharp\partial c=\partial f_\sharp c$）写すので、$f_*[z]=[f_\sharp z]$ が矛盾なく定まる。
\end{solution}

\begin{problem}[難易度 B]
円周の 2 つの単体分割として、頂点 $u_0,\dots,u_5$ と辺 $[u_i,u_{i+1}]$（添字は $\bmod 6$）からなる六角形 $P$、および $L=\partial[v_0,v_1,v_2]$ を考える。\\
(1) $f(u_i)=v_{i\bmod 3}$ で定まる単体写像 $f\colon P\to L$ について、$f_*\colon H_1(P)\to H_1(L)$ が $\Z\to\Z$ の 2 倍写像であることを示せ。\\
(2) $g(v_0)=v_0$、$g(v_1)=v_2$、$g(v_2)=v_1$ で定まる単体写像 $g\colon L\to L$ について、$g_*$ が $-1$ 倍であることを示せ。
\end{problem}

\begin{solution}
$P$ は連結グラフで $E-V+1=1$ なので $H_1(P)\cong\Z$ であり、円周の計算と同様に $\ker\partial_1$ は $z_P=\sum_{i=0}^{5}[u_i,u_{i+1}]$ で生成される（$\partial c=0$ から各辺の係数が等しいことが分かる）。$H_1(L)=\Z z$、$z=[v_0,v_1]+[v_1,v_2]+[v_2,v_0]$ である。

(1) 辺 $\{u_i,u_{i+1}\}$ は辺 $\{v_{i\bmod3},v_{(i+1)\bmod3}\}$ に写るので $f$ は単体写像で、辺上で退化しない。
\[
f_\sharp z_P=\sum_{i=0}^{5}[v_{i\bmod 3},v_{(i+1)\bmod3}]
=2\bigl([v_0,v_1]+[v_1,v_2]+[v_2,v_0]\bigr)=2z .
\]
よって生成元同士の対応で $f_*$ は $\Z\to\Z$、$1\mapsto 2$ である。

(2) $g_\sharp z=[v_0,v_2]+[v_2,v_1]+[v_1,v_0]=-\bigl([v_2,v_0]+[v_1,v_2]+[v_0,v_1]\bigr)=-z$。よって $g_*=-\id$。

(1) は円周を 2 周する写像、(2) は折り返し（鏡映）であり、それぞれ写像度 $2$、$-1$ に対応する。
\end{solution}

\begin{problem}[難易度 C]
$L$ を $K$ の部分複体とし、相対鎖群 $C_n(K,L)=C_n(K)/C_n(L)$ と誘導される境界作用素から相対ホモロジー $H_n(K,L)$ を定める。\\
(1) $n\ge1$ とし、$K=\Delta^n$（全ての面）、$L=\partial\Delta^n$ とする。$H_k(K,L)$ を計算せよ。\\
(2) $K$ が連結で $L\ne\emptyset$ ならば $H_0(K,L)=0$ であることを示せ。
\end{problem}

\begin{solution}
$\partial(C_n(L))\subset C_{n-1}(L)$ なので $\partial$ は商に誘導され、$\partial\circ\partial=0$ も保たれる。

(1) $K$ と $L$ の単体の違いは $n$ 単体 $\sigma=[v_0,\dots,v_n]$ のみである。よって
\[
C_k(K,L)=\begin{cases}\Z\bar\sigma & k=n,\\ 0 & k\ne n\end{cases}
\]
であり、相対鎖複体は次数 $n$ にだけ $\Z$ がある複体で、境界作用素はすべて $0$ である（$\partial\sigma\in C_{n-1}(L)$ なので $\partial\bar\sigma=0$）。したがって $H_n(K,L)\cong\Z$（$\bar\sigma$ で生成）、$H_k(K,L)=0$（$k\ne n$）。

$n\ge2$ のとき長完全列 $H_n(K)\to H_n(K,L)\to H_{n-1}(L)\to H_{n-1}(K)$ で両端が $0$ なので $H_n(K,L)\cong H_{n-1}(\partial\Delta^n)\cong\Z$ となり、連結準同型は $\bar\sigma\mapsto[\partial\sigma]$ である。これは上の結果と整合する。

(2) $H_0(K,L)=C_0(K)/\bigl(\operatorname{im}\partial_1+C_0(L)\bigr)$ である。$L$ の頂点 $p$ を 1 つとる。$K$ は連結なので、任意の頂点 $u$ は辺の列で $p$ と結べ、$u-p\in\operatorname{im}\partial_1$。よって $u=(u-p)+p\in\operatorname{im}\partial_1+C_0(L)$ となり、$C_0(K)$ 全体が割られるので $H_0(K,L)=0$。
\end{solution}

\begin{problem}[難易度 C]
頂点 $0,1,2,3,4$ と 5 つの 2 単体 $\sigma_i=[i,i+1,i+2]$（$i=0,\dots,4$、添字は $\bmod 5$）からなる単体複体 $M$ を考える。これは帯状に並べた三角形 $[0,1,2],[1,2,3],[2,3,4],[3,4,0],[4,0,1]$ の両端の辺 $\{0,1\}$ を、ひねって貼り合わせたものであり、メビウスの帯の単体分割である。$H_n(M)$ を計算せよ。
\begin{center}
\begin{tikzpicture}[scale=1.15]
  \coordinate (b0) at (0,0); \coordinate (t1) at (1,1.3); \coordinate (b2) at (2,0);
  \coordinate (t3) at (3,1.3); \coordinate (b4) at (4,0); \coordinate (t0) at (5,1.3);
  \coordinate (b1) at (6,0);
  \fill[blue!6] (b0) -- (b2) -- (b4) -- (b1) -- (t0) -- (t3) -- (t1) -- cycle;
  \draw (b0) -- (b2) -- (b4) -- (b1) (t1) -- (t3) -- (t0);
  \draw[thick,gray!70!black] (t1) -- (b2) -- (t3) -- (b4) -- (t0);
  \draw[very thick,red!70!black,mid arrow] (b0) -- (t1);
  \draw[very thick,red!70!black,mid arrow] (t0) -- (b1);
  \foreach \c/\l/\pos in {b0/0/below,t1/1/above,b2/2/below,t3/3/above,b4/4/below,t0/0/above,b1/1/below}
    \node[vtx,label=\pos:{$\l$}] at (\c) {};
  \foreach \x/\s/\y in {1/0/0.43,2/1/0.87,3/2/0.43,4/3/0.87,5/4/0.43} \node at (\x,\y) {\footnotesize $\sigma_{\s}$};
\end{tikzpicture}
\zucap{左端の辺 $0\to1$ と右端の辺 $0\to1$ を矢印に従って貼る（ひねりが入る）。灰色の辺 $e_1,\dots,e_4$ と赤い辺 $e_0$ をつなぐと中心線を1周する。上下の辺 $f_i=[i,i+2]$ が境界。}
\end{center}
\end{problem}

\begin{solution}
辺は $e_i=[i,i+1]$ と $f_i=[i,i+2]$（$i=0,\dots,4$、$\bmod 5$）の 10 本で、5 頂点の 2 元部分集合すべてである。$\partial\sigma_i=[i+1,i+2]-[i,i+2]+[i,i+1]=e_i+e_{i+1}-f_i$ なので、行を $e_0,\dots,e_4,f_0,\dots,f_4$、列を $\sigma_0,\dots,\sigma_4$ として
\[
\partial_2=\begin{pmatrix}
1&0&0&0&1\\ 1&1&0&0&0\\ 0&1&1&0&0\\ 0&0&1&1&0\\ 0&0&0&1&1\\
-1&0&0&0&0\\ 0&-1&0&0&0\\ 0&0&-1&0&0\\ 0&0&0&-1&0\\ 0&0&0&0&-1
\end{pmatrix}.
\]
$H_2$：$\partial(\sum a_i\sigma_i)$ の $f_i$ の係数は $-a_i$ なので、$\ker\partial_2=0$、$H_2(M)=0$。

$H_1$：$z=e_0+e_1+e_2+e_3+e_4$ はサイクルである。任意のサイクル $c=\sum x_ie_i+\sum y_if_i$ に対し $c+\sum y_i\partial\sigma_i$ は $e_i$ のみを含むサイクルであり、5 角形 $0\to1\to\dots\to4\to0$ 上のサイクルなので係数はすべて等しく、$mz$ の形である。よって $\Z\to H_1(M)$、$m\mapsto[mz]$ は全射。$mz=\partial(\sum a_i\sigma_i)$ なら $f_i$ の係数から $a_i=0$ となり $m=0$ なので単射。ゆえに $H_1(M)\cong\Z$（生成元は中心線 $z$）。

$H_0$：$M$ は連結なので $H_0(M)\cong\Z$。オイラー標数 $5-10+5=0=1-1+0$ とも整合する。なお境界の辺は $f_i$（1 つの三角形にのみ含まれる）で、$0,2,4,1,3$ を回る 1 つの円周をなす。
\end{solution}

\begin{problem}[難易度 D]
有限 2 次元単体複体 $K$ が次を満たすとする：各辺はちょうど 2 つの 2 単体の面であり、各頂点・辺はある 2 単体に含まれる。各 2 単体 $\sigma$ に向きを選び $s=\sum_\sigma\sigma$ とおく。辺 $e$ を含む 2 単体 $\sigma,\tau$ が $e$ に逆の向きを誘導するとき $e$ を整合的という。整合的な辺を共有する 2 単体同士を結んだグラフ $\Gamma$ が連結であるとする。\\
(a) すべての辺が整合的なら $H_2(K)\cong\Z$ であり、$H_1(K)$ はねじれを持たないことを示せ。\\
(b) 整合的でない辺が存在するなら $H_2(K)=0$ であり、$H_1(K)$ のねじれ部分群は $\Z/2$ であることを示せ。
\end{problem}

\begin{solution}
辺 $e$ を含む $\sigma,\tau$ について、$\partial\sigma,\partial\tau$ における $e$ の係数を $\epsilon_\sigma,\epsilon_\tau\in\{\pm1\}$ とすると、$d=\sum d_\sigma\sigma$ の境界の $e$ の係数は $\epsilon_\sigma d_\sigma+\epsilon_\tau d_\tau$ で、整合的なら $\epsilon_\sigma(d_\sigma-d_\tau)$、そうでなければ $\epsilon_\sigma(d_\sigma+d_\tau)$。特に $\partial s=2w$（$w$ は整合的でない辺の $\pm$ 和）で、$2\partial w=0$ より $w$ はサイクル。

(a) $\partial d=0$ なら隣接する $\sigma,\tau$ で $d_\sigma=d_\tau$、$\Gamma$ が連結なので $d=ks$。$\partial s=0$ より $H_2=\ker\partial_2=\Z s\cong\Z$。サイクル $c$ が $mc=\partial d$（$m\ge1$）なら、$d_\sigma\equiv d_\tau\pmod m$ が隣接で成り立ち $d=ks+md'$。$mc=m\,\partial d'$ で $C_1$ はねじれがないので $c=\partial d'$。よって $H_1$ はねじれなし。

(b) $\partial d=0$ なら $\Gamma$ に沿って $d_\sigma=k$（一定）。整合的でない辺では $2k=0$ より $k=0$。よって $H_2=0$。ねじれ元 $[c]$、$mc=\partial d$ をとると、同様に $d=ks+md'$ かつ $2k\equiv 0\pmod m$、$2k=ml$ とおけて
\[
mc=k\,\partial s+m\,\partial d'=m(lw+\partial d'),\qquad c=lw+\partial d'.
\]
よってねじれ部分群は $[w]$ で生成され、$2[w]=[\partial s]=0$。$w=\partial d$ なら $\partial(s-2d)=0$ より $s=2d$ となるが $s$ の係数は $1$ なので矛盾。ゆえに $[w]\ne0$ でねじれ部分群は $\Z/2$。
\end{solution}

\begin{problem}[難易度 C]
正方形を $3\times3$ の小正方形に分け、各小正方形を左上から右下への対角線で 2 つの三角形に分ける。頂点に
\[
\begin{matrix}0&1&2&0\\3&4&5&3\\6&7&8&6\\0&1&2&0\end{matrix}
\]
（上の行が正方形の上辺）とラベルを付けて同じラベルを同一視すると、トーラス $T^2$ の単体分割 $K$（9 頂点、27 辺、18 個の 2 単体）が得られる。前問を用いて $H_n(K)$ を計算せよ。
\begin{center}
\begin{tikzpicture}[scale=1.0]
  \fill[blue!5] (0,0) rectangle (3,3);
  \draw[gray] (0,0) grid (3,3);
  \foreach \i in {0,1,2} \foreach \j in {0,1,2} \draw[gray] (\i,\j+1) -- (\i+1,\j);
  \draw[very thick,red!70!black,mid arrow] (0,0) -- (3,0);
  \draw[very thick,red!70!black,mid arrow] (0,3) -- (3,3);
  \draw[very thick,blue!70!black,mid arrows] (0,0) -- (0,3);
  \draw[very thick,blue!70!black,mid arrows] (3,0) -- (3,3);
  \foreach \x/\y/\l in {0/3/0,1/3/1,2/3/2,3/3/0, 0/2/3,1/2/4,2/2/5,3/2/3,
                        0/1/6,1/1/7,2/1/8,3/1/6, 0/0/0,1/0/1,2/0/2,3/0/0}
    \node[fill=white,inner sep=1pt,font=\footnotesize] at (\x,\y) {$\l$};
\end{tikzpicture}
\zucap{トーラスの単体分割。上下（単矢印）・左右（二重矢印）の辺を同じ向きに貼る。どの三角形も3頂点のラベルが相異なり、2つの三角形が同じ3頂点をもつこともないので、単体複体になっている。}
\end{center}
\end{problem}

\begin{solution}
小正方形の左上・右上・左下・右下を $a,b,c,d$ とし、2 単体を $[a,c,d]$、$[a,d,b]$ と向き付ける（平面上で反時計回り）。辺は水平・垂直・対角それぞれ 9 本で計 27 本、三角形は 18 個であり、各辺はちょうど 2 つの三角形に含まれる。

整合性：正方形の内部の辺では、隣り合う 2 つの反時計回りの三角形が逆の向きを誘導する。上辺の辺は右上の三角形 $[a,d,b]$ から $-[a,b]$（右から左）、下辺の同じラベルの辺は左下の三角形 $[a,c,d]$ から $+[c,d]$（左から右）を受け取り、同じ向きで同一視されるので打ち消し合う。左右の辺も同様（左辺の $[a,c]$ は上から下、右辺の $[d,b]$ は下から上、同一視は同じ向き）。よってすべての辺が整合的で $\partial s=0$。$\Gamma$ は正方形の三角形の隣接関係を含み連結。

前問 (a) より $H_2(K)\cong\Z$（生成元 $s$）、$H_1(K)$ はねじれなし。$K$ は連結で $H_0(K)\cong\Z$。オイラー標数は $\chi(K)=9-27+18=0$ なので
\[
0=b_0-b_1+b_2=1-b_1+1,\qquad b_1=2 .
\]
有限生成アーベル群の基本定理より $H_1(K)\cong\Z^2$。まとめて
\[
H_0(K)\cong\Z,\quad H_1(K)\cong\Z^2,\quad H_2(K)\cong\Z,\quad H_n(K)=0\ (n\ge3).
\]
\end{solution}

\begin{problem}[難易度 D]
前問と同じ $3\times3$ の分割で、頂点のラベルを
\[
\begin{matrix}0&1&2&0\\3&4&5&6\\6&7&8&3\\0&1&2&0\end{matrix}
\]
に変える（右辺は左辺と逆向きに同一視される）と、クラインの壺の単体分割 $K$（9 頂点、27 辺、18 個の 2 単体）が得られる。$H_n(K)$ を計算せよ。
\begin{center}
\begin{tikzpicture}[scale=1.0]
  \fill[blue!5] (0,0) rectangle (3,3);
  \draw[gray] (0,0) grid (3,3);
  \foreach \i in {0,1,2} \foreach \j in {0,1,2} \draw[gray] (\i,\j+1) -- (\i+1,\j);
  \draw[very thick,red!70!black,mid arrow] (0,0) -- (3,0);
  \draw[very thick,red!70!black,mid arrow] (0,3) -- (3,3);
  \draw[very thick,blue!70!black,mid arrows] (0,0) -- (0,3);
  \draw[very thick,blue!70!black,mid arrows] (3,3) -- (3,0);
  \foreach \x/\y/\l in {0/3/0,1/3/1,2/3/2,3/3/0, 0/2/3,1/2/4,2/2/5,3/2/6,
                        0/1/6,1/1/7,2/1/8,3/1/3, 0/0/0,1/0/1,2/0/2,3/0/0}
    \node[fill=white,inner sep=1pt,font=\footnotesize] at (\x,\y) {$\l$};
\end{tikzpicture}
\zucap{クラインの壺の単体分割。前問との違いは右辺のラベル（$3$ と $6$ が入れ替わる）だけで、左右の辺は逆向きに貼られる（二重矢印の向きに注意）。}
\end{center}
\end{problem}

\begin{solution}
前問と同じく各三角形を反時計回りに $[a,c,d]$、$[a,d,b]$ と向き付け、$s$ をその和とする。各辺はちょうど 2 つの三角形に含まれる。内部の辺および上下の辺（同じ向きで同一視）は前問と同様に整合的であり、$\Gamma$ は連結である。

左辺の三角形 $[a,c,d]$ は辺 $[0,3],[3,6],[6,0]$ を係数 $+1$ で含む。右辺の三角形 $[a,d,b]$ は $[d,b]$ を係数 $+1$ で含み、右辺のラベルは上から $0,6,3,0$ なので $[6,0],[3,6],[0,3]$ となる。よってこれら 3 辺は整合的でなく
\[
\partial s=2w,\qquad w=[0,3]+[3,6]+[6,0].
\]
前問 (b) より $H_2(K)=0$、$H_1(K)$ のねじれ部分群は $[w]$ で生成される $\Z/2$。

$K$ は連結で $H_0\cong\Z$。$\chi(K)=9-27+18=0=1-b_1+0$ より $b_1=1$。基本定理より
\[
H_0(K)\cong\Z,\quad H_1(K)\cong\Z\oplus\Z/2,\quad H_n(K)=0\ (n\ge2).
\]
$w$ は左辺（＝右辺）に沿う閉曲線で、2 倍すると境界 $\partial s$ になる。
\end{solution}

\begin{problem}[難易度 D]
メビウスの帯の問題の複体 $M$（頂点 $0,\dots,4$、$\sigma_i=[i,i+1,i+2]$）に新しい頂点 $5$ を加え、境界の辺 $f_i=[i,i+2]$ 上の錐 $\tau_i=[5,i,i+2]$（$i=0,\dots,4$、$\bmod 5$）を貼ると、実射影平面 $\mathbb{RP}^2$ の単体分割 $K$（6 頂点、15 辺、10 個の 2 単体）が得られる。$H_n(K)$ を計算せよ。
\end{problem}

\begin{solution}
$e_i=[i,i+1]$、$g_j=[5,j]$ とおく。辺は $e_i,f_i,g_j$ の 15 本で、各辺はちょうど 2 つの 2 単体に含まれる（$e_i$ は $\sigma_{i-1},\sigma_i$、$f_i$ は $\sigma_i,\tau_i$、$g_j$ は $\tau_j,\tau_{j-2}$）。向きを $s=\sum_i(-1)^i(\sigma_i+\tau_i)$ と選ぶ。
\[
\partial\sigma_i=e_i+e_{i+1}-f_i,\qquad \partial\tau_i=f_i-g_{i+2}+g_i .
\]
$f_i$ の係数は $0$、$e_j$（$1\le j\le4$）の係数は $(-1)^j+(-1)^{j-1}=0$、$e_0$ の係数は $1+(-1)^4=2$。$g_j$ の係数は $(-1)^j-(-1)^{j-2\bmod5}$ で、$j=0,1,2,3,4$ に対し $2,-2,0,0,0$。よって
\[
\partial s=2w,\qquad w=e_0+g_0-g_1=[0,1]+[1,5]+[5,0].
\]
$e_1,\dots,e_4,f_0,\dots,f_4$ は整合的で、$\sigma_0,\dots,\sigma_4$ を鎖状に結び、各 $\tau_i$ を $\sigma_i$ に結ぶので $\Gamma$ は連結。$e_0$ は整合的でない。

先の閉曲面の問題 (b) より $H_2(K)=0$、$H_1(K)$ のねじれ部分群は $[w]$ で生成される $\Z/2$。$K$ は連結で $H_0\cong\Z$。$\chi(K)=6-15+10=1=1-b_1+0$ より $b_1=0$、よって $H_1(K)$ は有限で
\[
H_0(K)\cong\Z,\quad H_1(K)\cong\Z/2,\quad H_n(K)=0\ (n\ge2).
\]
メビウスの帯の中心線 $z$ は $K$ では $2z$ が境界となる。
\end{solution}

%================================================
\chapter{特異ホモロジー}
% 25問
%================================================

\begin{summary}
\begin{itemize}
  \item \textbf{特異鎖}：特異 $n$ 単体は連続写像 $\sigma\colon\Delta^n\to X$。$C_n(X)$ はそれらで生成される自由アーベル群で、$\partial\sigma=\sum_i(-1)^i\sigma|_{[v_0,\dots,\widehat{v_i},\dots,v_n]}$。
  \item \textbf{基本性質}：$H_0(X)\cong\bigoplus_{\text{弧状連結成分}}\Z$。一点のホモロジーは $H_0=\Z$、それ以外は $0$。\textbf{ホモトピー不変性}：$f\simeq g$ なら $f_*=g_*$。特にホモトピー同値写像は同型を誘導する。
  \item \textbf{長完全列}：鎖複体の短完全列 $0\to A_\bullet\to B_\bullet\to C_\bullet\to0$ から、連結準同型 $\partial$ を含む長完全列
  \[
  \cdots\to H_n(A)\to H_n(B)\to H_n(C)\xrightarrow{\partial}H_{n-1}(A)\to\cdots
  \]
  が得られる。特に対 $(X,A)$ の長完全列がある。\textbf{5項補題}もよく使う。
  \item \textbf{切除定理}：$\overline Z\subset\operatorname{int}A$ なら $H_n(X\setminus Z,A\setminus Z)\cong H_n(X,A)$。
  \item \textbf{良い対}：$A$ が閉で、$A$ のある近傍が $A$ に変位レトラクトするなら $H_n(X,A)\cong\tilde H_n(X/A)$。
  \item \textbf{マイヤー・ヴィートリス完全列}：$X=\operatorname{int}A\cup\operatorname{int}B$ なら
  \[
  \cdots\to H_n(A\cap B)\to H_n(A)\oplus H_n(B)\to H_n(X)\to H_{n-1}(A\cap B)\to\cdots
  \]
  \item \textbf{球面}：$\tilde H_k(S^n)\cong\Z$（$k=n$）、$0$（$k\neq n$）。
  \item \textbf{フレヴィッチの定理}：$X$ が弧状連結なら $h\colon\pi_1(X)\to H_1(X)$ は全射で、核は交換子群。すなわち $H_1(X)\cong\pi_1(X)^{\mathrm{ab}}$。
\end{itemize}
\end{summary}

\begin{problem}[難易度 A]
位相空間 $X$ の特異 $n$ 単体 $\sigma\colon \Delta^n\to X$ に対し、境界を
\[
\partial_n\sigma=\sum_{i=0}^{n}(-1)^i\,\sigma|_{[v_0,\dots,\hat{v}_i,\dots,v_n]}
\]
で定め、$C_n(X)$ 上に線形に拡張する。ここで $\sigma|_{[v_0,\dots,\hat{v}_i,\dots,v_n]}$ は $\sigma$ を第 $i$ 面に制限し、頂点の順序を保つ線形同型で $\Delta^{n-1}$ と同一視したものである。$\partial_{n-1}\circ\partial_n=0$ を示せ。
\end{problem}

\begin{solution}
特異 $n$ 単体 $\sigma$ について示せば十分である。定義を2回適用すると、第 $i$ 面からさらに頂点 $v_j$ を除く項が現れる。$v_i$ を除いた後の頂点列で $v_j$ が何番目にあるかに注意すると、$j<i$ なら $j$ 番目、$j>i$ なら $j-1$ 番目である。したがって
\begin{align*}
\partial\partial\sigma
&=\sum_{j<i}(-1)^i(-1)^j\,\sigma|_{[v_0,\dots,\hat{v}_j,\dots,\hat{v}_i,\dots,v_n]}\\
&\quad+\sum_{j>i}(-1)^i(-1)^{j-1}\,\sigma|_{[v_0,\dots,\hat{v}_i,\dots,\hat{v}_j,\dots,v_n]}
\end{align*}
となる。どちらの和も、2つの頂点 $v_p,v_q$（$p<q$）を除いた面 $\sigma|_{[\dots,\hat{v}_p,\dots,\hat{v}_q,\dots]}$ を各 $(p,q)$ についてちょうど1回ずつ含む。第1の和では $(j,i)=(p,q)$ で係数は $(-1)^{p+q}$、第2の和では $(i,j)=(p,q)$ で係数は $(-1)^{p+q-1}$ である。これらは打ち消し合うので $\partial\partial\sigma=0$ である。
\end{solution}

\begin{problem}[難易度 A]
連続写像 $f\colon X\to Y$ に対し $f_\sharp\colon C_n(X)\to C_n(Y)$ を $f_\sharp(\sigma)=f\circ\sigma$ で定める。$f_\sharp$ が鎖写像であること、したがって準同型 $f_*\colon H_n(X)\to H_n(Y)$ を誘導することを示せ。さらに $(g\circ f)_*=g_*\circ f_*$ および $(\id_X)_*=\id$ を示せ。
\end{problem}

\begin{solution}
特異単体 $\sigma\colon\Delta^n\to X$ に対し、面への制限と $f$ の合成は順序を交換できる：$f\circ(\sigma|_{[v_0,\dots,\hat{v}_i,\dots,v_n]})=(f\circ\sigma)|_{[v_0,\dots,\hat{v}_i,\dots,v_n]}$。よって
\[
f_\sharp\partial\sigma=\sum_i(-1)^i f\circ\sigma|_{[\dots,\hat{v}_i,\dots]}=\sum_i(-1)^i (f\circ\sigma)|_{[\dots,\hat{v}_i,\dots]}=\partial f_\sharp\sigma
\]
であり、線形性から $f_\sharp\partial=\partial f_\sharp$、すなわち $f_\sharp$ は鎖写像である。

一般に鎖写像 $\varphi$ は、$z$ がサイクルなら $\partial\varphi(z)=\varphi(\partial z)=0$ よりサイクルを、$\partial w$ を $\varphi(\partial w)=\partial\varphi(w)$ より境界を境界に写す。したがって $[z]\mapsto[\varphi(z)]$ は $H_n$ 上の well-defined な準同型である。これを $f_*$ とする。

鎖のレベルで $(g\circ f)_\sharp(\sigma)=g\circ f\circ\sigma=g_\sharp(f_\sharp(\sigma))$、$(\id_X)_\sharp(\sigma)=\sigma$ が成り立つ。ホモロジー類の代表元にこれを適用すれば $(g\circ f)_*=g_*\circ f_*$ と $(\id_X)_*=\id$ が得られる。
\end{solution}

\begin{problem}[難易度 A]
一点からなる空間 $P$ の特異ホモロジーが $H_0(P)\cong\Z$、$H_n(P)=0$（$n\ge 1$）であることを示せ。
\end{problem}

\begin{solution}
各 $n\ge 0$ に対し特異 $n$ 単体 $\Delta^n\to P$ はただ一つ（定値写像 $\sigma_n$）であるから $C_n(P)=\Z\sigma_n\cong\Z$ である。$\sigma_n$ のどの面も $\sigma_{n-1}$ であるから、$n\ge1$ に対し
\[
\partial_n\sigma_n=\sum_{i=0}^n(-1)^i\sigma_{n-1}=
\begin{cases}\sigma_{n-1} & (n\ \text{が偶数}),\\ 0 & (n\ \text{が奇数})\end{cases}
\]
である。よって鎖複体は
\[
\cdots\xrightarrow{\ \cong\ }\Z\xrightarrow{\ 0\ }\Z\xrightarrow{\ \cong\ }\Z\xrightarrow{\ 0\ }\Z\longrightarrow 0
\]
（右端が $C_0$、$\partial_1=0$、$\partial_2$ は同型、…）となる。$H_0=C_0/\operatorname{im}\partial_1=\Z$ である。$n$ が奇数なら $\ker\partial_n=\Z$、$\operatorname{im}\partial_{n+1}=\Z$（$n+1$ は偶数で $\partial_{n+1}$ は同型）なので $H_n=0$。$n\ge2$ が偶数なら $\partial_n$ は単射で $\ker\partial_n=0$ なので $H_n=0$ である。
\end{solution}

\begin{problem}[難易度 A]
$X$ の弧状連結成分を $\{X_\alpha\}_\alpha$ とするとき、包含写像が誘導する準同型により $\bigoplus_\alpha H_n(X_\alpha)\cong H_n(X)$ がすべての $n$ で成り立つことを示せ。
\end{problem}

\begin{solution}
$\Delta^n$ は弧状連結であるから、特異単体 $\sigma\colon\Delta^n\to X$ の像は弧状連結であり、ただ一つの $X_\alpha$ に含まれる。したがって $C_n(X)$ の基底は $X_\alpha$ ごとに分割され、
\[
C_n(X)=\bigoplus_\alpha C_n(X_\alpha)
\]
となる（右辺の各成分は包含写像による像と同一視している）。$\sigma$ が $X_\alpha$ に値をとればその面も $X_\alpha$ に値をとるので、$\partial$ はこの直和分解を保つ。すなわち鎖複体として $C_\bullet(X)=\bigoplus_\alpha C_\bullet(X_\alpha)$ である。

直和成分ごとに境界をとるので、$\ker\partial_n=\bigoplus_\alpha\ker\partial_n^{(\alpha)}$、$\operatorname{im}\partial_{n+1}=\bigoplus_\alpha\operatorname{im}\partial_{n+1}^{(\alpha)}$ が成り立つ。したがって
\[
H_n(X)=\bigoplus_\alpha \ker\partial^{(\alpha)}_n\big/\operatorname{im}\partial^{(\alpha)}_{n+1}=\bigoplus_\alpha H_n(X_\alpha)
\]
であり、この同型は包含写像の誘導する準同型の和で与えられる。
\end{solution}

\begin{problem}[難易度 B]
空でない弧状連結空間 $X$ に対し $H_0(X)\cong\Z$ を示せ。これと前問を用いて、一般の $X$ に対し $H_0(X)$ が $X$ の弧状連結成分全体を基底とする自由アーベル群であることを示せ。
\end{problem}

\begin{solution}
準同型 $\varepsilon\colon C_0(X)\to\Z$ を $\varepsilon(\sum_i n_ix_i)=\sum_i n_i$ で定める（特異0単体を点と同一視する）。$X\ne\emptyset$ より $\varepsilon$ は全射である。$\ker\varepsilon=\operatorname{im}\partial_1$ を示せば、$\partial_0=0$ より $H_0(X)=C_0(X)/\operatorname{im}\partial_1\cong\Z$ となる。

特異1単体 $\sigma$ に対し $\partial\sigma=\sigma(v_1)-\sigma(v_0)$ なので $\varepsilon\partial\sigma=0$、よって $\operatorname{im}\partial_1\subset\ker\varepsilon$。逆に $c=\sum_in_ix_i$ が $\sum_in_i=0$ を満たすとする。点 $x_0\in X$ を固定し、弧状連結性から $x_0$ から $x_i$ への道 $\gamma_i$ をとり特異1単体とみなせば
\[
\partial\Bigl(\sum_in_i\gamma_i\Bigr)=\sum_in_ix_i-\Bigl(\sum_in_i\Bigr)x_0=c
\]
となり、$c\in\operatorname{im}\partial_1$ である。

一般の $X$ については、前問より $H_0(X)\cong\bigoplus_\alpha H_0(X_\alpha)\cong\bigoplus_\alpha\Z$ である。基底は各成分の任意の1点の類 $[x_\alpha]$（$x_\alpha\in X_\alpha$）で与えられ、これは $x_\alpha$ の選び方によらない。
\end{solution}

\begin{problem}[難易度 B]
特異1単体 $\sigma\colon\Delta^1\to X$ を道 $s\mapsto\sigma((1-s)v_0+sv_1)$（$s\in[0,1]$）と同一視し、逆向きの道を $\bar\sigma$ と書く。次を示せ。
(1) 定値な特異1単体 $c_x$ は境界である。
(2) $\sigma+\bar\sigma$ は境界である。
\end{problem}

\begin{solution}
(1) 点 $x$ への定値写像 $\tau\colon\Delta^2\to X$ を考える。その3つの面はすべて $c_x$ なので $\partial\tau=c_x-c_x+c_x=c_x$ である。

(2) 線形写像 $p\colon\Delta^2\to\Delta^1$ を $p(v_0)=v_0$、$p(v_1)=v_1$、$p(v_2)=v_0$ で定め、$\tau=\sigma\circ p$ とおく。各面を調べると、
\begin{itemize}
\item 第0面 $[v_1,v_2]$ は $p$ により $v_1$ から $v_0$ へ線形に写るので $\tau|_{[v_1,v_2]}=\bar\sigma$、
\item 第1面 $[v_0,v_2]$ は $p$ により点 $v_0$ に潰れるので $\tau|_{[v_0,v_2]}=c_{x_0}$（$x_0=\sigma(v_0)$）、
\item 第2面 $[v_0,v_1]$ は恒等的に写るので $\tau|_{[v_0,v_1]}=\sigma$
\end{itemize}
である。よって $\partial\tau=\bar\sigma-c_{x_0}+\sigma$。(1) より $c_{x_0}=\partial\tau'$ となる $\tau'$ があるので、$\sigma+\bar\sigma=\partial(\tau+\tau')$ は境界である。
\end{solution}

\begin{problem}[難易度 A]
鎖写像 $f,g\colon C_\bullet\to D_\bullet$ が鎖ホモトピック、すなわち準同型 $P\colon C_n\to D_{n+1}$ で $f-g=\partial P+P\partial$ を満たすものが存在するとき、$f_*=g_*\colon H_n(C)\to H_n(D)$ を示せ。
\end{problem}

\begin{solution}
$z\in C_n$ をサイクル（$\partial z=0$）とする。仮定より
\[
f(z)-g(z)=\partial P(z)+P(\partial z)=\partial P(z)
\]
であるから、$f(z)-g(z)$ は $D_\bullet$ の境界である。よって $H_n(D)$ において $[f(z)]=[g(z)]$、すなわち $f_*[z]=g_*[z]$ が成り立つ。$[z]$ は任意なので $f_*=g_*$ である。

（注）特異ホモロジーのホモトピー不変性の証明では、ホモトピー $F\colon X\times[0,1]\to X$ から $\Delta^n\times[0,1]$ の単体分割を用いて「プリズム作用素」$P$ を構成し、$F_0$ と $F_1$ の誘導する鎖写像が鎖ホモトピックであることを示す。その最後の一歩がこの問題の主張である。
\end{solution}

\begin{problem}[難易度 A]
ホモトピー不変性（$f\simeq g$ ならば $f_*=g_*$）を用いて、次を示せ。
(1) $f\colon X\to Y$ がホモトピー同値写像ならば $f_*\colon H_n(X)\to H_n(Y)$ はすべての $n$ で同型である。
(2) 可縮な空間 $X$ では $H_0(X)\cong\Z$、$H_n(X)=0$（$n\ge1$）である。特に $\R^m$ や円板 $D^m$ はそうである。
\end{problem}

\begin{solution}
(1) ホモトピー逆写像 $g\colon Y\to X$ をとると $g\circ f\simeq\id_X$、$f\circ g\simeq\id_Y$ である。関手性とホモトピー不変性から
\[
g_*\circ f_*=(g\circ f)_*=(\id_X)_*=\id,\qquad f_*\circ g_*=(f\circ g)_*=(\id_Y)_*=\id
\]
となり、$f_*$ は同型で逆写像は $g_*$ である。

(2) $X$ が可縮であるとは、一点空間 $P$ への写像 $X\to P$ がホモトピー同値であることである。(1) より $H_n(X)\cong H_n(P)$ であり、一点のホモロジーの計算から結論を得る。$\R^m$ と $D^m$ は直線ホモトピー $H(x,t)=(1-t)x$ により原点に可縮であるから、これらにも適用できる。
\end{solution}

\begin{problem}[難易度 A]
連続写像 $f\colon X\to Y$ が定値写像にホモトピックならば、$n\ge1$ に対し $f_*=0\colon H_n(X)\to H_n(Y)$ であることを示せ。
\end{problem}

\begin{solution}
$f$ が点 $y_0\in Y$ への定値写像 $c$ にホモトピックであるとする。ホモトピー不変性より $f_*=c_*$ である。$c$ は一点空間 $P$ を経由して
\[
c\colon X\xrightarrow{\ q\ }P\xrightarrow{\ j\ }Y,\qquad j(P)=y_0
\]
と分解するので、関手性から $c_*=j_*\circ q_*$ である。$n\ge1$ では $H_n(P)=0$ であるから、$q_*\colon H_n(X)\to H_n(P)=0$ は零写像であり、したがって $c_*=0$、よって $f_*=0$ である。

（注）$n=0$ では一般に $f_*\ne0$ である。例えば $X=Y=P$ のとき $f_*=\id_{\Z}$ である。
\end{solution}

\begin{problem}[難易度 B]
空でない空間 $X$ に対し、添加写像 $\varepsilon\colon C_0(X)\to\Z$、$\varepsilon(\sum n_ix_i)=\sum n_i$ を付け加えた鎖複体
\[
\cdots\to C_1(X)\xrightarrow{\partial_1}C_0(X)\xrightarrow{\varepsilon}\Z\to0
\]
のホモロジーを簡約ホモロジー $\tilde H_n(X)$ という。$n\ge1$ では $\tilde H_n(X)=H_n(X)$、また $H_0(X)\cong\tilde H_0(X)\oplus\Z$ であることを示せ。さらに $X$ が可縮ならばすべての $n$ で $\tilde H_n(X)=0$ であることを示せ。
\end{problem}

\begin{solution}
$\varepsilon\partial_1=0$（$\partial\sigma=\sigma(v_1)-\sigma(v_0)$）であるから、これは鎖複体である。$n\ge1$ では、次数 $n$ のサイクルと境界は元の鎖複体と同じなので $\tilde H_n(X)=H_n(X)$ である。

次数0では $\tilde H_0(X)=\ker\varepsilon/\operatorname{im}\partial_1$、$H_0(X)=C_0(X)/\operatorname{im}\partial_1$ である。$\varepsilon$ は $\operatorname{im}\partial_1$ 上で $0$ なので $\bar\varepsilon\colon H_0(X)\to\Z$ を誘導し、完全列
\[
0\to\tilde H_0(X)\to H_0(X)\xrightarrow{\ \bar\varepsilon\ }\Z\to0
\]
を得る。$X\ne\emptyset$ なので点 $x_0$ をとると $\bar\varepsilon[x_0]=1$ であり、$\Z\to H_0(X)$、$1\mapsto[x_0]$ は $\bar\varepsilon$ の切断である。よってこの完全列は分裂し、$H_0(X)\cong\tilde H_0(X)\oplus\Z$ である。

$X$ が可縮ならば $X$ は弧状連結で $H_0(X)\cong\Z$ であり、$\bar\varepsilon[x_0]=1$ より $\bar\varepsilon$ は同型、よって $\tilde H_0(X)=0$ である。$n\ge1$ では $\tilde H_n(X)=H_n(X)=0$ である。
\end{solution}

\begin{problem}[難易度 C]
アーベル群の鎖複体の短完全列 $0\to A_\bullet\xrightarrow{i}B_\bullet\xrightarrow{p}C_\bullet\to0$ が与えられたとする。連結準同型 $\partial_*\colon H_n(C)\to H_{n-1}(A)$ を構成し、それが代表元などの選び方によらず well-defined な準同型であることを示せ。
\end{problem}

\begin{solution}
\textbf{構成。} $c\in C_n$ をサイクルとする。$p$ は全射なので $p(b)=c$ となる $b\in B_n$ がある。$p(\partial b)=\partial p(b)=\partial c=0$ なので、完全性から $\partial b=i(a)$ となる $a\in A_{n-1}$ がただ一つ存在する（$i$ は単射）。さらに $i(\partial a)=\partial i(a)=\partial\partial b=0$ と $i$ の単射性から $\partial a=0$、すなわち $a$ はサイクルである。そこで $\partial_*[c]=[a]$ と定める。

\textbf{$b$ の選び方によらないこと。} 別の $b'$ で $p(b')=c$ とすると $p(b'-b)=0$ より $b'=b+i(a'')$（$a''\in A_n$）と書ける。$\partial b'=i(a)+i(\partial a'')$ なので対応する元は $a+\partial a''$ であり、$[a+\partial a'']=[a]$ である。

\textbf{$c$ の選び方によらないこと。} $c'=c+\partial\gamma$（$\gamma\in C_{n+1}$）とし、$\gamma=p(\beta)$ となる $\beta\in B_{n+1}$ をとる。$b'=b+\partial\beta$ とおくと $p(b')=c+\partial p(\beta)=c'$ であり、$\partial b'=\partial b=i(a)$ だから同じ $a$ が得られる。

\textbf{準同型であること。} $c_1,c_2$ に対し $b_1,b_2,a_1,a_2$ を上のようにとれば、$c_1+c_2$ に対して $b_1+b_2$、$a_1+a_2$ を選べる。選び方によらないことから $\partial_*[c_1+c_2]=\partial_*[c_1]+\partial_*[c_2]$ である。
\end{solution}

\begin{problem}[難易度 B]
前問の記号のもとで、長い列
\[
\cdots\to H_n(B)\xrightarrow{p_*}H_n(C)\xrightarrow{\partial_*}H_{n-1}(A)\xrightarrow{i_*}H_{n-1}(B)\to\cdots
\]
が $H_n(C)$ と $H_{n-1}(A)$ において完全であることを示せ。
\end{problem}

\begin{solution}
\textbf{$H_n(C)$ での完全性。} $b$ がサイクルなら $c=p(b)$ に対して $b$ 自身を選べ、$\partial b=0=i(0)$ より $\partial_*p_*[b]=0$。よって $\operatorname{im}p_*\subset\ker\partial_*$。逆に $\partial_*[c]=0$ とし、構成に現れる $b,a$ をとると $a=\partial a'$（$a'\in A_n$）と書ける。このとき
\[
\partial\bigl(b-i(a')\bigr)=i(a)-i(\partial a')=0,\qquad p\bigl(b-i(a')\bigr)=c
\]
なので $b-i(a')$ はサイクルで $p_*[b-i(a')]=[c]$。よって $\ker\partial_*\subset\operatorname{im}p_*$。

\textbf{$H_{n-1}(A)$ での完全性。} 構成より $i(a)=\partial b$ は境界だから $i_*\partial_*[c]=0$、すなわち $\operatorname{im}\partial_*\subset\ker i_*$。逆に $a\in A_{n-1}$ をサイクルで $i_*[a]=0$、つまり $i(a)=\partial b$（$b\in B_n$）とする。$c=p(b)$ とおくと $\partial c=p(\partial b)=p(i(a))=0$ なので $c$ はサイクルであり、構成の定義から $\partial_*[c]=[a]$ である。よって $\ker i_*\subset\operatorname{im}\partial_*$。
\end{solution}

\begin{problem}[難易度 C]
アーベル群と準同型の可換図式
\[
\begin{array}{ccccccccc}
A_1 & \xrightarrow{\alpha_1} & A_2 & \xrightarrow{\alpha_2} & A_3 & \xrightarrow{\alpha_3} & A_4 & \xrightarrow{\alpha_4} & A_5\\
\big\downarrow{\scriptstyle f_1} & & \big\downarrow{\scriptstyle f_2} & & \big\downarrow{\scriptstyle f_3} & & \big\downarrow{\scriptstyle f_4} & & \big\downarrow{\scriptstyle f_5}\\
B_1 & \xrightarrow{\beta_1} & B_2 & \xrightarrow{\beta_2} & B_3 & \xrightarrow{\beta_3} & B_4 & \xrightarrow{\beta_4} & B_5
\end{array}
\]
において2つの行が完全であり、$f_1,f_2,f_4,f_5$ が同型であるとする。$f_3$ が同型であることを示せ（5項補題）。
\end{problem}

\begin{solution}
\textbf{単射性。} $x\in A_3$、$f_3(x)=0$ とする。$f_4\alpha_3(x)=\beta_3f_3(x)=0$ と $f_4$ の単射性から $\alpha_3(x)=0$、よって $x=\alpha_2(y)$（$y\in A_2$）。$\beta_2f_2(y)=f_3\alpha_2(y)=0$ より $f_2(y)=\beta_1(z)$（$z\in B_1$）であり、$f_1$ の全射性から $z=f_1(w)$。すると $f_2\alpha_1(w)=\beta_1f_1(w)=f_2(y)$ で、$f_2$ の単射性から $y=\alpha_1(w)$。したがって $x=\alpha_2\alpha_1(w)=0$。

\textbf{全射性。} $z\in B_3$ とする。$f_4$ の全射性から $\beta_3(z)=f_4(u)$（$u\in A_4$）。$f_5\alpha_4(u)=\beta_4f_4(u)=\beta_4\beta_3(z)=0$ と $f_5$ の単射性から $\alpha_4(u)=0$、よって $u=\alpha_3(x)$（$x\in A_3$）。すると
\[
\beta_3\bigl(z-f_3(x)\bigr)=f_4(u)-f_4\alpha_3(x)=0
\]
なので $z-f_3(x)=\beta_2(v)$（$v\in B_2$）、さらに $v=f_2(y)$（$y\in A_2$）と書ける。よって
\[
z=f_3(x)+\beta_2f_2(y)=f_3(x)+f_3\alpha_2(y)=f_3\bigl(x+\alpha_2(y)\bigr).
\]
（実際には単射性には $f_2,f_4$ の単射性と $f_1$ の全射性、全射性には $f_2,f_4$ の全射性と $f_5$ の単射性しか使っていない。）
\end{solution}

\begin{problem}[難易度 B]
$A\subset X$ とする。$H_0(X,A)$ は、$A$ と交わらない $X$ の弧状連結成分全体を基底とする自由アーベル群に同型であることを示せ。特に $H_0(X,A)=0$ であることと、$X$ のすべての弧状連結成分が $A$ と交わることは同値である。
\end{problem}

\begin{solution}
相対鎖複体は $C_n(X,A)=C_n(X)/C_n(A)$ であり、$C_{-1}=0$ なので相対0鎖はすべてサイクルである。相対境界は $\partial C_1(X)$ の像であるから
\[
H_0(X,A)=C_0(X)\big/\bigl(C_0(A)+\partial C_1(X)\bigr)
\]
である。まず $\partial C_1(X)$ で割ると、$H_0(X)=C_0(X)/\partial C_1(X)\cong\bigoplus_\alpha\Z e_\alpha$ であり、この同型で点 $x$ は $x$ を含む弧状連結成分 $X_\alpha$ に対応する基底 $e_\alpha$ に写る。したがって
\[
H_0(X,A)\cong\Bigl(\bigoplus_\alpha\Z e_\alpha\Bigr)\Big/N
\]
となる。ここで $N$ は $C_0(A)$ の像、すなわち $a\in A$ の像 $e_\alpha$ たちで生成される部分群であり、これは $X_\alpha\cap A\ne\emptyset$ となる $\alpha$ についての $e_\alpha$ 全体で生成される。よって商は $X_\alpha\cap A=\emptyset$ となる $\alpha$ の $e_\alpha$ を基底とする自由アーベル群である。後半の主張はこれから直ちに従う。
\end{solution}

\begin{problem}[難易度 B]
$x_0\in X$ とする。すべての $n$ に対し $H_n(X,\{x_0\})\cong\tilde H_n(X)$ であることを、対の長完全列を用いて示せ。
\end{problem}

\begin{solution}
対 $(X,\{x_0\})$ の長完全列
\[
\cdots\to H_n(\{x_0\})\to H_n(X)\xrightarrow{j_*}H_n(X,\{x_0\})\xrightarrow{\partial}H_{n-1}(\{x_0\})\to\cdots
\]
を考える。$n\ge2$ では両隣の $H_n(\{x_0\})$、$H_{n-1}(\{x_0\})$ が $0$ なので $H_n(X,\{x_0\})\cong H_n(X)=\tilde H_n(X)$ である。

低次の部分は
\[
0\to H_1(X)\xrightarrow{j_*}H_1(X,\{x_0\})\xrightarrow{\partial}H_0(\{x_0\})\xrightarrow{i_*}H_0(X)\xrightarrow{j_*}H_0(X,\{x_0\})\to0
\]
である。$H_0(\{x_0\})=\Z[x_0]$ であり、$i_*[x_0]=[x_0]\in H_0(X)$ は $\bar\varepsilon[x_0]=1$ を満たすので $i_*$ は単射である（$\bar\varepsilon$ は添加写像の誘導する $H_0(X)\to\Z$）。よって $\partial=0$ となり $H_1(X,\{x_0\})\cong H_1(X)=\tilde H_1(X)$。また
\[
H_0(X,\{x_0\})\cong H_0(X)/\Z[x_0]
\]
である。簡約ホモロジーについての分裂 $H_0(X)=\tilde H_0(X)\oplus\Z[x_0]$ により、右辺は $\tilde H_0(X)$ に同型である。
\end{solution}

\begin{problem}[難易度 B]
包含写像 $\iota\colon A\hookrightarrow X$ がホモトピー同値写像ならば、すべての $n$ で $H_n(X,A)=0$ であることを示せ。これを用いて $H_n(D^m,\{0\})=0$ を確かめよ。
\end{problem}

\begin{solution}
仮定とホモトピー不変性より $\iota_*\colon H_n(A)\to H_n(X)$ はすべての $n$ で同型である。対の長完全列
\[
H_n(A)\xrightarrow{\iota_*}H_n(X)\xrightarrow{j_*}H_n(X,A)\xrightarrow{\partial}H_{n-1}(A)\xrightarrow{\iota_*}H_{n-1}(X)
\]
において、左の $\iota_*$ が全射なので $\ker j_*=\operatorname{im}\iota_*=H_n(X)$ となり $j_*=0$ である。右の $\iota_*$ が単射なので $\operatorname{im}\partial=\ker\iota_*=0$、すなわち $\partial=0$ である。$H_n(X,A)$ における完全性から
\[
H_n(X,A)=\ker\partial=\operatorname{im}j_*=0
\]
を得る。

$D^m$ は原点 $\{0\}$ に可縮であり、包含 $\{0\}\hookrightarrow D^m$ はホモトピー同値（ホモトピー逆は定値写像 $D^m\to\{0\}$、$D^m$ 上の恒等写像との間のホモトピーは $H(x,t)=tx$）である。よって $H_n(D^m,\{0\})=0$ である。
\end{solution}

\begin{problem}[難易度 B]
$A\subset X$ がレトラクトである、すなわち $r\circ\iota=\id_A$ となる $r\colon X\to A$ が存在するとする（$\iota$ は包含）。次を示せ。
(1) $\iota_*\colon H_n(A)\to H_n(X)$ は単射、$r_*\colon H_n(X)\to H_n(A)$ は全射である。
(2) すべての $n$ で $H_n(X)\cong H_n(A)\oplus H_n(X,A)$ である。
\end{problem}

\begin{solution}
(1) 関手性より $r_*\circ\iota_*=(r\circ\iota)_*=\id$ である。よって $\iota_*$ は単射、$r_*$ は全射である。

(2) 対の長完全列
\[
\cdots\to H_n(A)\xrightarrow{\iota_*}H_n(X)\xrightarrow{j_*}H_n(X,A)\xrightarrow{\partial}H_{n-1}(A)\xrightarrow{\iota_*}H_{n-1}(X)\to\cdots
\]
において、すべての $\iota_*$ が単射なので $\operatorname{im}\partial=\ker\iota_*=0$、すなわち $\partial=0$ である。よって短完全列
\[
0\to H_n(A)\xrightarrow{\iota_*}H_n(X)\xrightarrow{j_*}H_n(X,A)\to0
\]
を得る。$\Phi=(r_*,j_*)\colon H_n(X)\to H_n(A)\oplus H_n(X,A)$ が同型であることを示す。

単射性：$r_*x=0$、$j_*x=0$ とすると $x=\iota_*a$ と書け、$a=r_*\iota_*a=r_*x=0$、よって $x=0$。

全射性：$(a,z)$ を与える。$j_*$ は全射なので $z=j_*x'$ となる $x'$ がある。$x=x'-\iota_*r_*x'+\iota_*a$ とおくと、$j_*\iota_*=0$ より $j_*x=z$、また $r_*\iota_*=\id$ より $r_*x=r_*x'-r_*x'+a=a$ である。
\end{solution}

\begin{problem}[難易度 B]
$X$ をハウスドルフ空間、$x\in X$、$U$ を $x$ の開近傍とする。切除定理を用いて、包含写像が同型
\[
H_n(U,U\setminus\{x\})\xrightarrow{\ \cong\ }H_n(X,X\setminus\{x\})
\]
を誘導することを示せ。特に $U\subset\R^m$ が開集合で $x\in U$ ならば $H_n(U,U\setminus\{x\})\cong H_n(\R^m,\R^m\setminus\{x\})$ である。
\end{problem}

\begin{solution}
切除定理は次の主張である：$Z\subset A\subset X$ で $Z$ の閉包が $A$ の内部に含まれるならば、包含写像 $(X\setminus Z,A\setminus Z)\hookrightarrow(X,A)$ はすべての $n$ でホモロジーの同型を誘導する。

$A=X\setminus\{x\}$、$Z=X\setminus U$ とおく。$x\in U$ より $Z\subset A$ である。$U$ は開集合なので $Z$ は閉集合で $\overline Z=Z$ である。$X$ はハウスドルフなので一点集合 $\{x\}$ は閉集合であり、$A$ は開集合、よって $A$ の内部は $A$ 自身である。したがって $\overline Z=Z\subset A=\operatorname{int}A$ が成り立ち、切除定理が使える。さらに
\[
X\setminus Z=U,\qquad A\setminus Z=U\setminus\{x\}
\]
であるから、主張の同型が得られる。

後半は、ハウスドルフ空間 $X=\R^m$ と $x$ の開近傍 $U$ に前半を適用すれば直ちに従う。
\end{solution}

\begin{problem}[難易度 C]
$n\ge0$ に対し $\tilde H_k(S^n)\cong\Z$（$k=n$）、$\tilde H_k(S^n)=0$（$k\ne n$）であることを、マイヤー・ヴィートリス完全列を用いて $n$ に関する帰納法で示せ。
\end{problem}

\begin{solution}
$n=0$ のとき $S^0$ は2点からなり、$H_0(S^0)\cong\Z^2$、$H_k(S^0)=0$（$k\ge1$）である（弧状連結成分への分解と一点のホモロジー）。$H_0\cong\tilde H_0\oplus\Z$ より $\tilde H_0(S^0)\cong\Z$ である。

$n\ge1$ とし、$n-1$ で主張が成り立つとする。北極 $N$、南極 $S$ に対し $U=S^n\setminus\{N\}$、$V=S^n\setminus\{S\}$ とおく。
\begin{center}
\begin{tikzpicture}[scale=0.9]
  \foreach \xs/\nm/\hn/\hs in {0/U/1/0, 3.6/V/0/1, 7.2/U\cap V/1/1}{
    \begin{scope}[xshift=\xs cm]
      \fill[blue!8] (0,0) circle (1.2);
      \draw (0,0) circle (1.2);
      \draw[dashed,gray] (1.2,0) arc (0:180:1.2 and 0.3);
      \draw (-1.2,0) arc (180:360:1.2 and 0.3);
      \ifnum\hn=1 \node[hvtx] at (0,1.2) {}; \else \node[vtx] at (0,1.2) {}; \fi
      \ifnum\hs=1 \node[hvtx] at (0,-1.2) {}; \else \node[vtx] at (0,-1.2) {}; \fi
      \node[above] at (0,1.25) {\footnotesize $N$};
      \node[below] at (0,-1.25) {\footnotesize $S$};
      \node at (0,-2.0) {$\nm$};
    \end{scope}
  }
  \begin{scope}[xshift=7.2cm]
    \draw[very thick,red!70!black] (-1.2,0) arc (180:360:1.2 and 0.3);
    \foreach \x in {-0.6,0,0.6} {
      \draw[-{Stealth[length=1.6mm]},gray] (\x,0.9) -- (\x,0.4);
      \draw[-{Stealth[length=1.6mm]},gray] (\x,-0.9) -- (\x,-0.4);
    }
  \end{scope}
\end{tikzpicture}
\zucap{$U$, $V$ はそれぞれ1点を除いた球面で $\R^n$ と同相（可縮）。$U\cap V$ は両極を除いた球面で、緯線に沿って赤道 $S^{n-1}$（赤）に変位レトラクトする。}
\end{center}これらは開集合で $U\cup V=S^n$、立体射影により $U,V\cong\R^n$ は可縮である。また $U\cap V\cong\R^n\setminus\{0\}$ であり、これは $x\mapsto x/\lvert x\rvert$ により $S^{n-1}$ に変位レトラクトするので $U\cap V\simeq S^{n-1}$ である。$U\cap V\ne\emptyset$ なので簡約ホモロジーのマイヤー・ヴィートリス完全列
\[
\tilde H_k(U)\oplus\tilde H_k(V)\to\tilde H_k(S^n)\xrightarrow{\ \partial\ }\tilde H_{k-1}(U\cap V)\to\tilde H_{k-1}(U)\oplus\tilde H_{k-1}(V)
\]
が成り立つ。両端は可縮性より $0$ なので、すべての $k$ で
\[
\tilde H_k(S^n)\cong\tilde H_{k-1}(U\cap V)\cong\tilde H_{k-1}(S^{n-1})
\]
である。帰納法の仮定より右辺は $k-1=n-1$、すなわち $k=n$ のとき $\Z$、それ以外で $0$ である。
\end{solution}

\begin{problem}[難易度 B]
$n\ge1$ とする。次を計算せよ。
(1) $H_k(D^n,S^{n-1})$。また $(D^n,S^{n-1})$ が良い対であり $D^n/S^{n-1}\cong S^n$ であることを用いて、結果が $\tilde H_k(S^n)$ と一致することを確かめよ。
(2) $S^{n-1}$ を $S^n$ の赤道とするときの $H_k(S^n,S^{n-1})$。
\end{problem}

\begin{solution}
空でない $A\subset X$ に対し、簡約ホモロジーによる対の長完全列
\[
\cdots\to\tilde H_k(A)\to\tilde H_k(X)\to H_k(X,A)\xrightarrow{\partial}\tilde H_{k-1}(A)\to\tilde H_{k-1}(X)\to\cdots
\]
が成り立つことを用いる。

(1) $D^n$ は可縮なので $\tilde H_*(D^n)=0$ であり、$H_k(D^n,S^{n-1})\cong\tilde H_{k-1}(S^{n-1})$ である。前問より、これは $k=n$ のとき $\Z$、それ以外で $0$ である。一方、良い対 $(X,A)$ では商写像が同型 $H_k(X,A)\cong\tilde H_k(X/A)$ を誘導するので、$H_k(D^n,S^{n-1})\cong\tilde H_k(S^n)$ であり、これも $k=n$ で $\Z$、それ以外で $0$ である。両者は一致する。

(2) 長完全列
\[
\tilde H_k(S^{n-1})\to\tilde H_k(S^n)\to H_k(S^n,S^{n-1})\to\tilde H_{k-1}(S^{n-1})\to\tilde H_{k-1}(S^n)
\]
を考える。$k\ne n$ のとき $\tilde H_k(S^n)=0$、$\tilde H_{k-1}(S^{n-1})=0$ なので $H_k(S^n,S^{n-1})=0$ である。$k=n$ のときは $\tilde H_n(S^{n-1})=0$、$\tilde H_{n-1}(S^n)=0$ なので短完全列
\[
0\to\Z\to H_n(S^n,S^{n-1})\to\Z\to0
\]
を得る。$\Z$ は自由なのでこれは分裂し、$H_n(S^n,S^{n-1})\cong\Z^2$ である。
\end{solution}

\begin{problem}[難易度 C]
次を示せ。
(1) $m\ne n$ ならば $S^m$ と $S^n$ はホモトピー同値でない。また $S^n$ は可縮でない。
(2) 空でない開集合 $U\subset\R^m$ と $V\subset\R^n$（$m,n\ge1$）が同相ならば $m=n$ である。特に $m\ne n$ なら $\R^m$ と $\R^n$ は同相でない。
\end{problem}

\begin{solution}
(1) ホモトピー同値な空間のホモロジーは同型であり、簡約ホモロジーも同型である（$H_0\cong\tilde H_0\oplus\Z$ の分解は誘導写像と両立する）。$\tilde H_m(S^m)\cong\Z$ だが $m\ne n$ なら $\tilde H_m(S^n)=0$ なので、$S^m\not\simeq S^n$。可縮空間では $\tilde H_*=0$ だが $\tilde H_n(S^n)\ne0$ なので $S^n$ は可縮でない。

(2) $x\in\R^m$ とする。$\R^m$ は可縮、$\R^m\setminus\{x\}\simeq S^{m-1}$ なので、簡約ホモロジーの対の長完全列から
\[
H_k(\R^m,\R^m\setminus\{x\})\cong\tilde H_{k-1}(\R^m\setminus\{x\})\cong\tilde H_{k-1}(S^{m-1})
\]
であり、これは $k=m$ のとき $\Z$、それ以外で $0$ である。切除定理による局所ホモロジーの同型より、$x\in U$ に対し $H_k(U,U\setminus\{x\})$ も同じである。

同相写像 $h\colon U\to V$ と $x\in U$ に対し、$h$ は対の同相 $(U,U\setminus\{x\})\to(V,V\setminus\{h(x)\})$ を与えるので
\[
H_k(U,U\setminus\{x\})\cong H_k(V,V\setminus\{h(x)\})
\]
である。左辺は $k=m$ でのみ非零、右辺は $k=n$ でのみ非零であるから $m=n$ である。
\end{solution}

\begin{problem}[難易度 B]
$n\ge1$ とする。$S^{n-1}$ は $D^n$ のレトラクトでないことを示せ。これを用いて、任意の連続写像 $f\colon D^n\to D^n$ が不動点をもつこと（ブラウワーの不動点定理）を示せ。
\end{problem}

\begin{solution}
レトラクション $r\colon D^n\to S^{n-1}$（$r\circ\iota=\id$）が存在したとする。簡約ホモロジー $\tilde H_{n-1}$ をとると
\[
\tilde H_{n-1}(S^{n-1})\xrightarrow{\iota_*}\tilde H_{n-1}(D^n)\xrightarrow{r_*}\tilde H_{n-1}(S^{n-1})
\]
の合成は恒等写像である。しかし左右は $\Z$、中央は $D^n$ の可縮性から $0$ であり、恒等写像 $\Z\to\Z$ が $0$ を経由することになり矛盾する（$n=1$ でも $\tilde H_0(S^0)\cong\Z$ なので同じ議論が成り立つ）。

次に $f$ が不動点をもたないとする。各 $x$ で $u=x-f(x)\ne0$ であり、$f(x)$ から $x$ へ向かう半直線と $S^{n-1}$ の交点を $r(x)=x+tu$（$t\ge0$）とする。$\lvert x+tu\rvert^2=1$ を解いて
\[
t=\frac{-x\cdot u+\sqrt{(x\cdot u)^2+\lvert u\rvert^2(1-\lvert x\rvert^2)}}{\lvert u\rvert^2}
\]
であり、根号内は非負で $t\ge0$、また $t$ は $x$ について連続なので $r$ は連続である。$\lvert x\rvert=1$ ならば $x\cdot u=1-x\cdot f(x)\ge0$ より $t=0$、すなわち $r(x)=x$ である。よって $r$ はレトラクション $D^n\to S^{n-1}$ となり、前半に矛盾する。
\end{solution}

\begin{problem}[難易度 C]
空でない空間 $X$ の懸垂を $SX=X\times[0,1]/{\sim}$（$X\times\{0\}$ と $X\times\{1\}$ をそれぞれ一点に潰したもの）とする。マイヤー・ヴィートリス完全列を用いて、すべての $n$ で $\tilde H_{n+1}(SX)\cong\tilde H_n(X)$ であることを示せ。
\end{problem}

\begin{solution}
商写像を $q\colon X\times[0,1]\to SX$ とし、$U=q(X\times[0,2/3))$、$V=q(X\times(1/3,1])$ とおく。$X\times[0,2/3)$ は $q$ に関して飽和な開集合なので $U$ は開集合であり、$V$ も同様である。$U\cup V=SX$ である。

$U$ は錐 $X\times[0,2/3)/X\times\{0\}$ であり、$([x,s],\tau)\mapsto[x,(1-\tau)s]$ により頂点に可縮である（$[0,1]$ が局所コンパクト・ハウスドルフなので $q\times\id_{[0,1]}$ も商写像となり、このホモトピーは連続である）。$V$ も同様に可縮である。また $U\cap V=q(X\times(1/3,2/3))\cong X\times(1/3,2/3)\simeq X$ である。

$X\ne\emptyset$ より $U\cap V\ne\emptyset$ なので、簡約ホモロジーのマイヤー・ヴィートリス完全列
\[
\tilde H_{n+1}(U)\oplus\tilde H_{n+1}(V)\to\tilde H_{n+1}(SX)\xrightarrow{\partial}\tilde H_n(U\cap V)\to\tilde H_n(U)\oplus\tilde H_n(V)
\]
が使える。両端は $0$ なので $\partial$ は同型であり、
\[
\tilde H_{n+1}(SX)\cong\tilde H_n(U\cap V)\cong\tilde H_n(X)
\]
を得る。（$SS^{n}\cong S^{n+1}$ であるから、これは球面のホモロジーの別証明も与える。）
\end{solution}

\begin{problem}[難易度 D]
$x_0\in X$ とし、道 $f\colon[0,1]\to X$ を特異1単体 $\sigma((1-s)v_0+sv_1)=f(s)$ と同一視する。次を示せ。
(1) $f(1)=g(0)$ となる道 $f,g$ に対し、$f\cdot g-f-g$ は境界である（$f\cdot g$ は道の積）。
(2) 端点を止めてホモトピックな道 $f\simeq g$ に対し、$f-g$ は境界である。
(3) $h([f])=[f]$ により well-defined な準同型 $h\colon\pi_1(X,x_0)\to H_1(X)$ が定まる（フレヴィッチ準同型）。
\end{problem}

\begin{solution}
(1) $\sigma\colon\Delta^2\to X$ を重心座標で $\sigma(t_0,t_1,t_2)=(f\cdot g)(t_1/2+t_2)$ と定める。第2面（$t_2=0$、$t_1=s$）では $(f\cdot g)(s/2)=f(s)$、第0面（$t_0=0$、$t_1=1-s$、$t_2=s$）では $(f\cdot g)((1+s)/2)=g(s)$、第1面（$t_1=0$、$t_2=s$）では $(f\cdot g)(s)$ である。よって $\partial\sigma=g-f\cdot g+f$。

(2) まず端点を止めて定値道 $c_x$ にホモトピックなループ $\ell$ が境界であることを示す。ホモトピー $H\colon[0,1]^2\to X$ を $H(s,0)=\ell(s)$、$H(s,1)=H(0,t)=H(1,t)=x$ とし、$\tau(t_0,t_1,t_2)=H\bigl(t_1/(t_0+t_1),\,t_2\bigr)$（$t_2<1$）、$\tau(v_2)=x$ と定める。$H$ の一様連続性と $H(\cdot,1)\equiv x$ から $\tau$ は $v_2$ でも連続である。第2面は $\ell$、第0面・第1面は $c_x$ なので $\partial\tau=c_x-c_x+\ell=\ell$。
次に $f\simeq g$ なら $f\cdot\bar g\simeq c_{f(0)}$（端点を止めて）なので $f\cdot\bar g$ は境界である。(1) と、$g+\bar g$ が境界であることより、ホモロジー類として
\[
0=[f\cdot\bar g]=[f]+[\bar g]=[f]-[g].
\]

(3) ループ $f$ は $\partial f=x_0-x_0=0$ よりサイクルである。(2) より $[f]$ はホモトピー類のみで決まり、(1) より $h([f][g])=[f\cdot g]=[f]+[g]$ である。
\end{solution}

\begin{problem}[難易度 B]
フレヴィッチの定理（$X$ が弧状連結ならば、フレヴィッチ準同型 $h$ は全射でその核は交換子部分群であり、$H_1(X)\cong\pi_1(X,x_0)^{\mathrm{ab}}$）を用いて、次の空間の $H_1$ を計算せよ：(1) $S^1\vee S^1$、(2) トーラス $T^2$、(3) 実射影平面 $\R P^2$、(4) クラインの壺 $K$。
\end{problem}

\begin{solution}
いずれも弧状連結なので、基本群のアーベル化を計算すればよい。アーベル化は、生成元と関係式による表示 $\langle S\mid R\rangle$ に対し、$S$ を基底とする自由アーベル群を $R$ を加法的に書いたもので割ったものである。

(1) $\pi_1(S^1\vee S^1)=\langle a,b\mid\ \rangle$（自由群）なので $H_1\cong\Z^2$。

(2) $\pi_1(T^2)=\langle a,b\mid aba^{-1}b^{-1}\rangle\cong\Z^2$ はすでにアーベル群なので $H_1(T^2)\cong\Z^2$。

(3) $\pi_1(\R P^2)\cong\Z/2$ はアーベル群なので $H_1(\R P^2)\cong\Z/2$。

(4) $\pi_1(K)=\langle a,b\mid abab^{-1}\rangle$ である。関係式を加法的に書くと $a+b+a-b=2a=0$ となるので
\[
H_1(K)\cong(\Z a\oplus\Z b)/\langle 2a\rangle\cong\Z/2\oplus\Z.
\]
特に $H_1$ はトーラスとクラインの壺を区別する。
\end{solution}

%================================================
\chapter{ホモロジーの計算と応用}
% 20問
%================================================

\begin{summary}
\begin{itemize}
  \item \textbf{写像度}：$f\colon S^n\to S^n$ に対し、$f_*\colon H_n(S^n)\cong\Z\to\Z$ は整数 $\deg f$ 倍。$\deg(g\circ f)=\deg g\deg f$、ホモトピー不変。鏡映は $-1$、対蹠写像は $(-1)^{n+1}$。不動点をもたない $f$ は対蹠写像とホモトピック。
  \item \textbf{局所写像度}：$y$ の逆像が有限集合 $\{x_i\}$ なら $\deg f=\sum_i\deg f|_{x_i}$。
  \item \textbf{胞体ホモロジー}：$C_n^{\mathrm{CW}}=H_n(X^n,X^{n-1})\cong\Z^{(n\text{ 胞体の数})}$。境界は $d_n(e^n_\alpha)=\sum_\beta d_{\alpha\beta}e^{n-1}_\beta$ で、$d_{\alpha\beta}$ は「接着写像 $S^{n-1}\to X^{n-1}$ と、$e^{n-1}_\beta$ 以外をつぶす写像 $X^{n-1}\to S^{n-1}$」の合成の写像度。$H_n^{\mathrm{CW}}(X)\cong H_n(X)$。
  \item \textbf{標準的なCW構造}：$S^n=e^0\cup e^n$；$\mathbb{RP}^n=e^0\cup e^1\cup\dots\cup e^n$（$d_k=1+(-1)^k$）；$\mathbb{CP}^n=e^0\cup e^2\cup\dots\cup e^{2n}$；向き付け可能な閉曲面 $\Sigma_g$ は 0胞体1個、1胞体 $2g$ 個、2胞体1個（語 $[a_1,b_1]\cdots[a_g,b_g]$）。向き付け不能な $N_h$ は語 $a_1^2\cdots a_h^2$。
  \item \textbf{オイラー標数}：$\chi(X)=\sum(-1)^n(\text{$n$ 胞体の数})=\sum(-1)^n\operatorname{rank}H_n(X)$。$n$ 葉の被覆 $\tilde X\to X$ では $\chi(\tilde X)=n\chi(X)$。
  \item \textbf{レフシェッツ数}：$L(f)=\sum_k(-1)^k\operatorname{tr}\bigl(f_*\colon H_k(X)/\text{ねじれ}\to H_k(X)/\text{ねじれ}\bigr)$。有限CW複体（およびそのレトラクト）上で $L(f)\neq0$ なら $f$ は不動点をもつ。
  \item \textbf{局所ホモロジー}：$H_k(\R^n,\R^n\setminus\{x\})\cong\tilde H_{k-1}(S^{n-1})$。次元の不変性や領域不変性に用いる。
\end{itemize}
\end{summary}

\begin{problem}[難易度 A]
$n\geq 0$ とし、連続写像 $f\colon S^n\to S^n$ の写像度 $\operatorname{deg} f$ を、$f_*\colon \tilde H_n(S^n)\cong\Z\to\tilde H_n(S^n)\cong\Z$ がかける整数倍として定める。$S^n\subset\R^{n+1}$ の超平面に関する鏡映 $r(x_0,x_1,\dots,x_n)=(-x_0,x_1,\dots,x_n)$ について $\operatorname{deg} r=-1$ を示せ。さらに、原点を通る任意の超平面に関する鏡映の写像度も $-1$ であることを示せ。
\end{problem}

\begin{solution}
$n$ に関する帰納法で示す。鏡映を $r_n\colon S^n\to S^n$ と書く。$n=0$ のとき $S^0=\{p,q\}$（$p=(1)$、$q=(-1)$）であり、$\tilde H_0(S^0)=\ker(\varepsilon\colon H_0(S^0)\to\Z)$ は $p-q$ で生成される。$r_0$ は $p$ と $q$ を入れ替えるので $(r_0)_*(p-q)=q-p$ となり、$\operatorname{deg} r_0=-1$ である。

$n\geq 1$ とし、$U=\{x\in S^n : x_n>-1/2\}$、$V=\{x\in S^n : x_n<1/2\}$ とおく。$U,V$ は開円板と同相で可縮であり、$U\cap V=\{\lvert x_n\rvert<1/2\}$ は赤道 $S^{n-1}=\{x_n=0\}$ を変位レトラクトとしてもつ。被約ホモロジーのマイヤー・ヴィートリス完全列から、連結準同型
\[
\partial\colon \tilde H_n(S^n)\to\tilde H_{n-1}(U\cap V)
\]
は同型である（両隣の項 $\tilde H_n(U)\oplus\tilde H_n(V)$、$\tilde H_{n-1}(U)\oplus\tilde H_{n-1}(V)$ が $0$ だから）。$r_n$ は第 $0$ 座標しか動かさないので $r_n(U)=U$、$r_n(V)=V$ であり、マイヤー・ヴィートリス完全列の自然性から $\partial\circ(r_n)_*=(r_n|_{U\cap V})_*\circ\partial$ が成り立つ。また包含 $i\colon S^{n-1}\to U\cap V$ はホモトピー同値で $r_n\circ i=i\circ r_{n-1}$ をみたす。したがって $\tilde H_n(S^n)\cong\tilde H_{n-1}(S^{n-1})$ という同型のもとで $(r_n)_*$ と $(r_{n-1})_*$ が対応し、$\operatorname{deg} r_n=\operatorname{deg} r_{n-1}=-1$ である。

一般の単位ベクトル $v$ に対し、$v^{\perp}$ に関する鏡映 $r_v(x)=x-2(x\cdot v)v$ を考える。$Ae_0=v$ となる直交行列 $A\in O(n+1)$ をとれば $r_v=A\circ r\circ A^{-1}$ である。写像度は合成に関して乗法的なので
\[
\operatorname{deg} r_v=\operatorname{deg} A\cdot\operatorname{deg} r\cdot\operatorname{deg}(A^{-1})=\operatorname{deg} r\cdot\operatorname{deg}(A\circ A^{-1})=-1
\]
となる。
\end{solution}

\begin{problem}[難易度 B]
対蹠写像 $a\colon S^n\to S^n$、$a(x)=-x$ の写像度が $(-1)^{n+1}$ であることを示せ。また、$a$ が恒等写像にホモトピックであるための必要十分条件は $n$ が奇数であることを示せ。
\end{problem}

\begin{solution}
$i=0,1,\dots,n$ に対し、第 $i$ 座標の符号を反転する鏡映を $r^{(i)}$ とすると $a=r^{(0)}\circ r^{(1)}\circ\dots\circ r^{(n)}$ である。前問より各 $r^{(i)}$ の写像度は $-1$ であり、写像度は合成に関して乗法的だから $\operatorname{deg} a=(-1)^{n+1}$ である。

写像度はホモトピー不変である。$n$ が偶数なら $\operatorname{deg} a=-1\neq 1=\operatorname{deg}\id$ なので、$a$ は $\id$ にホモトピックでない。

$n$ が奇数のとき、$n+1=2m$ として $\R^{n+1}=\mathbb{C}^m$ と同一視し、
\[
F(x,t)=e^{\pi i t}x\quad(x\in S^{2m-1}\subset\mathbb{C}^m,\ 0\leq t\leq 1)
\]
とおく。$\lvert e^{\pi i t}x\rvert=\lvert x\rvert=1$ なので $F$ は $S^n$ 上のホモトピーであり、$F(x,0)=x$、$F(x,1)=-x$ である。よって $a\simeq\id$ である。
\end{solution}

\begin{problem}[難易度 B]
$f\colon S^n\to S^n$ を連続写像とする。(1) $f$ が不動点をもたなければ $f$ は対蹠写像 $a$ にホモトピックであり、$\operatorname{deg} f=(-1)^{n+1}$ であることを示せ。(2) すべての $x$ で $f(x)\neq -x$ ならば $\operatorname{deg} f=1$ であることを示せ。(3) $n$ が偶数ならば、任意の連続写像 $f\colon S^n\to S^n$ に対し $f(x)=x$ または $f(x)=-x$ となる点 $x$ が存在することを示せ。
\end{problem}

\begin{solution}
(1) $F(x,t)=\dfrac{(1-t)f(x)-tx}{\lvert (1-t)f(x)-tx\rvert}$ とおく。分母が $0$ になるとすると $(1-t)f(x)=tx$ であり、両辺のノルムをとって $1-t=t$、すなわち $t=1/2$ を得るので $f(x)=x$ となり仮定に反する。よって $F$ は $S^n\times[0,1]$ 上連続で、$F(\cdot,0)=f$、$F(\cdot,1)=a$ を結ぶホモトピーである。前問より $\operatorname{deg} f=\operatorname{deg} a=(-1)^{n+1}$ となる。

(2) 同様に $G(x,t)=\dfrac{(1-t)f(x)+tx}{\lvert (1-t)f(x)+tx\rvert}$ を考える。分母が $0$ なら $t=1/2$ かつ $f(x)=-x$ となり仮定に反するので、$G$ は $f$ と $\id$ を結ぶホモトピーである。よって $\operatorname{deg} f=1$ である。

(3) $f(x)\neq x$ かつ $f(x)\neq -x$ がすべての $x$ で成り立つと仮定する。(1) より $\operatorname{deg} f=(-1)^{n+1}=-1$、(2) より $\operatorname{deg} f=1$ となり矛盾する。
\end{solution}

\begin{problem}[難易度 C]
$S^n\subset\R^{n+1}$ 上の連続な接ベクトル場とは、連続写像 $v\colon S^n\to\R^{n+1}$ で各 $x$ について $v(x)\cdot x=0$ をみたすものである。至るところ $0$ にならない接ベクトル場が存在するための必要十分条件は $n$ が奇数であることを示せ（毛玉の定理）。
\end{problem}

\begin{solution}
$n$ が奇数のとき、$n+1=2m$ として
\[
v(x_1,x_2,\dots,x_{2m-1},x_{2m})=(-x_2,x_1,\dots,-x_{2m},x_{2m-1})
\]
とおけば $v(x)\cdot x=0$ かつ $\lvert v(x)\rvert=\lvert x\rvert=1$ であり、至るところ $0$ でない接ベクトル場である。

逆に $v$ が至るところ $0$ でない接ベクトル場であるとし、$w(x)=v(x)/\lvert v(x)\rvert$ とおく。$w(x)$ は $x$ に直交する単位ベクトルなので、
\[
F(x,t)=\cos(\pi t)\,x+\sin(\pi t)\,w(x)
\]
は $\lvert F(x,t)\rvert^2=\cos^2(\pi t)+\sin^2(\pi t)=1$ をみたし、$S^n$ 上で $F(\cdot,0)=\id$ と $F(\cdot,1)=a$ を結ぶホモトピーとなる。よって $1=\operatorname{deg}\id=\operatorname{deg} a=(-1)^{n+1}$ であり、$n$ は奇数である。
\end{solution}

\begin{problem}[難易度 A]
$n\geq 1$ とする。任意の連続写像 $f\colon D^n\to D^n$ は不動点をもつことを示せ（ブラウアーの不動点定理）。ただし $\tilde H_{n-1}(S^{n-1})\cong\Z$ と $D^n$ の可縮性は既知としてよい。
\end{problem}

\begin{solution}
$f$ が不動点をもたないと仮定し、$u(x)=x-f(x)\neq 0$ とおく。$f(x)$ から $x$ に向かう半直線が $S^{n-1}$ と交わる点を $r(x)=x+t(x)u(x)$（$t(x)\geq 0$）とする。$\lvert x+tu\rvert^2=1$ を $t$ について解くと、$b=x\cdot u$ として
\[
t(x)=\frac{-b+\sqrt{b^2+\lvert u\rvert^2(1-\lvert x\rvert^2)}}{\lvert u\rvert^2}
\]
であり、$u\neq 0$ より $t$ は連続、したがって $r\colon D^n\to S^{n-1}$ は連続である。$x\in S^{n-1}$ なら $b=1-x\cdot f(x)\geq 0$ なので $t(x)=(-b+\lvert b\rvert)/\lvert u\rvert^2=0$、すなわち $r(x)=x$ である。よって包含 $i\colon S^{n-1}\to D^n$ に対し $r\circ i=\id$ が成り立つ。

被約ホモロジーをとると $r_*\circ i_*=\id\colon \tilde H_{n-1}(S^{n-1})\to\tilde H_{n-1}(S^{n-1})$ である。しかし $i_*$ は $\tilde H_{n-1}(D^n)=0$ を経由するので $\id=0$ となり、$\tilde H_{n-1}(S^{n-1})\cong\Z$ に矛盾する（$n=1$ でも $\tilde H_0(S^0)\cong\Z$ なので同じ議論が使える）。
\end{solution}

\begin{problem}[難易度 B]
空でない空間 $X$ の懸垂 $SX=X\times[0,1]/(X\times\{0\},\,X\times\{1\})$（$X\times\{0\}$ と $X\times\{1\}$ をそれぞれ $1$ 点につぶした空間）について、すべての $k$ で $\tilde H_{k+1}(SX)\cong\tilde H_k(X)$ であることをマイヤー・ヴィートリス完全列を用いて示せ。また、これから $\tilde H_k(S^n)$ を求めよ。
\end{problem}

\begin{solution}
$X\times\{0\}$、$X\times\{1\}$ の像をそれぞれ $s,n$ とし、$U=SX\setminus\{s\}$、$V=SX\setminus\{n\}$ とおく。$1$ 点は閉集合なので $U,V$ は開集合で $U\cup V=SX$ である。$U$ は $X\times(0,1]$ の像であり、$([x,u],\tau)\mapsto[x,u+\tau(1-u)]$ によって $n$ に変位レトラクトする（$[0,1]$ が局所コンパクトなので商写像と $[0,1]$ の直積は商写像となり、このホモトピーは連続である）。よって $U$ は可縮であり、同様に $V$ も可縮である。また $U\cap V$ は $X\times(0,1)$ と同相で $X$ とホモトピー同値である。

被約ホモロジーのマイヤー・ヴィートリス完全列（$U\cap V\neq\emptyset$）
\[
\tilde H_{k+1}(U)\oplus\tilde H_{k+1}(V)\to\tilde H_{k+1}(SX)\xrightarrow{\ \partial\ }\tilde H_k(U\cap V)\to\tilde H_k(U)\oplus\tilde H_k(V)
\]
で両端が $0$ なので $\partial$ は同型であり、$\tilde H_{k+1}(SX)\cong\tilde H_k(U\cap V)\cong\tilde H_k(X)$ を得る。

$SS^{n-1}\cong S^n$（$S^{n-1}\times[0,1]\to S^n$、$(x,u)\mapsto(\sin(\pi u)\,x,\,-\cos(\pi u))$ が誘導する同相）と $\tilde H_k(S^0)=\Z$（$k=0$）、$0$（$k\neq0$）から、帰納的に
\[
\tilde H_k(S^n)\cong\begin{cases}\Z & (k=n)\\ 0 & (k\neq n)\end{cases}
\]
である。
\end{solution}

\begin{problem}[難易度 B]
$X,Y$ を空間とし、基点 $x_0\in X$、$y_0\in Y$ がそれぞれ $\{x_0\}$、$\{y_0\}$ を変位レトラクトにもつ開近傍 $A\subset X$、$B\subset Y$ をもつとする。楔和 $X\vee Y$ について $\tilde H_k(X\vee Y)\cong\tilde H_k(X)\oplus\tilde H_k(Y)$ を示せ。これを用いて $S^1\vee S^1\vee S^2$ のホモロジー群を計算せよ。
\end{problem}

\begin{solution}
$X\vee Y$ において $U=X\vee B$、$V=A\vee Y$ とおく。$U$ の $X$ および $Y$ への逆像はそれぞれ $X$、$B$ で開集合だから $U$ は開集合であり、同様に $V$ も開集合、$U\cup V=X\vee Y$ である。$B$ を基点に変位レトラクトするホモトピーを $X$ 上では恒等的に延長すれば、$U$ は $X$ を変位レトラクトにもつ。同様に $V\simeq Y$、$U\cap V=A\vee B\simeq\{\mathrm{pt}\}$ である。

被約マイヤー・ヴィートリス完全列
\[
\tilde H_k(U\cap V)\to\tilde H_k(U)\oplus\tilde H_k(V)\to\tilde H_k(X\vee Y)\to\tilde H_{k-1}(U\cap V)
\]
で両端が $0$ だから、$\tilde H_k(X\vee Y)\cong\tilde H_k(U)\oplus\tilde H_k(V)\cong\tilde H_k(X)\oplus\tilde H_k(Y)$ である。

球面の基点は開円板を近傍にもつので仮定をみたす。$S^1\vee S^1\vee S^2=(S^1\vee S^1)\vee S^2$ に $2$ 回適用し、$H_0$ は弧状連結性から $\Z$ であることを用いると、
\[
H_0\cong\Z,\quad H_1\cong\Z^2,\quad H_2\cong\Z,\quad H_k=0\ (k\geq 3)
\]
を得る。
\end{solution}

\begin{problem}[難易度 B]
トーラス $T=S^1\times S^1$ のホモロジー群をマイヤー・ヴィートリス完全列を用いて計算せよ。
\end{problem}

\begin{solution}
$S^1$ を $2$ つの開弧 $I_1,I_2$（それぞれ可縮で、$I_1\cap I_2$ は $2$ つの開弧の非交和）で覆い、$U=I_1\times S^1$、$V=I_2\times S^1$ とおく。
\begin{center}
\begin{tikzpicture}[scale=1.1]
  \fill[blue!60,fill opacity=0.25] (0.3,0) rectangle (2.7,3);
  \fill[red!60,fill opacity=0.25] (0,0) rectangle (1.2,3);
  \fill[red!60,fill opacity=0.25] (1.8,0) rectangle (3,3);
  \draw[thick,mid arrow] (0,0) -- (3,0);
  \draw[thick,mid arrow] (0,3) -- (3,3);
  \draw[thick,mid arrows] (0,0) -- (0,3);
  \draw[thick,mid arrows] (3,0) -- (3,3);
  \draw[dashed] (0.3,0) -- (0.3,3) (2.7,0) -- (2.7,3) (1.2,0) -- (1.2,3) (1.8,0) -- (1.8,3);
  \node at (0.75,1.5) {$W_1$};
  \node at (2.25,1.5) {$W_2$};
  \draw[decorate,decoration={brace,amplitude=4pt}] (0.3,3.15) -- (2.7,3.15) node[midway,above=3pt] {$U=I_1\times S^1$};
  \draw[decorate,decoration={brace,amplitude=4pt,mirror}] (0,-0.15) -- (1.2,-0.15);
  \draw[decorate,decoration={brace,amplitude=4pt,mirror}] (1.8,-0.15) -- (3,-0.15);
  \node at (1.5,-0.6) {$V=I_2\times S^1$（左右の帯は貼り合わさって1つの帯）};
\end{tikzpicture}
\zucap{トーラスを正方形で表し、横方向を第1因子とする。青が $U$、赤が $V$、重なり（紫）が $U\cap V=W_1\sqcup W_2$。どれも縦方向の円周 $\{\mathrm{pt}\}\times S^1$ を芯とする円環である。}
\end{center}$U\simeq V\simeq S^1$ であり、$U\cap V$ は $2$ つの円環 $W_1\sqcup W_2$ で、各 $W_j\simeq S^1$ である。どの円環の $H_1$ の生成元も $\{\mathrm{pt}\}\times S^1$ のクラスであり、$H_1(U)$、$H_1(V)$ の生成元に移る。$\Phi(c)=(i_*c,-j_*c)$ という符号の規約のもとで、マイヤー・ヴィートリス完全列は
\begin{align*}
0&\to H_2(T)\to H_1(U\cap V)\cong\Z^2\xrightarrow{\ \Phi_1\ }\Z^2\to H_1(T)\\
&\to H_0(U\cap V)\cong\Z^2\xrightarrow{\ \Phi_0\ }\Z^2\to H_0(T)\to 0
\end{align*}
となり（$H_2(U)=H_2(V)=0$）、$\Phi_1(a,b)=(a+b,-(a+b))$、$\Phi_0(a,b)=(a+b,-(a+b))$ である。

$\ker\Phi_1=\{(a,-a)\}\cong\Z$ なので $H_2(T)\cong\Z$ である。$\operatorname{im}\Phi_1=\{(m,-m)\}$ であり、$(x,y)\mapsto x+y$ により $\operatorname{coker}\Phi_1\cong\Z$ となる。$\ker\Phi_0\cong\Z$ なので短完全列
\[
0\to\operatorname{coker}\Phi_1\cong\Z\to H_1(T)\to\ker\Phi_0\cong\Z\to 0
\]
を得る。右端が自由加群だからこれは分裂し $H_1(T)\cong\Z^2$ である。最後に $H_0(T)\cong\operatorname{coker}\Phi_0\cong\Z$、また $k\geq 3$ で $H_k(T)=0$ である（$U\cap V$、$U$、$V$ の $H_k$ $(k\geq 2)$ がすべて $0$ だから）。
\end{solution}

\begin{problem}[難易度 B]
$X$ を $k$ 胞体を $c_k$ 個もつ有限CW複体とし、オイラー標数を $\chi(X)=\sum_k(-1)^k c_k$ で定める。
\[
\chi(X)=\sum_k(-1)^k\operatorname{rank} H_k(X)
\]
を示せ。特に $\chi(X)$ はCW構造のとり方によらない。
\end{problem}

\begin{solution}
胞体鎖複体 $C_k=H_k(X^k,X^{k-1})\cong\Z^{c_k}$ と境界写像 $d_k\colon C_k\to C_{k-1}$ を考え、$Z_k=\ker d_k$、$B_k=\operatorname{im} d_{k+1}$ とおく。胞体ホモロジーと特異ホモロジーの同型により $H_k(X)\cong Z_k/B_k$ である。

有限生成アーベル群の短完全列に対して階数は加法的である。短完全列 $0\to Z_k\to C_k\xrightarrow{d_k} B_{k-1}\to 0$ と $0\to B_k\to Z_k\to H_k(X)\to 0$ から
\[
c_k=\operatorname{rank} Z_k+\operatorname{rank} B_{k-1},\qquad
\operatorname{rank} H_k(X)=\operatorname{rank} Z_k-\operatorname{rank} B_k
\]
を得る。したがって
\[
\sum_k(-1)^k c_k-\sum_k(-1)^k\operatorname{rank} H_k(X)=\sum_k(-1)^k\bigl(\operatorname{rank} B_{k-1}+\operatorname{rank} B_k\bigr)
\]
となり、右辺は $k$ と $k+1$ の項が打ち消し合って $0$ である（有限和）。右辺はホモロジー群だけで決まるので、$\chi(X)$ はCW構造によらない。

例えば $S^n$ $(n\geq1)$ は $0$ 胞体と $n$ 胞体を $1$ つずつもつので $\chi(S^n)=1+(-1)^n$ であり、ホモロジー $H_0\cong H_n\cong\Z$ から計算した値と一致する。
\end{solution}

\begin{problem}[難易度 D]
実射影空間 $\mathbb{RP}^n$ は各次元 $0\leq k\leq n$ にちょうど $1$ つの $k$ 胞体 $e^k$ をもち、$e^k$ の接着写像は二重被覆 $\varphi_k\colon S^{k-1}\to\mathbb{RP}^{k-1}$ である。胞体境界写像 $d_k$ が $1+(-1)^k$ 倍であることを示し、$H_*(\mathbb{RP}^n)$ を計算せよ。
\end{problem}

\begin{solution}
胞体境界公式により、$d_k(e^k)=\operatorname{deg}(q\circ\varphi_k)\,e^{k-1}$ である。ここで $q\colon\mathbb{RP}^{k-1}\to\mathbb{RP}^{k-1}/\mathbb{RP}^{k-2}\cong S^{k-1}$ は商写像である。$q\circ\varphi_k$ は赤道 $S^{k-2}$ を $1$ 点につぶし、上半球と下半球の内部をそれぞれ同相に $S^{k-1}$ から $1$ 点を除いた集合に写す。下半球上の写像は上半球上の写像に対蹠写像 $a$ を合成したものである。局所写像度の和として写像度を計算すると、上半球の寄与を $1$ となるように向きを選べば下半球の寄与は $\operatorname{deg} a=(-1)^k$（$a$ は $S^{k-1}$ 上の対蹠写像）なので、
\[
d_k=1+(-1)^k=\begin{cases}2 & (k\text{ 偶数})\\ 0 & (k\text{ 奇数})\end{cases}
\]
となる。$k=1$ でも $d_1=0$ であり、これは $0$ 胞体が $1$ つであることと整合する。

胞体鎖複体は $0\leftarrow\Z\xleftarrow{0}\Z\xleftarrow{2}\Z\xleftarrow{0}\Z\xleftarrow{2}\cdots\leftarrow\Z$（$C_k=\Z$、$0\le k\le n$）となる。
\begin{itemize}
\item $k=0$：$H_0=\Z/\operatorname{im} d_1=\Z$。
\item $0<k<n$、$k$ 奇数：$\ker d_k=\Z$、$\operatorname{im} d_{k+1}=2\Z$ より $H_k=\Z/2$。
\item $0<k\leq n$、$k$ 偶数：$d_k=2$ は単射なので $H_k=0$。
\item $k=n$ 奇数：$d_n=0$ で $C_{n+1}=0$ より $H_n=\Z$。
\end{itemize}
$k>n$ では $H_k=0$ である。
\end{solution}

\begin{problem}[難易度 B]
複素射影空間 $\mathbb{CP}^n$ が、各偶数次元 $0,2,\dots,2n$ にちょうど $1$ つずつ胞体をもつCW構造をもつことを示し、$H_*(\mathbb{CP}^n)$ と $\chi(\mathbb{CP}^n)$ を求めよ。
\end{problem}

\begin{solution}
$\mathbb{CP}^{n-1}=\{[z_0:\dots:z_n] : z_n=0\}\subset\mathbb{CP}^n$ とみなす。$D^{2n}=\{z\in\mathbb{C}^n:\lvert z\rvert\leq1\}$ に対して
\[
\Phi\colon D^{2n}\to\mathbb{CP}^n,\qquad \Phi(z)=\bigl[z_0:\dots:z_{n-1}:\sqrt{1-\lvert z\rvert^2}\bigr]
\]
とおく。$z_n\neq0$ の点は $z_n$ が正の実数で $\lvert(z_0,\dots,z_n)\rvert=1$ となる代表元をただ $1$ つもつので、$\Phi$ は開円板を $\mathbb{CP}^n\setminus\mathbb{CP}^{n-1}$ に同相に写す。また $\Phi(\partial D^{2n})\subset\mathbb{CP}^{n-1}$ であり、$\partial D^{2n}=S^{2n-1}\to\mathbb{CP}^{n-1}$ は商写像である。$\mathbb{CP}^n$ はコンパクト・ハウスドルフなので $\mathbb{CP}^n=\mathbb{CP}^{n-1}\cup_{\Phi}e^{2n}$ である。帰納的に $\mathbb{CP}^0=\mathrm{pt}$ から、各次元 $0,2,\dots,2n$ に $1$ つずつ胞体をもつCW構造を得る。

胞体鎖複体は $C_{2k}=\Z$（$0\leq k\leq n$）、奇数次では $0$ である。境界写像 $d_j\colon C_j\to C_{j-1}$ は始域か終域の一方が必ず $0$ なので零写像である。したがって
\[
H_j(\mathbb{CP}^n)\cong\begin{cases}\Z & (j=0,2,\dots,2n)\\ 0 & (\text{それ以外})\end{cases}
\]
であり、$\chi(\mathbb{CP}^n)=n+1$ である。
\end{solution}

\begin{problem}[難易度 B]
$g\geq1$ とする。種数 $g$ の向き付け可能閉曲面 $\Sigma_g$ は、$0$ 胞体 $1$ 個、$1$ 胞体 $a_1,b_1,\dots,a_g,b_g$、$2$ 胞体 $1$ 個からなり、$2$ 胞体は語 $a_1b_1a_1^{-1}b_1^{-1}\cdots a_gb_ga_g^{-1}b_g^{-1}$ に沿って接着される。胞体ホモロジーを用いて $H_*(\Sigma_g)$ と $\chi(\Sigma_g)$ を求めよ。
\end{problem}

\begin{solution}
\begin{center}
\begin{tikzpicture}[scale=1.0]
  \fill[blue!6] (247.5:1.6) \foreach \k in {1,...,7} { -- ({247.5+45*\k}:1.6) } -- cycle;
  \foreach \k/\lab/\s in {0/a_1/1,1/b_1/1,2/a_1/-1,3/b_1/-1,4/a_2/1,5/b_2/1,6/a_2/-1,7/b_2/-1}{
    \pgfmathsetmacro\angA{247.5+45*\k}
    \pgfmathsetmacro\angB{247.5+45*(\k+1)}
    \ifnum\s=1
      \draw[thick,mid arrow] (\angA:1.6) -- (\angB:1.6);
    \else
      \draw[thick,mid arrow] (\angB:1.6) -- (\angA:1.6);
    \fi
    \node at ({(\angA+\angB)/2}:1.95) {$\lab$};
    \node[vtx] at (\angA:1.6) {};
  }
\end{tikzpicture}
\qquad
\begin{tikzpicture}[scale=0.8]
  % 種数2の曲面の略図
  \draw[thick] (0,0) .. controls (0,1.3) and (2.2,1.3) .. (3,0.6) .. controls (3.8,1.3) and (6,1.3) .. (6,0)
               .. controls (6,-1.3) and (3.8,-1.3) .. (3,-0.6) .. controls (2.2,-1.3) and (0,-1.3) .. (0,0);
  \foreach \c in {1.4,4.6}{
    \draw[thick] (\c-0.55,0.08) .. controls (\c-0.2,-0.22) and (\c+0.2,-0.22) .. (\c+0.55,0.08);
    \draw[thick] (\c-0.42,-0.04) .. controls (\c-0.15,0.18) and (\c+0.15,0.18) .. (\c+0.42,-0.04);
  }
  \draw[red!70!black,thick] (1.4,0) ellipse (0.95 and 0.5);
  \draw[blue!70!black,thick] (1.4,-0.14) .. controls (1.2,-0.4) and (1.2,-0.7) .. (1.4,-0.92);
  \draw[blue!70!black,thick,dashed] (1.4,-0.14) .. controls (1.6,-0.4) and (1.6,-0.7) .. (1.4,-0.92);
  \node[red!70!black] at (1.4,0.68) {\footnotesize $a_1$};
  \node[blue!70!black] at (1.75,-0.75) {\footnotesize $b_1$};
  \draw[red!70!black,thick] (4.6,0) ellipse (0.95 and 0.5);
  \draw[blue!70!black,thick] (4.6,-0.14) .. controls (4.4,-0.4) and (4.4,-0.7) .. (4.6,-0.92);
  \draw[blue!70!black,thick,dashed] (4.6,-0.14) .. controls (4.8,-0.4) and (4.8,-0.7) .. (4.6,-0.92);
  \node[red!70!black] at (4.6,0.68) {\footnotesize $a_2$};
  \node[blue!70!black] at (4.95,-0.75) {\footnotesize $b_2$};
\end{tikzpicture}
\zucap{$g=2$ の場合。八角形の辺を語 $a_1b_1a_1^{-1}b_1^{-1}a_2b_2a_2^{-1}b_2^{-1}$ に従って貼る（左）と、8頂点はすべて1点になり種数2の曲面（右）ができる。1胞体 $a_i,b_i$ は曲面上の閉曲線になる。}
\end{center}
$0$ 胞体が $v$ ただ $1$ つなので、各 $1$ 胞体の境界は $v-v=0$ であり $d_1=0$ である。

$2$ 胞体の境界は胞体境界公式で計算する。接着写像 $S^1\to X^1$ を $X^1/(X^1\setminus a_j)\cong S^1$ に合成した写像は、語の中で $a_j^{\pm1}$ が現れるごとに円周を $\pm1$ 回まわり、他の部分は $1$ 点に写す。よって $a_j$ の係数は語における $a_j$ の指数和に等しい。各 $a_j$、$b_j$ は語の中に $+1$ と $-1$ で $1$ 回ずつ現れるので指数和は $0$ であり、$d_2=0$ となる。

胞体鎖複体 $0\to\Z\xrightarrow{0}\Z^{2g}\xrightarrow{0}\Z\to0$ の境界写像がすべて $0$ なので
\[
H_0(\Sigma_g)\cong\Z,\quad H_1(\Sigma_g)\cong\Z^{2g},\quad H_2(\Sigma_g)\cong\Z,\quad H_k(\Sigma_g)=0\ (k\geq3)
\]
である。また $\chi(\Sigma_g)=1-2g+1=2-2g$ である。
\end{solution}

\begin{problem}[難易度 C]
$h\geq1$ とする。向き付け不可能閉曲面 $N_h$ は、$0$ 胞体 $1$ 個、$1$ 胞体 $a_1,\dots,a_h$、語 $a_1a_1a_2a_2\cdots a_ha_h$ に沿って接着された $2$ 胞体 $1$ 個からなる。$H_*(N_h)$ を計算し、$N_h$ がどの $\Sigma_g$ ともホモトピー同値でないこと、また $h\neq h'$ なら $N_h$ と $N_{h'}$ がホモトピー同値でないことを示せ。
\end{problem}

\begin{solution}
前問と同様に $d_1=0$ である。語において各 $a_j$ の指数和は $2$ なので
\[
d_2(e^2)=2(a_1+a_2+\dots+a_h)
\]
である。$d_2\colon\Z\to\Z^h$ は単射だから $H_2(N_h)=\ker d_2=0$ である。

$H_1(N_h)=\Z^h/\langle 2(a_1+\dots+a_h)\rangle$ を求める。$\Z^h$ の基底を $a_1,\dots,a_{h-1},\ c=a_1+\dots+a_h$ にとり替えると（変換行列は行列式 $1$ の整数行列）、
\[
H_1(N_h)\cong\Z a_1\oplus\dots\oplus\Z a_{h-1}\oplus\Z c/\langle 2c\rangle\cong\Z^{h-1}\oplus\Z/2
\]
を得る。また $H_0(N_h)\cong\Z$ である。$h=1$ のとき $N_1=\mathbb{RP}^2$ であり、$H_1\cong\Z/2$ は $\mathbb{RP}^n$ の胞体ホモロジーの計算結果と一致する。

ホモトピー同値な空間のホモロジー群は同型である。$H_2(\Sigma_g)\cong\Z$（$g=0$ のときは $\Sigma_0=S^2$）、$H_2(N_h)=0$ なので $N_h\not\simeq\Sigma_g$ である。また $H_1(N_h)$ の階数は $h-1$ なので、$h\neq h'$ なら $N_h\not\simeq N_{h'}$ である。
\end{solution}

\begin{problem}[難易度 C]
$p\colon\tilde X\to X$ を $n$ 重被覆、$X$ を有限CW複体とする。$\tilde X$ は各 $k$ について $k$ 胞体を $X$ の $n$ 倍の個数もつCW構造をもち、$\chi(\tilde X)=n\,\chi(X)$ となることを示せ。これを用いて、$g,h\geq 2$ のとき $\Sigma_g$ が $\Sigma_h$ の $n$ 重被覆空間であるならば $g-1=n(h-1)$ であることを示せ。
\end{problem}

\begin{solution}
$X$ の各胞体の特性写像 $\Phi\colon D^k\to X$ を考える。$D^k$ は単連結かつ局所弧状連結なので、持ち上げの基準により、$\Phi(0)$ の上の $n$ 個の点 $\tilde x_1,\dots,\tilde x_n$ のそれぞれについて $\tilde\Phi_i(0)=\tilde x_i$ となる持ち上げ $\tilde\Phi_i\colon D^k\to\tilde X$ が一意に存在する。開胞体 $e=\Phi(\operatorname{int} D^k)$ は単連結なので $p^{-1}(e)$ は $n$ 個の開集合 $\tilde\Phi_i(\operatorname{int}D^k)$ の非交和で、それぞれ $p$ により $e$ に同相に写る。$\tilde\Phi_i(\partial D^k)$ は $p^{-1}(X^{k-1})$ に含まれる。こうして $\tilde X$ は $\tilde\Phi_i$ を特性写像とする胞体分割をもち、閉包有限性と弱位相の条件も $p$ が局所同相であることから確かめられる。よって $\tilde X$ の $k$ 胞体の個数は $n c_k$ であり、
\[
\chi(\tilde X)=\sum_k(-1)^k n c_k=n\,\chi(X)
\]
である。

$\Sigma_g\to\Sigma_h$ が $n$ 重被覆なら、前々問の $\chi(\Sigma_g)=2-2g$ より $2-2g=n(2-2h)$、すなわち $g-1=n(h-1)$ を得る。例えば $\Sigma_3$ は $\Sigma_2$ の $2$ 重被覆になりうるが、$\Sigma_4$ は $\Sigma_3$ の被覆空間にはならない（$3=n\cdot2$ をみたす整数 $n$ はない）。
\end{solution}

\begin{problem}[難易度 C]
正六角形の辺に、反時計回りに語を読んでラベルと向きを付け、同じラベルの辺を矢印の向きがそろうように貼り合わせて得られる閉曲面を考える。
\begin{center}
\begin{tikzpicture}[scale=1.0]
  \foreach \xs/\word in {0/{0/a/1,1/b/1,2/c/1,3/a/-1,4/b/-1,5/c/-1}, 5.5/{0/a/1,1/b/1,2/c/1,3/a/1,4/b/1,5/c/1}}{
    \begin{scope}[xshift=\xs cm]
      \fill[blue!6] (240:1.5) \foreach \k in {1,...,5} { -- ({240+60*\k}:1.5) } -- cycle;
      \foreach \k/\lab/\s in \word {
        \pgfmathsetmacro\angA{240+60*\k}
        \pgfmathsetmacro\angB{240+60*(\k+1)}
        \ifnum\s=1
          \draw[thick,mid arrow] (\angA:1.5) -- (\angB:1.5);
        \else
          \draw[thick,mid arrow] (\angB:1.5) -- (\angA:1.5);
        \fi
        \node at ({(\angA+\angB)/2}:1.62) [anchor={(\angA+\angB)/2+180}] {$\lab$};
      }
    \end{scope}
  }
  \node at (0,-2.4) {(1) $abca^{-1}b^{-1}c^{-1}$};
  \node at (5.5,-2.4) {(2) $abcabc$};
\end{tikzpicture}
\end{center}
それぞれについて、頂点の同値類を求めて胞体ホモロジー $H_*$ とオイラー標数を計算し、閉曲面の分類定理（閉曲面は $\Sigma_g$ か $N_h$ のいずれかに同相）を用いて、どの曲面かを決定せよ。
\end{problem}

\begin{solution}
六角形の頂点を、下の辺の左端から反時計回りに $P_0,P_1,\dots,P_5$ とし、$k$ 番目の辺を $P_k$ から $P_{k+1}$（$P_6=P_0$）への辺とする。ラベル $x$ の辺が指数 $+1$ ならその辺の $x$ は $P_k\to P_{k+1}$ 向き、指数 $-1$ なら $P_{k+1}\to P_k$ 向きである。同じラベルの辺を貼ると、始点どうし・終点どうしが同一視される。

(1) $a$：$P_0\to P_1$ と $P_4\to P_3$ より $P_0\sim P_4$、$P_1\sim P_3$。$b$：$P_1\to P_2$ と $P_5\to P_4$ より $P_1\sim P_5$、$P_2\sim P_4$。$c$：$P_2\to P_3$ と $P_0\to P_5$ より $P_2\sim P_0$、$P_3\sim P_5$。よって頂点は
\[
v=\{P_0,P_2,P_4\},\qquad w=\{P_1,P_3,P_5\}
\]
の2つである。
\begin{center}
\begin{tikzpicture}[scale=0.9]
  \fill[blue!6] (240:1.5) \foreach \k in {1,...,5} { -- ({240+60*\k}:1.5) } -- cycle;
  \foreach \k/\lab/\s in {0/a/1,1/b/1,2/c/1,3/a/-1,4/b/-1,5/c/-1} {
    \pgfmathsetmacro\angA{240+60*\k}
    \pgfmathsetmacro\angB{240+60*(\k+1)}
    \ifnum\s=1 \draw[thick,mid arrow] (\angA:1.5) -- (\angB:1.5);
    \else \draw[thick,mid arrow] (\angB:1.5) -- (\angA:1.5); \fi
    \node at ({(\angA+\angB)/2}:1.62) [anchor={(\angA+\angB)/2+180}] {$\lab$};
  }
  \foreach \k in {0,2,4} \node[vtx,inner sep=1.8pt] at ({240+60*\k}:1.5) {};
  \foreach \k in {1,3,5} \node[hvtx,inner sep=1.8pt] at ({240+60*\k}:1.5) {};
  \node[left] at (240:1.6) {\footnotesize $P_0$};
  \node[right] at (300:1.6) {\footnotesize $P_1$};
\end{tikzpicture}
\qquad
\begin{tikzpicture}[scale=0.9]
  \fill[blue!6] (240:1.5) \foreach \k in {1,...,5} { -- ({240+60*\k}:1.5) } -- cycle;
  \foreach \k/\lab in {0/a,1/b,2/c,3/a,4/b,5/c} {
    \pgfmathsetmacro\angA{240+60*\k}
    \pgfmathsetmacro\angB{240+60*(\k+1)}
    \draw[thick,mid arrow] (\angA:1.5) -- (\angB:1.5);
    \node at ({(\angA+\angB)/2}:1.62) [anchor={(\angA+\angB)/2+180}] {$\lab$};
  }
  \foreach \k in {0,3} \node[vtx,inner sep=1.8pt] at ({240+60*\k}:1.5) {};
  \foreach \k in {1,4} \node[hvtx,inner sep=1.8pt] at ({240+60*\k}:1.5) {};
  \foreach \k in {2,5} \node[draw,fill=red!60,rectangle,inner sep=2pt] at ({240+60*\k}:1.5) {};
  \node[left] at (240:1.6) {\footnotesize $P_0$};
  \node[right] at (300:1.6) {\footnotesize $P_1$};
\end{tikzpicture}
\zucap{同一視される頂点を同じ記号で示した。左（(1)）は2類、右（(2)）は3類。}
\end{center}
胞体鎖複体は $C_2=\Z e$、$C_1=\Z a\oplus\Z b\oplus\Z c$、$C_0=\Z v\oplus\Z w$。$a$ は $P_0$（$v$）から $P_1$（$w$）へ、$b$ は $P_1$（$w$）から $P_2$（$v$）へ、$c$ は $P_2$（$v$）から $P_3$（$w$）へ向かうので
\[
d_1a=w-v,\qquad d_1b=v-w,\qquad d_1c=w-v .
\]
$d_2e$ における各ラベルの係数は語の中での指数和なので $d_2e=a+b+c-a-b-c=0$。よって $H_2\cong\Z$。$\operatorname{im}d_1=\Z(w-v)$ なので $H_0\cong\Z$、$\ker d_1$ は階数 $3-1=2$ で、$\operatorname{im}d_2=0$ より $H_1\cong\Z^2$（基底 $a+b$, $b+c$）。オイラー標数は $\chi=2-3+1=0$。$H_2\cong\Z$ より向き付け可能で、$\chi(\Sigma_g)=2-2g=0$ から $g=1$。よって (1) は\textbf{トーラス}である。

(2) $a$：$P_0\to P_1$ と $P_3\to P_4$ より $P_0\sim P_3$、$P_1\sim P_4$。$b$：$P_1\to P_2$ と $P_4\to P_5$ より $P_1\sim P_4$、$P_2\sim P_5$。$c$：$P_2\to P_3$ と $P_5\to P_0$ より $P_2\sim P_5$、$P_3\sim P_0$。よって頂点は $u_0=\{P_0,P_3\}$、$u_1=\{P_1,P_4\}$、$u_2=\{P_2,P_5\}$ の3つで、
\[
d_1a=u_1-u_0,\qquad d_1b=u_2-u_1,\qquad d_1c=u_0-u_2,\qquad d_2e=2a+2b+2c .
\]
$d_2$ は単射なので $H_2=0$。$\operatorname{im}d_1$ は $u_1-u_0,u_2-u_1$ で生成される階数2の群で $H_0\cong\Z$。$\ker d_1=\Z(a+b+c)$（$d_1(xa+yb+zc)=(z-x)u_0+(x-y)u_1+(y-z)u_2$ が $0$ となるのは $x=y=z$ のとき）なので
\[
H_1=\Z(a+b+c)/\Z\,2(a+b+c)\cong\Z/2 .
\]
$\chi=3-3+1=1$。$H_2=0$ より向き付け不可能で、$\chi(N_h)=2-h=1$ から $h=1$。よって (2) は\textbf{実射影平面} $\mathbb{RP}^2$ である。

\begin{note}
頂点の同一視は、図で「ある頂点から出発し、隣の辺を貼り合わせ先へ移る」ことを繰り返して追跡するとよい。(2) は六角形の中心に関する対蹠点を同一視しているので、円板の境界の対蹠点を同一視する $\mathbb{RP}^2$ の標準的な構成 $D^2/(x\sim-x,\ x\in\partial D^2)$ そのものである。
\end{note}
\end{solution}

\begin{problem}[難易度 B]
2次元球面 $S^2\subset\R^3$ に、北極 $N$ と南極 $S$ を結ぶ直径（線分）$D$ を付け加えた空間 $X=S^2\cup D$ を考える。
\begin{center}
\begin{tikzpicture}[scale=1.0]
  \fill[blue!6] (0,0) circle (1.3);
  \draw (0,0) circle (1.3);
  \draw[dashed,gray] (1.3,0) arc (0:180:1.3 and 0.32);
  \draw (-1.3,0) arc (180:360:1.3 and 0.32);
  \draw[very thick,red!70!black] (0,1.3) -- (0,-1.3);
  \node[vtx] at (0,1.3) {}; \node[above] at (0,1.3) {$N$};
  \node[vtx] at (0,-1.3) {}; \node[below] at (0,-1.3) {$S$};
  \node[red!70!black,right] at (0,0.45) {$D$};
\end{tikzpicture}
\end{center}
$X$ が $S^2\vee S^1$ とホモトピー同値であることを示し、$H_*(X)$ を求めよ。ただし「CW対 $(X,A)$ で $A$ が可縮ならば $X\to X/A$ はホモトピー同値」であること、および楔和のホモロジー $\tilde H_n(Y\vee Z)\cong\tilde H_n(Y)\oplus\tilde H_n(Z)$ は用いてよい。
\end{problem}

\begin{solution}
$S^2$ 上で $N$ と $S$ を結ぶ経線（大円の半分）$A$ をとる。$X$ は $0$ 胞体 $N,S$、$1$ 胞体 $D, A$ と、もう1本の経線 $A'$ を加えて $S^2$ を2つの $2$ 胞体に分けたCW複体とみなせ、$A$（に端点を含めたもの）は可縮な部分複体である。
\begin{center}
\begin{tikzpicture}[scale=0.85]
  \fill[blue!6] (0,0) circle (1.3);
  \draw (0,0) circle (1.3);
  \draw[very thick,green!50!black] (0,1.3) arc (90:270:0.55 and 1.3);
  \draw[very thick,red!70!black] (0,1.3) -- (0,-1.3);
  \node[vtx] at (0,1.3) {}; \node[vtx] at (0,-1.3) {};
  \node[green!50!black] at (-0.85,0) {$A$};
  \node[red!70!black,right] at (0,0.4) {$D$};
  \draw[-{Stealth[length=2.2mm]}] (1.8,0) -- (3.2,0) node[midway,above] {\footnotesize $A$ をつぶす};
  \begin{scope}[xshift=5cm]
    \fill[blue!6] (0,0.3) circle (1.0);
    \draw (0,0.3) circle (1.0);
    \draw[very thick,red!70!black] (0,-0.7) .. controls (-1.6,-2.2) and (1.6,-2.2) .. (0,-0.7);
    \node[vtx] at (0,-0.7) {};
    \node[red!70!black] at (1.0,-1.6) {$D/\partial D\cong S^1$};
    \node at (1.4,0.9) {$S^2/A\cong S^2$};
  \end{scope}
\end{tikzpicture}
\zucap{経線 $A$（緑）を1点につぶすと、球面は球面のまま、直径 $D$ は両端が1点になって円周になる。}
\end{center}
よって $X\simeq X/A$ である。$S^2/A$ は $S^2$ と同相である（$A$ は球面上の弧で、たとえば $S^2\setminus A$ は開円板と同相であり、$S^2/A$ はその1点コンパクト化となるので $S^2$）。一方 $D$ の像は、両端点 $N,S$ がともに $A$ の像である1点に写るので $D/\{N,S\}\cong S^1$ となる。$S^2/A$ と $D$ の像は、$A$ の像である1点でのみ交わるので
\[
X\simeq X/A\cong S^2\vee S^1 .
\]
楔和のホモロジーより
\[
H_0(X)\cong\Z,\quad H_1(X)\cong\Z,\quad H_2(X)\cong\Z,\quad H_n(X)=0\ (n\ge3).
\]
$H_1$ の生成元は、$D$ を $N$ から $S$ へたどり、球面上の経線で $N$ に戻るループのクラスである。

\begin{note}
「球面に穴をあけずに取っ手（直径）をつけると、$H_1$ が1つ増える」という直観どおりの結果である。逆に、直径を球面の外側に付けた空間（$S^2$ に $N$ と $S$ を結ぶ外側の弧を付けたもの）も同じ議論で $S^2\vee S^1$ とホモトピー同値になる。直径が球面の内部にあるか外部にあるかはホモトピー型に影響しない。
\end{note}
\end{solution}

\begin{problem}[難易度 C]
$p\geq2$、$q$ を $p$ と互いに素な整数とし、$\zeta=e^{2\pi i/p}$ とする。レンズ空間 $L(p,q)$ は $S^3\subset\mathbb{C}^2$ を $(z,w)\mapsto(\zeta z,\zeta^q w)$ で生成される $\Z/p$ の作用で割った空間である。$L(p,q)$ は各次元 $0\leq k\leq3$ に $1$ つずつ胞体 $e^k$ をもつCW構造をもち、$2$ 胞体の接着写像 $S^1\to L^{(1)}\cong S^1$ の写像度は $p$ であることが知られている。これを用いて $H_*(L(p,q))$ を計算せよ。
\end{problem}

\begin{solution}
胞体鎖複体は $C_k=\Z e^k$（$0\leq k\leq 3$）である。$0$ 胞体が $1$ つなので $d_1=0$ である。$1$ 骨格は円周 $S^1$ で、$X^1/X^0=X^1$ だから、胞体境界公式より $d_2(e^2)=p\,e^1$ である。

$d_3$ を求める。$d_2\circ d_3=0$ であり、$d_2\colon\Z\to\Z$ は $p$ 倍で単射なので $d_3=0$ でなければならない。したがって胞体鎖複体は
\[
0\to\Z\xrightarrow{\ 0\ }\Z\xrightarrow{\ p\ }\Z\xrightarrow{\ 0\ }\Z\to 0
\]
（左から次数 $3,2,1,0$）である。これより
\[
H_0\cong\Z,\quad H_1\cong\Z/p,\quad H_2=\ker d_2=0,\quad H_3=\ker d_3=\Z
\]
であり、$k\geq4$ では $H_k=0$ である。なお $\chi(L(p,q))=1-1+1-1=0$ で、$\operatorname{rank}$ の交代和 $1-0+0-1=0$ と一致する。ホモロジー群は $q$ によらないことにも注意する。
\end{solution}

\begin{problem}[難易度 B]
$X$ を有限CW複体、$f\colon X\to X$ を連続写像とし、レフシェッツ数を
\[
L(f)=\sum_k(-1)^k\operatorname{tr}\bigl(f_*\colon H_k(X)/T_k\to H_k(X)/T_k\bigr)
\]
（$T_k$ は $H_k(X)$ のねじれ部分群）で定める。レフシェッツの不動点定理「$L(f)\neq0$ ならば $f$ は不動点をもつ」は既知とする。(1) $L(\id_X)=\chi(X)$ を示し、$\chi(X)\neq0$ なら $\id_X$ にホモトピックな任意の写像が不動点をもつことを示せ。(2) $n\geq1$、$f\colon S^n\to S^n$ の写像度を $d$ とするとき $L(f)$ を求め、$d\neq(-1)^{n+1}$ ならば $f$ は不動点をもつことを示せ。
\end{problem}

\begin{solution}
(1) $H_k(X)/T_k$ は階数 $\operatorname{rank} H_k(X)$ の自由アーベル群であり、恒等写像のトレースはその階数に等しい。オイラー標数の階数による表示（既出の問題）より
\[
L(\id_X)=\sum_k(-1)^k\operatorname{rank} H_k(X)=\chi(X)
\]
である。$f\simeq\id_X$ なら $f_*=(\id_X)_*$ なので $L(f)=L(\id_X)=\chi(X)\neq0$ となり、$f$ は不動点をもつ。

(2) $S^n$ のホモロジーは $H_0\cong\Z$、$H_n\cong\Z$ で、他は $0$ である（ねじれはない）。$S^n$ は弧状連結なので、任意の $1$ 点のクラスが $H_0$ を生成し $f$ はそれを $1$ 点のクラスに写す。よって $f_*$ は $H_0$ 上で恒等写像でトレースは $1$ である。$H_n$ 上では $d$ 倍なのでトレースは $d$ である。したがって
\[
L(f)=1+(-1)^n d
\]
である。$L(f)=0\iff d=-(-1)^n=(-1)^{n+1}$ だから、$d\neq(-1)^{n+1}$ ならば $L(f)\neq0$ となり $f$ は不動点をもつ。これは不動点をもたない写像の写像度に関する既出の結果と同じ結論である。
\end{solution}

\begin{problem}[難易度 C]
$k\geq1$ とする。任意の連続写像 $f\colon\mathbb{RP}^{2k}\to\mathbb{RP}^{2k}$ は不動点をもつことを示せ。一方、$\mathbb{RP}^{2k+1}$ には不動点をもたない連続写像が存在することを示せ。（前問のレフシェッツ数の定義と不動点定理を用いてよい。）
\end{problem}

\begin{solution}
$\mathbb{RP}^n$ のホモロジーの計算より $H_0(\mathbb{RP}^{2k})\cong\Z$、$0<j<2k$ の奇数 $j$ で $H_j\cong\Z/2$、その他の $j>0$ で $H_j=0$ である。したがって $j>0$ で $H_j/T_j=0$ であり、レフシェッツ数には $H_0$ の項のみが寄与する。$\mathbb{RP}^{2k}$ は弧状連結なので $f_*$ は $H_0\cong\Z$ 上で恒等写像である。よって $L(f)=1\neq0$ であり、$f$ は不動点をもつ。

$\mathbb{RP}^{2k+1}$ については、$\R^{2k+2}$ の線形写像
\[
J(x_1,x_2,\dots,x_{2k+1},x_{2k+2})=(-x_2,x_1,\dots,-x_{2k+2},x_{2k+1})
\]
を考える。$J$ は可逆な線形写像なので $f([x])=[Jx]$ は well-defined な連続写像 $\mathbb{RP}^{2k+1}\to\mathbb{RP}^{2k+1}$ を定める。$f([x])=[x]$ とすると、ある実数 $\lambda$ で $Jx=\lambda x$（$x\neq0$）となる。$J^2=-\id$ より $-x=J^2x=\lambda^2x$ となり、$\lambda^2=-1$ が導かれて矛盾する。よって $f$ は不動点をもたない。
\end{solution}

\begin{problem}[難易度 B]
$m,n\geq0$、$m\neq n$ ならば $\R^m$ と $\R^n$ は同相でないことを示せ。さらに、空でない開集合 $U\subset\R^m$ と $V\subset\R^n$ が同相ならば $m=n$ であることを示せ。
\end{problem}

\begin{solution}
まず $n\geq1$、$x\in\R^n$ に対して局所ホモロジー $H_k(\R^n,\R^n\setminus\{x\})$ を計算する。$\R^n$ は可縮、$\R^n\setminus\{x\}$ は $S^{n-1}$ にホモトピー同値なので、対の被約ホモロジー長完全列
\[
\tilde H_k(\R^n)\to H_k(\R^n,\R^n\setminus\{x\})\xrightarrow{\ \partial\ }\tilde H_{k-1}(\R^n\setminus\{x\})\to\tilde H_{k-1}(\R^n)
\]
で両端が $0$ だから、$H_k(\R^n,\R^n\setminus\{x\})\cong\tilde H_{k-1}(S^{n-1})$ であり、これは $k=n$ のとき $\Z$、それ以外で $0$ である。

$h\colon\R^m\to\R^n$ が同相なら、$h$ は対の同相 $(\R^m,\R^m\setminus\{x\})\to(\R^n,\R^n\setminus\{h(x)\})$ を与えるので局所ホモロジーは同型である。$m,n\geq1$ なら上の計算から $m=n$ となる。$m=0$ のとき $\R^0$ は $1$ 点で、$n\geq1$ の $\R^n$ とは濃度が異なるので同相でない。

開集合の場合、$x\in U$ に対し閉集合 $\R^m\setminus U$ は開集合 $\R^m\setminus\{x\}$ の内部に含まれるので、切除により
\[
H_k(U,U\setminus\{x\})\cong H_k(\R^m,\R^m\setminus\{x\})
\]
である。同相 $U\cong V$ は局所ホモロジーの同型を与えるので、上と同様に $m=n$ を得る（$m=0$ または $n=0$ のときは、空でない開集合が $1$ 点か無限集合かで区別できる）。
\end{solution}

\begin{problem}[難易度 D]
次の事実を既知とする：「$h\colon D^k\to S^n$ が埋め込みならば、すべての $i$ で $\tilde H_i(S^n\setminus h(D^k))=0$。」これを用いて、$0\leq k<n$ と埋め込み $h\colon S^k\to S^n$ に対し
\[
\tilde H_i(S^n\setminus h(S^k))\cong\begin{cases}\Z & (i=n-k-1)\\ 0 & (\text{それ以外})\end{cases}
\]
を示せ。特に $S^n$ に埋め込まれた $S^{n-1}$ の補集合はちょうど $2$ つの連結成分をもつことを示せ（ジョルダン・ブラウワーの分離定理）。
\end{problem}

\begin{solution}
$k$ に関する帰納法で示す。$k=0$ のとき $h(S^0)$ は $2$ 点で、$S^n$ から $2$ 点を除いた空間は（立体射影により）$\R^n\setminus\{0\}\simeq S^{n-1}$ と同相なので主張が成り立つ。

$k\geq1$ とし、$S^k=D_+\cup D_-$（上半球と下半球、$D_\pm\cong D^k$、$D_+\cap D_-=S^{k-1}$）と分ける。$A=S^n\setminus h(D_+)$、$B=S^n\setminus h(D_-)$ とおくと、$h(D_\pm)$ はコンパクトで閉集合だから $A,B$ は開集合であり、
\[
A\cap B=S^n\setminus h(S^k),\qquad A\cup B=S^n\setminus h(S^{k-1})
\]
である。$k<n$ なので $S^k$ と $S^n$ のホモロジーは異なり、$h(S^k)\neq S^n$、すなわち $A\cap B\neq\emptyset$ である。既知の事実より $\tilde H_*(A)=\tilde H_*(B)=0$ なので、被約マイヤー・ヴィートリス完全列から
\[
\tilde H_i(S^n\setminus h(S^k))=\tilde H_i(A\cap B)\cong\tilde H_{i+1}(A\cup B)=\tilde H_{i+1}(S^n\setminus h(S^{k-1}))
\]
を得る。$h|_{S^{k-1}}$ は埋め込みなので帰納法の仮定から、右辺は $i+1=n-(k-1)-1$、すなわち $i=n-k-1$ のときのみ $\Z$ で、それ以外は $0$ である。

$k=n-1$ のとき $\tilde H_0(S^n\setminus h(S^{n-1}))\cong\Z$ なので、補集合はちょうど $2$ つの弧状連結成分をもつ。補集合は $S^n$ の開集合で局所弧状連結だから、弧状連結成分と連結成分は一致する。
\end{solution}

\begin{problem}[難易度 D]
$U\subset\R^n$ を開集合、$f\colon U\to\R^n$ を単射連続写像とするとき、$f(U)$ は $\R^n$ の開集合であり $f\colon U\to f(U)$ は同相であることを示せ（領域不変性定理）。前問とその中の既知の事実を用いてよい。
\end{problem}

\begin{solution}
$\R^n\subset S^n=\R^n\cup\{\infty\}$ とみなす。$x\in U$ を任意にとり、$x$ を中心とする閉球体 $D\subset U$ とその境界球面 $S=\partial D$ をとる。$D$ はコンパクトで $S^n$ はハウスドルフなので、$f|_D$、$f|_S$ は埋め込みである。

前問より $S^n\setminus f(S)$ はちょうど $2$ つの連結成分をもつ。一方
\[
S^n\setminus f(S)=\bigl(S^n\setminus f(D)\bigr)\sqcup f(D\setminus S)
\]
（$f$ の単射性より非交和）である。既知の事実より $\tilde H_0(S^n\setminus f(D))=0$ であり、$\infty\notin f(D)$ よりこれは空でないので連結である。また $f(D\setminus S)$ は連結集合 $D\setminus S$ の連続像で連結である。連結集合はどれか $1$ つの連結成分に含まれる。もし両者が同じ成分に含まれれば、その和である $S^n\setminus f(S)$ 全体が $1$ つの成分となり矛盾する。よって $2$ つの集合はそれぞれちょうど $1$ つの連結成分である。

$S^n\setminus f(S)$ は局所連結な開集合なので、その連結成分は開集合である。したがって $f(D\setminus S)$ は $S^n$ の開集合、よって $\R^n$ の開集合で $f(x)$ を含む。$x$ は任意なので $f(U)$ は開集合である。同じ議論を $U$ の任意の開部分集合に適用すれば $f$ は開写像となるので、$f\colon U\to f(U)$ は同相である。
\end{solution}

%================================================
\chapter{コホモロジー}
% 20問
%================================================

\begin{summary}
\begin{itemize}
  \item \textbf{定義}：$C^n(X;G)=\operatorname{Hom}(C_n(X),G)$、余境界は $(\delta\varphi)(\sigma)=\varphi(\partial\sigma)$ で $\delta\circ\delta=0$。$H^n(X;G)=\ker\delta/\operatorname{im}\delta$。
  \item \textbf{$H^0$}：$H^0(X;G)$ は弧状連結成分上で定数となる関数全体であり、$\prod_{\text{弧状連結成分}}G$ と同型（直和ではなく直積）。
  \item \textbf{関手性}：$f\colon X\to Y$ は $f^*\colon H^n(Y;G)\to H^n(X;G)$ を誘導し（向きが逆になる）、ホモトピー不変。
  \item \textbf{完全列}：対の長完全列
  \[
  \cdots\to H^n(X,A)\to H^n(X)\to H^n(A)\xrightarrow{\delta}H^{n+1}(X,A)\to\cdots
  \]
  と、切除定理、マイヤー・ヴィートリス完全列がホモロジーと同様に成り立つ（矢印の向きは逆）。
  \item \textbf{胞体コホモロジー}：胞体鎖複体に $\operatorname{Hom}(-,G)$ を施したもので計算できる。余境界の行列は胞体境界の行列の転置。
  \item \textbf{$H^1$}：$H^1(X;G)\cong\operatorname{Hom}(H_1(X),G)$（普遍係数定理の特別な場合。詳しくは第8章）。
  \item \textbf{ド・ラームコホモロジー}：滑らかな多様体 $M$ 上の微分形式の複体 $(\Omega^\bullet(M),d)$ に対し $H^k_{\mathrm{dR}}(M)=\{\text{閉形式}\}/\{\text{完全形式}\}$。\textbf{ポアンカレの補題}：星形開集合では $k\ge1$ で $H^k_{\mathrm{dR}}=0$。ド・ラームの定理 $H^k_{\mathrm{dR}}(M)\cong H^k(M;\R)$ も成り立つ。$\R^2\setminus\{0\}$ 上の $d\theta=\dfrac{-y\,dx+x\,dy}{x^2+y^2}$ は閉だが完全でない。
\end{itemize}
\end{summary}

\begin{problem}[難易度 A]
$G$ をアーベル群とする。特異コチェイン複体を $C^n(X;G)=\operatorname{Hom}(C_n(X),G)$、余境界を $(\delta\varphi)(c)=\varphi(\partial c)$ で定める。$1$ 点空間 $P$ について $H^n(P;G)$ を計算せよ。
\end{problem}

\begin{solution}
各 $n\geq0$ に対し $P$ への特異 $n$ 単体は定値写像 $\sigma_n$ ただ $1$ つなので $C_n(P)=\Z\sigma_n$ である。$\sigma_n$ の $n+1$ 個の面はすべて $\sigma_{n-1}$ だから、$n\geq1$ で
\[
\partial\sigma_n=\sum_{i=0}^n(-1)^i\sigma_{n-1}=\begin{cases}\sigma_{n-1} & (n\text{ 偶数})\\ 0 & (n\text{ 奇数})\end{cases}
\]
である。$\varphi\mapsto\varphi(\sigma_n)$ により $C^n(P;G)\cong G$ と同一視すると、$\delta\colon C^n\to C^{n+1}$ は $\partial_{n+1}$ の双対なので、$n+1$ が偶数（$n$ 奇数）のとき恒等写像、$n$ が偶数のとき $0$ である。よってコチェイン複体は
\[
G\xrightarrow{\ 0\ }G\xrightarrow{\ \id\ }G\xrightarrow{\ 0\ }G\xrightarrow{\ \id\ }\cdots
\]
（左から次数 $0,1,2,3,\dots$）となる。
\begin{itemize}
\item $n=0$：$H^0=\ker(0)=G$。
\item $n$ 奇数：$\delta^n=\id$ は単射なので $H^n=0$。
\item $n\geq2$ 偶数：$\ker\delta^n=G$、$\operatorname{im}\delta^{n-1}=G$ なので $H^n=0$。
\end{itemize}
したがって $H^0(P;G)\cong G$、$n\geq1$ で $H^n(P;G)=0$ である。
\end{solution}

\begin{problem}[難易度 B]
空間 $X$ とアーベル群 $G$ に対し、$H^0(X;G)$ は $X$ の各弧状連結成分上で定数となる関数 $X\to G$ 全体のなす群と同型であることを示せ。特に $X$ の弧状連結成分の集合を $\pi_0(X)$ とすると $H^0(X;G)\cong\prod_{\pi_0(X)}G$ である。$X$ が可算無限個の点からなる離散空間のとき、$H_0(X)$ と $H^0(X;\Z)$ を比較せよ。
\end{problem}

\begin{solution}
特異 $0$ 単体は点と同一視できるので、$C^0(X;G)$ は関数 $\varphi\colon X\to G$ 全体である。$C^{-1}=0$ なので $H^0(X;G)=\ker\delta^0$ である。特異 $1$ 単体 $\sigma\colon\Delta^1\to X$ は道であり、$\partial\sigma=\sigma(e_1)-\sigma(e_0)$ だから
\[
(\delta\varphi)(\sigma)=\varphi(\sigma(e_1))-\varphi(\sigma(e_0))
\]
である。よって $\delta\varphi=0$ であることは、任意の道の始点と終点で $\varphi$ の値が等しいこと、すなわち $\varphi$ が各弧状連結成分上で定数であることと同値である。そのような関数は各成分での値で決まり、値は自由に選べるので $H^0(X;G)\cong\prod_{\pi_0(X)}G$ である。

$X$ が可算無限離散空間なら $\pi_0(X)=X$ であり、
\[
H_0(X)\cong\bigoplus_{X}\Z,\qquad H^0(X;\Z)\cong\prod_{X}\Z
\]
である。前者は可算集合、後者は非可算集合（整数列全体）なので、両者は同型でない。ホモロジーでは直和、コホモロジーでは直積が現れることに注意する。
\end{solution}

\begin{problem}[難易度 A]
連続写像 $f\colon X\to Y$ は $f^{\sharp}\colon C^n(Y;G)\to C^n(X;G)$、$f^{\sharp}\varphi=\varphi\circ f_{\sharp}$ を誘導する。$f^{\sharp}$ がコチェイン写像であることを示し、さらに $f\simeq g$ ならば $f^*=g^*\colon H^n(Y;G)\to H^n(X;G)$ であることを示せ。ただしホモトピックな写像の誘導するチェイン写像がチェインホモトピックであることは既知とする。また、可縮な空間 $X$ について $H^n(X;G)$ を求めよ。
\end{problem}

\begin{solution}
$f_{\sharp}$ はチェイン写像なので、$\varphi\in C^n(Y;G)$ と $c\in C_{n+1}(X)$ に対し
\[
(\delta f^{\sharp}\varphi)(c)=\varphi(f_{\sharp}\partial c)=\varphi(\partial f_{\sharp}c)=(\delta\varphi)(f_{\sharp}c)=(f^{\sharp}\delta\varphi)(c)
\]
となり、$\delta f^{\sharp}=f^{\sharp}\delta$ である。よって $f^*$ が well-defined に定まる。

$f\simeq g$ のとき、$P\colon C_n(X)\to C_{n+1}(Y)$ で $g_{\sharp}-f_{\sharp}=\partial P+P\partial$ をみたすものがある。$\varphi\in C^n(Y;G)$ に対し $P^*\varphi\in C^{n-1}(X;G)$ を $(P^*\varphi)(c)=\varphi(Pc)$ で定めると、$c\in C_n(X)$ について
\[
(g^{\sharp}\varphi-f^{\sharp}\varphi)(c)=\varphi(\partial Pc)+\varphi(P\partial c)=(P^*\delta\varphi)(c)+(\delta P^*\varphi)(c)
\]
である。よって $g^{\sharp}-f^{\sharp}=P^*\delta+\delta P^*$ であり、$\varphi$ がコサイクルなら $g^{\sharp}\varphi-f^{\sharp}\varphi=\delta(P^*\varphi)$ はコバウンダリーである。したがって $f^*=g^*$ となる。

$X$ が可縮なら、$1$ 点 $P$ への写像 $c\colon X\to P$ と包含 $i\colon P\to X$ は $c\circ i=\id_P$、$i\circ c\simeq\id_X$ をみたすので、$c^*\colon H^n(P;G)\to H^n(X;G)$ は同型である。前々問より $H^0(X;G)\cong G$、$n\geq1$ で $H^n(X;G)=0$ である。
\end{solution}

\begin{problem}[難易度 A]
$X=\bigsqcup_\alpha X_\alpha$ を空間の非交和とするとき、$H^n(X;G)\cong\prod_\alpha H^n(X_\alpha;G)$ を示せ。
\end{problem}

\begin{solution}
$\Delta^n$ は弧状連結なので、特異単体 $\sigma\colon\Delta^n\to X$ の像はただ $1$ つの $X_\alpha$ に含まれる。よって $C_n(X)=\bigoplus_\alpha C_n(X_\alpha)$ であり、$\partial$ はこの直和分解を保つ。$\operatorname{Hom}(-,G)$ は直和を直積に変えるので
\[
C^n(X;G)=\operatorname{Hom}\Bigl(\bigoplus_\alpha C_n(X_\alpha),G\Bigr)\cong\prod_\alpha C^n(X_\alpha;G)
\]
であり、この同型のもとで $\delta$ は成分ごとの $\delta$ の積に対応する。コチェイン複体の直積については、成分ごとに考えて
\[
\ker\Bigl(\prod_\alpha\delta_\alpha\Bigr)=\prod_\alpha\ker\delta_\alpha,\qquad \operatorname{im}\Bigl(\prod_\alpha\delta_\alpha\Bigr)=\prod_\alpha\operatorname{im}\delta_\alpha
\]
が成り立つ。したがって $H^n(X;G)\cong\prod_\alpha\ker\delta_\alpha/\operatorname{im}\delta_\alpha=\prod_\alpha H^n(X_\alpha;G)$ である。
\end{solution}

\begin{problem}[難易度 B]
$A\subset X$ とし、$C^n(X,A;G)=\operatorname{Hom}(C_n(X,A),G)$ と定める。
\[
0\to C^n(X,A;G)\xrightarrow{\ j^{\sharp}\ }C^n(X;G)\xrightarrow{\ i^{\sharp}\ }C^n(A;G)\to0
\]
がコチェイン複体の短完全列であることを示し、それから得られる長完全列の連結準同型 $\delta\colon H^n(A;G)\to H^{n+1}(X,A;G)$ を具体的に記述せよ。
\end{problem}

\begin{solution}
$C_n(X)$ は $X$ の特異 $n$ 単体を基底とする自由加群で、$C_n(A)$ は $A$ に像が含まれる単体で生成される部分加群である。$C_n(X,A)=C_n(X)/C_n(A)$ は残りの単体を基底とする自由加群なので、短完全列
\[
0\to C_n(A)\xrightarrow{\ i_{\sharp}\ }C_n(X)\xrightarrow{\ j_{\sharp}\ }C_n(X,A)\to0
\]
は分裂する。分裂短完全列に $\operatorname{Hom}(-,G)$ を施すと再び（分裂）短完全列になるので、問題の列は完全である。具体的には、$C^n(X,A;G)$ は $A$ 内の単体上で $0$ となるコチェインの全体と同一視され、$i^{\sharp}$ は制限であり、その全射性は $A$ 上のコチェインを $A$ に含まれない単体で $0$ と延長すればよいことからわかる。$i_{\sharp},j_{\sharp}$ がチェイン写像なので $i^{\sharp},j^{\sharp}$ はコチェイン写像である。

連結準同型は次のように与えられる。$\varphi\in C^n(A;G)$ をコサイクルとし、$\tilde\varphi\in C^n(X;G)$ を $i^{\sharp}\tilde\varphi=\varphi$ となる延長とする。$i^{\sharp}\delta\tilde\varphi=\delta\varphi=0$ なので、$\delta\tilde\varphi$ は $A$ 内の単体上で $0$ であり、$C^{n+1}(X,A;G)$ の元とみなせる。これはコサイクルで、
\[
\delta[\varphi]=[\delta\tilde\varphi]\in H^{n+1}(X,A;G)
\]
と定める。延長の取り方を変えると $\tilde\varphi$ は $C^n(X,A;G)$ の元 $\psi$ だけずれ、$\delta\tilde\varphi$ は $\delta\psi$ だけずれるので、クラスは変わらない。こうして長完全列
\[
\cdots\to H^n(X,A;G)\to H^n(X;G)\to H^n(A;G)\xrightarrow{\ \delta\ }H^{n+1}(X,A;G)\to\cdots
\]
を得る。
\end{solution}

\begin{problem}[難易度 B]
(1) 弧状連結な空間 $X$ と $x_0\in X$ に対し、$n\geq1$ で $H^n(X,\{x_0\};G)\cong H^n(X;G)$、また $H^0(X,\{x_0\};G)=0$ であることを示せ。(2) $S^0=\{p,q\}\subset S^1$ とするとき、対 $(S^1,S^0)$ の長完全列を用いて $H^n(S^1,S^0;\Z)$ を計算せよ。ただし $H^0(S^1;\Z)\cong H^1(S^1;\Z)\cong\Z$、$H^n(S^1;\Z)=0$（$n\geq2$）は既知とする。
\end{problem}

\begin{solution}
(1) 対の長完全列
\[
H^{n-1}(\{x_0\};G)\to H^n(X,\{x_0\};G)\to H^n(X;G)\to H^n(\{x_0\};G)
\]
において、$n\geq2$ なら両端は $0$ なので同型である。低次の部分は
\[
0\to H^0(X,x_0)\to H^0(X)\xrightarrow{\ r\ }H^0(x_0)\to H^1(X,x_0)\to H^1(X)\to 0
\]
である。$H^0$ を弧状連結成分上定数の関数とみなすと、$H^0(X)\cong G$ は定数関数全体で、$r$ は $x_0$ での値をとる写像だから同型である。よって $H^0(X,x_0)=\ker r=0$、また $r$ が全射なので $H^1(X,x_0)\to H^1(X)$ は単射かつ全射となり同型である。

(2) 長完全列の低次部分は
\begin{align*}
0&\to H^0(S^1,S^0)\to H^0(S^1)\cong\Z\xrightarrow{\ r\ }H^0(S^0)\cong\Z^2\\
&\to H^1(S^1,S^0)\to H^1(S^1)\cong\Z\to H^1(S^0)=0
\end{align*}
である。$r$ は定数関数 $c$ を $(c,c)$ に写す対角写像で単射なので $H^0(S^1,S^0)=0$ である。$\operatorname{coker} r=\Z^2/\{(c,c)\}\cong\Z$（$(a,b)\mapsto a-b$）だから短完全列
\[
0\to\Z\to H^1(S^1,S^0)\to\Z\to0
\]
を得て、右端が自由なので分裂し $H^1(S^1,S^0;\Z)\cong\Z^2$ である。$n\geq2$ では $H^{n-1}(S^0)=0$、$H^n(S^1)=0$ より $H^n(S^1,S^0;\Z)=0$ である。
\end{solution}

\begin{problem}[難易度 B]
コホモロジーのマイヤー・ヴィートリス完全列を用いて、$n\geq1$ のとき
\[
H^k(S^n;G)\cong\begin{cases}G & (k=0,n)\\ 0 & (\text{それ以外})\end{cases}
\]
を示せ。
\end{problem}

\begin{solution}
$U=S^n\setminus\{\text{南極}\}$、$V=S^n\setminus\{\text{北極}\}$ とすると $U,V$ は可縮、$U\cap V\simeq S^{n-1}$ である。マイヤー・ヴィートリス完全列
\[
H^{k-1}(U)\oplus H^{k-1}(V)\xrightarrow{\ \rho\ }H^{k-1}(U\cap V)\xrightarrow{\ \delta\ }H^k(S^n)\to H^k(U)\oplus H^k(V)
\]
において、$\rho(a,b)=a|_{U\cap V}-b|_{U\cap V}$ である。$k\geq2$ では両端が $0$ なので $H^k(S^n;G)\cong H^{k-1}(S^{n-1};G)$ となる。$S^n$ は弧状連結なので $H^0(S^n;G)\cong G$ である。

$k=1$ のとき、$H^1(U)\oplus H^1(V)=0$ より $H^1(S^n)\cong\operatorname{coker}\rho$（$\rho\colon G\oplus G\to H^0(U\cap V)$）である。
\begin{itemize}
\item $n\geq2$：$U\cap V$ は弧状連結で $H^0(U\cap V)=G$、$\rho(a,b)=a-b$ は全射なので $H^1(S^n)=0$。
\item $n=1$：$U\cap V$ は $2$ つの弧の非交和で $H^0(U\cap V)=G^2$、$\rho(a,b)=(a-b,a-b)$ なので $\operatorname{coker}\rho\cong G$（$(x,y)\mapsto x-y$）。よって $H^1(S^1)\cong G$。
\end{itemize}
$S^1$ では $k\geq2$ で $H^k(S^1)\cong H^{k-1}(S^0)=0$ である。以上と $k\geq2$ での同型 $H^k(S^n)\cong H^{k-1}(S^{n-1})$ から、$n$ に関する帰納法で主張が従う。
\end{solution}

\begin{problem}[難易度 B]
$X,Y$ を弧状連結な空間とし、基点がそれぞれ自身を変位レトラクトにもつ開近傍をもつとする。$n\geq1$ に対して $H^n(X\vee Y;G)\cong H^n(X;G)\oplus H^n(Y;G)$ を示せ。これを用いて $S^2\vee S^4$ のコホモロジー群が $\mathbb{CP}^2$ のコホモロジー群と同型であることを確かめよ（$\mathbb{CP}^2$ については胞体分割を用いてよい）。
\end{problem}

\begin{solution}
基点の開近傍を $A\subset X$、$B\subset Y$ とし、$U=X\vee B$、$V=A\vee Y$ とおく。$U,V$ は $X\vee Y$ の開集合で、$U\simeq X$、$V\simeq Y$、$U\cap V=A\vee B$ は可縮である。マイヤー・ヴィートリス完全列
\begin{align*}
H^{n-1}(U)\oplus H^{n-1}(V)&\xrightarrow{\ \rho\ }H^{n-1}(U\cap V)\to H^n(X\vee Y)\\
&\to H^n(U)\oplus H^n(V)\to H^n(U\cap V)
\end{align*}
を考える。$n\geq2$ なら $H^{n-1}(U\cap V)=H^n(U\cap V)=0$ なので同型を得る。$n=1$ のとき $H^1(U\cap V)=0$ であり、$U,V,U\cap V$ はすべて弧状連結なので $\rho\colon G\oplus G\to G$、$\rho(a,b)=a-b$ は全射である。よって $H^1(X\vee Y)\to H^1(U)\oplus H^1(V)$ は単射かつ全射で同型である。いずれの場合も $H^n(X\vee Y)\cong H^n(X)\oplus H^n(Y)$ である。

前問より $S^2\vee S^4$ のコホモロジーは
\[
H^0\cong\Z,\quad H^2\cong\Z,\quad H^4\cong\Z,\quad \text{その他は }0
\]
である（$H^0$ は弧状連結性による）。一方 $\mathbb{CP}^2$ は $0,2,4$ 次元に $1$ つずつ胞体をもち、胞体コチェイン複体の余境界はすべて $0$（隣り合う次数の一方が $0$）なので、同じ群が得られる。両者はコホモロジー群では区別できないが、カップ積を用いると区別できることが知られている。
\end{solution}

\begin{problem}[難易度 C]
コホモロジーのマイヤー・ヴィートリス完全列を用いて $H^n(\mathbb{RP}^2;G)$ を計算せよ。$G=\Z$ と $G=\Z/2$ の場合を具体的に書け。ただし、$H_0$ が自由な空間に対する自然な同型 $H^1(X;G)\cong\operatorname{Hom}(H_1(X),G)$（普遍係数定理の特別な場合）は引用してよい。
\end{problem}

\begin{solution}
$\mathbb{RP}^2$ を $D^2$ の境界の対蹠点を同一視した空間とみなし、$p$ を中心の点とする。$U=\mathbb{RP}^2\setminus\{p\}$、$V$ を $p$ の小さな開円板とする。$U$ は径方向の変位レトラクトにより $\partial D^2/\!\sim\ =\mathbb{RP}^1\cong S^1$ にホモトピー同値、$V$ は可縮、$U\cap V$ は穴あき円板で $S^1$ にホモトピー同値である。$U\cap V$ の中の円周を $U$ に含め、$\mathbb{RP}^1$ にレトラクトすると、それは $S^1\to\partial D^2\to\mathbb{RP}^1$ という二重被覆になるので、$H_1(U\cap V)\cong\Z\to H_1(U)\cong\Z$ は $\pm2$ 倍である。したがって $H^1(U;G)\to H^1(U\cap V;G)$ は $\operatorname{Hom}(\pm2,G)$、すなわち $G$ 上の $\pm2$ 倍写像である。

マイヤー・ヴィートリス完全列
\begin{align*}
0&\to H^0(\mathbb{RP}^2)\to H^0(U)\oplus H^0(V)\xrightarrow{\ \rho\ }H^0(U\cap V)\to H^1(\mathbb{RP}^2)\\
&\to H^1(U)\oplus H^1(V)\xrightarrow{\ \rho_1\ }H^1(U\cap V)\to H^2(\mathbb{RP}^2)\to H^2(U)\oplus H^2(V)
\end{align*}
において、$\rho(a,b)=a-b\colon G\oplus G\to G$ は全射、$H^1(V)=0$、$\rho_1$ は $G$ 上の $\pm2$ 倍、$H^2(U)\cong H^2(S^1)=0$、$H^2(V)=0$ である。よって
\[
H^0(\mathbb{RP}^2;G)\cong G,\quad H^1(\mathbb{RP}^2;G)\cong\ker(2\colon G\to G),\quad H^2(\mathbb{RP}^2;G)\cong G/2G
\]
であり、$n\geq3$ では $H^n=0$ である。$G=\Z$ なら $H^0=\Z$、$H^1=0$、$H^2=\Z/2$。$G=\Z/2$ なら $2$ 倍は零写像なので $H^0=H^1=H^2=\Z/2$ である。
\end{solution}

\begin{problem}[難易度 B]
CW複体 $X$ の胞体コチェイン複体を $C^k_{\mathrm{CW}}(X;G)=\operatorname{Hom}(C^{\mathrm{CW}}_k(X),G)$、余境界を胞体境界写像 $d_{k+1}$ の双対 $\delta^k\varphi=\varphi\circ d_{k+1}$ で定める。そのコホモロジーが特異コホモロジー $H^k(X;G)$ と同型であることは既知とする。これを用いて、任意のアーベル群 $G$ についてトーラス $T=S^1\times S^1$ のコホモロジー群 $H^k(T;G)$ を計算せよ。$G=\Z$、$\Z/2$ の場合を書け。
\end{problem}

\begin{solution}
$T$ は $0$ 胞体 $v$、$1$ 胞体 $a,b$、語 $aba^{-1}b^{-1}$ に沿って接着された $2$ 胞体 $e$ からなるCW構造をもつ。$0$ 胞体が $1$ つなので $d_1=0$ であり、$2$ 胞体の境界の $a,b$ の係数はそれぞれ語における指数和 $1-1=0$ なので $d_2=0$ である。

したがって胞体コチェイン複体は
\[
0\to\operatorname{Hom}(\Z,G)\cong G\xrightarrow{\ 0\ }\operatorname{Hom}(\Z^2,G)\cong G^2\xrightarrow{\ 0\ }\operatorname{Hom}(\Z,G)\cong G\to0
\]
であり、余境界はすべて $0$ である。よって
\[
H^0(T;G)\cong G,\quad H^1(T;G)\cong G^2,\quad H^2(T;G)\cong G,\quad H^k(T;G)=0\ (k\geq3)
\]
である。$G=\Z$ なら $\Z,\Z^2,\Z$、$G=\Z/2$ なら $\Z/2,(\Z/2)^2,\Z/2$ となる。
\end{solution}

\begin{problem}[難易度 B]
$\mathbb{RP}^n$ の胞体鎖複体の境界写像が $d_k=1+(-1)^k$ 倍（$1\leq k\leq n$）であることを用いて、$H^k(\mathbb{RP}^n;\Z)$ を計算せよ。また、ホモロジー群 $H_k(\mathbb{RP}^n)$ と比較してねじれ部分群の現れる次数がどう変わるかを述べよ。
\end{problem}

\begin{solution}
$C^k_{\mathrm{CW}}=\operatorname{Hom}(\Z e^k,\Z)\cong\Z$（$0\leq k\leq n$）と同一視すると、$\delta^k$ は $d_{k+1}$ の双対なので $1+(-1)^{k+1}$ 倍、すなわち $k$ が奇数なら $2$ 倍、偶数なら $0$ 倍である（$\delta^n=0$）。コチェイン複体は
\[
\Z\xrightarrow{\ 0\ }\Z\xrightarrow{\ 2\ }\Z\xrightarrow{\ 0\ }\Z\xrightarrow{\ 2\ }\cdots\to\Z
\]
（左から次数 $0,1,2,\dots,n$）である。
\begin{itemize}
\item $k=0$：$H^0=\ker\delta^0=\Z$。
\item $k$ 奇数、$k<n$：$\delta^k=2$ は単射なので $H^k=0$。
\item $k$ 偶数、$0<k\leq n$：$\ker\delta^k=\Z$（$k<n$ なら $\delta^k=0$、$k=n$ なら $\delta^n=0$）、$\operatorname{im}\delta^{k-1}=2\Z$ なので $H^k=\Z/2$。
\item $k=n$ 奇数：$\ker\delta^n=\Z$、$\operatorname{im}\delta^{n-1}=0$ なので $H^n=\Z$。
\end{itemize}
$k>n$ では $0$ である。

ホモロジーでは $\Z/2$ は $0<k<n$ の奇数次に現れたが、コホモロジーでは $0<k\leq n$ の偶数次に現れる。すなわちねじれ部分群は次数が $1$ つ上にずれる。これは普遍係数定理から予想される現象である。
\end{solution}

\begin{problem}[難易度 B]
$H^k(\mathbb{RP}^n;\Z/2)$ を計算せよ。また、無限次元実射影空間 $\mathbb{RP}^\infty$（各次元に $1$ つずつ胞体をもち、$d_k=1+(-1)^k$）について $H^k(\mathbb{RP}^\infty;\Z/2)$ と $H^k(\mathbb{RP}^\infty;\Z)$ を求めよ。
\end{problem}

\begin{solution}
$C^k_{\mathrm{CW}}(\mathbb{RP}^n;\Z/2)\cong\Z/2$（$0\le k\le n$）であり、$\delta^k$ は $0$ 倍または $2$ 倍である。$\Z/2$ では $2=0$ なので余境界はすべて零写像である。よって
\[
H^k(\mathbb{RP}^n;\Z/2)\cong\begin{cases}\Z/2 & (0\leq k\leq n)\\ 0 & (k>n)\end{cases}
\]
である。$\Z$ 係数では見えなかった奇数次の情報が、$\Z/2$ 係数ではすべての次数に現れる。

$\mathbb{RP}^\infty$ でも同様に、$\Z/2$ 係数では余境界がすべて $0$ なので、すべての $k\geq0$ で $H^k(\mathbb{RP}^\infty;\Z/2)\cong\Z/2$ である。$\Z$ 係数では前問の計算で「最上次元」が存在しない場合に当たり、
\[
H^0\cong\Z,\qquad H^k\cong\Z/2\ (k\text{ 偶数},\ k\geq2),\qquad H^k=0\ (k\text{ 奇数})
\]
となる。（$\mathbb{RP}^\infty$ は無限CW複体であるが、胞体コホモロジーと特異コホモロジーの同型はこの場合も成り立つ。）
\end{solution}

\begin{problem}[難易度 A]
$\mathbb{CP}^n$ のコホモロジー群 $H^k(\mathbb{CP}^n;G)$ を任意のアーベル群 $G$ について計算せよ。
\end{problem}

\begin{solution}
$\mathbb{CP}^n$ は次元 $0,2,\dots,2n$ に $1$ つずつ胞体をもつCW構造をもつ。胞体コチェイン複体は
\[
C^k_{\mathrm{CW}}(\mathbb{CP}^n;G)\cong\begin{cases}G & (k=0,2,\dots,2n)\\ 0 & (\text{それ以外})\end{cases}
\]
であり、余境界 $\delta^k\colon C^k\to C^{k+1}$ は始域か終域の一方が必ず $0$ なので零写像である。したがって
\[
H^k(\mathbb{CP}^n;G)\cong\begin{cases}G & (k=0,2,\dots,2n)\\ 0 & (\text{それ以外})\end{cases}
\]
である。
\end{solution}

\begin{problem}[難易度 B]
$h\geq1$ とし、向き付け不可能閉曲面 $N_h$ を $0$ 胞体 $1$ 個、$1$ 胞体 $a_1,\dots,a_h$、語 $a_1a_1\cdots a_ha_h$ に沿って接着された $2$ 胞体 $e$ からなるCW複体とする。$H^k(N_h;\Z)$ と $H^k(N_h;\Z/2)$ を計算せよ。特にクラインの壺 $K=N_2$ の場合を書け。
\end{problem}

\begin{solution}
胞体境界は $d_1=0$、$d_2(e)=2(a_1+\dots+a_h)$ である。$\varphi\in C^1=\operatorname{Hom}(\Z^h,G)$ を値 $\varphi(a_i)=g_i$ で表すと、
\[
\delta^1\varphi=\varphi\circ d_2,\qquad (\delta^1\varphi)(e)=2(g_1+\dots+g_h)
\]
であり、$\delta^0=0$ である。

$G=\Z$ の場合：$H^0=\Z$。$H^1=\ker\delta^1=\{(g_i): g_1+\dots+g_h=0\}$ であり、これは $g_1,\dots,g_{h-1}$ を自由に選べば $g_h$ が決まるので $\Z^{h-1}$ と同型である。$\operatorname{im}\delta^1=2\Z$（$g_1=1$、他を $0$ とすれば $2$ が実現され、像は常に偶数）なので $H^2=\Z/2\Z=\Z/2$ である。

$G=\Z/2$ の場合：$2=0$ なので $\delta^1=0$ であり、
\[
H^0\cong\Z/2,\quad H^1\cong(\Z/2)^h,\quad H^2\cong\Z/2
\]
である。

クラインの壺 $K=N_2$ では、$\Z$ 係数で $H^0=\Z$、$H^1=\Z$、$H^2=\Z/2$、$\Z/2$ 係数で $\Z/2,(\Z/2)^2,\Z/2$ となる。いずれの場合も $k\geq3$ で $H^k=0$ である。
\end{solution}

\begin{problem}[難易度 C]
レンズ空間 $L(p,q)$（$p\geq2$）の胞体鎖複体が
\[
0\to\Z\xrightarrow{\ 0\ }\Z\xrightarrow{\ p\ }\Z\xrightarrow{\ 0\ }\Z\to0
\]
（左から次数 $3,2,1,0$）であることを用いて、$H^k(L(p,q);\Z)$ と $H^k(L(p,q);\Z/p)$ を計算せよ。さらに、$p$ を割らない素数 $\ell$ について $H^k(L(p,q);\Z/\ell)$ を求めよ。
\end{problem}

\begin{solution}
各 $C^k_{\mathrm{CW}}(L;G)\cong G$（$0\leq k\leq3$）と同一視すると、$\delta^k$ は $d_{k+1}$ の双対なので
\[
G\xrightarrow{\ \delta^0=0\ }G\xrightarrow{\ \delta^1=p\ }G\xrightarrow{\ \delta^2=0\ }G
\]
となる。したがって
\[
H^0\cong G,\quad H^1\cong\ker(p\colon G\to G),\quad H^2\cong G/pG,\quad H^3\cong G
\]
である。

$G=\Z$：$p$ 倍は単射なので $H^0=\Z$、$H^1=0$、$H^2=\Z/p$、$H^3=\Z$。

$G=\Z/p$：$p$ 倍は零写像なので $H^k\cong\Z/p$（$k=0,1,2,3$）。

$G=\Z/\ell$（$\ell\nmid p$）：$p$ は $\Z/\ell$ で可逆なので $p$ 倍は同型であり、$H^0=H^3=\Z/\ell$、$H^1=H^2=0$。これは $3$ 次元球面 $S^3$ と同じであり、$\Z/\ell$ 係数では $L(p,q)$ と $S^3$ は区別できない。
\end{solution}

\begin{problem}[難易度 C]
$X$ を有限CW複体、$F$ を体とする。
\[
\chi(X)=\sum_k(-1)^k\dim_F H^k(X;F)
\]
を示せ。これを用いて $\mathbb{RP}^n$ と $L(p,q)$ のオイラー標数を $\Z/2$ 係数、$\Z/p$ 係数の計算から求め、$\mathbb{Q}$ 係数での計算と一致することを確かめよ。
\end{problem}

\begin{solution}
$k$ 胞体の個数を $c_k$ とすると $C^k=C^k_{\mathrm{CW}}(X;F)=\operatorname{Hom}(\Z^{c_k},F)\cong F^{c_k}$ は $c_k$ 次元 $F$ ベクトル空間で、$\delta^k$ は $F$ 線形写像である。$Z^k=\ker\delta^k$、$B^k=\operatorname{im}\delta^{k-1}$ とおくと、次元定理から
\[
c_k=\dim Z^k+\dim B^{k+1},\qquad \dim H^k(X;F)=\dim Z^k-\dim B^k
\]
である。よって
\[
\sum_k(-1)^k c_k-\sum_k(-1)^k\dim H^k(X;F)=\sum_k(-1)^k\bigl(\dim B^{k+1}+\dim B^k\bigr)=0
\]
（隣り合う項が打ち消し合う）となる。

$\mathbb{RP}^n$：$\Z/2$ 係数では $0\leq k\leq n$ ですべて $\Z/2$ なので $\chi=\sum_{k=0}^n(-1)^k$、すなわち $n$ 偶数で $1$、奇数で $0$ である。$\mathbb{Q}$ 係数でも胞体コチェイン複体は $\mathbb{Q}\xrightarrow{0}\mathbb{Q}\xrightarrow{2}\mathbb{Q}\xrightarrow{0}\cdots$ で $2$ が可逆なので、$H^0=\mathbb{Q}$、$n$ が奇数のときのみ $H^n=\mathbb{Q}$ となり、同じ値を得る。

$L(p,q)$：$\Z/p$ 係数ではすべて $\Z/p$ なので $\chi=1-1+1-1=0$。$\mathbb{Q}$ 係数では $p$ が可逆なので $H^0=H^3=\mathbb{Q}$、$H^1=H^2=0$ となり $\chi=1-1=0$ で一致する。
\end{solution}

\begin{problem}[難易度 A]
$M$ を滑らかな多様体とし、ド・ラームコホモロジーを $H^k_{\mathrm{dR}}(M)=\ker(d\colon\Omega^k(M)\to\Omega^{k+1}(M))/\operatorname{im}(d\colon\Omega^{k-1}(M)\to\Omega^k(M))$ で定める。$H^0_{\mathrm{dR}}(M)$ は $M$ 上の局所定数関数全体と一致し、$M$ の連結成分の集合を $\pi_0(M)$ とすると $H^0_{\mathrm{dR}}(M)\cong\prod_{\pi_0(M)}\R$ であることを示せ。特異コホモロジー $H^0(M;\R)$ と比較せよ。
\end{problem}

\begin{solution}
$\Omega^{-1}(M)=0$ なので $H^0_{\mathrm{dR}}(M)=\{f\in C^\infty(M):df=0\}$ である。局所座標 $(x_1,\dots,x_n)$ で $df=\sum_i(\partial f/\partial x_i)dx_i$ なので、$df=0$ はすべての偏導関数が $0$ であることと同値である。

$df=0$ なら、各点の周りの座標近傍として凸な開球をとると、その上で平均値の定理により $f$ は定数である。よって $f$ は局所定数である。逆に局所定数関数は滑らかで $df=0$ をみたす。

局所定数関数の各値の逆像は開集合で、これらは互いに交わらず $M$ を覆うので、$f$ は各連結成分上で定数である。逆に各連結成分上の値を任意に指定すれば局所定数関数が得られる（多様体は局所連結なので連結成分は開集合）。よって $H^0_{\mathrm{dR}}(M)\cong\prod_{\pi_0(M)}\R$ である。

多様体は局所弧状連結なので連結成分と弧状連結成分は一致し、$H^0(M;\R)\cong\prod_{\pi_0(M)}\R$ である。したがって $H^0_{\mathrm{dR}}(M)\cong H^0(M;\R)$ であり、これはド・ラームの定理 $H^k_{\mathrm{dR}}(M)\cong H^k(M;\R)$ の $k=0$ の場合である。
\end{solution}

\begin{problem}[難易度 B]
ポアンカレの補題は「$\R^n$（より一般に星形開集合）上の閉形式は完全形式である。すなわち $k\geq1$ で $H^k_{\mathrm{dR}}(\R^n)=0$」と主張する。(1) $n=1$ の場合を示せ。(2) $\R^2$ 上の閉 $1$ 形式 $\omega=P\,dx+Q\,dy$（$\partial P/\partial y=\partial Q/\partial x$）に対し、
\[
f(x,y)=\int_0^1\bigl(xP(tx,ty)+yQ(tx,ty)\bigr)\,dt
\]
が $df=\omega$ をみたすことを示せ。(3) $H^*_{\mathrm{dR}}(\R^n)$ と $H^*(\R^n;\R)$ が一致することを確かめよ。
\end{problem}

\begin{solution}
(1) $\R$ 上では $\Omega^2(\R)=0$ なので任意の $1$ 形式 $\omega=g(x)\,dx$ は閉である。$f(x)=\int_0^x g(s)\,ds$ とおけば、微分積分学の基本定理より $f$ は滑らかで $df=g\,dx=\omega$ となる。よって $H^1_{\mathrm{dR}}(\R)=0$ であり、$k\geq2$ では $\Omega^k(\R)=0$ である。

(2) 積分記号下の微分により（被積分関数は $(x,y,t)$ について滑らか）
\[
\frac{\partial f}{\partial x}=\int_0^1\Bigl(P(tx,ty)+tx\,P_x(tx,ty)+ty\,Q_x(tx,ty)\Bigr)dt
\]
である。閉条件 $Q_x=P_y$ を用いると、被積分関数は
\[
P(tx,ty)+t\bigl(xP_x(tx,ty)+yP_y(tx,ty)\bigr)=\frac{d}{dt}\bigl(tP(tx,ty)\bigr)
\]
に等しい。よって $\partial f/\partial x=\bigl[tP(tx,ty)\bigr]_{t=0}^{t=1}=P(x,y)$ である。$x$ と $y$、$P$ と $Q$ の役割を入れ替えて同様に $\partial f/\partial y=Q(x,y)$ を得るので $df=\omega$ である。

(3) ポアンカレの補題と $H^0_{\mathrm{dR}}(\R^n)=\R$（$\R^n$ は連結）より、$H^0_{\mathrm{dR}}(\R^n)\cong\R$、$k\geq1$ で $H^k_{\mathrm{dR}}(\R^n)=0$ である。一方 $\R^n$ は可縮なので $H^0(\R^n;\R)\cong\R$、$k\geq1$ で $H^k(\R^n;\R)=0$ であり、両者は一致する。
\end{solution}

\begin{problem}[難易度 C]
$S^1$ 上の $1$ 形式で、$\pi\colon\R\to S^1$、$\pi(t)=(\cos t,\sin t)$ による引き戻しが $dt$ となるものを $d\theta$ と書く（$d\theta$ は完全形式とは限らない記号である）。積分写像 $I\colon\Omega^1(S^1)\to\R$、$I(\omega)=\int_{S^1}\omega$ を用いて、$H^1_{\mathrm{dR}}(S^1)\cong\R$ を示せ。また $H^0_{\mathrm{dR}}(S^1)$ を求め、特異コホモロジー $H^*(S^1;\R)$ と比較せよ。
\end{problem}

\begin{solution}
任意の $\omega\in\Omega^1(S^1)$ に対し $\pi^*\omega=g(t)\,dt$ と書け、$g$ は周期 $2\pi$ の滑らかな関数である。逆に周期 $2\pi$ の滑らかな $g$ に対し $\omega=g\,d\theta$ が定まる。$I(\omega)=\int_0^{2\pi}g(t)\,dt$ である。

$I$ は全射である：$I(c\,d\theta)=2\pi c$。$\Omega^2(S^1)=0$ なのですべての $1$ 形式は閉であり、$H^1_{\mathrm{dR}}(S^1)=\Omega^1(S^1)/d\Omega^0(S^1)$ である。そこで $\ker I=d\Omega^0(S^1)$ を示す。

$\omega=dh$（$h\in C^\infty(S^1)$）なら $\pi^*\omega=d(h\circ\pi)$ であり、$I(\omega)=h(\pi(2\pi))-h(\pi(0))=0$ である。逆に $I(\omega)=0$ のとき $F(t)=\int_0^t g(s)\,ds$ とおくと、
\[
F(t+2\pi)-F(t)=\int_t^{t+2\pi}g(s)\,ds=\int_0^{2\pi}g(s)\,ds=0
\]
（$g$ の周期性）なので $F$ は周期 $2\pi$ であり、$F=h\circ\pi$ となる滑らかな $h\colon S^1\to\R$ が存在する（$\pi$ は局所微分同相）。$\pi^*dh=dF=g\,dt=\pi^*\omega$ であり、$\pi$ が局所微分同相なので $dh=\omega$ である。

よって $I$ は同型 $H^1_{\mathrm{dR}}(S^1)\cong\R$ を誘導し、その生成元は $[d\theta]$ である。また $S^1$ は連結なので $H^0_{\mathrm{dR}}(S^1)\cong\R$ である。特異コホモロジーでも $H^0(S^1;\R)\cong H^1(S^1;\R)\cong\R$ であり、ド・ラームの定理と整合する。
\end{solution}

\begin{problem}[難易度 C]
$X=\R^2\setminus\{0\}$ 上の $1$ 形式
\[
\omega=\frac{-y\,dx+x\,dy}{x^2+y^2}
\]
が閉形式であるが完全形式でないことを示せ。さらに、閉 $1$ 形式 $\eta$ に単位円周上の積分 $\int_{S^1}\eta$ を対応させる写像が同型 $H^1_{\mathrm{dR}}(X)\cong\R$ を誘導することを示し、$H^1(X;\R)$ と比較せよ。
\end{problem}

\begin{solution}
$P=-y/(x^2+y^2)$、$Q=x/(x^2+y^2)$ とすると
\[
\frac{\partial Q}{\partial x}=\frac{y^2-x^2}{(x^2+y^2)^2}=\frac{\partial P}{\partial y}
\]
なので $d\omega=(Q_x-P_y)\,dx\wedge dy=0$ である。単位円周 $\gamma(t)=(\cos t,\sin t)$ に沿って $\gamma^*\omega=(\sin^2t+\cos^2t)\,dt=dt$ だから $\int_{S^1}\omega=2\pi$ である。もし $\omega=df$ なら $\int_{S^1}df=f(\gamma(2\pi))-f(\gamma(0))=0$ となり矛盾するので、$\omega$ は完全でない。

積分写像 $J([\eta])=\int_{S^1}\eta$ は、完全形式の閉曲線上の積分が $0$ であることから well-defined であり、$J([\omega])=2\pi$ より全射である。単射性を示す。$\eta$ を閉 $1$ 形式で $\int_{S^1}\eta=0$ とする。極座標 $\Phi\colon(0,\infty)\times\R\to X$、$\Phi(r,t)=(r\cos t,r\sin t)$ は局所微分同相で、$\Phi^*\eta=a\,dr+b\,dt$ は凸集合 $(0,\infty)\times\R$ 上の閉形式である（$a,b$ は $t$ について周期 $2\pi$）。凸集合上のポアンカレの補題により $dF=\Phi^*\eta$ となる $F$ が存在する。$D(r)=F(r,t+2\pi)-F(r,t)=\int_t^{t+2\pi}b(r,s)\,ds$ は $b$ の周期性より $t$ によらず、
\[
D'(r)=\frac{\partial F}{\partial r}(r,t+2\pi)-\frac{\partial F}{\partial r}(r,t)=a(r,t+2\pi)-a(r,t)=0
\]
なので定数である。$D(1)=\int_0^{2\pi}b(1,s)\,ds=\int_{S^1}\eta=0$ だから $F$ は $t$ について周期 $2\pi$ であり、$F=f\circ\Phi$ となる滑らかな $f\colon X\to\R$ がある。$\Phi^*df=\Phi^*\eta$ と $\Phi$ の局所微分同相性から $df=\eta$ である。よって $J$ は単射で、$H^1_{\mathrm{dR}}(X)\cong\R$ である。

$X$ は $S^1$ にホモトピー同値なので $H^1(X;\R)\cong H^1(S^1;\R)\cong\R$ であり、ド・ラームの定理と整合する。
\end{solution}

%================================================
\chapter{カップ積とコホモロジー環}
% 15問
%================================================

\begin{summary}
\begin{itemize}
  \item \textbf{カップ積}：$\varphi\in C^k$, $\psi\in C^l$ に対し
  \[
  (\varphi\smile\psi)(\sigma)=\varphi\bigl(\sigma|_{[v_0,\dots,v_k]}\bigr)\,\psi\bigl(\sigma|_{[v_k,\dots,v_{k+l}]}\bigr).
  \]
  $\delta(\varphi\smile\psi)=\delta\varphi\smile\psi+(-1)^k\varphi\smile\delta\psi$ より、コホモロジー上の積が定まる。
  \item \textbf{性質}：結合的で、$1\in H^0$ が単位元。自然性 $f^*(a\smile b)=f^*a\smile f^*b$。次数付き可換性 $a\smile b=(-1)^{kl}\,b\smile a$（$a\in H^k$, $b\in H^l$）。
  \item \textbf{代表的なコホモロジー環}：
  \begin{align*}
  &H^*(S^n)=\Z[x]/(x^2),\\
  &H^*(T^n)=\Lambda_\Z[x_1,\dots,x_n]\ (\lvert x_i\rvert=1),\\
  &H^*(\mathbb{CP}^n)=\Z[x]/(x^{n+1})\ (\lvert x\rvert=2),\\
  &H^*(\mathbb{RP}^n;\Z/2)=\Z/2[\alpha]/(\alpha^{n+1})\ (\lvert\alpha\rvert=1).
  \end{align*}
  $\Sigma_g$ では $H^1$ の基底 $a_i,b_i$ について $a_i\smile b_i=-b_i\smile a_i$ が $H^2$ の生成元で、それ以外の積は $0$。
  \item \textbf{楔和と懸垂}：$\tilde H^*(X\vee Y)\cong\tilde H^*(X)\times\tilde H^*(Y)$（環として直積。$X$ 側と $Y$ 側の積は $0$）。懸垂 $SX$ では正の次数の元どうしの積はすべて $0$。
  \item \textbf{キャップ積}：$\sigma\frown\varphi=\varphi\bigl(\sigma|_{[v_0,\dots,v_k]}\bigr)\,\sigma|_{[v_k,\dots,v_n]}$ により $H_n\times H^k\to H_{n-k}$ が定まる。
\end{itemize}
\end{summary}

\begin{problem}[難易度 B]
$R$ を可換環とし、$X$ を位相空間とする。$\varphi\in C^p(X;R)$、$\psi\in C^q(X;R)$ のカップ積 $\varphi\smile\psi\in C^{p+q}(X;R)$ を、特異単体 $\sigma\colon\Delta^{p+q}\to X$ に対して
\[
(\varphi\smile\psi)(\sigma)=\varphi(\sigma|[v_0,\dots,v_p])\,\psi(\sigma|[v_p,\dots,v_{p+q}])
\]
で定める。ただし余境界は $(\delta\varphi)(\tau)=\sum_i(-1)^i\varphi(\tau|[v_0,\dots,\widehat{v_i},\dots,v_{p+1}])$ とする。このとき
\[
\delta(\varphi\smile\psi)=\delta\varphi\smile\psi+(-1)^p\varphi\smile\delta\psi
\]
を示せ。さらに、これからカップ積がコホモロジー上の双線形写像 $H^p(X;R)\times H^q(X;R)\to H^{p+q}(X;R)$ を誘導することを示せ。
\end{problem}

\begin{solution}
$\sigma\colon\Delta^{p+q+1}\to X$ を特異単体とする。定義より
\begin{align*}
(\delta\varphi\smile\psi)(\sigma)&=\sum_{i=0}^{p+1}(-1)^i\varphi(\sigma|[v_0,\dots,\widehat{v_i},\dots,v_{p+1}])\,\psi(\sigma|[v_{p+1},\dots,v_{p+q+1}]),\\
(-1)^p(\varphi\smile\delta\psi)(\sigma)&=\sum_{i=p}^{p+q+1}(-1)^i\varphi(\sigma|[v_0,\dots,v_p])\,\psi(\sigma|[v_p,\dots,\widehat{v_i},\dots,v_{p+q+1}])
\end{align*}
である。第2式では $\delta\psi$ の第 $j$ 項の符号 $(-1)^{j}$（$j=i-p$）に $(-1)^p$ を掛けて $(-1)^i$ とした。第1式の $i=p+1$ の項と第2式の $i=p$ の項はともに $\varphi(\sigma|[v_0,\dots,v_p])\psi(\sigma|[v_{p+1},\dots,v_{p+q+1}])$ の定数倍で、符号がそれぞれ $(-1)^{p+1}$、$(-1)^p$ なので打ち消し合う。残りの項の和は
\[
\sum_{i=0}^{p+q+1}(-1)^i(\varphi\smile\psi)(\sigma|[v_0,\dots,\widehat{v_i},\dots,v_{p+q+1}])=\delta(\varphi\smile\psi)(\sigma)
\]
に他ならない（$i\le p$ のとき面は前半部分に、$i\ge p+1$ のとき後半部分に入る）。これで公式が示された。

この公式から、$\varphi,\psi$ がともにコサイクルなら $\delta(\varphi\smile\psi)=0$ で、$\varphi\smile\psi$ はコサイクルである。また $\varphi$ がコサイクルで $\psi=\delta\eta$ なら $\varphi\smile\delta\eta=(-1)^p\delta(\varphi\smile\eta)$ はコバウンダリであり、同様に $\delta\eta\smile\psi=\delta(\eta\smile\psi)$ もコバウンダリである。よって $[\varphi]\smile[\psi]=[\varphi\smile\psi]$ は代表元の取り方によらず定まる。双線形性はコチェインレベルの定義式から明らかである。
\end{solution}

\begin{problem}[難易度 B]
カップ積について次を示せ。
(1) コチェインレベルで結合的である：$(\varphi\smile\psi)\smile\chi=\varphi\smile(\psi\smile\chi)$。
(2) すべての特異0単体に値 $1$ をとる0コチェイン $1$ はコサイクルであり、$H^*(X;R)$ の単位元を与える。
(3) 連続写像 $f\colon X\to Y$ に対し $f^*(\alpha\smile\beta)=f^*\alpha\smile f^*\beta$ が成り立つ。特に、ホモトピー同値写像はコホモロジー環の同型を誘導する。
\end{problem}

\begin{solution}
(1) $\varphi\in C^p$、$\psi\in C^q$、$\chi\in C^r$ とし、$\sigma\colon\Delta^{p+q+r}\to X$ とすると、両辺とも
\[
\varphi(\sigma|[v_0,\dots,v_p])\,\psi(\sigma|[v_p,\dots,v_{p+q}])\,\chi(\sigma|[v_{p+q},\dots,v_{p+q+r}])
\]
に等しい。よって結合的である。

(2) 1単体 $\tau$ に対し $(\delta 1)(\tau)=1(\tau|[v_1])-1(\tau|[v_0])=1-1=0$ なので $1$ はコサイクルである。また $(1\smile\varphi)(\sigma)=1(\sigma|[v_0])\varphi(\sigma)=\varphi(\sigma)$、$(\varphi\smile1)(\sigma)=\varphi(\sigma)1(\sigma|[v_p])=\varphi(\sigma)$ なので $[1]$ は単位元である。

(3) $f^\sharp\varphi=\varphi\circ f_\sharp$、すなわち $(f^\sharp\varphi)(\sigma)=\varphi(f\circ\sigma)$ である。$(f\circ\sigma)|[v_0,\dots,v_p]=f\circ(\sigma|[v_0,\dots,v_p])$ などに注意すると
\[
(f^\sharp\varphi\smile f^\sharp\psi)(\sigma)=\varphi(f\circ\sigma|[v_0,\dots,v_p])\,\psi(f\circ\sigma|[v_p,\dots,v_{p+q}])=f^\sharp(\varphi\smile\psi)(\sigma)
\]
となる。コホモロジーに移れば $f^*(\alpha\smile\beta)=f^*\alpha\smile f^*\beta$ である。また $f^\sharp1=1$ なので $f^*$ は単位元を保つ環準同型である。$f$ がホモトピー同値写像ならば $f^*$ は加法群として同型で、環準同型でもあるから環同型である。
\end{solution}

\begin{problem}[難易度 B]
(1) $n\ge1$ のとき $H^*(S^n;\Z)\cong\Z[x]/(x^2)$（$\deg x=n$）を示せ。
(2) $X=X_1\sqcup X_2$（非交和）のとき、包含写像 $i_k\colon X_k\to X$ が誘導する $(i_1^*,i_2^*)\colon H^*(X;R)\to H^*(X_1;R)\times H^*(X_2;R)$ は環同型であることを示せ。これを用いて $H^*(S^0;\Z)$ を求めよ。
\end{problem}

\begin{solution}
(1) $H^k(S^n;\Z)$ は $k=0,n$ のとき $\Z$、それ以外で $0$ である。$H^0$ の生成元は単位元 $1$ で、$H^n$ の生成元を $x$ とする。$x\smile x\in H^{2n}(S^n;\Z)=0$ なので $x^2=0$ であり、$1\smile x=x$。したがって $\Z[x]\to H^*(S^n;\Z)$、$x\mapsto x$ は全射環準同型で、核は $(x^2)$ である。よって $H^*(S^n;\Z)\cong\Z[x]/(x^2)$。

(2) 単体 $\Delta^k$ は弧状連結なので、特異単体 $\sigma\colon\Delta^k\to X$ の像は $X_1,X_2$ のどちらか一方に含まれる。よって $C_k(X)=C_k(X_1)\oplus C_k(X_2)$ であり、$C^k(X;R)=C^k(X_1;R)\times C^k(X_2;R)$ はコチェイン複体の同型である。したがって $(i_1^*,i_2^*)$ は加法群の同型である。各 $i_k^*$ は前問 (3) により環準同型なので、$(i_1^*,i_2^*)$ は直積環への環準同型であり、全単射だから環同型である。

$S^0$ は2点の非交和で、点のコホモロジー環は $R$（次数0）なので $H^*(S^0;\Z)\cong\Z\times\Z$（次数0に集中）である。
\end{solution}

\begin{problem}[難易度 A]
カップ積の次数付き可換性 $\alpha\smile\beta=(-1)^{pq}\beta\smile\alpha$（$\alpha\in H^p(X;R)$、$\beta\in H^q(X;R)$）を認めて、次を示せ。
(1) $p$ が奇数ならば $2(\alpha\smile\alpha)=0$ である。特に $H^{2p}(X;R)$ に位数2の元がなければ $\alpha\smile\alpha=0$ である。
(2) 偶数次の元全体 $\bigoplus_k H^{2k}(X;R)$ は可換な部分環である。
(3) $R=\Z/2$ のとき $H^*(X;\Z/2)$ は可換環である。
\end{problem}

\begin{solution}
(1) $\beta=\alpha$、$q=p$ とおくと $(-1)^{p^2}=-1$ なので $\alpha\smile\alpha=-\alpha\smile\alpha$、すなわち $2(\alpha\smile\alpha)=0$ である。$H^{2p}(X;R)$ に位数2の元がなければ、$2\gamma=0$ となる $\gamma\in H^{2p}(X;R)$ は $\gamma=0$ に限るから $\alpha\smile\alpha=0$ である。たとえば $R=\Z$ で $H^{2p}(X;\Z)$ が自由加群のときや、$R=\mathbb{Q}$ のときにこれが成り立つ。

(2) 偶数次の元の積は偶数次であり、単位元 $1$ は次数0なので部分環である。$p,q$ がともに偶数なら $pq$ は偶数で $(-1)^{pq}=1$ だから、斉次元同士は可換であり、双線形性により一般の元同士も可換である。

(3) $\Z/2$ では $-1=1$ なので $(-1)^{pq}=1$ であり、斉次元同士は可換である。双線形性により任意の元同士が可換になる。
\end{solution}

\begin{problem}[難易度 C]
トーラス $T^2$ を、正方形 $[0,1]^2$ の対辺を平行移動で同一視して得られる Δ 複体とみる。頂点は1つ $v$、辺は下辺（$(0,0)\to(1,0)$）の $a$、左辺（$(0,0)\to(0,1)$）の $b$、対角線（$(0,0)\to(1,1)$）の $c$ であり、2単体は $U=[(0,0),(1,0),(1,1)]$ と $L=[(0,0),(0,1),(1,1)]$ である。単体的コチェインでカップ積を計算し、$H^*(T^2;\Z)\cong\Lambda_\Z[\alpha,\beta]$（$\deg\alpha=\deg\beta=1$ の外積代数）を示せ。
\begin{center}
\begin{tikzpicture}[scale=2.3]
  \fill[blue!5] (0,0) rectangle (1,1);
  \draw[thick,red!70!black,mid arrow] (0,0) -- (1,0);
  \draw[thick,red!70!black,mid arrow] (0,1) -- (1,1);
  \draw[thick,blue!70!black,mid arrows] (0,0) -- (0,1);
  \draw[thick,blue!70!black,mid arrows] (1,0) -- (1,1);
  \draw[thick,mid arrow] (0,0) -- (1,1);
  \node[below] at (0.5,0) {$a$}; \node[above] at (0.5,1) {$a$};
  \node[left] at (0,0.5) {$b$}; \node[right] at (1,0.5) {$b$};
  \node at (0.45,0.58) {$c$};
  \node at (0.72,0.25) {$U$}; \node at (0.25,0.72) {$L$};
  \foreach \p in {(0,0),(1,0),(0,1),(1,1)} \node[vtx] at \p {};
  \node[below left] at (0,0) {\footnotesize $v$};
\end{tikzpicture}
\zucap{トーラスの Δ 複体。4頂点はすべて $v$。$U$ の頂点順序は $(0,0),(1,0),(1,1)$、$L$ は $(0,0),(0,1),(1,1)$。カップ積 $(\varphi\smile\psi)(U)=\varphi(\text{前の辺})\psi(\text{後の辺})$ は、$U$ では「$a$ の次に $b$」、$L$ では「$b$ の次に $a$」を読む。}
\end{center}
\end{problem}

\begin{solution}
Δ 複体の単体的コホモロジーは特異コホモロジーと自然に同型であり、カップ積も同じ公式（頂点の前半と後半への制限）で計算できる。$U$ の辺は $[v_0v_1]=a$、$[v_1v_2]=b$、$[v_0v_2]=c$、$L$ の辺は $[v_0v_1]=b$、$[v_1v_2]=a$、$[v_0v_2]=c$ なので
\[
\partial U=b-c+a,\qquad \partial L=a-c+b
\]
である。頂点が1つなので $\delta\colon C^0\to C^1$ は $0$ である。1コチェイン $\varphi$ がコサイクルである条件は $\varphi(a)+\varphi(b)-\varphi(c)=0$ だから、$H^1\cong\Z^2$ で、基底として
\[
\alpha(a)=1,\ \alpha(b)=0,\ \alpha(c)=1,\qquad \beta(a)=0,\ \beta(b)=1,\ \beta(c)=1
\]
がとれる。2コチェインは $(\psi(U),\psi(L))\in\Z^2$ で決まり、コバウンダリ全体は $\{(t,t)\}$ なので、$\psi\mapsto\psi(U)-\psi(L)$ が同型 $H^2\cong\Z$ を与える。

カップ積を計算すると、$(\alpha\smile\beta)(U)=\alpha(a)\beta(b)=1$、$(\alpha\smile\beta)(L)=\alpha(b)\beta(a)=0$ なので $\alpha\smile\beta$ は $H^2$ の生成元である。また $(\alpha\smile\alpha)(U)=\alpha(a)\alpha(b)=0$、$(\alpha\smile\alpha)(L)=\alpha(b)\alpha(a)=0$ より $\alpha^2=0$、同様に $\beta^2=0$ である。さらに $(\beta\smile\alpha)(U)=0$、$(\beta\smile\alpha)(L)=\beta(b)\alpha(a)=1$ より $\beta\smile\alpha=-\alpha\smile\beta$ である（次数付き可換性と整合する）。以上より $H^*(T^2;\Z)$ は $1,\alpha,\beta,\alpha\beta$ を基底とし、関係 $\alpha^2=\beta^2=0$、$\beta\alpha=-\alpha\beta$ をもつ外積代数 $\Lambda_\Z[\alpha,\beta]$ である。
\end{solution}

\begin{problem}[難易度 C]
射影平面 $\mathbb{RP}^2$ を、正方形の頂点を $P=(0,0)$、$Q=(1,0)$、$R=(1,1)$、$S=(0,1)$ として、辺 $PQ$ と $RS$ を（$P\mapsto R$、$Q\mapsto S$）、辺 $PS$ と $RQ$ を（$P\mapsto R$、$S\mapsto Q$）同一視した空間とみる。辺 $a=[P,Q]=[R,S]$、$b=[P,S]=[R,Q]$、$c=[Q,S]$ と2単体 $U=[P,Q,S]$、$L=[R,Q,S]$ からなる Δ 複体を用いて、$H^*(\mathbb{RP}^2;\Z/2)\cong\Z/2[\alpha]/(\alpha^3)$ を示せ。
\begin{center}
\begin{tikzpicture}[scale=2.3]
  \fill[blue!5] (0,0) rectangle (1,1);
  \draw[thick,red!70!black,mid arrow] (0,0) -- (1,0);
  \draw[thick,red!70!black,mid arrow] (1,1) -- (0,1);
  \draw[thick,blue!70!black,mid arrows] (0,0) -- (0,1);
  \draw[thick,blue!70!black,mid arrows] (1,1) -- (1,0);
  \draw[thick,mid arrow] (1,0) -- (0,1);
  \node[below] at (0.5,0) {$a$}; \node[above] at (0.5,1) {$a$};
  \node[left] at (0,0.5) {$b$}; \node[right] at (1,0.5) {$b$};
  \node at (0.58,0.58) {$c$};
  \node at (0.28,0.28) {$U$}; \node at (0.72,0.72) {$L$};
  \node[vtx] at (0,0) {}; \node[vtx] at (1,1) {};
  \node[hvtx] at (1,0) {}; \node[hvtx] at (0,1) {};
  \node[below left] at (0,0) {\footnotesize $P$}; \node[below right] at (1,0) {\footnotesize $Q$};
  \node[above right] at (1,1) {\footnotesize $R$}; \node[above left] at (0,1) {\footnotesize $S$};
\end{tikzpicture}
\zucap{射影平面の Δ 複体。対辺を逆向きに貼る（中心に関する点対称）。黒丸 $P=R$、白丸 $Q=S$ の2頂点。}
\end{center}
\end{problem}

\begin{solution}
頂点は $v=P=R$ と $w=Q=S$ の2つである。係数は $\Z/2$ なので符号は無視してよい。$U$ の面は $[Q,S]=c$、$[P,S]=b$、$[P,Q]=a$、$L$ の面は $[Q,S]=c$、$[R,S]=a$、$[R,Q]=b$ なので、$\partial U=\partial L=a+b+c$ であり、$\partial a=\partial b=v+w$、$\partial c=0$ である。

1コチェイン $\varphi$ がコサイクルである条件は $\varphi(a)+\varphi(b)+\varphi(c)=0$ で、これは2次元の空間をなす。0コチェイン $f$ の余境界は $\delta f(a)=\delta f(b)=f(v)+f(w)$、$\delta f(c)=0$ なので、コバウンダリは1次元である。よって $H^1(\mathbb{RP}^2;\Z/2)\cong\Z/2$ で、生成元として $\alpha(a)=1$、$\alpha(b)=0$、$\alpha(c)=1$ がとれる（$\alpha(c)\neq0$ なのでコバウンダリではない）。2コチェインは $(\psi(U),\psi(L))$ で決まり、コバウンダリは $(t,t)$ の形なので、$\psi\mapsto\psi(U)+\psi(L)$ が同型 $H^2(\mathbb{RP}^2;\Z/2)\cong\Z/2$ を与える。

$U$ の頂点順序は $P,Q,S$、$L$ の頂点順序は $R,Q,S$ だから
\[
(\alpha\smile\alpha)(U)=\alpha([P,Q])\alpha([Q,S])=\alpha(a)\alpha(c)=1,\qquad (\alpha\smile\alpha)(L)=\alpha(b)\alpha(c)=0
\]
となり、$(\alpha\smile\alpha)(U)+(\alpha\smile\alpha)(L)=1$。よって $\alpha^2$ は $H^2$ の生成元である。$\alpha^3\in H^3=0$ であるから、$H^*(\mathbb{RP}^2;\Z/2)=\Z/2\{1,\alpha,\alpha^2\}\cong\Z/2[\alpha]/(\alpha^3)$ である。
\end{solution}

\begin{problem}[難易度 C]
$\sigma\colon\Delta^n\to X$ と $\varphi\in C^p(X;R)$（$p\le n$）に対し、キャップ積を $\sigma\frown\varphi=\varphi(\sigma|[v_0,\dots,v_p])\,\sigma|[v_p,\dots,v_n]\in C_{n-p}(X;R)$ で定め、線形に拡張する。
(1) $\psi\in C^{n-p}(X;R)$ に対し $\psi(\sigma\frown\varphi)=(\varphi\smile\psi)(\sigma)$ を示せ。
(2) キャップ積がホモロジー上の写像 $H_n\times H^p\to H_{n-p}$ を誘導することを認めて、本章のトーラスの問題と同じ Δ 複体（辺 $a,b,c$、2単体 $U,L$、コホモロジー類 $\alpha,\beta$ は同じ記号）において、基本類 $[T]=[U-L]$ に対し $[T]\frown\alpha$ と $[T]\frown\beta$ を求めよ。
\end{problem}

\begin{solution}
(1) 定義と線形性より
\[
\psi(\sigma\frown\varphi)=\varphi(\sigma|[v_0,\dots,v_p])\,\psi(\sigma|[v_p,\dots,v_n])=(\varphi\smile\psi)(\sigma)
\]
である。コホモロジーに移ると $\langle\psi,x\frown\varphi\rangle=\langle\varphi\smile\psi,x\rangle$ が成り立つ。

(2) $U$ の辺は $[v_0v_1]=a$、$[v_1v_2]=b$、$L$ の辺は $[v_0v_1]=b$、$[v_1v_2]=a$ であった。また $\alpha(a)=1$、$\alpha(b)=0$、$\beta(a)=0$、$\beta(b)=1$ である。よって
\begin{align*}
U\frown\alpha&=\alpha(a)\,b=b, & L\frown\alpha&=\alpha(b)\,a=0,\\
U\frown\beta&=\beta(a)\,b=0, & L\frown\beta&=\beta(b)\,a=a
\end{align*}
であり、$[T]\frown\alpha=[b]$、$[T]\frown\beta=-[a]$ となる。ここで $H_1(T^2)$ は $[a],[b]$ を基底とする自由加群である（$\partial U=a+b-c$ より $[c]=[a]+[b]$）。検算として (1) を用いると $\langle\beta,[T]\frown\alpha\rangle=\beta(b)=1=\langle\alpha\smile\beta,[T]\rangle$ となり、前問の計算と一致する。このように $[T]\frown\cdot\colon H^1\to H_1$ は同型であり、これはポアンカレ双対性の一例である。
\end{solution}

\begin{problem}[難易度 B]
$X,Y$ を弧状連結な空間とし、基点 $x_0\in X$、$y_0\in Y$ はそれぞれある近傍の変形レトラクトであるとする。
(1) $k\ge1$ に対し、$\alpha\in H^k(X;R)$、$\beta\in H^l(Y;R)$（$l\ge1$）を潰す写像 $r_X\colon X\vee Y\to X$、$r_Y\colon X\vee Y\to Y$ で引き戻した類の積は $r_X^*\alpha\smile r_Y^*\beta=0$ であることを示せ。これから $\widetilde H^*(X\vee Y;R)\cong\widetilde H^*(X;R)\times\widetilde H^*(Y;R)$（直積環）を説明せよ。
(2) $S^2\vee S^1\vee S^1$ と $T^2$ は同型なコホモロジー群をもつが、ホモトピー同値でないことを示せ。
\end{problem}

\begin{solution}
(1) 仮定から、$k\ge1$ で $(i_X^*,i_Y^*)\colon H^k(X\vee Y;R)\to H^k(X;R)\oplus H^k(Y;R)$ は同型である（$i_X,i_Y$ は包含写像）。$r_X\circ i_X=\id$、$r_X\circ i_Y$ は定値写像なので、$i_X^*r_X^*\alpha=\alpha$、$i_Y^*r_X^*\alpha=0$ である（定値写像は正次数で $0$ を誘導する）。同様に $i_X^*r_Y^*\beta=0$ である。よって自然性より
\[
i_X^*(r_X^*\alpha\smile r_Y^*\beta)=\alpha\smile0=0,\qquad i_Y^*(r_X^*\alpha\smile r_Y^*\beta)=0\smile\beta=0
\]
となり、$k+l\ge1$ での単射性から $r_X^*\alpha\smile r_Y^*\beta=0$ である。また $r_X^*$ は環準同型なので $r_X^*\alpha\smile r_X^*\alpha'=r_X^*(\alpha\smile\alpha')$。したがって正次数部分において、同型 $(\alpha,\beta)\mapsto r_X^*\alpha+r_Y^*\beta$ は積を成分ごとに保ち、$\widetilde H^*(X\vee Y;R)$ は直積環 $\widetilde H^*(X;R)\times\widetilde H^*(Y;R)$ と同型である。

(2) どちらも $H^0=\Z$、$H^1=\Z^2$、$H^2=\Z$ である。$W=S^2\vee S^1\vee S^1$ において、$H^1(W)$ は2つの $S^1$ の生成元の引き戻し $x_1,x_2$ で生成される。(1) より $x_1x_2=x_2x_1=0$ であり、$x_i^2$ は $H^2(S^1)=0$ の元の引き戻しなので $0$ である。よって積 $H^1(W)\times H^1(W)\to H^2(W)$ は恒等的に $0$ である。一方、$T^2$ では $\alpha\smile\beta$ が $H^2$ の生成元である。ホモトピー同値写像はコホモロジー環の同型を誘導するので、両者はホモトピー同値でない。
\end{solution}

\begin{problem}[難易度 B]
$X$ の懸垂を $SX=X\times[-1,1]/(X\times\{1\},\,X\times\{-1\})$ とする。開集合 $A,B\subset Z$ に対する相対カップ積 $H^k(Z,A;R)\times H^l(Z,B;R)\to H^{k+l}(Z,A\cup B;R)$ が存在し、$H^*(Z;R)$ への写像と可換であることを認めて、$SX$ では正次数の任意の2元のカップ積が $0$ になることを示せ。これを用いて $T^2$ も $\mathbb{RP}^2$ も懸垂とホモトピー同値でないことを示せ。
\end{problem}

\begin{solution}
$A$ を $X\times(-1/2,1]$ の像、$B$ を $X\times[-1,1/2)$ の像とすると、$A,B$ は $SX$ の開集合で $A\cup B=SX$ であり、それぞれ錐の頂点に変形レトラクトするので可縮である。

$\alpha\in H^k(SX;R)$、$k\ge1$ とする。対 $(SX,A)$ の長完全列
\[
H^k(SX,A;R)\xrightarrow{j^*}H^k(SX;R)\to H^k(A;R)=0
\]
より $\alpha=j^*\tilde\alpha$ となる $\tilde\alpha\in H^k(SX,A;R)$ がある。同様に $\beta\in H^l(SX;R)$、$l\ge1$ に対し $\beta=j^*\tilde\beta$、$\tilde\beta\in H^l(SX,B;R)$ がとれる。相対カップ積との可換性から
\[
\alpha\smile\beta=j^*(\tilde\alpha\smile\tilde\beta),\qquad \tilde\alpha\smile\tilde\beta\in H^{k+l}(SX,A\cup B;R)=H^{k+l}(SX,SX;R)=0
\]
なので $\alpha\smile\beta=0$ である。

コホモロジー環はホモトピー不変なので、懸垂とホモトピー同値な空間でも正次数の積はすべて消える。しかし $T^2$ では $\alpha\smile\beta\neq0$（$\Z$ 係数）、$\mathbb{RP}^2$ では $\alpha\smile\alpha\neq0$（$\Z/2$ 係数）であった。よってどちらも懸垂とホモトピー同値でない。
\end{solution}

\begin{problem}[難易度 C]
$g\ge1$ とし、$\Sigma_g$ を種数 $g$ の向き付け可能閉曲面とする。次の事実を認める：$\Sigma_g$ から $g$ 個のトーラスの楔和 $W=T_1\vee\dots\vee T_g$ への「つまみ写像」$q$ があって、$q^*\colon H^1(W;\Z)\to H^1(\Sigma_g;\Z)$ は同型であり、各 $i$ について合成 $H^2(T_i;\Z)\xrightarrow{r_i^*}H^2(W;\Z)\xrightarrow{q^*}H^2(\Sigma_g;\Z)\cong\Z$ は同型である（$r_i\colon W\to T_i$ は他の成分を潰す写像）。このとき $H^1(\Sigma_g;\Z)$ の基底 $a_1,b_1,\dots,a_g,b_g$ と $H^2$ の生成元 $\omega$ で
\[
a_i\smile b_i=\omega=-b_i\smile a_i,\qquad a_i\smile a_j=b_i\smile b_j=0,\qquad a_i\smile b_j=0\ (i\neq j)
\]
となるものが存在することを示せ。
\end{problem}

\begin{solution}
$H^2(\Sigma_g;\Z)$ の生成元 $\omega$ を1つ固定する。トーラス $T_i$ の $H^1$ の基底 $\alpha_i,\beta_i$ を、$\alpha_i\smile\beta_i$ が $H^2(T_i)$ の生成元となるようにとる（トーラスの問題）。仮定より $q^*r_i^*(\alpha_i\smile\beta_i)=\varepsilon_i\omega$、$\varepsilon_i=\pm1$ である。そこで
\[
a_i=q^*r_i^*\alpha_i,\qquad b_i=\varepsilon_i\,q^*r_i^*\beta_i
\]
とおく。楔和の問題から $H^1(W)$ は $r_i^*\alpha_i,r_i^*\beta_i$（$i=1,\dots,g$）を基底とし、$q^*$ は $H^1$ で同型なので、$a_i,b_i$ は $H^1(\Sigma_g)$ の基底である。

自然性より $a_i\smile b_i=\varepsilon_iq^*r_i^*(\alpha_i\smile\beta_i)=\varepsilon_i^2\omega=\omega$ であり、次数付き可換性より $b_i\smile a_i=-\omega$。また $a_i\smile a_i=q^*r_i^*(\alpha_i^2)=0$、$b_i\smile b_i=0$ である。$i\neq j$ のとき、$a_i\smile a_j$、$a_i\smile b_j$、$b_i\smile b_j$ はいずれも $q^*(r_i^*x\smile r_j^*y)$ の形（符号を除く）で、楔和の問題により $r_i^*x\smile r_j^*y=0$ である。以上で主張が示された。特に $H^1\times H^1\to H^2\cong\Z$ は非退化な交代形式（シンプレクティック形式）である。
\end{solution}

\begin{problem}[難易度 C]
$A=(a_{ij})$ を整数成分の $2\times2$ 行列とし、$f_A\colon T^2=\R^2/\Z^2\to T^2$ を $A$ が誘導する写像とする。$f_A^*\colon H^2(T^2;\Z)\to H^2(T^2;\Z)$ は $\det A$ 倍写像であることを、カップ積を用いて示せ。ただし $H^1(X;\Z)\cong\operatorname{Hom}(H_1(X),\Z)$（自然な同型）を用いてよい。
\end{problem}

\begin{solution}
$e_1,e_2$ を $\R^2$ の標準基底とし、$\gamma_j(t)=te_j$（$0\le t\le1$）の像を $T^2$ のループとみる。$H_1(T^2)\cong\Z^2$ は $[\gamma_1],[\gamma_2]$ を基底とする。$f_A\circ\gamma_j$ は $\R^2$ における $0$ から $Ae_j=\sum_ia_{ij}e_i$ への線分の射影であり、これは端点を止めて $e_1$ 方向に $a_{1j}$、$e_2$ 方向に $a_{2j}$ 進む折れ線の射影とホモトピックである。よって $(f_A)_*[\gamma_j]=\sum_ia_{ij}[\gamma_i]$ である。

$x_1,x_2\in H^1(T^2;\Z)$ を $[\gamma_1],[\gamma_2]$ の双対基底とすると、$f_A^*x_i=x_i\circ(f_A)_*$ であり $(f_A^*x_i)([\gamma_j])=a_{ij}$ なので、$f_A^*x_i=a_{i1}x_1+a_{i2}x_2$ である。トーラスの問題の $\alpha,\beta$ は $x_1,x_2$ に一致し（$a=\gamma_1$、$b=\gamma_2$）、$x_1\smile x_2$ は $H^2$ の生成元である。$x_1^2=x_2^2=0$、$x_2x_1=-x_1x_2$ を用いると
\begin{align*}
f_A^*(x_1\smile x_2)&=(a_{11}x_1+a_{12}x_2)\smile(a_{21}x_1+a_{22}x_2)\\
&=a_{11}a_{22}\,x_1x_2+a_{12}a_{21}\,x_2x_1=(a_{11}a_{22}-a_{12}a_{21})\,x_1x_2
\end{align*}
となる。よって $f_A^*$ は $H^2$ 上で $\det A$ 倍である。すなわち $f_A$ の写像度は $\det A$ である。
\end{solution}

\begin{problem}[難易度 B]
$H^*(\mathbb{CP}^n;\Z)\cong\Z[x]/(x^{n+1})$（$\deg x=2$）を認める。
(1) $\mathbb{CP}^2$ と $S^2\vee S^4$ は同型なホモロジー群をもつが、ホモトピー同値でないことを示せ。
(2) 写像 $f\colon\mathbb{CP}^n\to\mathbb{CP}^n$ に対し $f^*x=kx$ とすると、$f$ の写像度は $k^n$ であることを示せ。特に $n$ が偶数のとき、写像度 $-1$ の自己写像は存在しないことを示せ。
\end{problem}

\begin{solution}
(1) どちらもホモロジー群は次数 $0,2,4$ で $\Z$、それ以外で $0$ である。$\mathbb{CP}^2$ では $H^2$ の生成元 $x$ の平方 $x^2$ が $H^4$ の生成元である。一方 $S^2\vee S^4$ では、$H^2$ の生成元は $y=r^*\sigma_2$（$r$ は $S^4$ を潰す写像、$\sigma_2$ は $H^2(S^2)$ の生成元）であり、$y^2=r^*(\sigma_2^2)=0$（$H^4(S^2)=0$）である。$H^2$ の任意の元 $my$ の平方も $0$ である。環同型は $H^2$ の生成元を生成元に移すので、両者のコホモロジー環は同型でなく、ホモトピー同値でない。

(2) $f^*$ は環準同型なので $f^*(x^n)=(f^*x)^n=k^nx^n$。$x^n$ は $H^{2n}(\mathbb{CP}^n;\Z)\cong\Z$ の生成元だから、$f^*$ は最高次で $k^n$ 倍であり、写像度は $k^n$ である（ホモロジーの $f_*$ も $H^{2n}=\operatorname{Hom}(H_{2n},\Z)$ を通じて同じ倍率になる）。$n$ が偶数なら $k^n\ge0$ なので写像度 $-1$ は実現できない。たとえば $\mathbb{CP}^2$ には向きを逆にするホモトピー同値写像が存在しない。なお複素共役は $x\mapsto -x$ を誘導し、写像度 $(-1)^n$ をもつ。
\end{solution}

\begin{problem}[難易度 B]
$H^*(\mathbb{RP}^n;\Z/2)\cong\Z/2[\alpha]/(\alpha^{n+1})$（$\deg\alpha=1$）を認める。$1\le m<n$ とする。
(1) 任意の写像 $f\colon\mathbb{RP}^n\to\mathbb{RP}^m$ は $H^1(\,\cdot\,;\Z/2)$ 上で $0$ を誘導することを示せ。
(2) 標準的な包含 $\mathbb{RP}^m\subset\mathbb{RP}^n$ に対し、$\mathbb{RP}^m$ は $\mathbb{RP}^n$ のレトラクトでないことを示せ。ただし包含写像 $i$ は $i^*\alpha_n=\alpha_m$ をみたす。
\end{problem}

\begin{solution}
$\alpha_m,\alpha_n$ をそれぞれの $H^1(\,\cdot\,;\Z/2)\cong\Z/2$ の生成元とする。

(1) $f^*\alpha_m=c\,\alpha_n$（$c\in\Z/2$）とおく。$\alpha_m^{m+1}=0$ なので
\[
0=f^*(\alpha_m^{m+1})=(f^*\alpha_m)^{m+1}=c^{m+1}\alpha_n^{m+1}=c\,\alpha_n^{m+1}
\]
である。$m+1\le n$ より $\alpha_n^{m+1}\neq0$ だから $c=0$、すなわち $f^*=0$ である。

(2) レトラクション $r\colon\mathbb{RP}^n\to\mathbb{RP}^m$（$r\circ i=\id$）が存在したとすると、$i^*r^*=\id$ なので $i^*r^*\alpha_m=\alpha_m\neq0$、よって $r^*\alpha_m\neq0$ である。これは (1) に矛盾する。したがってレトラクトでない。なお $m=0$ のときは1点なのでレトラクトになり、仮定 $m\ge1$ は必要である。また $m\ge2$ のとき $i_*\colon H_1(\mathbb{RP}^m)\to H_1(\mathbb{RP}^n)$ は $\Z/2$ の同型なので、$H_1$ を見るだけではこの結論は得られないことに注意する。
\end{solution}

\begin{problem}[難易度 C]
$m,n\ge1$ とする。$H^*(S^m\times S^n;\Z)$ は $a=p_1^*\sigma_m$、$b=p_2^*\sigma_n$（$\sigma_k$ は $H^k(S^k)$ の生成元、$p_i$ は射影）で生成され、$a^2=b^2=0$ で $a\smile b$ が $H^{m+n}$ の生成元であることを認める（第8章のキュネット公式）。同様に $H^*(T^n;\Z)\cong\Lambda_\Z[x_1,\dots,x_n]$ を認める。
(1) 任意の写像 $f\colon S^{m+n}\to S^m\times S^n$ は $H^{m+n}$ 上で $0$ を誘導することを示せ。
(2) $n\ge2$ のとき、任意の写像 $S^n\to T^n$ の写像度は $0$ であることを示せ。
(3) $S^m\times S^n$ と $S^m\vee S^n\vee S^{m+n}$ はホモトピー同値でないことを示せ。
\end{problem}

\begin{solution}
(1) $0<m<m+n$ なので $H^m(S^{m+n};\Z)=0$ であり、$f^*a=0$ である。よって $f^*(a\smile b)=f^*a\smile f^*b=0$。$a\smile b$ は $H^{m+n}(S^m\times S^n)$ の生成元なので $f^*=0$ である。

(2) $H^n(T^n;\Z)\cong\Z$ は $x_1\smile\cdots\smile x_n$ で生成される。$g\colon S^n\to T^n$ に対し、$n\ge2$ より $H^1(S^n)=0$ なので $g^*x_i=0$、よって $g^*(x_1\cdots x_n)=0$ であり、写像度は $0$ である。同じ議論で、$k\ge2$ なら任意の写像 $S^k\to T^n$ は正次数のコホモロジーに $0$ を誘導する（$H^k(T^n)$ は1次の元の積で生成されるから）。

(3) どちらもホモロジー群は同型である（次数 $0,m,n,m+n$ に $\Z$、$m=n$ なら次数 $m$ に $\Z^2$）。$S^m\times S^n$ では積 $H^m\times H^n\to H^{m+n}$ は $a\smile b\neq0$ で非自明である。一方 $W=S^m\vee S^n\vee S^{m+n}$ では、$H^m(W)$、$H^n(W)$ の元は $S^m$、$S^n$ の成分からの引き戻しの和である。異なる成分からの積は楔和の問題により $0$、同じ成分 $S^k$ からの積は $H^{m+n}(S^k)=0$（$k<m+n$）の元の引き戻しで $0$ である。よって $H^m(W)\times H^n(W)\to H^{m+n}(W)$ は零写像で、両者のコホモロジー環は同型でない。
\end{solution}

\begin{problem}[難易度 D]
$n$ を正の偶数とする。$S^n\times S^n$ のコホモロジー環について前問の記述を認めて、連続写像 $\mu\colon S^n\times S^n\to S^n$ で、基点 $e$ に対し $\mu(\,\cdot\,,e)$ と $\mu(e,\,\cdot\,)$ がともに恒等写像とホモトピックとなるもの（$H$ 空間の積）は存在しないことを示せ。
\end{problem}

\begin{solution}
$\sigma$ を $H^n(S^n;\Z)$ の生成元、$a=p_1^*\sigma$、$b=p_2^*\sigma$ とする。$H^n(S^n\times S^n)$ は $a,b$ を基底とする自由加群なので $\mu^*\sigma=\lambda a+\lambda'b$ と書ける。$i_1(x)=(x,e)$、$i_2(x)=(e,x)$ とおくと、$p_1\circ i_1=\id$、$p_2\circ i_1$ は定値なので $i_1^*a=\sigma$、$i_1^*b=0$ である。仮定 $\mu\circ i_1\simeq\id$ より
\[
\sigma=(\mu\circ i_1)^*\sigma=i_1^*(\lambda a+\lambda'b)=\lambda\sigma
\]
となり $\lambda=1$。同様に $i_2$ を用いて $\lambda'=1$ である。よって $\mu^*\sigma=a+b$。

$\sigma^2\in H^{2n}(S^n)=0$ なので $\mu^*(\sigma^2)=0$ である。一方、$n$ が偶数なので $ba=(-1)^{n^2}ab=ab$ であり、$a^2=b^2=0$ から
\[
(\mu^*\sigma)^2=(a+b)^2=a^2+ab+ba+b^2=2ab
\]
となる。$ab$ は $H^{2n}(S^n\times S^n)\cong\Z$ の生成元なので $2ab\neq0$ であり、$\mu^*(\sigma^2)=(\mu^*\sigma)^2$ に矛盾する。よってそのような $\mu$ は存在しない。なお $n$ が奇数なら $ba=-ab$ で $(a+b)^2=0$ となり矛盾は生じない（実際 $S^1,S^3,S^7$ は $H$ 空間である）。
\end{solution}

%================================================
\chapter{普遍係数定理とKunneth公式}
% 15問
%================================================

\begin{summary}
\begin{itemize}
  \item \textbf{Tor と Ext}：アーベル群 $A$ の自由分解 $0\to F_1\to F_0\to A\to0$ をとり、
  \begin{align*}
  \operatorname{Tor}(A,G)&=\ker(F_1\otimes G\to F_0\otimes G),\\
  \operatorname{Ext}(A,G)&=\operatorname{coker}(\operatorname{Hom}(F_0,G)\to\operatorname{Hom}(F_1,G)).
  \end{align*}
  \item \textbf{計算規則}：直和について加法的。$\operatorname{Tor}(\Z,G)=\operatorname{Ext}(\Z,G)=0$、$\operatorname{Tor}(\Z/n,G)=\{g\in G\mid ng=0\}$、$\operatorname{Ext}(\Z/n,G)=G/nG$。特に $\operatorname{Tor}(\Z/m,\Z/n)\cong\operatorname{Ext}(\Z/m,\Z/n)\cong\Z/\gcd(m,n)$。
  \item \textbf{普遍係数定理（ホモロジー）}：分裂する（ただし自然でない）短完全列
  \[
  0\to H_n(X)\otimes G\to H_n(X;G)\to\operatorname{Tor}(H_{n-1}(X),G)\to0.
  \]
  \item \textbf{普遍係数定理（コホモロジー）}：分裂する（自然でない）短完全列
  \[
  0\to\operatorname{Ext}(H_{n-1}(X),G)\to H^n(X;G)\to\operatorname{Hom}(H_n(X),G)\to0.
  \]
  $H_n(X)\cong\Z^{b_n}\oplus T_n$（$T_n$ は有限）なら $H^n(X;\Z)\cong\Z^{b_n}\oplus T_{n-1}$（ねじれが1次ずれる）。体 $F$ 係数では $H^n(X;F)\cong\operatorname{Hom}_F(H_n(X;F),F)$。
  \item \textbf{キュネット公式}：分裂する短完全列
  \begin{align*}
  0\to\bigoplus_{i+j=n}H_i(X)\otimes H_j(Y)&\to H_n(X\times Y)\\
  &\to\bigoplus_{i+j=n-1}\operatorname{Tor}(H_i(X),H_j(Y))\to0.
  \end{align*}
  体 $F$ 係数で各 $H^i(Y;F)$ が有限次元なら、クロス積 $a\times b=p_1^*a\smile p_2^*b$ により環同型 $H^*(X\times Y;F)\cong H^*(X;F)\otimes H^*(Y;F)$。
  \item \textbf{ポアンカレ多項式}：$P_X(t)=\sum_i b_i t^i$（$b_i$ はベッチ数）。$P_{X\times Y}=P_XP_Y$、$\chi(X)=P_X(-1)$。
\end{itemize}
\end{summary}

\begin{problem}[難易度 B]
アーベル群 $A$ の自由分解 $0\to F_1\xrightarrow{\partial}F_0\to A\to0$ に対し、$\operatorname{Tor}(A,G)=\ker(\partial\otimes1\colon F_1\otimes G\to F_0\otimes G)$、$\operatorname{Ext}(A,G)=\operatorname{coker}(\partial^*\colon\operatorname{Hom}(F_0,G)\to\operatorname{Hom}(F_1,G))$ と定める（分解の取り方によらない）。$m\ge1$ とし、$G[m]=\{g\in G\mid mg=0\}$ とおく。
(1) $\operatorname{Tor}(\Z/m,G)\cong G[m]$、$\operatorname{Ext}(\Z/m,G)\cong G/mG$ を示せ。
(2) $d=\gcd(m,n)$ とするとき、$\operatorname{Hom}(\Z/m,\Z/n)$、$\Z/m\otimes\Z/n$、$\operatorname{Tor}(\Z/m,\Z/n)$、$\operatorname{Ext}(\Z/m,\Z/n)$ はすべて $\Z/d$ と同型であることを示せ。
(3) $G=\Z$ および $G=\mathbb{Q}$ のとき、$\operatorname{Tor}(\Z/m,G)$ と $\operatorname{Ext}(\Z/m,G)$ を求めよ。
\end{problem}

\begin{solution}
(1) 自由分解 $0\to\Z\xrightarrow{m}\Z\to\Z/m\to0$ を用いる。$\Z\otimes G\cong G$ なので $\partial\otimes1$ は $G\xrightarrow{m}G$ となり、$\operatorname{Tor}(\Z/m,G)=\ker(m)=G[m]$、また $\Z/m\otimes G\cong\operatorname{coker}(m)=G/mG$ である。同様に $\operatorname{Hom}(\Z,G)\cong G$ なので $\partial^*$ も $G\xrightarrow{m}G$ となり、$\operatorname{Hom}(\Z/m,G)\cong G[m]$、$\operatorname{Ext}(\Z/m,G)\cong G/mG$ である。

(2) $G=\Z/n$ とする。$mx\equiv0\pmod n$ は $(m/d)x\equiv0\pmod{n/d}$ と同値で、$\gcd(m/d,n/d)=1$ より $x\equiv0\pmod{n/d}$ と同値である。よって $G[m]=(n/d)\Z/n\Z\cong\Z/d$。また $mG$ は $\gcd(m,n)=d$ が生成する部分群 $d\Z/n\Z$ なので $G/mG\cong\Z/d$。(1) の計算から4つとも $\Z/d$ と同型である。

(3) $G=\Z$ のとき $\Z[m]=0$、$\Z/m\Z$ より $\operatorname{Tor}(\Z/m,\Z)=0$、$\operatorname{Ext}(\Z/m,\Z)\cong\Z/m$ である。$G=\mathbb{Q}$ のとき、$\mathbb{Q}$ はねじれをもたないので $\mathbb{Q}[m]=0$、また任意の $r\in\mathbb{Q}$ は $r=m(r/m)$ なので $m\mathbb{Q}=\mathbb{Q}$ である。よって $\operatorname{Tor}(\Z/m,\mathbb{Q})=\operatorname{Ext}(\Z/m,\mathbb{Q})=0$ である。
\end{solution}

\begin{problem}[難易度 B]
$A=\Z^r\oplus\Z/m_1\oplus\dots\oplus\Z/m_s$ を有限生成アーベル群とし、$p$ を素数とする。$\operatorname{Hom},\otimes,\operatorname{Tor},\operatorname{Ext}$ が第1変数の有限直和を直和に移すことを用いて、次を示せ。
(1) $\operatorname{Hom}(A,\Z)\cong\Z^r$、$\operatorname{Ext}(A,\Z)\cong\Z/m_1\oplus\dots\oplus\Z/m_s$。
(2) $\dim_{\Z/p}(A\otimes\Z/p)-\dim_{\Z/p}\operatorname{Tor}(A,\Z/p)=r$。
(3) $A=\Z^2\oplus\Z/4\oplus\Z/6$ のとき、$p=2,3,5$ について $A\otimes\Z/p$ と $\operatorname{Tor}(A,\Z/p)$ を求めよ。
\end{problem}

\begin{solution}
$\Z$ は自由分解 $0\to0\to\Z\to\Z\to0$ をもつので、任意の $G$ について $\operatorname{Tor}(\Z,G)=\operatorname{Ext}(\Z,G)=0$ であり、$\operatorname{Hom}(\Z,G)\cong\Z\otimes G\cong G$ である。

(1) $\operatorname{Hom}(\Z/m,\Z)=0$ なので $\operatorname{Hom}(A,\Z)\cong\operatorname{Hom}(\Z,\Z)^r=\Z^r$。前問より $\operatorname{Ext}(\Z/m_i,\Z)\cong\Z/m_i$ なので $\operatorname{Ext}(A,\Z)\cong\bigoplus_i\Z/m_i$、すなわち $A$ のねじれ部分と同型である。

(2) 前問 (2) より $\Z/m_i\otimes\Z/p\cong\operatorname{Tor}(\Z/m_i,\Z/p)\cong\Z/\gcd(m_i,p)$ である。これは $p\mid m_i$ なら $\Z/p$、そうでなければ $0$ である。$c$ を $p\mid m_i$ となる $i$ の個数とすると
\[
A\otimes\Z/p\cong(\Z/p)^{r+c},\qquad \operatorname{Tor}(A,\Z/p)\cong(\Z/p)^{c}
\]
となり、次元の差は $r$ である。

(3) $r=2$、$m_1=4$、$m_2=6$ である。$p=2$ では両方の $m_i$ が $2$ で割れるので $c=2$ で、$A\otimes\Z/2\cong(\Z/2)^4$、$\operatorname{Tor}(A,\Z/2)\cong(\Z/2)^2$。$p=3$ では $c=1$ で、$A\otimes\Z/3\cong(\Z/3)^3$、$\operatorname{Tor}(A,\Z/3)\cong\Z/3$。$p=5$ では $c=0$ で、$A\otimes\Z/5\cong(\Z/5)^2$、$\operatorname{Tor}(A,\Z/5)=0$ である。
\end{solution}

\begin{problem}[難易度 A]
$X$ の整数係数ホモロジー群がすべて有限生成であるとし、$H_k(X)$ の階数（ベッチ数）を $b_k$ とする。普遍係数定理を用いて
\[
H_k(X;\mathbb{Q})\cong H_k(X)\otimes\mathbb{Q}\cong\mathbb{Q}^{b_k},\qquad H^k(X;\mathbb{Q})\cong\operatorname{Hom}(H_k(X),\mathbb{Q})\cong\mathbb{Q}^{b_k}
\]
を示せ。
\end{problem}

\begin{solution}
ホモロジーの普遍係数定理より、分裂完全列
\[
0\to H_k(X)\otimes\mathbb{Q}\to H_k(X;\mathbb{Q})\to\operatorname{Tor}(H_{k-1}(X),\mathbb{Q})\to0
\]
がある。$H_{k-1}(X)$ は有限生成なので $\Z$ と $\Z/m$ の直和であり、本章最初の問題の (3) と $\operatorname{Tor}(\Z,\mathbb{Q})=0$ から $\operatorname{Tor}(H_{k-1}(X),\mathbb{Q})=0$ である。よって $H_k(X;\mathbb{Q})\cong H_k(X)\otimes\mathbb{Q}$。さらに $\Z\otimes\mathbb{Q}\cong\mathbb{Q}$、$\Z/m\otimes\mathbb{Q}\cong\mathbb{Q}/m\mathbb{Q}=0$ なので、$H_k(X)\otimes\mathbb{Q}\cong\mathbb{Q}^{b_k}$ である。

コホモロジーの普遍係数定理より
\[
0\to\operatorname{Ext}(H_{k-1}(X),\mathbb{Q})\to H^k(X;\mathbb{Q})\to\operatorname{Hom}(H_k(X),\mathbb{Q})\to0
\]
は分裂完全である。同様に $\operatorname{Ext}(\Z/m,\mathbb{Q})=\mathbb{Q}/m\mathbb{Q}=0$、$\operatorname{Ext}(\Z,\mathbb{Q})=0$ より Ext 項は消える。また $\operatorname{Hom}(\Z/m,\mathbb{Q})=\mathbb{Q}[m]=0$、$\operatorname{Hom}(\Z,\mathbb{Q})=\mathbb{Q}$ なので $H^k(X;\mathbb{Q})\cong\mathbb{Q}^{b_k}$ である。このように有理係数ではねじれの情報がすべて失われ、ベッチ数だけが残る。
\end{solution}

\begin{problem}[難易度 A]
$\mathbb{RP}^n$ の整数係数ホモロジーは
\[
H_k(\mathbb{RP}^n)\cong\begin{cases}\Z & k=0,\ \text{または } k=n \text{ が奇数},\\ \Z/2 & k \text{ が奇数で } 0<k<n,\\ 0 & \text{それ以外}\end{cases}
\]
である。普遍係数定理を用いて $H_k(\mathbb{RP}^n;\Z/2)$ と $H_k(\mathbb{RP}^n;\mathbb{Q})$ を計算せよ。
\end{problem}

\begin{solution}
普遍係数定理 $H_k(X;G)\cong H_k(X)\otimes G\oplus\operatorname{Tor}(H_{k-1}(X),G)$ を用いる。$G=\Z/2$ のとき、$\Z\otimes\Z/2=\Z/2$、$\Z/2\otimes\Z/2=\Z/2$、$\operatorname{Tor}(\Z,\Z/2)=0$、$\operatorname{Tor}(\Z/2,\Z/2)=\Z/2$ である。$0\le k\le n$ で場合分けする。
\begin{itemize}
\item $k=0$：$H_0\otimes\Z/2=\Z/2$ なので $\Z/2$。
\item $k$ 奇数、$0<k<n$：$H_k\otimes\Z/2=\Z/2$、$H_{k-1}$ は $\Z$（$k=1$）または $0$ なので Tor 項は $0$。よって $\Z/2$。
\item $k=n$ 奇数：$H_n\otimes\Z/2=\Z/2$、$H_{n-1}$ は $0$ または $\Z$（$n=1$）で Tor 項は $0$。よって $\Z/2$。
\item $k$ 偶数、$0<k\le n$：$H_k=0$ で、$H_{k-1}=\Z/2$ なので Tor 項が $\Z/2$。よって $\Z/2$。
\end{itemize}
$k>n$ では両項とも $0$ である（$k=n+1$ のとき $H_n$ は $\Z$ か $0$ で Tor 項は $0$）。したがって $H_k(\mathbb{RP}^n;\Z/2)\cong\Z/2$（$0\le k\le n$）、それ以外は $0$ である。これは胞体鎖複体の境界写像 $0,2$ が mod 2 ですべて $0$ になることとも整合する。

$G=\mathbb{Q}$ のときは前問より $H_k(\mathbb{RP}^n;\mathbb{Q})\cong H_k(\mathbb{RP}^n)\otimes\mathbb{Q}$ なので、$k=0$ で $\mathbb{Q}$、$n$ が奇数なら $k=n$ でも $\mathbb{Q}$、それ以外は $0$ である。つまり $n$ が偶数なら $\mathbb{RP}^n$ は有理ホモロジーでは1点と区別できない。
\end{solution}

\begin{problem}[難易度 B]
前問のホモロジー群を用い、コホモロジーの普遍係数定理によって $H^k(\mathbb{RP}^n;\Z)$、$H^k(\mathbb{RP}^n;\Z/2)$、$H^k(\mathbb{RP}^n;\mathbb{Q})$ を計算せよ。
\end{problem}

\begin{solution}
$H^k(X;G)\cong\operatorname{Hom}(H_k(X),G)\oplus\operatorname{Ext}(H_{k-1}(X),G)$ を用いる。

$G=\Z$：$\operatorname{Hom}(\Z/2,\Z)=0$、$\operatorname{Ext}(\Z/2,\Z)=\Z/2$、$\operatorname{Ext}(\Z,\Z)=0$ である。$k=0$ では $\Z$。$k$ 奇数 $0<k<n$ では $\operatorname{Hom}(\Z/2,\Z)=0$ で、$H_{k-1}$ は自由なので Ext 項も $0$。$k$ 偶数 $0<k\le n$ では $\operatorname{Hom}(0,\Z)=0$、$\operatorname{Ext}(H_{k-1},\Z)=\operatorname{Ext}(\Z/2,\Z)=\Z/2$。$k=n$ 奇数では $\operatorname{Hom}(\Z,\Z)=\Z$ で Ext 項は $0$。よって
\[
H^k(\mathbb{RP}^n;\Z)\cong\begin{cases}\Z & k=0,\ \text{または } k=n \text{ が奇数},\\ \Z/2 & k \text{ が偶数で } 0<k\le n,\\ 0 & \text{それ以外}\end{cases}
\]
である。ホモロジーでは奇数次にあった $\Z/2$ が、コホモロジーでは1つ上の偶数次に現れる（ねじれのずれ）。

$G=\Z/2$：$\operatorname{Hom}(\Z,\Z/2)=\operatorname{Hom}(\Z/2,\Z/2)=\Z/2$、$\operatorname{Ext}(\Z/2,\Z/2)=\Z/2$、$\operatorname{Ext}(\Z,\Z/2)=0$。$0\le k\le n$ の各 $k$ で、$H_k\neq0$ なら Hom 項が $\Z/2$ で Ext 項は $0$（$H_{k-1}$ は $\Z$ か $0$）、$H_k=0$（$k$ 偶数 $>0$）なら Hom 項は $0$ で Ext 項が $\Z/2$ である。よって $H^k(\mathbb{RP}^n;\Z/2)\cong\Z/2$（$0\le k\le n$）、それ以外は $0$。

$G=\mathbb{Q}$：有理係数の問題より $H^k(\mathbb{RP}^n;\mathbb{Q})\cong\operatorname{Hom}(H_k,\mathbb{Q})$ で、$k=0$ と（$n$ 奇数のとき）$k=n$ で $\mathbb{Q}$、それ以外は $0$ である。
\end{solution}

\begin{problem}[難易度 B]
(1) $X$ のホモロジー群が有限生成で、$H_k(X)\cong\Z^{b_k}\oplus T_k$（$T_k$ は有限群）と表されるとき、$H^k(X;\Z)\cong\Z^{b_k}\oplus T_{k-1}$ を示せ（$T_{-1}=0$）。
(2) クラインの壺 $K$ のホモロジーは $H_0=\Z$、$H_1=\Z\oplus\Z/2$、$H_2=0$ である。$H^*(K;\Z)$ と $H^*(K;\Z/2)$ を計算せよ。
\end{problem}

\begin{solution}
(1) コホモロジーの普遍係数定理より $H^k(X;\Z)\cong\operatorname{Hom}(H_k(X),\Z)\oplus\operatorname{Ext}(H_{k-1}(X),\Z)$ である。本章2番目の問題の (1) より $\operatorname{Hom}(\Z^{b_k}\oplus T_k,\Z)\cong\Z^{b_k}$、$\operatorname{Ext}(\Z^{b_{k-1}}\oplus T_{k-1},\Z)\cong T_{k-1}$ である。よって $H^k(X;\Z)\cong\Z^{b_k}\oplus T_{k-1}$。すなわち、自由部分は同じ次数に、ねじれ部分は1つ上の次数に現れる。

(2) (1) より $H^0(K;\Z)=\Z$、$H^1(K;\Z)\cong\Z^{1}\oplus T_0=\Z$、$H^2(K;\Z)\cong\Z^0\oplus T_1=\Z/2$ であり、$k\ge3$ では $0$ である。

$\Z/2$ 係数では $H^k(K;\Z/2)\cong\operatorname{Hom}(H_k,\Z/2)\oplus\operatorname{Ext}(H_{k-1},\Z/2)$ より
\begin{align*}
H^0(K;\Z/2)&\cong\operatorname{Hom}(\Z,\Z/2)=\Z/2,\\
H^1(K;\Z/2)&\cong\operatorname{Hom}(\Z\oplus\Z/2,\Z/2)\oplus\operatorname{Ext}(\Z,\Z/2)=(\Z/2)^2,\\
H^2(K;\Z/2)&\cong\operatorname{Hom}(0,\Z/2)\oplus\operatorname{Ext}(\Z\oplus\Z/2,\Z/2)=\Z/2
\end{align*}
である。検算としてオイラー標数はどちらの係数でも $1-1+0=1-2+1=0$ となる。
\end{solution}

\begin{problem}[難易度 B]
$n\ge1$、$m\ge2$ とし、$M=S^n\cup_m e^{n+1}$ を写像度 $m$ の写像 $S^n\to S^n$ で $(n+1)$ 胞体を貼り付けた空間（ムーア空間）とする。$H_0(M)=\Z$、$H_n(M)=\Z/m$、それ以外の $H_k(M)=0$ を用いて、$H^*(M;\Z)$ および素数 $p$ に対する $H^*(M;\Z/p)$ を計算せよ。
\end{problem}

\begin{solution}
コホモロジーの普遍係数定理を用いる。$\Z$ 係数では、$H^0=\Z$ であり、
\[
H^n(M;\Z)\cong\operatorname{Hom}(\Z/m,\Z)\oplus\operatorname{Ext}(H_{n-1}(M),\Z)=0
\]
である（$H_{n-1}(M)$ は $n\ge2$ なら $0$、$n=1$ なら $\Z$ で、いずれも Ext は $0$）。また
\[
H^{n+1}(M;\Z)\cong\operatorname{Hom}(0,\Z)\oplus\operatorname{Ext}(\Z/m,\Z)\cong\Z/m
\]
で、その他の次数は $0$ である。

$\Z/p$ 係数では、$H^0=\Z/p$ で、本章最初の問題の (2) より
\begin{align*}
H^n(M;\Z/p)&\cong\operatorname{Hom}(\Z/m,\Z/p)\cong\Z/\gcd(m,p),\\
H^{n+1}(M;\Z/p)&\cong\operatorname{Ext}(\Z/m,\Z/p)\cong\Z/\gcd(m,p)
\end{align*}
である。よって $p\mid m$ なら $H^n(M;\Z/p)\cong H^{n+1}(M;\Z/p)\cong\Z/p$ で、$p\nmid m$ なら $M$ の $\Z/p$ 係数コホモロジーは1点のそれと同じである。特に $p\mid m$ のとき、$H^n(M;\Z)=0$ であるにもかかわらず $H^n(M;\Z/p)\neq0$ であり、係数を変えると新たな情報が得られることがわかる。
\end{solution}

\begin{problem}[難易度 B]
$X$ を有限 CW 複体、$F$ を体とする。$X$ のオイラー標数 $\chi(X)=\sum_k(-1)^kb_k$（$b_k$ は $H_k(X)$ の階数）が
\[
\chi(X)=\sum_k(-1)^k\dim_FH_k(X;F)
\]
をみたすことを、普遍係数定理を用いて示せ。
\end{problem}

\begin{solution}
$H_k(X)\cong\Z^{b_k}\oplus\bigoplus_i\Z/m_{k,i}$ と書く。$X$ は有限 CW 複体なので、これらは有限生成で、十分大きい $k$ と $k<0$ では $0$ である。普遍係数定理より
\[
H_k(X;F)\cong H_k(X)\otimes F\oplus\operatorname{Tor}(H_{k-1}(X),F)
\]
である。本章最初の問題の (1) より $\Z/m\otimes F\cong F/mF$、$\operatorname{Tor}(\Z/m,F)\cong F[m]$ である。$F$ の標数が $m$ を割り切るなら $m\cdot1=0$ なので $F/mF=F[m]=F$、割り切らない（標数 $0$ を含む）なら $m$ は $F$ で可逆なので $F/mF=F[m]=0$ である。

そこで $c_k$ を、$F$ の標数が $m_{k,i}$ を割り切るような $i$ の個数とすると（標数 $0$ なら $c_k=0$）、
\[
\dim_FH_k(X;F)=b_k+c_k+c_{k-1}
\]
となる。交代和をとると $c_k$ の項は $(-1)^kc_k$ と $(-1)^{k+1}c_k$ として2回現れて打ち消し合う（$c_{-1}=0$ で、十分大きい $k$ で $c_k=0$）。よって $\sum_k(-1)^k\dim_FH_k(X;F)=\sum_k(-1)^kb_k=\chi(X)$ である。
\end{solution}

\begin{problem}[難易度 D]
コホモロジーの普遍係数定理の完全列
\[
0\to\operatorname{Ext}(H_{n-1}(X),G)\to H^n(X;G)\xrightarrow{h}\operatorname{Hom}(H_n(X),G)\to0
\]
は自然であり分裂するが、分裂は自然にとれないことを示せ。すなわち、すべての空間 $X$ について $h\circ s_X=\id$ をみたし、任意の写像 $f\colon X\to Y$ について $f^*\circ s_Y=s_X\circ\operatorname{Hom}(f_*,G)$ となる準同型 $s_X\colon\operatorname{Hom}(H_n(X),G)\to H^n(X;G)$ の族は存在しないことを示せ。（$n=2$、$G=\Z/2$、商写像 $f\colon\mathbb{RP}^2\to\mathbb{RP}^2/\mathbb{RP}^1\cong S^2$ を考えよ。）
\end{problem}

\begin{solution}
$\mathbb{RP}^2$ は胞体 $e^0,e^1,e^2$ をもち、$f$ は $\mathbb{RP}^1$ を1点に潰す胞体写像で、$e^2$ の内部を $S^2=e^0\cup e^2$ の $e^2$ の内部に同相に写す。$\Z/2$ 係数の胞体コチェイン複体では、$\mathbb{RP}^2$ の $\delta\colon C^1\to C^2$ は $2$ 倍なので $0$ となり、$H^2(\mathbb{RP}^2;\Z/2)=C^2\cong\Z/2$ である。$S^2$ でも $H^2(S^2;\Z/2)=C^2\cong\Z/2$ であり、$f^\sharp\colon C^2(S^2)\to C^2(\mathbb{RP}^2)$ は同型だから、$f^*\colon H^2(S^2;\Z/2)\to H^2(\mathbb{RP}^2;\Z/2)$ は同型（特に $0$ でない）である。

自然な分裂 $s$ が存在したと仮定する。$H_1(S^2)=0$ なので $S^2$ では Ext 項が消え、$h$ は同型、したがって $s_{S^2}=h^{-1}$ も同型である。一方 $H_2(\mathbb{RP}^2)=0$ なので $\operatorname{Hom}(H_2(\mathbb{RP}^2),\Z/2)=0$ である。任意の $x\in H^2(S^2;\Z/2)$ に対し $x=s_{S^2}(h(x))$ なので、自然性から
\[
f^*x=f^*s_{S^2}(h(x))=s_{\mathbb{RP}^2}\bigl(\operatorname{Hom}(f_*,\Z/2)(h(x))\bigr)=s_{\mathbb{RP}^2}(0)=0
\]
となる。これは $f^*$ が同型であることに矛盾する。よって分裂は自然にとれない。同じ写像 $f$ を用いると、ホモロジーの普遍係数定理の分裂も自然でないことがわかる（$f_*\colon H_2(\mathbb{RP}^2;\Z/2)\to H_2(S^2;\Z/2)$ は同型だが、Tor 項から $\otimes$ 項への写像になる）。
\end{solution}

\begin{problem}[難易度 B]
ホモロジーのキュネット公式
\[
H_n(X\times Y)\cong\bigoplus_{i+j=n}H_i(X)\otimes H_j(Y)\ \oplus\bigoplus_{i+j=n-1}\operatorname{Tor}(H_i(X),H_j(Y))
\]
（$X,Y$ は CW 複体）を用いて、次を示せ。
(1) $m,n\ge1$ のとき $H_*(S^m\times S^n)$ を求め、有理係数のポアンカレ多項式 $P(t)=\sum_k\dim_{\mathbb{Q}}H_k(\,\cdot\,;\mathbb{Q})\,t^k$ が $(1+t^m)(1+t^n)$ であることを示せ。
(2) $H_k(T^n)\cong\Z^{\binom{n}{k}}$ を示し、$T^n$ のポアンカレ多項式とオイラー標数を求めよ。
\end{problem}

\begin{solution}
(1) $H_*(S^m)$、$H_*(S^n)$ は自由なので Tor 項はすべて $0$ である。$H_i(S^m)\otimes H_j(S^n)$ が $0$ でないのは $(i,j)=(0,0),(m,0),(0,n),(m,n)$ のときで、いずれも $\Z$ である。よって $H_k(S^m\times S^n)$ は $k=0,m,n,m+n$ で $\Z$（$m=n$ なら $H_m\cong\Z^2$）、それ以外で $0$ である。有理係数では $\dim H_k(\,\cdot\,;\mathbb{Q})$ は階数に等しいので
\[
P_{S^m\times S^n}(t)=1+t^m+t^n+t^{m+n}=(1+t^m)(1+t^n)
\]
である。

(2) $n$ に関する帰納法で示す。$n=1$ は $T^1=S^1$ で正しい。$T^n=T^{n-1}\times S^1$ とし、$H_*(T^{n-1})$ が自由であると仮定すると、Tor 項は消えて
\[
H_k(T^n)\cong H_k(T^{n-1})\otimes\Z\oplus H_{k-1}(T^{n-1})\otimes\Z\cong\Z^{\binom{n-1}{k}+\binom{n-1}{k-1}}=\Z^{\binom{n}{k}}
\]
となり、再び自由である。よって $P_{T^n}(t)=\sum_k\binom nk t^k=(1+t)^n$ であり、$\chi(T^n)=P_{T^n}(-1)=0$（$n\ge1$）である。
\end{solution}

\begin{problem}[難易度 C]
(1) キュネット公式を用いて $H_*(\mathbb{RP}^2\times\mathbb{RP}^2;\Z)$ を計算せよ。
(2) その結果と普遍係数定理から $H_*(\mathbb{RP}^2\times\mathbb{RP}^2;\Z/2)$ の次元を求め、体係数のキュネット公式による値 $(1+t+t^2)^2$ と一致することを確かめよ。
(3) より一般に、$M_m=S^1\cup_me^2$（$H_1(M_m)=\Z/m$、$H_2(M_m)=0$）に対し、$H_*(M_m\times M_n;\Z)$ を求めよ。
\end{problem}

\begin{solution}
(1) $H_*(\mathbb{RP}^2)$ は $H_0=\Z$、$H_1=\Z/2$、$H_2=0$ である。テンソル積部分は
\[
H_0=\Z,\quad H_1=(\Z\otimes\Z/2)\oplus(\Z/2\otimes\Z)=(\Z/2)^2,\quad H_2=\Z/2\otimes\Z/2=\Z/2
\]
を与え（$H_0\otimes H_2$ などは $0$）、次数3以上のテンソル積部分は $0$ である。Tor 部分は $i+j=n-1$ で $\operatorname{Tor}(H_1,H_1)=\operatorname{Tor}(\Z/2,\Z/2)=\Z/2$ のみが $0$ でなく、これは $n=3$ に寄与する。よって
\[
H_0\cong\Z,\quad H_1\cong(\Z/2)^2,\quad H_2\cong\Z/2,\quad H_3\cong\Z/2,\quad H_k=0\ (k\ge4)
\]
である。

(2) $\dim H_k(\,\cdot\,;\Z/2)=\dim(H_k\otimes\Z/2)+\dim\operatorname{Tor}(H_{k-1},\Z/2)$ より、$k=0,1,2,3,4$ の順に $1$、$2+0$、$1+2$、$1+1$、$0+1$、すなわち $1,2,3,2,1$ となる。一方 $H_*(\mathbb{RP}^2;\Z/2)$ は各次数 $0,1,2$ で $\Z/2$ なので、体係数のキュネット公式から $(1+t+t^2)^2=1+2t+3t^2+2t^3+t^4$ となり、一致する。

(3) $d=\gcd(m,n)$ とする。(1) と同様にテンソル積部分は $H_0=\Z$、$H_1=\Z/m\oplus\Z/n$、$H_2=\Z/m\otimes\Z/n\cong\Z/d$ を与え、Tor 部分は $\operatorname{Tor}(\Z/m,\Z/n)\cong\Z/d$ が $H_3$ に寄与する。よって
\[
H_*(M_m\times M_n)\cong\Z,\ \Z/m\oplus\Z/n,\ \Z/d,\ \Z/d,\ 0,\dots
\]
である。$m=n=2$ が (1) である。特に $d=1$ なら $H_2=H_3=0$ となる。
\end{solution}

\begin{problem}[難易度 B]
$F$ を体とし、有限 CW 複体 $X$ に対し $P_X(t)=\sum_k\dim_FH_k(X;F)\,t^k$ とおく。体係数のキュネット公式 $H_n(X\times Y;F)\cong\bigoplus_{i+j=n}H_i(X;F)\otimes_FH_j(Y;F)$ を用いて、$P_{X\times Y}=P_XP_Y$ および $\chi(X\times Y)=\chi(X)\chi(Y)$ を示せ。また $\Sigma_g\times S^2$ の有理係数、$\mathbb{RP}^4\times\mathbb{RP}^2$ の $\Z/2$ 係数および有理係数のポアンカレ多項式を求めよ。
\end{problem}

\begin{solution}
$\dim_F(V\otimes_FW)=\dim V\cdot\dim W$ なので
\[
\dim_FH_n(X\times Y;F)=\sum_{i+j=n}\dim_FH_i(X;F)\dim_FH_j(Y;F)
\]
であり、これは $P_X(t)P_Y(t)$ の $t^n$ の係数に等しい。よって $P_{X\times Y}=P_XP_Y$。オイラー標数の問題より $\chi(X)=P_X(-1)$（どの体でも）なので、$t=-1$ を代入して $\chi(X\times Y)=\chi(X)\chi(Y)$ を得る。

$\Sigma_g$ は $H_*=\Z,\Z^{2g},\Z$ なので $P_{\Sigma_g}=1+2gt+t^2$、よって
\[
P_{\Sigma_g\times S^2}(t)=(1+2gt+t^2)(1+t^2)=1+2gt+2t^2+2gt^3+t^4
\]
であり、$\chi=2(2-2g)$ である。$\Z/2$ 係数では $P_{\mathbb{RP}^n}=1+t+\dots+t^n$ なので
\[
P_{\mathbb{RP}^4\times\mathbb{RP}^2}(t)=(1+t+t^2+t^3+t^4)(1+t+t^2)=1+2t+3t^2+3t^3+3t^4+2t^5+t^6
\]
である。有理係数では $\mathbb{RP}^4,\mathbb{RP}^2$ ともに $P=1$ なので積も $1$ である。どちらでも $P(-1)=1=\chi(\mathbb{RP}^4)\chi(\mathbb{RP}^2)$ となり整合する。
\end{solution}

\begin{problem}[難易度 C]
体 $F$ 上のコホモロジーのキュネット公式：有限 CW 複体 $X,Y$ に対しクロス積 $a\times b=p_1^*a\smile p_2^*b$ が環同型 $H^*(X;F)\otimes_FH^*(Y;F)\cong H^*(X\times Y;F)$ を与える（積は $(a\otimes b)(c\otimes d)=(-1)^{\lvert b\rvert\lvert c\rvert}ac\otimes bd$）ことを認める。$H^*(\mathbb{RP}^k;\Z/2)\cong\Z/2[\alpha]/(\alpha^{k+1})$ を用いて
\[
H^*(\mathbb{RP}^m\times\mathbb{RP}^n;\Z/2)\cong\Z/2[\alpha,\beta]/(\alpha^{m+1},\beta^{n+1})
\]
を示せ。さらに $(\alpha+\beta)^{m+n}\neq0$ となるための必要十分条件は $\binom{m+n}{m}$ が奇数であることを示せ。
\end{problem}

\begin{solution}
$\alpha=p_1^*\alpha_m=\alpha_m\times1$、$\beta=p_2^*\alpha_n=1\times\alpha_n$ とおく。キュネット公式より $H^*(\mathbb{RP}^m\times\mathbb{RP}^n;\Z/2)$ は
\[
\alpha_m^i\times\alpha_n^j=p_1^*\alpha_m^i\smile p_2^*\alpha_n^j=\alpha^i\beta^j\qquad(0\le i\le m,\ 0\le j\le n)
\]
を基底とする。$\Z/2$ 係数なので符号は現れず、環は可換で、$\alpha^{m+1}=p_1^*(\alpha_m^{m+1})=0$、$\beta^{n+1}=0$ である。したがって $\Z/2[\alpha,\beta]\to H^*$ は全射で $(\alpha^{m+1},\beta^{n+1})$ を核に含み、商 $\Z/2[\alpha,\beta]/(\alpha^{m+1},\beta^{n+1})$ は単項式 $\alpha^i\beta^j$（$i\le m$、$j\le n$）を基底とするので、誘導される写像は基底を基底に写す同型である。

可換環なので二項定理が使えて
\[
(\alpha+\beta)^{m+n}=\sum_{i=0}^{m+n}\binom{m+n}{i}\alpha^i\beta^{m+n-i}
\]
である。$i>m$ なら $\alpha^i=0$、$i<m$ なら $m+n-i>n$ で $\beta^{m+n-i}=0$ なので、残るのは $i=m$ の項 $\binom{m+n}{m}\alpha^m\beta^n$ のみである。$\alpha^m\beta^n$ は $H^{m+n}$ の基底なので、$(\alpha+\beta)^{m+n}\neq0$ と $\binom{m+n}{m}$ が奇数であることは同値である。例えば $m=n=1$ では $(\alpha+\beta)^2=0$、$m=1$、$n=2$ では $(\alpha+\beta)^3=\alpha\beta^2\neq0$ である。
\end{solution}

\begin{problem}[難易度 C]
$Y$ のすべての整数係数コホモロジー群が有限生成自由であるとき、クロス積が環同型 $H^*(X;\Z)\otimes H^*(Y;\Z)\cong H^*(X\times Y;\Z)$ を与える（積の符号は前問と同じ）ことを認める。これを用いて次を示せ。
(1) $m,n\ge1$ のとき $H^*(S^m\times S^n;\Z)$ は $a=p_1^*\sigma_m$、$b=p_2^*\sigma_n$ で生成され、$a^2=b^2=0$、$ba=(-1)^{mn}ab$ で、$ab$ が $H^{m+n}$ の生成元である。
(2) $H^*(T^n;\Z)$ は $x_i=p_i^*\sigma$（$p_i\colon T^n\to S^1$ は第 $i$ 成分への射影、$\sigma$ は $H^1(S^1)$ の生成元）を生成元とする外積代数 $\Lambda_\Z[x_1,\dots,x_n]$ である。
\end{problem}

\begin{solution}
(1) $H^*(S^n;\Z)$ は $1,\sigma_n$ を基底とする自由加群である。よって $H^*(S^m\times S^n)$ は
\[
1\times1=1,\quad \sigma_m\times1=a,\quad 1\times\sigma_n=b,\quad \sigma_m\times\sigma_n=p_1^*\sigma_m\smile p_2^*\sigma_n=ab
\]
を基底とする。$a^2=p_1^*(\sigma_m^2)=0$、$b^2=p_2^*(\sigma_n^2)=0$ であり、次数付き可換性から $ba=(-1)^{mn}ab$ である。特に $ab$ は $H^{m+n}\cong\Z$ の生成元である。

(2) $n$ に関する帰納法で、$H^*(T^n)$ が単項式 $x_I=x_{i_1}\cdots x_{i_k}$（$I=\{i_1<\dots<i_k\}\subset\{1,\dots,n\}$）を基底とする自由加群で、$x_i^2=0$、$x_ix_j=-x_jx_i$ をみたすことを示す。$n=1$ は $H^*(S^1)=\Z\{1,\sigma\}$ で正しい。$T^n=T^{n-1}\times S^1$ とし、$q\colon T^n\to T^{n-1}$ を最初の $n-1$ 成分への射影とする。$H^*(S^1)$ は自由なので仮定が使え、$H^*(T^n)$ は $x_I\times1=q^*x_I$ と $x_I\times\sigma=q^*x_I\smile x_n$ を基底とする。$q^*x_i=x_i$（$i<n$）なので、これらはちょうど $x_J$（$J\subset\{1,\dots,n\}$）である。関係式については、各 $x_i$ が次数1なので次数付き可換性から $x_ix_j=-x_jx_i$、また $x_i^2=p_i^*(\sigma^2)=0$ である。よって $H^*(T^n)\cong\Lambda_\Z[x_1,\dots,x_n]$ で、$H^k(T^n)$ の階数は $\binom nk$、$x_1\cdots x_n$ が $H^n$ の生成元である。
\end{solution}

\begin{problem}[難易度 C]
$\mathbb{CP}^3$、$S^2\times S^4$、$S^2\vee S^4\vee S^6$ はいずれも同型な整数係数ホモロジー群（したがって同じポアンカレ多項式）をもつが、どの2つもホモトピー同値でないことを示せ。ただし $H^*(\mathbb{CP}^3;\Z)\cong\Z[x]/(x^4)$（$\deg x=2$）を用いてよい。
\end{problem}

\begin{solution}
3つの空間はいずれも次数 $0,2,4,6$ で $\Z$、それ以外で $0$ のホモロジーをもち（$S^2\times S^4$ はキュネット公式による）、ポアンカレ多項式はすべて $1+t^2+t^4+t^6$ である。コホモロジー群も普遍係数定理により同じ次数に $\Z$ である。$H^2$ の生成元を $u$、$H^4$ の生成元を $v$ とし、次の2つの性質を調べる：(A) $u^2\neq0$、(B) $u\smile v\neq0$。これらは生成元の符号の選び方によらず、環同型で保たれる。

$\mathbb{CP}^3$ では $u=x$、$v=x^2$ で、$u^2=x^2\neq0$、$uv=x^3\neq0$ なので (A)(B) ともに成り立つ。

$S^2\times S^4$ では前問より $u=a$、$v=b$ で、$u^2=a^2=0$、$uv=ab\neq0$ なので (A) は不成立、(B) は成立する。

$S^2\vee S^4\vee S^6$ では、$u,v$ はそれぞれ $S^2,S^4$ の成分からの引き戻しである。$u^2$ は $H^4(S^2)=0$ の元の引き戻しなので $0$ であり、$uv$ は異なる成分からの積なので楔和の性質（第7章）により $0$ である。よって (A)(B) ともに不成立である。

以上より3つのコホモロジー環は互いに同型でなく、どの2つもホモトピー同値でない。ホモロジー群やポアンカレ多項式では区別できない空間が、カップ積によって区別される例である。
\end{solution}

%================================================
\chapter{双対性と多様体}
% 15問
%================================================

\begin{summary}
\begin{itemize}
  \item \textbf{局所ホモロジー}：$n$ 次元多様体 $M$ の各点で $H_n(M,M\setminus\{x\})\cong\Z$。その生成元 $\mu_x$ の、局所的に整合する選び方が\textbf{向き}。連結な $M$ は向きの2重被覆 $\tilde M\to M$ をもち、$M$ が向き付け可能 $\iff$ $\tilde M$ が非連結。特に $\pi_1(M)$ が指数2の部分群をもたなければ向き付け可能。
  \item \textbf{最高次のホモロジー}：閉連結 $n$ 次元多様体では、$H_n(M;\Z)\cong\Z$（向き付け可能）または $0$（向き付け不能）。$H_n(M;\Z/2)\cong\Z/2$ はつねに成り立ち、$k>n$ で $H_k(M)=0$。生成元が基本類 $[M]$。
  \item \textbf{ポアンカレ双対性}：閉で $R$ 向き付け可能な $n$ 次元多様体 $M$ について
  \[
  D\colon H^k(M;R)\to H_{n-k}(M;R),\qquad a\mapsto[M]\frown a
  \]
  は同型。$\Z/2$ 係数ならすべての閉多様体で成り立つ。
  \item \textbf{帰結}：体係数のベッチ数は $b_k=b_{n-k}$。カップ積の対 $H^k(M;F)\times H^{n-k}(M;F)\to F$ は非退化。奇数次元の閉多様体では $\chi(M)=0$。
  \item \textbf{交叉形式}：$n=4k$ のとき $Q(a,b)=\langle a\smile b,[M]\rangle$ は $H^{2k}(M;\Z)/\text{ねじれ}$ 上の対称なユニモジュラ形式（行列式 $\pm1$）。
  \item \textbf{レフシェッツ双対性}：コンパクトで向き付け可能な境界付き多様体では $H^k(M;R)\cong H_{n-k}(M,\partial M;R)$。
  \item \textbf{アレクサンダー双対性}：$K\subset S^n$ がコンパクトで局所可縮なら $\tilde H_i(S^n\setminus K)\cong\tilde H^{n-i-1}(K)$。
\end{itemize}
\end{summary}

\begin{problem}[難易度 A]
$M$ を $n$ 次元多様体（境界なし）、$x\in M$ とする。局所ホモロジー群について
\[
H_k(M, M-\{x\};\Z)\cong
\begin{cases} \Z & (k=n),\\ 0 & (k\neq n)\end{cases}
\]
が成り立つことを示せ。
\end{problem}

\begin{solution}
$x$ の開近傍 $U$ で $\R^n$ と同相なものをとり、同相で $x$ が原点 $0$ に対応するとしてよい。
$Z=M-U$ とおく。$U$ は開集合なので $Z$ は閉集合であり、$M$ はハウスドルフなので $M-\{x\}$ は開集合である。
したがって
\[
\overline{Z}=Z\subset M-\{x\}=\operatorname{int}(M-\{x\})
\]
となり、切除定理が使える。$Z$ を切除すると
\[
H_k(M,M-\{x\})\cong H_k(U,U-\{x\})\cong H_k(\R^n,\R^n-\{0\})
\]
を得る。$\R^n$ は可縮なので、対 $(\R^n,\R^n-\{0\})$ の被約ホモロジー長完全列から
連結準同型 $H_k(\R^n,\R^n-\{0\})\to \tilde H_{k-1}(\R^n-\{0\})$ は同型である。
さらに $\R^n-\{0\}$ は $S^{n-1}$ にホモトピー同値（$v\mapsto v/\lvert v\rvert$ が変位レトラクション）なので
\[
H_k(M,M-\{x\})\cong \tilde H_{k-1}(S^{n-1})\cong
\begin{cases} \Z & (k=n),\\ 0 & (k\neq n)\end{cases}
\]
となる。
\end{solution}

\begin{problem}[難易度 B]
$\mathbb{H}^n=\{(x_1,\dots,x_n)\in\R^n : x_n\geq 0\}$ とし、$x\in\partial\mathbb{H}^n=\{x_n=0\}$ とする。
すべての $k$ に対して $H_k(\mathbb{H}^n,\mathbb{H}^n-\{x\};\Z)=0$ であることを示せ。
これと前問を用いて、境界付き多様体の内点と境界点は局所ホモロジーで区別されること
（したがって境界は同相写像で保たれること）を説明せよ。
\end{problem}

\begin{solution}
$p=x+e_n$（$e_n=(0,\dots,0,1)$）とおく。$y\in\mathbb{H}^n-\{x\}$ と $t\in[0,1]$ に対し
$(1-t)y+tp$ の第 $n$ 座標は $(1-t)y_n+t$ であり、これは $t>0$ なら正なので $x$ と異なる。
$t=0$ のときは $y\neq x$ である。よって線分 $[y,p]$ は $\mathbb{H}^n-\{x\}$ に含まれ、
$\mathbb{H}^n-\{x\}$ は $p$ に関して星型、特に可縮である。$\mathbb{H}^n$ も凸なので可縮である。
対 $(\mathbb{H}^n,\mathbb{H}^n-\{x\})$ の被約ホモロジー長完全列
\[
\tilde H_k(\mathbb{H}^n-\{x\})\to\tilde H_k(\mathbb{H}^n)\to H_k(\mathbb{H}^n,\mathbb{H}^n-\{x\})\to\tilde H_{k-1}(\mathbb{H}^n-\{x\})
\]
の両端と左から2番目が $0$ なので、$H_k(\mathbb{H}^n,\mathbb{H}^n-\{x\})=0$ である。

境界付き多様体 $M$ の点 $x$ に対し、前問と同じ切除の議論により $H_n(M,M-\{x\})$ は
$x$ の座標近傍 $U$ だけで決まる。$x$ が内点なら $U\cong\R^n$ ととれて $H_n(M,M-\{x\})\cong\Z$、
境界点なら $U\cong\mathbb{H}^n$（$x\mapsto\partial\mathbb{H}^n$ の点）ととれて
$H_n(M,M-\{x\})=0$ となる。同相写像 $f\colon M\to N$ は対の同相
$(M,M-\{x\})\to(N,N-\{f(x)\})$ を与えるので局所ホモロジーを保ち、内点を内点に、
境界点を境界点に写す。
\end{solution}

\begin{problem}[難易度 B]
$n\geq 1$ とする。実射影空間 $\mathbb{RP}^n$ の胞体ホモロジーを用いて $H_n(\mathbb{RP}^n;\Z)$ を求め、
$\mathbb{RP}^n$ が向き付け可能であるための必要十分条件は $n$ が奇数であることを示せ。
ただし「閉連結 $n$ 次元多様体 $M$ が向き付け可能であることと $H_n(M;\Z)\cong\Z$ であることは同値であり、
向き付け不可能なら $H_n(M;\Z)=0$」という定理は用いてよい。
\end{problem}

\begin{solution}
$\mathbb{RP}^n$ は各次元 $0\leq k\leq n$ に胞体を1つずつもつ CW 複体であり、
胞体鎖複体は $C_k=\Z$（$0\leq k\leq n$）で、境界写像は
\[
d_k\colon C_k\to C_{k-1},\qquad d_k=1+(-1)^k
\]
すなわち $k$ が偶数なら $2$ 倍写像、奇数なら $0$ 写像である。
$C_{n+1}=0$ なので $H_n(\mathbb{RP}^n)=\ker d_n$ である。
$n$ が奇数なら $d_n=0$ なので $H_n(\mathbb{RP}^n)\cong\Z$、
$n$ が偶数なら $d_n$ は $\Z\to\Z$ の $2$ 倍写像で単射なので $H_n(\mathbb{RP}^n)=0$ である。

$\mathbb{RP}^n$ はコンパクトで境界のない連結 $n$ 次元多様体（閉連結多様体）であるから、
上の定理により、向き付け可能 $\iff H_n(\mathbb{RP}^n;\Z)\cong\Z\iff n$ は奇数、となる。
\end{solution}

\begin{problem}[難易度 B]
単連結な（連結）多様体は向き付け可能であることを示せ。
ただし、連結多様体 $M$ の向き付け二重被覆 $\tilde M\to M$ について、
「$\tilde M$ が連結であることと $M$ が向き付け不可能であることは同値」という事実は用いてよい。
\end{problem}

\begin{solution}
$M$ を単連結な連結 $n$ 次元多様体とし、向き付け不可能であると仮定して矛盾を導く。
向き付け二重被覆 $p\colon\tilde M\to M$ は2重の被覆写像であり、仮定より $\tilde M$ は連結である。
多様体は局所弧状連結なので、連結性と弧状連結性は一致し、被覆空間論が適用できる。
$\tilde x\in\tilde M$、$x=p(\tilde x)$ とすると、弧状連結な被覆空間のシート数は部分群の指数
\[
[\pi_1(M,x):p_*\pi_1(\tilde M,\tilde x)]
\]
に等しい。$\pi_1(M,x)=1$ なのでこの指数は $1$ であり、シート数が $2$ であることに矛盾する。
よって $M$ は向き付け可能である。
\end{solution}

\begin{problem}[難易度 C]
$M$ を閉連結な向き付け可能 $n$ 次元多様体とし、$b_k=\rank H_k(M;\Z)$ をベッチ数、
$T_k$ を $H_k(M;\Z)$ のねじれ部分群とする。ポアンカレ双対性を用いて次を示せ。
\begin{enumerate}
\item $b_k=b_{n-k}$。
\item $T_k\cong T_{n-k-1}$。
\end{enumerate}
ただし閉多様体のホモロジー群は有限生成であることを用いてよい。
\end{problem}

\begin{solution}
(1) $M$ は $\Z$ 向き付け可能なので $\mathbb{Q}$ 向き付け可能でもあり、
$\mathbb{Q}$ 係数のポアンカレ双対性 $H^{n-k}(M;\mathbb{Q})\cong H_k(M;\mathbb{Q})$ が成り立つ。
体係数の普遍係数定理より $H^{n-k}(M;\mathbb{Q})\cong\operatorname{Hom}_{\mathbb{Q}}(H_{n-k}(M;\mathbb{Q}),\mathbb{Q})$ であり、
有限次元ベクトル空間の双対は同じ次元をもつ。また $H_k(M;\mathbb{Q})\cong H_k(M;\Z)\otimes\mathbb{Q}$ の次元は $b_k$ である。よって
\[
b_k=\dim H_k(M;\mathbb{Q})=\dim H^{n-k}(M;\mathbb{Q})=\dim H_{n-k}(M;\mathbb{Q})=b_{n-k}.
\]

(2) $\Z$ 係数のポアンカレ双対性より $H_k(M;\Z)\cong H^{n-k}(M;\Z)$ である。普遍係数定理から
\[
H^{n-k}(M;\Z)\cong\operatorname{Hom}(H_{n-k}(M),\Z)\oplus\operatorname{Ext}(H_{n-k-1}(M),\Z).
\]
第1項は自由、第2項は $\operatorname{Ext}(\Z,\Z)=0$、$\operatorname{Ext}(\Z/m,\Z)\cong\Z/m$ より $T_{n-k-1}$ と同型な有限群である。
したがって $H^{n-k}(M;\Z)$ のねじれ部分群は $T_{n-k-1}$ と同型であり、同型で対応する
$H_k(M;\Z)$ のねじれ部分群 $T_k$ もこれと同型である。
\end{solution}

\begin{problem}[難易度 C]
$M$ を奇数次元 $n$ の閉多様体とする（向き付け可能とは仮定しない）。$\chi(M)=0$ を示せ。
ただし閉多様体のホモロジー群は有限生成で、$k>n$ で $0$ になることを用いてよい。
\end{problem}

\begin{solution}
まず、任意の素数 $p$ に対して
$\chi(M)=\sum_k(-1)^k\dim_{\Z/p}H_k(M;\Z/p)$ であることを示す。
$H_k(M;\Z)$ の階数を $b_k$、その巡回分解に現れる $p$ べき位数の巡回因子の個数を $t_k$ とすると、
普遍係数定理 $H_k(M;\Z/p)\cong H_k(M)\otimes\Z/p\oplus\operatorname{Tor}(H_{k-1}(M),\Z/p)$ より
\[
\dim H_k(M;\Z/p)=b_k+t_k+t_{k-1}
\]
である。交代和をとると $t$ の項は打ち消し合い（$t_{-1}=0$、$k$ が大きければ $t_k=0$）、
$\sum(-1)^k\dim H_k(M;\Z/p)=\sum(-1)^kb_k=\chi(M)$ となる。

$p=2$ とする。任意の多様体は $\Z/2$ 向き付け可能なので、$\Z/2$ 係数のポアンカレ双対性
$H_k(M;\Z/2)\cong H^{n-k}(M;\Z/2)$ が成り立ち、体係数の普遍係数定理から
$\dim H^{n-k}(M;\Z/2)=\dim H_{n-k}(M;\Z/2)$ である。$d_k=\dim H_k(M;\Z/2)$ とおくと $d_k=d_{n-k}$ であり、
$n$ が奇数なので $(-1)^{n-k}=-(-1)^k$ である。よって
\[
\chi(M)=\sum_{k=0}^n(-1)^kd_k=\sum_{k=0}^n(-1)^{n-k}d_{n-k}=-\sum_{k=0}^n(-1)^{k}d_{k}=-\chi(M)
\]
（2つ目の等号は $k\mapsto n-k$ の添字の付け替え）となり、$\chi(M)=0$ である。
\end{solution}

\begin{problem}[難易度 D]
$M$ を閉連結 $n$ 次元多様体とする。次を示せ。
\begin{enumerate}
\item $M$ が向き付け可能なら $H_{n-1}(M;\Z)$ は自由アーベル群である。
\item $M$ が向き付け不可能なら $H_{n-1}(M;\Z)$ のねじれ部分群は $2$ べき位数の自明でない巡回群である。
\end{enumerate}
ただし次の事実は用いてよい：閉多様体のホモロジーは有限生成である。
環 $R$ について、$M$ が $R$ 向き付け可能なら $H_n(M;R)\cong R$、
向き付け不可能なら $H_n(M;R)\cong\{r\in R: 2r=0\}$ である。
\end{problem}

\begin{solution}
素数 $p$ に対し、普遍係数定理より
\[
H_n(M;\Z/p)\cong H_n(M;\Z)\otimes\Z/p\;\oplus\;\operatorname{Tor}(H_{n-1}(M;\Z),\Z/p).
\]
また $\operatorname{Tor}(\Z,\Z/p)=0$、$\operatorname{Tor}(\Z/m,\Z/p)\cong\Z/\gcd(m,p)$ である。

(1) 向き付け可能なら $H_n(M;\Z)\cong\Z$ であり、$\Z$ 向き付けから $\Z/p$ 向き付けが誘導されるので
$H_n(M;\Z/p)\cong\Z/p$。上式は $\Z/p\cong\Z/p\oplus\operatorname{Tor}(H_{n-1},\Z/p)$ となり、
位数を比べて $\operatorname{Tor}(H_{n-1},\Z/p)=0$ である。
これがすべての素数 $p$ で成り立つので、有限生成アーベル群 $H_{n-1}$ の巡回分解には
有限巡回因子が現れず、$H_{n-1}$ は自由である。

(2) 向き付け不可能なら $H_n(M;\Z)=\{r\in\Z:2r=0\}=0$ である。
$p$ が奇素数なら $\{r\in\Z/p:2r=0\}=0$ なので $H_n(M;\Z/p)=0$、よって
$\operatorname{Tor}(H_{n-1},\Z/p)=0$ で、$H_{n-1}$ に奇素数べき位数の巡回因子はない。
$p=2$ なら $H_n(M;\Z/2)\cong\Z/2$ なので $\operatorname{Tor}(H_{n-1},\Z/2)\cong\Z/2$ であり、
$2$ べき位数の巡回因子はちょうど1個である。したがってねじれ部分群は $\Z/2^m$（$m\geq1$）の形である。
（実際には $\Z/2$ であることも示せるが、ここでは立ち入らない。）
\end{solution}

\begin{problem}[難易度 C]
$\mathbb{CP}^n$ の整数係数コホモロジー群は $0\leq k\leq n$ で $H^{2k}(\mathbb{CP}^n;\Z)\cong\Z$、それ以外で $0$ である。
ポアンカレ双対性を用いて、環として
\[
H^*(\mathbb{CP}^n;\Z)\cong\Z[a]/(a^{n+1}),\qquad \deg a=2
\]
であることを示せ。
\end{problem}

\begin{solution}
$n$ に関する帰納法で「$a\in H^2(\mathbb{CP}^n)$ を生成元とすると、$0\leq k\leq n$ で $a^k$ が $H^{2k}(\mathbb{CP}^n)$ を生成する」ことを示す。
$n=1$ は明らか。$n\geq2$ とする。$\mathbb{CP}^n$ は $\mathbb{CP}^{n-1}$ に $2n$ 胞体を1つ貼って得られるので、
包含 $i\colon\mathbb{CP}^{n-1}\to\mathbb{CP}^n$ は $k\leq 2n-2$ で同型 $i^*\colon H^k(\mathbb{CP}^n)\to H^k(\mathbb{CP}^{n-1})$ を誘導する
（胞体コホモロジーで $2n-1$ 骨格以下が一致するため）。
$i^*a$ は $H^2(\mathbb{CP}^{n-1})$ の生成元であり、帰納法の仮定より $(i^*a)^k=i^*(a^k)$ は $k\leq n-1$ で生成元である。
$i^*$ が同型なので $a^k$ も $H^{2k}(\mathbb{CP}^n)$ の生成元である（$k\leq n-1$）。

残るは $a^n$ である。$\mathbb{CP}^n$ は閉連結で向き付け可能（単連結）な $2n$ 次元多様体であり、ポアンカレ双対性より
カップ積のペアリング
\[
H^2(\mathbb{CP}^n)\times H^{2n-2}(\mathbb{CP}^n)\to\Z,\qquad(\varphi,\psi)\mapsto\langle\varphi\smile\psi,[\mathbb{CP}^n]\rangle
\]
は（ねじれがないので）非退化、すなわち表現行列の行列式が $\pm1$ である。
両群は $a$、$a^{n-1}$ で生成される $\Z$ なので $\langle a\smile a^{n-1},[\mathbb{CP}^n]\rangle=\pm1$ であり、
$a^n$ は $H^{2n}(\mathbb{CP}^n)\cong\Z$ の生成元である。
また $a^{n+1}\in H^{2n+2}=0$ である。よって $\Z[a]/(a^{n+1})\to H^*(\mathbb{CP}^n)$、$a\mapsto a$ は環同型である。
\end{solution}

\begin{problem}[難易度 C]
$n\geq1$ とする。$\Z/2$ 係数のポアンカレ双対性を用いて、環として
\[
H^*(\mathbb{RP}^n;\Z/2)\cong\Z/2[a]/(a^{n+1}),\qquad\deg a=1
\]
であることを示せ。
\end{problem}

\begin{solution}
$\mathbb{RP}^n$ の胞体鎖複体の境界写像は $0$ か $2$ 倍なので、$\Z/2$ 係数ではすべて $0$ になり、
$0\leq k\leq n$ で $H^k(\mathbb{RP}^n;\Z/2)\cong\Z/2$、それ以外で $0$ である。
$n$ に関する帰納法で「$a$ を $H^1$ の生成元とすると $a^k$（$k\leq n$）は $H^k$ の生成元」を示す。
$n=1$ は明らか。$n\geq2$ とする。$\mathbb{RP}^n$ は $\mathbb{RP}^{n-1}$ に $n$ 胞体を貼ったものなので、
包含 $i$ は $k\leq n-1$ で同型 $i^*\colon H^k(\mathbb{RP}^n;\Z/2)\to H^k(\mathbb{RP}^{n-1};\Z/2)$ を誘導する。
帰納法の仮定と $i^*(a^k)=(i^*a)^k$ から、$k\leq n-1$ で $a^k\neq0$ である。

$\mathbb{RP}^n$ は閉連結多様体であり、任意の多様体は $\Z/2$ 向き付け可能なので、$\Z/2$ 係数の基本類が存在し、
体 $\Z/2$ 係数のポアンカレ双対性より
\[
H^1(\mathbb{RP}^n;\Z/2)\times H^{n-1}(\mathbb{RP}^n;\Z/2)\to\Z/2,\qquad(\varphi,\psi)\mapsto\langle\varphi\smile\psi,[\mathbb{RP}^n]\rangle
\]
は非退化である。両者は $\Z/2$ で $a$、$a^{n-1}$ が非零元なので、非退化性から $\langle a\smile a^{n-1},[\mathbb{RP}^n]\rangle=1$、
特に $a^n\neq0$ である。$H^{n+1}=0$ より $a^{n+1}=0$。以上から主張の環同型を得る。
\end{solution}

\begin{problem}[難易度 C]
$n\geq1$ とし、$f\colon\mathbb{CP}^{2n}\to\mathbb{CP}^{2n}$ を連続写像とする。
$\deg f\geq0$ であることを示せ。特に $\mathbb{CP}^{2n}$ には向きを反転するホモトピー同値写像
（したがって向きを反転する同相写像）が存在しないことを示せ。
\end{problem}

\begin{solution}
$m=2n$ とおき、$a\in H^2(\mathbb{CP}^m;\Z)$ を生成元とする。$\mathbb{CP}^n$ のコホモロジー環の問題より $a^m$ は $H^{2m}(\mathbb{CP}^m)\cong\Z$ の生成元である。
$f^*a\in H^2\cong\Z$ なので $f^*a=ka$（$k\in\Z$）と書ける。カップ積の自然性より
\[
f^*(a^m)=(f^*a)^m=k^ma^m.
\]
一方、普遍係数定理により $H^{2m}(\mathbb{CP}^m)\cong\operatorname{Hom}(H_{2m}(\mathbb{CP}^m),\Z)$ であり、$f^*$ は $f_*$ の双対である。
$f_*[\mathbb{CP}^m]=(\deg f)[\mathbb{CP}^m]$ なので $f^*$ は $H^{2m}$ 上で $\deg f$ 倍写像である。よって
\[
\deg f=k^m=k^{2n}\geq0.
\]
向きを反転するホモトピー同値写像は次数 $-1$ をもつ（ホモトピー同値写像の次数は $\pm1$ で、向き反転は $-1$ を意味する）が、
これは $\deg f\geq0$ に反する。同相写像はホモトピー同値写像なので後半も従う。
\end{solution}

\begin{problem}[難易度 C]
閉連結な向き付けられた $4$ 次元多様体 $M$ の交叉形式を
\[
Q_M\colon H^2(M;\Z)/T\times H^2(M;\Z)/T\to\Z,\qquad Q_M(\alpha,\beta)=\langle\alpha\smile\beta,[M]\rangle
\]
で定める（$T$ はねじれ部分群）。$Q_M$ が対称でユニモジュラー（表現行列の行列式が $\pm1$）であることを説明し、
$M=\mathbb{CP}^2$ と $M=S^2\times S^2$ の交叉形式の表現行列を求めよ。
また $S^2\times S^2$ では任意の $x$ について $Q(x,x)$ が偶数であることを示せ。
\end{problem}

\begin{solution}
次数2の元について $\alpha\smile\beta=(-1)^{2\cdot2}\beta\smile\alpha=\beta\smile\alpha$ なので $Q_M$ は対称である。
ポアンカレ双対性の帰結として、閉向き付け可能 $n$ 次元多様体ではカップ積ペアリング
$H^k/T\times H^{n-k}/T\to\Z$ が非退化（随伴写像 $H^k/T\to\operatorname{Hom}(H^{n-k}/T,\Z)$ が同型）である。
$n=4,k=2$ の場合、これは $Q_M$ の表現行列が $\Z$ 上可逆、すなわち行列式 $\pm1$ であることを意味する。

$\mathbb{CP}^2$：$H^2=\Z a$ で、$a^2$ は $H^4$ の生成元である。$[\mathbb{CP}^2]$ の向きを $\langle a^2,[\mathbb{CP}^2]\rangle=1$ となるように選べば、表現行列は $(1)$ である。

$S^2\times S^2$：$u\in H^2(S^2)$ を生成元、$p_1,p_2$ を射影とし、$\alpha=p_1^*u$、$\beta=p_2^*u$ とおく。
キュネット公式より $H^2(S^2\times S^2)=\Z\alpha\oplus\Z\beta$、$H^4\cong\Z$ は $u\times u=\alpha\smile\beta$ で生成される。
$\alpha^2=p_1^*(u^2)=0$（$H^4(S^2)=0$）、同様に $\beta^2=0$ である。
$\langle\alpha\smile\beta,[M]\rangle=1$ となる向きを選べば、基底 $\alpha,\beta$ に関する表現行列は
\[
\begin{pmatrix}0&1\\1&0\end{pmatrix}
\]
であり、行列式は $-1$ である。$x=s\alpha+t\beta$ に対して $Q(x,x)=2st$ は偶数である。
（$\mathbb{CP}^2$ では $Q(a,a)=1$ が奇数である点が対照的である。）
\end{solution}

\begin{problem}[難易度 C]
$M$ を閉連結な向き付け可能 $(4k+2)$ 次元多様体とする。
ベッチ数 $b_{2k+1}(M)$ は偶数であることを示せ。さらに $\chi(M)$ も偶数であることを示せ。
\end{problem}

\begin{solution}
$m=2k+1$、$V=H^m(M;\R)$ とおく。$\R$ 係数のポアンカレ双対性により、双線形形式
\[
\omega\colon V\times V\to\R,\qquad\omega(\varphi,\psi)=\langle\varphi\smile\psi,[M]\rangle
\]
は非退化である（体係数では $H^m\times H^{n-m}\to\R$ が非退化で、ここでは $n-m=m$）。
また $\varphi\smile\psi=(-1)^{m\cdot m}\psi\smile\varphi=-\psi\smile\varphi$ なので $\omega$ は交代的（歪対称）である。
$\omega$ の表現行列を $A$、$d=\dim V$ とすると $A^{\mathsf T}=-A$ より
\[
\det A=\det A^{\mathsf T}=\det(-A)=(-1)^d\det A.
\]
非退化より $\det A\neq0$ なので $(-1)^d=1$、すなわち $d$ は偶数である。
普遍係数定理より $d=\dim H^m(M;\R)=\dim H_m(M;\R)=b_m$ なので $b_{2k+1}$ は偶数である。

次に、ベッチ数の対称性（本章の前出の問題）より $b_i=b_{4k+2-i}$ である。$i$ と $4k+2-i$ は偶奇が等しいので
\[
\chi(M)=\sum_{i=0}^{4k+2}(-1)^ib_i=2\sum_{i=0}^{2k}(-1)^ib_i-b_{2k+1}
\]
となり、$b_{2k+1}$ が偶数だから $\chi(M)$ も偶数である。例えば種数 $g$ の閉曲面では $b_1=2g$ である。
\end{solution}

\begin{problem}[難易度 B]
アレクサンダー双対性
（$K\subset S^n$ がコンパクトで局所可縮な空でない真部分集合なら $\tilde H_i(S^n-K;\Z)\cong\tilde H^{n-i-1}(K;\Z)$）を用いて、
次のホモロジー群を計算せよ。
\begin{enumerate}
\item $K\subset S^3$ が $S^1$ と同相な部分集合（結び目）のとき、$H_*(S^3-K;\Z)$。
\item $L\subset S^3$ が2つの互いに交わらない $S^1$ の和（2成分の絡み目）のとき、$H_*(S^3-L;\Z)$。
\end{enumerate}
\end{problem}

\begin{solution}
$S^1$ や有限個の円周の非交和はコンパクトな多様体なので局所可縮であり、双対性が使える。
$n=3$ のとき $\tilde H_i(S^3-K)\cong\tilde H^{2-i}(K)$ である。

(1) $\tilde H^j(S^1)$ は $j=1$ で $\Z$、それ以外で $0$ である。$2-i=1$ となるのは $i=1$ のみなので
\[
H_0(S^3-K)\cong\Z,\quad H_1(S^3-K)\cong\Z,\quad H_i(S^3-K)=0\ (i\geq2).
\]
（$H_0$ は被約ホモロジーが $0$ であることから $\Z$、すなわち補空間は弧状連結である。）
結び目がどのように結ばれていても、補空間のホモロジーは変わらないことに注意する
（結び目の区別には基本群などより精密な不変量が必要である）。

(2) $\tilde H^0(S^1\sqcup S^1)\cong\Z$、$H^1(S^1\sqcup S^1)\cong\Z^2$、それ以外は $0$ である。
$2-i=0$ より $i=2$、$2-i=1$ より $i=1$ なので
\[
H_0(S^3-L)\cong\Z,\quad H_1(S^3-L)\cong\Z^2,\quad H_2(S^3-L)\cong\Z,\quad H_i(S^3-L)=0\ (i\geq3).
\]
\end{solution}

\begin{problem}[難易度 D]
$n\geq2$ とする。アレクサンダー双対性を用いて次を示せ。
\begin{enumerate}
\item $h\colon D^k\to S^n$ が埋め込み（像への同相写像）なら、$S^n-h(D^k)$ の被約ホモロジーはすべて $0$ である。特に $S^n-h(D^k)$ は弧状連結である。
\item （ジョルダン・ブラウアーの分離定理）$h\colon S^{n-1}\to S^n$ が埋め込みなら、$S^n-h(S^{n-1})$ はちょうど2つの連結成分をもち、各成分は1点と同じホモロジーをもつ。
\end{enumerate}
\end{problem}

\begin{solution}
$h(D^k)$、$h(S^{n-1})$ はそれぞれ $D^k$、$S^{n-1}$ と同相なのでコンパクトかつ局所可縮であり、
$S^n$ の空でない真部分集合である（次元の理由、あるいは下の計算自体からわかる）。

(1) $D^k$ は可縮なので $\tilde H^j(h(D^k))=0$（すべての $j$）。アレクサンダー双対性より
$\tilde H_i(S^n-h(D^k))\cong\tilde H^{n-i-1}(h(D^k))=0$ である。$\tilde H_0=0$ は弧状連結性を意味する。

(2) $U=S^n-h(S^{n-1})$ とおく。$\tilde H_i(U)\cong\tilde H^{n-i-1}(S^{n-1})$ であり、右辺は $n-i-1=n-1$、すなわち $i=0$ のときのみ $\Z$ である。
よって $\tilde H_0(U)\cong\Z$、つまり $H_0(U)\cong\Z^2$ で、$U$ はちょうど2つの弧状連結成分 $U_1,U_2$ をもつ。
$U$ は $S^n$ の開集合なので局所弧状連結であり、連結成分と弧状連結成分は一致する。
また $i\geq1$ では
\[
0=H_i(U)\cong H_i(U_1)\oplus H_i(U_2)
\]
なので $H_i(U_1)=H_i(U_2)=0$ である。$H_0(U_j)\cong\Z$ と合わせて、各成分は1点と同じホモロジーをもつ。
（なお $n=3$ のアレクサンダーの角付き球面のように、成分が球体と同相とは限らない。）
\end{solution}

\begin{problem}[難易度 D]
$M$ をコンパクトで向き付け可能な境界付き $n$ 次元多様体とする。レフシェッツ双対性
$H^k(M,\partial M;\mathbb{Q})\cong H_{n-k}(M;\mathbb{Q})$ を用いて
\[
\chi(\partial M)=\bigl(1-(-1)^n\bigr)\chi(M)
\]
を示せ。これを用いて、$\mathbb{CP}^2$ はコンパクトな向き付け可能 $5$ 次元多様体の境界にならないことを示せ。
ただし $M$、$\partial M$ のホモロジーは有限生成であることを用いてよい。
\end{problem}

\begin{solution}
有理数係数で考え、$\chi(X,A)=\sum_k(-1)^k\dim H_k(X,A;\mathbb{Q})$ とおく（奇数次元多様体のオイラー標数の問題と同様に、$\mathbb{Q}$ 係数の次元の交代和は整数係数の階数の交代和に一致する）。
体係数の普遍係数定理より $\dim H^k(M,\partial M;\mathbb{Q})=\dim H_k(M,\partial M;\mathbb{Q})$ なので、レフシェッツ双対性から
\[
\dim H_k(M,\partial M;\mathbb{Q})=\dim H_{n-k}(M;\mathbb{Q}).
\]
よって
\begin{align*}
\chi(M,\partial M)&=\sum_k(-1)^k\dim H_{n-k}(M;\mathbb{Q})\\
&=(-1)^n\sum_j(-1)^j\dim H_j(M;\mathbb{Q})=(-1)^n\chi(M).
\end{align*}
対 $(M,\partial M)$ のホモロジー長完全列は有限次元ベクトル空間の有界な完全列なので、次元の交代和は $0$ であり
\[
\chi(\partial M)-\chi(M)+\chi(M,\partial M)=0
\]
を得る。したがって $\chi(\partial M)=\chi(M)-(-1)^n\chi(M)=(1-(-1)^n)\chi(M)$ である。

$\mathbb{CP}^2=\partial M$（$M$ はコンパクト向き付け可能な5次元多様体）と仮定すると、
$n=5$ より $\chi(\mathbb{CP}^2)=2\chi(M)$ は偶数でなければならない。
しかし $\chi(\mathbb{CP}^2)=1+1+1=3$ は奇数であり矛盾する。
（$n$ が偶数のときの式 $\chi(\partial M)=0$ は、$\partial M$ が奇数次元閉多様体であることからも従い、奇数次元閉多様体のオイラー標数が $0$ であることと整合する。）
\end{solution}

%================================================
\chapter{ファイバー束・スペクトル系列への導入}
% 15問
%================================================

\begin{summary}
\begin{itemize}
  \item \textbf{ファイバー束}：$p\colon E\to B$ で、各点の近傍 $U$ について $p^{-1}(U)\cong U\times F$（ファイバーを保つ同相）となるもの。例：被覆空間、メビウスの帯、ホップ束 $S^1\to S^3\to S^2$、$S^1\to S^{2n+1}\to\mathbb{CP}^n$。
  \item \textbf{ファイブレーション}：任意の空間に対してホモトピー持ち上げ性質をもつ写像。底空間がパラコンパクトなファイバー束はファイブレーション。パス・ループファイブレーション $\Omega X\to PX\to X$ では $PX$ が可縮。
  \item \textbf{高次ホモトピー群}：$\pi_n(X,x_0)=[(I^n,\partial I^n),(X,x_0)]$。$n\ge2$ で可換群。$k<n$ で $\pi_k(S^n)=0$、$\pi_n(S^n)\cong\Z$。被覆写像は $n\ge2$ で $\pi_n$ の同型を誘導する。
  \item \textbf{ホモトピー長完全列}：ファイブレーション $F\to E\to B$ に対し
  \[
  \cdots\to\pi_n(F)\to\pi_n(E)\to\pi_n(B)\to\pi_{n-1}(F)\to\cdots\to\pi_0(E)\to0.
  \]
  $S^\infty$ は可縮であり、$\mathbb{CP}^\infty$ は $\pi_2=\Z$ でほかの $\pi_n$ が $0$。
  \item \textbf{セール・スペクトル系列}：$B$ が単連結なファイブレーション $F\to E\to B$ に対し、$E_2^{p,q}=H^p(B;H^q(F))\Rightarrow H^{p+q}(E)$。微分は $d_r\colon E_r^{p,q}\to E_r^{p+r,q-r+1}$ で、カップ積についてライプニッツ則 $d_r(xy)=d_r(x)y+(-1)^{p+q}x\,d_r(y)$ をみたす。
  \item \textbf{ギシン完全列}：ファイバー $S^{k-1}$ の向き付け可能な球面束 $E\to B$ に対し、オイラー類 $e\in H^k(B)$ を用いて
  \[
  \cdots\to H^i(B)\xrightarrow{\smile e}H^{i+k}(B)\xrightarrow{p^*}H^{i+k}(E)\to H^{i+1}(B)\to\cdots
  \]
  \item \textbf{ルレイ・ヒルシュの定理}：$c_j\in H^{k_j}(E;R)$ の各ファイバーへの制限が $H^*(F;R)$ の（有限生成自由）基底を与えるなら、$H^*(E;R)\cong H^*(B;R)\otimes_R H^*(F;R)$（加群として）。
\end{itemize}
\end{summary}

\begin{problem}[難易度 A]
$B$ を連結な空間、$p\colon E\to B$ を被覆写像とする。$p$ は離散空間をファイバーとするファイバー束であることを示せ。
また $p\colon\R\to S^1$、$p(t)=e^{2\pi it}$ について、局所自明化を具体的に書け。
\end{problem}

\begin{solution}
各 $b\in B$ に対し、$b$ の均等被覆開近傍 $U$ をとると $p^{-1}(U)=\bigsqcup_{\alpha\in A}V_\alpha$（$V_\alpha$ は互いに交わらない開集合で、
$p|_{V_\alpha}\colon V_\alpha\to U$ は同相）と書ける。$A$ に離散位相を入れ、
\[
\phi\colon U\times A\to p^{-1}(U),\qquad\phi(u,\alpha)=(p|_{V_\alpha})^{-1}(u)
\]
と定めると、$\phi$ は同相で $p\circ\phi$ は第1射影である。よって $U$ 上で $p$ は自明である。
ファイバーの濃度 $\lvert p^{-1}(b)\rvert$ は均等被覆近傍上で一定なので $B$ 上の局所定数関数であり、
$B$ が連結だから一定である。よって全ファイバーは同一の離散空間 $F$ と同相で、$p$ はファイバー $F$ のファイバー束である。

$\R\to S^1$ の場合、$U_1=S^1-\{1\}$、$U_2=S^1-\{-1\}$ とおく。
$\theta_1\colon U_1\to(0,1)$、$\theta_2\colon U_2\to(-\tfrac12,\tfrac12)$ を $e^{2\pi i\theta_j(z)}=z$ となる連続な偏角とし、
\[
\phi_j\colon U_j\times\Z\to p^{-1}(U_j),\qquad\phi_j(z,m)=\theta_j(z)+m
\]
とおけば、これらが局所自明化（ファイバー $\Z$）を与える。
\end{solution}

\begin{problem}[難易度 A]
次を示せ。
\begin{enumerate}
\item 射影 $\mathrm{pr}\colon B\times F\to B$ は任意の空間に関してホモトピー持ち上げ性質をもつ（ファイブレーションである）。
\item $p\colon E\to B$ がファイブレーションで、$B$ が弧状連結、$E\neq\emptyset$ ならば $p$ は全射である。
\end{enumerate}
\end{problem}

\begin{solution}
(1) 空間 $X$、ホモトピー $G\colon X\times[0,1]\to B$ と $G_0=G(-,0)$ の持ち上げ $\tilde g_0\colon X\to B\times F$ が与えられたとする。
$\mathrm{pr}\circ\tilde g_0=G_0$ より $\tilde g_0(x)=(G(x,0),f(x))$（$f\colon X\to F$ は連続）と書ける。
\[
\tilde G(x,t)=(G(x,t),f(x))
\]
とおくと $\tilde G$ は連続で、$\mathrm{pr}\circ\tilde G=G$、$\tilde G(-,0)=\tilde g_0$ を満たす。よって持ち上げが存在する。

(2) $e\in E$ をとり $b=p(e)$ とする。任意の $b'\in B$ に対し、$b$ から $b'$ への道 $\gamma\colon[0,1]\to B$ がある。
$X$ を1点とするホモトピー持ち上げ性質を $\gamma$ と初期値 $e$ に適用すると、$\tilde\gamma(0)=e$、$p\circ\tilde\gamma=\gamma$ となる道 $\tilde\gamma$ が得られる。
すると $p(\tilde\gamma(1))=\gamma(1)=b'$ なので $b'\in p(E)$ であり、$p$ は全射である。
\end{solution}

\begin{problem}[難易度 B]
メビウスの帯を
\[
E=[0,1]\times[-1,1]\big/(0,t)\sim(1,-t),\qquad p\colon E\to S^1,\ p([s,t])=e^{2\pi is}
\]
で定める。$p$ がファイバー $[-1,1]$ のファイバー束であることを示し、それが自明束（$S^1\times[-1,1]$ とファイバーを保って同相）でないことを示せ。
\end{problem}

\begin{solution}
$U_1=p(\{0<s<1\})=S^1-\{1\}$ 上では $p^{-1}(U_1)=(0,1)\times[-1,1]$ であり自明である。
$U_2=p(\{s\neq\tfrac12\})=S^1-\{-1\}$ 上では、$s\in(\tfrac12,1]$ の部分を $s-1\in(-\tfrac12,0]$ にずらし、そのとき $t$ を $-t$ に置き換える写像
\[
[s,t]\mapsto\begin{cases}(s,t)&(0\leq s<\tfrac12)\\(s-1,-t)&(\tfrac12<s\leq1)\end{cases}
\]
が同値関係と整合し、同相 $p^{-1}(U_2)\cong(-\tfrac12,\tfrac12)\times[-1,1]$ を与える。よって $p$ はファイバー束である。

\begin{center}
\begin{tikzpicture}[scale=1.1]
  \fill[blue!6] (0,-1) rectangle (4,1);
  \draw[very thick,green!50!black] (0,1) -- (4,1);
  \draw[very thick,orange!90!black] (0,-1) -- (4,-1);
  \draw[thick,red!70!black,mid arrow] (0,-1) -- (0,1);
  \draw[thick,red!70!black,mid arrow] (4,1) -- (4,-1);
  \draw[thick,gray] (1.3,-1) -- (1.3,1) node[above] {\footnotesize $p^{-1}(e^{2\pi is})$};
  \node[hvtx,inner sep=1.8pt] at (4,1) {}; \node[hvtx,inner sep=1.8pt] at (0,-1) {};
  \node[vtx,inner sep=1.6pt] at (4,-1) {}; \node[vtx,inner sep=1.6pt] at (0,1) {};
  \draw[-{Stealth[length=2.2mm]}] (2,-1.35) -- (2,-1.95) node[midway,right] {$p$};
  \draw[thick] (0,-2.3) -- (4,-2.3);
  \node[vtx] at (0,-2.3) {}; \node[vtx] at (4,-2.3) {};
  \node[right] at (4.1,-2.3) {$[0,1]/(0\sim1)\cong S^1$};
\end{tikzpicture}
\zucap{メビウスの帯 $E$ と射影 $p$（縦線分がファイバー）。ファイバーの端点の集合は、上辺（緑）の右端 $(1,1)$ が下辺（橙）の左端 $(0,-1)$ に（白丸）、下辺の右端 $(1,-1)$ が上辺の左端 $(0,1)$ に（黒丸）つながるので、緑→橙→緑と全体で1つの円周になる。自明束 $S^1\times[-1,1]$ では上辺・下辺がそれぞれ別の円周になる。}
\end{center}

自明でないことを示す。区間 $[-1,1]$ の端点は「取り除いても連結なままの点」として特徴づけられるので、区間の自己同相は端点を端点に写す。
ファイバーを保つ同相 $\Phi\colon E\to S^1\times[-1,1]$ があったとすると、$\Phi$ は各ファイバーの端点の集合
\[
\partial_fE=\{[s,\pm1]\}\subset E
\]
を $S^1\times\{\pm1\}$ に同相に写す。後者は2つの連結成分をもつ。しかし $\partial_fE$ は2本の線分 $[0,1]\times\{1\}$、$[0,1]\times\{-1\}$ を
$(1,1)\sim(0,-1)$、$(1,-1)\sim(0,1)$ で繋いだものであり、1つの円周で連結である。これは矛盾であり、$p$ は自明束ではない。
\end{solution}

\begin{problem}[難易度 B]
$S^{2n+1}=\{z\in\mathbb{C}^{n+1}:\lVert z\rVert=1\}$ とし、$p\colon S^{2n+1}\to\mathbb{CP}^n$、$p(z)=[z]$ とする。
$p$ がファイバー $S^1$ のファイバー束であることを示せ。（$n=1$ のとき $\mathbb{CP}^1\cong S^2$ であり、これがホップ束 $S^3\to S^2$ である。）
\end{problem}

\begin{solution}
$[z]=[w]$ と $w=\lambda z$（$\lambda\in\mathbb{C}^*$）は同値であり、$z,w\in S^{2n+1}$ なら $\lvert\lambda\rvert=1$ である。
よって $p^{-1}([z])=\{\lambda z:\lvert\lambda\rvert=1\}\cong S^1$ である。
$U_j=\{[z]: z_j\neq0\}$（$j=0,\dots,n$）は $\mathbb{CP}^n$ の開被覆である。
\[
\phi_j\colon U_j\times S^1\to p^{-1}(U_j),\qquad\phi_j([z],\lambda)=\lambda\frac{\lvert z_j\rvert}{z_j}\frac{z}{\lVert z\rVert}
\]
と定める。代表元 $z$ を $\mu z$（$\mu\in\mathbb{C}^*$）に替えると右辺は $\frac{\lvert\mu\rvert\mu}{\mu\lvert\mu\rvert}=1$ 倍しか変わらないので well-defined であり、
連続で $p\circ\phi_j([z],\lambda)=[z]$ を満たす。逆写像は
\[
\psi_j\colon p^{-1}(U_j)\to U_j\times S^1,\qquad\psi_j(w)=\Bigl([w],\frac{w_j}{\lvert w_j\rvert}\Bigr)
\]
で与えられる。実際 $w\in S^{2n+1}$ に対し $z=w$ を代表元に選べば $\phi_j(\psi_j(w))=\frac{w_j}{\lvert w_j\rvert}\frac{\lvert w_j\rvert}{w_j}w=w$、
また $\phi_j([z],\lambda)$ の第 $j$ 成分は $\lambda\lvert z_j\rvert/\lVert z\rVert$ なので、その偏角は $\lambda$ であり $\psi_j\circ\phi_j=\id$ である。
よって $\phi_j$ は局所自明化である。
\end{solution}

\begin{problem}[難易度 B]
$n\geq2$ のとき、ホモトピー群 $\pi_n(X,x_0)$ は可換群であることを示せ。
\end{problem}

\begin{solution}
$\pi_n(X,x_0)$ の元を写像 $f\colon(I^n,\partial I^n)\to(X,x_0)$ のホモトピー類で表す。$i=1,2$ に対し、第 $i$ 座標で連結した和を
\[
(f+_1g)(s)=\begin{cases}f(2s_1,s_2,\dots,s_n)&(s_1\leq\tfrac12)\\g(2s_1-1,s_2,\dots,s_n)&(s_1\geq\tfrac12)\end{cases}
\]
と定め、$+_2$ も第2座標で同様に定める（ここで $n\geq2$ を使う）。どちらもホモトピー類上の群演算を定め、定値写像 $c$ の類が単位元である。
定義に代入すれば、写像として
\[
(f+_1g)+_2(h+_1k)=(f+_2h)+_1(g+_2k)
\]
が成り立つ（両辺とも $I^n$ を4つの小正方形に分け、左下・右下・左上・右上にそれぞれ $f,g,h,k$ を置いた写像である）。
これを用いると、ホモトピー類について
\begin{align*}
[f]+_1[g]&=([f]+_2[c])+_1([c]+_2[g])=([f]+_1[c])+_2([c]+_1[g])=[f]+_2[g],\\
[f]+_1[g]&=([c]+_2[f])+_1([g]+_2[c])=([c]+_1[g])+_2([f]+_1[c])=[g]+_2[f]
\end{align*}
を得る。よって $[f]+_1[g]=[f]+_2[g]=[g]+_2[f]=[g]+_1[f]$ となり、2つの演算は一致して可換である。
\end{solution}

\begin{problem}[難易度 B]
$p\colon\tilde X\to X$ を被覆写像とし、$\tilde x_0\in\tilde X$、$x_0=p(\tilde x_0)$ とする。
ファイブレーションのホモトピー長完全列を用いて、$n\geq2$ で $p_*\colon\pi_n(\tilde X,\tilde x_0)\to\pi_n(X,x_0)$ が同型であることを示せ。
これを用いて $\pi_n(S^1)=0$、$\pi_n(\mathbb{RP}^m)\cong\pi_n(S^m)$（$n\geq2$）を示せ。
\end{problem}

\begin{solution}
被覆写像は任意の空間に関してホモトピー持ち上げ性質をもつのでファイブレーションであり、ファイバーは離散空間 $F=p^{-1}(x_0)$ である。長完全列
\[
\pi_n(F,\tilde x_0)\to\pi_n(\tilde X,\tilde x_0)\xrightarrow{p_*}\pi_n(X,x_0)\to\pi_{n-1}(F,\tilde x_0)
\]
を考える。$F$ は離散なので、$k\geq1$ に対し連結な $S^k$ からの基点付き写像は定値写像に限り、$\pi_k(F,\tilde x_0)=0$ である。
$n\geq2$ なら $n-1\geq1$ なので両端が $0$ となり、$p_*$ は同型である。

$p\colon\R\to S^1$ は被覆で $\R$ は可縮なので、$n\geq2$ で $\pi_n(S^1)\cong\pi_n(\R)=0$ である。
$m\geq1$ に対し商写像 $S^m\to\mathbb{RP}^m$ は2重被覆なので、$n\geq2$ で $\pi_n(S^m)\cong\pi_n(\mathbb{RP}^m)$ である。
（例えば $\pi_2(\mathbb{RP}^2)\cong\pi_2(S^2)$。）
\end{solution}

\begin{problem}[難易度 B]
$F\to E\xrightarrow{p}B$ をセールファイブレーション（$B$ は弧状連結）とし、$E$ が可縮であるとする。
$n\geq2$ で $\pi_n(B)\cong\pi_{n-1}(F)$ であることを示せ。これを次の2つに適用せよ。
\begin{enumerate}
\item ファイバー束 $S^1\to S^\infty\to\mathbb{CP}^\infty$（$S^\infty$ は可縮）から $\pi_*(\mathbb{CP}^\infty)$ を求めよ。
\item パス・ループファイブレーション $\Omega X\to PX\to X$ から $\pi_n(\Omega X)\cong\pi_{n+1}(X)$（$n\geq1$）を導け。
\end{enumerate}
\end{problem}

\begin{solution}
長完全列 $\pi_n(E)\to\pi_n(B)\xrightarrow{\partial}\pi_{n-1}(F)\to\pi_{n-1}(E)$ において、$E$ は可縮なので
$n\geq2$ では両端が $0$ であり（$n-1\geq1$）、$\partial$ は同型である。

(1) $n\geq2$ で $\pi_n(\mathbb{CP}^\infty)\cong\pi_{n-1}(S^1)$ であり、これは $n=2$ で $\Z$、$n\geq3$ で $0$ である。
また $\pi_1(\mathbb{CP}^\infty)$ については、完全列 $\pi_1(S^\infty)\to\pi_1(\mathbb{CP}^\infty)\to\pi_0(S^1)$ において
左端が自明、右端が1点なので $\pi_1(\mathbb{CP}^\infty)=0$ である。（$\mathbb{CP}^\infty$ の1骨格が1点であることからもわかる。）
よって $\mathbb{CP}^\infty$ の非自明なホモトピー群は $\pi_2\cong\Z$ のみである。

(2) $PX$ は $x_0$ を始点とする道の空間、$p(\gamma)=\gamma(1)$ で、ファイバーは $\Omega X$ である（$X$ は弧状連結とする）。
$PX$ は $H(\gamma,u)(t)=\gamma((1-u)t)$ により定値道へ縮約するので可縮である。よって $n\geq1$ に対し $n+1\geq2$ だから
$\pi_{n+1}(X)\cong\pi_n(\Omega X)$ である。
\end{solution}

\begin{problem}[難易度 B]
ホップ束 $S^1\to S^3\xrightarrow{p}S^2$ のホモトピー長完全列を用いて次を示せ。
\begin{enumerate}
\item $\pi_2(S^2)\cong\Z$。
\item $n\geq3$ で $p_*\colon\pi_n(S^3)\to\pi_n(S^2)$ は同型。特に $\pi_3(S^2)\cong\Z$ であり、ホップ写像 $p$ の類が生成元である。
\end{enumerate}
ただし $i<k$ で $\pi_i(S^k)=0$、$\pi_k(S^k)\cong\Z$（$[\id]$ が生成元）であることは用いてよい。
\end{problem}

\begin{solution}
長完全列の一部
\[
\pi_n(S^1)\to\pi_n(S^3)\xrightarrow{p_*}\pi_n(S^2)\xrightarrow{\partial}\pi_{n-1}(S^1)\to\pi_{n-1}(S^3)
\]
を用いる。$k\geq2$ で $\pi_k(S^1)=0$ である（被覆 $\R\to S^1$ による）。

(1) $n=2$ とすると
\[
0=\pi_2(S^3)\to\pi_2(S^2)\xrightarrow{\partial}\pi_1(S^1)\to\pi_1(S^3)=0
\]
であるから $\partial$ は同型で、$\pi_2(S^2)\cong\pi_1(S^1)\cong\Z$ である。

(2) $n\geq3$ なら $\pi_n(S^1)=0$ かつ $\pi_{n-1}(S^1)=0$（$n-1\geq2$）なので $p_*$ は同型である。
$n=3$ とすると $\pi_3(S^2)\cong\pi_3(S^3)\cong\Z$ であり、生成元 $[\id_{S^3}]$ の像 $p_*[\id]=[p]$ が $\pi_3(S^2)$ の生成元である。
特にホップ写像は定値写像とホモトピックでない。
\end{solution}

\begin{problem}[難易度 B]
$n\geq1$ とする。ファイバー束 $S^1\to S^{2n+1}\to\mathbb{CP}^n$ を用いて、
$\pi_1(\mathbb{CP}^n)=0$、$\pi_2(\mathbb{CP}^n)\cong\Z$、$3\leq k\leq 2n$ で $\pi_k(\mathbb{CP}^n)=0$、$\pi_{2n+1}(\mathbb{CP}^n)\cong\Z$ を示せ。
\end{problem}

\begin{solution}
長完全列
\[
\pi_k(S^1)\to\pi_k(S^{2n+1})\to\pi_k(\mathbb{CP}^n)\xrightarrow{\partial}\pi_{k-1}(S^1)\to\pi_{k-1}(S^{2n+1})
\]
を用いる。$2n+1\geq3$ なので $\pi_1(S^{2n+1})=\pi_2(S^{2n+1})=0$ である。

$k=1$：$\pi_1(S^{2n+1})\to\pi_1(\mathbb{CP}^n)\xrightarrow{\partial}\pi_0(S^1)$ の完全性より、左の写像の像は $\partial$ の核に一致する。
$\pi_0(S^1)$ は1点なので $\partial$ の核は $\pi_1(\mathbb{CP}^n)$ 全体であり、左端は自明なので $\pi_1(\mathbb{CP}^n)=0$。

$k=2$：$0=\pi_2(S^{2n+1})\to\pi_2(\mathbb{CP}^n)\to\pi_1(S^1)\to\pi_1(S^{2n+1})=0$ より $\pi_2(\mathbb{CP}^n)\cong\Z$。

$k\geq3$：$\pi_k(S^1)=\pi_{k-1}(S^1)=0$ より $\pi_k(\mathbb{CP}^n)\cong\pi_k(S^{2n+1})$ である。
これは $3\leq k\leq2n$ で $0$（$k<2n+1$）、$k=2n+1$ で $\Z$ である。
\end{solution}

\begin{problem}[難易度 C]
$S^2$ と $S^3\times\mathbb{CP}^\infty$ は、すべての $k\geq1$ で同型なホモトピー群をもつが、ホモトピー同値ではないことを示せ。
ただし $\pi_k(X\times Y)\cong\pi_k(X)\times\pi_k(Y)$ は用いてよい。
\end{problem}

\begin{solution}
既出の問題より $\pi_*(\mathbb{CP}^\infty)$ は $\pi_2\cong\Z$ のみが非自明である。よって
\[
\pi_k(S^3\times\mathbb{CP}^\infty)\cong\pi_k(S^3)\times\pi_k(\mathbb{CP}^\infty)\cong
\begin{cases}0&(k=1)\\\Z&(k=2)\\\pi_k(S^3)&(k\geq3)\end{cases}
\]
である（$\pi_1(S^3)=\pi_2(S^3)=0$）。一方 $\pi_1(S^2)=0$、$\pi_2(S^2)\cong\Z$ であり、ホップ束により $k\geq3$ で $\pi_k(S^2)\cong\pi_k(S^3)$ である。
よって全ての $k$ でホモトピー群は同型である。

しかしホモロジーは異なる。$\mathbb{CP}^\infty$ のホモロジーは偶数次で $\Z$、奇数次で $0$ であり、すべて自由なのでキュネット公式より
\[
H_3(S^3\times\mathbb{CP}^\infty)\cong\bigoplus_{i+j=3}H_i(S^3)\otimes H_j(\mathbb{CP}^\infty)=H_3(S^3)\otimes H_0(\mathbb{CP}^\infty)\cong\Z
\]
である。一方 $H_3(S^2)=0$ なので、両者はホモトピー同値でない。
（ホワイトヘッドの定理は、ホモトピー群の同型が「ある写像によって」誘導されることを要求する点に注意する。）
\end{solution}

\begin{problem}[難易度 C]
ホップ束 $S^1\to S^3\to S^2$ のセール・スペクトル系列（整数係数のコホモロジー）を考える。
$E_2$ 項を書き下し、$H^*(S^3)$ が既知であることから微分 $d_2\colon E_2^{0,1}\to E_2^{2,0}$ が同型であることを示せ。
また $E_\infty$ 項が $H^*(S^3)$ と整合することを確かめよ。
\end{problem}

\begin{solution}
$S^2$ は単連結なので局所係数は自明であり、$E_2^{p,q}=H^p(S^2;H^q(S^1;\Z))$ である。
$H^q(S^1)$ は $q=0,1$ で $\Z$、$H^p(S^2;\Z)$ は $p=0,2$ で $\Z$ なので
\[
E_2^{p,q}\cong\Z\quad((p,q)\in\{(0,0),(2,0),(0,1),(2,1)\}),\qquad\text{それ以外は }0
\]
である。$d_r\colon E_r^{p,q}\to E_r^{p+r,q-r+1}$ のうち非零の可能性があるのは $d_2\colon E_2^{0,1}\to E_2^{2,0}$ のみであり、
$r\geq3$ では始点か終点が $0$ なので $E_3=E_\infty$ である。

$E_\infty^{0,1}=\ker d_2$、$E_\infty^{2,0}=\operatorname{coker}d_2$ である。$H^1(S^3)$ は $E_\infty^{0,1},E_\infty^{1,0}$ を、
$H^2(S^3)$ は $E_\infty^{2,0},E_\infty^{1,1},E_\infty^{0,2}$ を段階商にもつフィルトレーションをもつ。
$H^1(S^3)=H^2(S^3)=0$ より $\ker d_2=0$、$\operatorname{coker}d_2=0$ であり、$d_2\colon\Z\to\Z$ は同型である。

このとき $E_\infty$ で生き残るのは $E_\infty^{0,0}\cong\Z$ と $E_\infty^{2,1}\cong\Z$ のみで、それぞれ全次数 $0,3$ にある。
よって $H^0\cong\Z$、$H^3\cong\Z$、他は $0$ となり $H^*(S^3)$ と整合する。
（対照的に自明束 $S^2\times S^1$ では $d_2=0$ で、$E_\infty=E_2$ がキュネット公式を再現する。）
\end{solution}

\begin{problem}[難易度 C]
ギシン完全列を用いて、環として $H^*(\mathbb{CP}^n;\Z)\cong\Z[e]/(e^{n+1})$（$\deg e=2$）であることを示せ（$n\geq1$）。
ただし $\mathbb{CP}^n$ が $2n$ 次元 CW 複体であることは用いてよい。
\end{problem}

\begin{solution}
$S^1\to S^{2n+1}\xrightarrow{p}\mathbb{CP}^n$ は向き付け可能な円周束（底空間が単連結）であり、オイラー類 $e\in H^2(\mathbb{CP}^n)$ をもつ。
$B=\mathbb{CP}^n$、$E=S^{2n+1}$ と書くと、ギシン完全列は
\[
\cdots\to H^{i-1}(E)\to H^{i-2}(B)\xrightarrow{\smile e}H^i(B)\xrightarrow{p^*}H^i(E)\to H^{i-1}(B)\to\cdots
\]
である（負の次数のコホモロジーは $0$）。$H^j(S^{2n+1})=0$（$1\leq j\leq2n$）に注意する。

$i=1$：$0=H^{-1}(\mathbb{CP}^n)\to H^1(\mathbb{CP}^n)\to H^1(S^{2n+1})=0$ より $H^1(\mathbb{CP}^n)=0$。

$2\leq i\leq2n$：$H^{i-1}(S^{2n+1})=0$ かつ $H^i(S^{2n+1})=0$ なので、$\smile e\colon H^{i-2}(\mathbb{CP}^n)\to H^i(\mathbb{CP}^n)$ は同型である。

$H^0(\mathbb{CP}^n)=\Z\cdot1$ から始めて、$0\leq k\leq n$ で $H^{2k}(\mathbb{CP}^n)=\Z e^k$、また $H^1=0$ から奇数次 $\leq2n-1$ はすべて $0$ である。
$2n$ 次元 CW 複体なので $i>2n$ で $H^i=0$、特に $e^{n+1}=0$ である。
よって $\Z[e]/(e^{n+1})\to H^*(\mathbb{CP}^n)$ は環同型であり、特に $e$ は $H^2$ の生成元である。
\end{solution}

\begin{problem}[難易度 D]
$n\geq2$ とする。パス・ループファイブレーション $\Omega S^n\to PS^n\to S^n$ のセール・スペクトル系列を用いて
\[
H^q(\Omega S^n;\Z)\cong\begin{cases}\Z&((n-1)\mid q,\ q\geq0)\\0&(\text{それ以外})\end{cases}
\]
を示せ。
\end{problem}

\begin{solution}
$S^n$ は単連結（$n\geq2$）なので $E_2^{p,q}=H^p(S^n;H^q(\Omega S^n))$ であり、これは $p=0,n$ のときのみ $H^q(\Omega S^n)$ に同型で、他は $0$ である。
非零の列は $p=0,n$ の2本だけなので、非零になりうる微分は
\[
d_n\colon E_n^{0,q}\to E_n^{n,q-n+1}
\]
のみであり、$E_2=E_n$、$E_{n+1}=E_\infty$ である。$PS^n$ は可縮なので、$E_\infty^{p,q}=0$（$(p,q)\neq(0,0)$）である。

$q\geq1$ に対し $E_\infty^{0,q}=\ker d_n=0$ である。また $q'\geq0$ に対し $E_\infty^{n,q'}=\operatorname{coker}(d_n\colon E_n^{0,q'+n-1}\to E_n^{n,q'})=0$ である。
$q=q'+n-1\geq n-1\geq1$ なので、両者を合わせて
\[
d_n\colon H^q(\Omega S^n)\xrightarrow{\ \cong\ }H^{q-n+1}(\Omega S^n)\quad(q\geq n-1)
\]
を得る。さらに $1\leq q\leq n-2$ では $d_n$ の行き先 $E^{n,q-n+1}$ は $q-n+1<0$ より $0$ なので、$\ker d_n=H^q(\Omega S^n)$ が $0$ でなければならない。
$\Omega S^n$ は弧状連結（$\pi_0(\Omega S^n)=\pi_1(S^n)=0$）なので $H^0\cong\Z$ である。
以上から $q$ に関する帰納法で、$H^q\cong H^{q-(n-1)}$ を繰り返し使えば、$q$ が $n-1$ の倍数なら $\Z$、そうでなければ $0$ となる。
\end{solution}

\begin{problem}[難易度 D]
クラインの壺 $K=[0,1]\times S^1/(0,z)\sim(1,\bar z)$ と、$p([s,z])=e^{2\pi is}$ で定まるファイバー $S^1$ のファイバー束 $p\colon K\to S^1$ を考える。
\begin{enumerate}
\item $\Z/2$ 係数ではルレイ・ヒルシュの定理の仮定が満たされることを示し、$H^*(K;\Z/2)$ の各次元を求めよ。
\item $\Z$ 係数では仮定が満たされず、実際に $H^*(K;\Z)\not\cong H^*(S^1;\Z)\otimes H^*(S^1;\Z)$ であることを示せ。
\end{enumerate}
\end{problem}

\begin{solution}
ルレイ・ヒルシュの定理は「ファイバー $F$ の $H^*(F;R)$ が各次数で有限生成自由 $R$ 加群で、各ファイバーへの制限が $H^*(F;R)$ の基底となる類 $c_j\in H^*(E;R)$ があれば、
$H^*(E;R)\cong H^*(B;R)\otimes_RH^*(F;R)$（加法的）」を主張する。底空間は弧状連結なので、1つのファイバーで確かめればよい。

$K$ は 0 胞体 $v$（$(s,z)=(0,1)$ の類）、1 胞体 $a=\{[s,1]\}$、$b=\{[0,z]\}$（基点上のファイバー $F$）、2 胞体 $e$ をもつ CW 複体で、
$e$ の境界の語は $ab^{-1}a^{-1}b^{-1}$ となるので、胞体境界は $d(e)=-2b$、$d(a)=d(b)=0$ である。
よって $H_1(K;\Z)=\Z a\oplus\Z/2\,b$、$H_2(K;\Z)=0$ である。

(1) $\Z/2$ 係数では $d(e)=0$ なので $H_1(K;\Z/2)=(\Z/2)^2$ で、$i_*[F]=b\neq0$ である。
体係数の普遍係数定理により $i^*\colon H^1(K;\Z/2)\to H^1(F;\Z/2)$ は $i_*$ の双対なので全射であり、
$i^*c$ が生成元となる $c$ がとれる。$1\in H^0(K)$ と合わせて仮定が満たされ、
$\dim H^k(K;\Z/2)$ は $H^*(S^1;\Z/2)\otimes H^*(S^1;\Z/2)$ と同じく $k=0,1,2$ で $1,2,1$ である。

(2) $H^1(K;\Z)\cong\operatorname{Hom}(H_1(K),\Z)$ であり、$\varphi$ の制限は $\varphi(b)$ で決まる。$b$ は位数2なので $\varphi(b)=0$、
すなわち $i^*\colon H^1(K;\Z)\to H^1(F;\Z)\cong\Z$ は零写像で、仮定は満たされない。
実際、普遍係数定理より $H^2(K;\Z)\cong\operatorname{Hom}(H_2(K),\Z)\oplus\operatorname{Ext}(H_1(K),\Z)\cong\Z/2$ であり、
$H^*(S^1)\otimes H^*(S^1)$ の2次の部分 $\Z$ とは異なる。
\end{solution}

\begin{problem}[難易度 D]
ファイバー束 $S^1\to S^\infty\to\mathbb{CP}^\infty$（$S^\infty$ は可縮）のセール・スペクトル系列とその乗法構造を用いて、
環として $H^*(\mathbb{CP}^\infty;\Z)\cong\Z[x]$（$\deg x=2$）であることを示せ。
\end{problem}

\begin{solution}
$B=\mathbb{CP}^\infty$ は単連結（前出）であり、$H^*(S^1)$ は自由なので $E_2^{p,q}\cong H^p(B)\otimes H^q(S^1)$ で、
これは $q=0,1$ の2行のみ非零である。$u\in H^1(S^1)$ を生成元とすると $E_2^{p,1}=H^p(B)\cdot u$ である。
非零になりうる微分は $d_2\colon E_2^{p,1}\to E_2^{p+2,0}$ のみなので $E_3=E_\infty$ であり、
$S^\infty$ が可縮だから $(p,q)\neq(0,0)$ で $E_\infty^{p,q}=0$ である。したがって
\begin{itemize}
\item $E_2^{1,0}=H^1(B)$ には何も入ってこない（$E_2^{-1,1}=0$）ので、$H^1(B)=E_\infty^{1,0}=0$。
\item $p\geq0$ で $d_2\colon E_2^{p,1}\to E_2^{p+2,0}$ は単射かつ全射、すなわち同型。
\end{itemize}
$x=d_2(u)\in E_2^{2,0}=H^2(B)$ とおくと、$p=0$ の場合から $x$ は $H^2(B)\cong H^0(B)\cong\Z$ の生成元である。

乗法構造を使う。$y\in H^p(B)=E_2^{p,0}$ に対し $d_2y\in E_2^{p+2,-1}=0$ なので、ライプニッツ則より
\[
d_2(y\cdot u)=d_2(y)\cdot u+(-1)^py\cdot d_2(u)=(-1)^py\smile x.
\]
$d_2\colon E_2^{p,1}=H^p(B)u\to H^{p+2}(B)$ が同型なので、$\smile x\colon H^p(B)\to H^{p+2}(B)$ は同型である。
$H^0=\Z$、$H^1=0$ から始めて、$H^{2k}(B)=\Z x^k$、$H^{2k+1}(B)=0$ を得る。
よって $\Z[x]\to H^*(\mathbb{CP}^\infty)$ は環同型である。
\end{solution}

\printsolutions

\backmatter

\chapter*{参考文献}
\begin{enumerate}
  \item A. Hatcher, Algebraic Topology.
  \item J. M. Munkres, Elements of Algebraic Topology.
  \item E. H. Spanier, Algebraic Topology.
\end{enumerate}

\end{document}